\section{Overview of
{\em Ricci Flow with Surgery on
Three-Manifolds}  \cite{Perelman2}} \label{overview2}

The paper \cite{Perelman2} is concerned with the Ricci flow on compact oriented
$3$-manifolds.  The main difference with respect to \cite{Perelman}
is that singularity formation is allowed, so the paper deals with
a ``Ricci flow with surgery''.

The main part of the paper is concerned with
setting up the surgery procedure and showing that it is
well-defined, in the sense that surgery times do not
accumulate.  In addition, the long-time behavior of a
Ricci flow with surgery is analyzed.

The paper can be divided into three main parts.
Sections II.1-II.3 contain preparatory material about ancient
solutions, the so-called standard solution and the geometry
at the first singular time.  Sections II.4-II.5 set up the
surgery procedure and prove that it is well-defined.
Sections II.6-II.8 analyze the long-time behavior.

\subsection{II.1-II.3}

Section II.1 continues the analysis of three-dimensional
$\kappa$-solutions from I.11.  From I.11, any $\kappa$-solution
contains an ``asymptotic soliton'', a gradient shrinking
soliton that arises as a rescaled limit of the $\kappa$-solution
as $t \rightarrow - \infty$. It is shown that any such
gradient shrinking soliton must be a shrinking round cylinder $\R \times S^2$,
its $\Z_2$-quotient $\R \times_{\Z_2} S^2$
or a finite quotient of the
round shrinking $S^3$. Using this, one
obtains a finer description of the $\kappa$-solutions. In
particular, any compact $\kappa$-solution must be isometric
to a finite quotient of the round shrinking $S^3$, or
diffeomorphic to $S^3$ or $\R P^3$. It is shown that there is
a universal number $\kappa_0 > 0$ so that any $\kappa$-solution
is a finite quotient of the round shrinking $S^3$ or is a
$\kappa_0$-solution. This implies universal
derivative bounds on the scalar curvature of a $\kappa$-solution.

Section II.2 defines and analyzes the Ricci flow of the
so-called standard solution.  This is a Ricci flow on $\R^3$ whose
initial metric is a capped-off half cylinder.  The surgery
procedure will amount to gluing in a truncated copy of the
time-zero slice of the standard solution. Hence one
needs to understand the Ricci flow on the standard solution
itself.  It is shown that the Ricci flow of the standard
solution exists on a maximal time interval $[0,1)$, and the
solution goes singular everywhere as $t \rightarrow 1$.

The geometry of the solution at the first singular time $T$
(assuming that there is one) is considered in II.3.  Put
$\Omega \: = \: \{x \in M \: : \: \limsup_{t \rightarrow T^-}
|\Rm(x,t)| < \infty \}$. Then $\Omega$ is an open subset of $M$,
and $x \in M-\Omega$ if and only if $\lim_{t \rightarrow T^-}
R(x,t) = \infty$. If $\Omega = \emptyset$ then for $t$ slightly less than $T$,
the manifold $(M, g(t))$ consists of nothing but high-scalar-curvature
regions. Using Theorem I.12.1, one shows that
$M$ is diffeomorphic to
$S^1 \times S^2$, $\R P^3 \# \R P^3$ or
a finite isometric quotient of $S^3$.

If $\Omega \neq \emptyset$ then there is a well-defined limit
metric $\overline{g}$ on $\Omega$, with scalar curvature
function $\overline{R}$. The set $\Omega$
could {\it a priori} have an infinite number of connected components,
for example if an infinite number of distinct $2$-spheres simultaneously
shrink to points at time $T$. For small $\rho > 0$,
put $\Omega_\rho \: = \: \{x \in \Omega \: : \: \overline{R}(x) \le \rho^{-2} \}$,
a compact subset of $M$.
The connected components of $\Omega$ can be divided into those
that intersect $\Omega_\rho$ and those that do not.  If a connected
component does not intersect $\Omega_\rho$ then it is a
``capped $\epsilon$-horn'' (consisting of a hornlike end capped off by a ball
or a copy of $\R P^3 - B^3$) or a ``double $\epsilon$-horn'' (with two
hornlike ends). If a connected component of $\Omega$ does intersect
$\Omega_\rho$ then it has a finite number of ends, each being
an $\epsilon$-horn.

Topologically, the surgery procedure of II.4 will amount to taking each
connected component of $\Omega$ that intersects $\Omega_\rho$,
truncating each of its $\epsilon$-horns and gluing a $3$-ball onto
each truncated horn.
The connected components of $\Omega$ that do not intersect
$\Omega_\rho$ are thrown away.
Call the new manifold $M^\prime$. At a time
$t$ slightly less than $T$, the region $M - \Omega_\rho$ consists
of high-scalar-curvature regions. Using the characterization of such
regions in I.12.1, one shows that $M$ can be reconstructed from
$M^\prime$ by taking the connected sum of its connected components,
along possibly with a finite number of $S^1 \times S^2$ and
$\R P^3$ factors.

\subsection{II.4-II.5}

Section II.4 defines the surgery procedure.
A Ricci flow with surgery consists of a sequence of smooth $3$-dimensional Ricci flows on
adjacent time intervals with the property that for any two adjacent intervals, there is a
a compact $3$-dimensional submanifold-with-boundary that is common to the
final slice of the first time interval and the initial slice of the second time interval.

There are two {\it a priori} assumptions on a Ricci flow with surgery, the
pinching assumption and the
canonical neighborhood assumption.
The pinching assumption is a form of Hamilton-Ivey pinching. The canonical
neighborhood assumption says that every spacetime
point $(x,t)$ with $R(x,t) \ge r(t)^{-2}$ has a neighborhood which,
after rescaling, is $\epsilon$-close to one of the neighborhoods that occur
in a $\kappa$-solution or in a time slice of the standard solution. Here
$\epsilon$ is a small but universal constant and $r(\cdot)$ is a decreasing
function, which is to be specified.

One wishes to define a Ricci flow with
surgery starting from any compact oriented $3$-manifold, say with a
normalized initial metric.  There are
various parameters that will enter into the definition : the above canonical neighborhood scale
$r(\cdot)$,
a nonincreasing function $\delta(\cdot)$ that decays to zero, the truncation
scale $\rho(t) = \delta(t) r(t)$ and the surgery scale $h$. In order
to show that one can construct the Ricci flow with surgery, it turns out that one wants to perform
the surgery only on necks with a radius that is very small compared to the
canonical neighborhood scale; this is the role of the parameter $\delta(\cdot)$.

Suppose that the Ricci flow with surgery is defined at times less than $T$, with
the {\it a priori} assumptions satisfied, and goes
singular at time $T$. Define the open subset $\Omega \subset M$ as before and
construct the compact subset $\Omega_\rho \subset M$
using $\rho = \rho(T)$. Any connected component $N$ of $\Omega$ that
intersects $\Omega_\rho$ has a finite number of ends, each of which is an
$\epsilon$-horn.  This means that each point in the horn is in the center of
an $\epsilon$-neck, i.e. has a neighborhood
that, after rescaling, is $\epsilon$-close to a cylinder $\left[ - \frac{1}{\epsilon},
\frac{1}{\epsilon} \right] \times S^2$. In II.4.3 it is shown that as one goes down the
end of the horn, there is a self-improvement phenomenon; for any
$\delta > 0$, one can find $h < \delta \rho$ so that if a point $x$ in the horn has
$\overline{R}(x) \ge h^{-2}$ then it is actually in the center of a $\delta$-neck.

With $\delta = \delta(T)$, let
$h$ be the corresponding number. One then
cuts off the $\epsilon$-horn at a $2$-sphere in the center of such a
$\delta$-neck and glues in a copy of a rescaled truncated standard solution.
One does this for each $\epsilon$-horn in $N$ and each connected component
$N$ that intersects $\Omega_\rho$, and throws away the connected components
of $\Omega$ that do not intersect $\Omega_\rho$. One lets the new manifold evolve
under the Ricci flow. If one encounters another singularity then one again
performs surgery.  Based on an estimate on the volume change under a surgery,
one concludes that a finite number of surgeries occur in any finite time
interval. (However,
one is not able to conclude from volume arguments that there is a finite
number of surgeries altogether.)

The preceding discussion was predicated on the condition that the
{\it a priori} assumptions hold for all times.  For the Ricci flow
before the first surgery time, the pinching condition
follows from the Hamilton-Ivey result.  One shows
that surgery can be performed so that it does not make the pinching any worse.
Then the pinching condition will hold up to the second surgery time, etc.
The main issue is to show that one can choose the parameters
$r(\cdot)$ and $\delta(\cdot)$ so that one knows {\it a priori} that the canonical
neighborhood assumption, with parameter $r(\cdot)$, will hold for the
Ricci flow with surgery. (For any singularity time $T$, one needs to know that the
canonical neighborhood assumption holds for $t \in [0, T)$ in order to do the
surgery at time $T$.)

As a preliminary step, in Lemma II.4.5 it is shown that after one glues in a standard
solution, the result will still look similar to a standard solution, for as long of
a time interval as one could expect, unless the entire region gets removed by some
exterior surgery.

The result of II.5 is that the time-dependent
parameters $r(\cdot)$ and $\delta(\cdot)$ can be
chosen so as to ensure that the {\it a priori} assumptions hold.
The normalization of the initial metric implies that there is a time
interval $[0, C]$, for a universal constant $C$, on which the
Ricci flow is smooth and has explicitly bounded curvature.
On this time interval the canonical neighborhood assumption
holds vacuously, if $r(\cdot) \big|_{[0,C]}$ is sufficiently small.
To handle later times, the strategy is to divide $[\epsilon, \infty)$ into
a countable sequence of finite time intervals and proceed by induction.
In II.5 the intervals $\{[2^{j-1} \epsilon, 2^j \epsilon]\}_{j=1}^\infty$ are used,
although the precise choice of intervals is immaterial.

We recall from I.12.1 that
in the case of smooth flows, the proof of the canonical neighborhood assumption
used the fact that one has $\kappa$-noncollapsing. It is not immediate that the
method of proof of I.12.1 extends to a Ricci flow with surgery.
(It is exactly for this reason that one takes $\delta(\cdot)$ to be a
time-dependent function which can be forced to be very small.)

Hence one needs to prove $\kappa$-noncollapsing and
the canonical neighborhoood assumption together. The main proposition of
II.5 says that there are decreasing sequences $r_j$, $\kappa_j$ and
$\overline{\delta}_j$ so that if $\delta(\cdot)$ is a function with
$\delta(\cdot) \big|_{[2^{j-1} \epsilon, 2^j \epsilon]} \le \overline{\delta}_j$
for each $j > 0$ then any Ricci flow with surgery, defined with the
parameters $r(\cdot)$ and $\delta(\cdot)$, is $\kappa_j$-noncollapsed on
the time interval $[2^{j-1} \epsilon, 2^j \epsilon]$ at scales less than
$\epsilon$ and satisfies the
canonical neighborhood assumption there. Here we take
$r(\cdot) \big|_{[2^{j-1} \epsilon, 2^j \epsilon)} = r_j$.

The proof of the proposition is by induction.  Suppose that it is true for
$1 \le j \le i$. In the induction step,
besides defining the parameters $r_{i+1}$, $\kappa_{i+1}$ and
$\overline{\delta}_{i+1}$, one redefines $\overline{\delta}_i$. As one
only redefines $\overline{\delta}$ in the previous interval, there is no
circularity.

The first step of the proof, Lemma II.5.2, consists of showing
that there is some $\kappa > 0$ so that for any $r$, one can
find $\overline{\delta} = \overline{\delta}(r) > 0$ with the following
property. Suppose that $g(\cdot)$ is a Ricci flow
with surgery defined on $[0, T)$, with $T \in [2^i \epsilon, 2^{i+1} \epsilon]$,
that satisfies the proposition on $[0, 2^i \epsilon]$.
Suppose that it also satisfies the canonical neighborhood assumption with
parameter $r$ on $[2^i \epsilon, T)$, and is constructed using a function
$\delta(\cdot)$ that satisfies $\delta(t) \: \le \: \overline{\delta}$ on
$[2^{i-1} \epsilon, T)$. Then it is
$\kappa$-noncollapsed at all scales less than $\epsilon$.

The proof of this
lemma is along the lines of the $\kappa$-noncollapsing result of I.7, with
some important modifications.  One again considers the ${\mathcal L}$-length of curves
$\gamma(\tau)$ starting from the point at which one wishes to prove the
noncollapsing. One wants to find a spacetime point
$(\overline{x}, \overline{t})$, with
$\overline{t} \in [2^{i-1} \epsilon, 2^i \epsilon]$, at which one has
an explicit upper bound on $l$. In I.7, the analogous statement came from
a differential inequality for $l$. In order to use this differential
equality in the present case, one needs to know that any curve $\gamma(\tau)$
that is competitive to be a minimizer for $L(\overline{x}, \overline{t})$
will avoid the surgery regions.
Choosing $\overline{\delta}$ small enough, one can ensure that the
surgeries in the time interval $[2^{i-1} \epsilon, T)$ are done on very
long thin necks. Using Lemma II.4.5, one shows that a curve $\gamma(\tau)$
passing near such a surgery region obtains a large value of ${\mathcal L}$,
thereby making it noncompetitive as a minimizer for
$L(\overline{x}, \overline{t})$. (This is the underlying
reason that the surgery parameter
$\delta(\cdot)$ is chosen in a time-dependent way.)
One also chooses the point
$(\overline{x}, \overline{t})$ so that there is a small parabolic neighborhood
around it with a bound on its geometry. One can then run the argument of
I.7 to prove $\kappa$-noncollapsing in the time interval
$\left[ 2^i \epsilon, T \right)$.

The proof of the main proposition of II.5 is now by contradiction.  Suppose
that it is not true.  Then for some sequences $r^\al \rightarrow 0$ and
$\overline{\delta}^\al \rightarrow 0$, for each $\al$ there is a counterexample to the
proposition with $r_{i+1} \le r^\al$ and $\delta_i, \delta_{i+1} \le
\overline{\delta}^\al$. That is, there is some spacetime point in the interval
$\left[ 2^i \epsilon, 2^{i+1} \epsilon \right]$ at which the
canonical neighborhood assumption fails.  Take a first such point
$(x^\al, t^\al)$.
By Lemma II.5.2, one has $\kappa$-noncollapsing up to the time of this
first counterexample. Using this noncollapsing, one can consider
taking rescaled limits.  If there are no surgeries in an
appropriate-sized backward spacetime region around $(x^\al, t^\al)$
then one can extract a convergent subsequence as $\al \rightarrow \infty$
and construct, as in the proof of Theorem I.12.1, a limit $\kappa$-solution,
thereby giving a contradiction. If there are nearby interfering surgeries
then one argues, using Lemma II.4.5, that the point $(x^\al, t^\al)$ is in fact in a
canonical neighborhood, again giving a contradiction.

Having constructed the Ricci flow with surgery, if the initial
manifold is simply-connected then according to
\cite{Colding-Minicozzi,Colding-Minicozzi2,Perelman3},
there is a finite extinction time.
One then concludes that the Poincar\'e Conjecture holds.


\subsection{II.6-II.8}

Sections II.6 and II.8 analyze the large-time behavior of a Ricci flow with
surgery.

Section II.6 establishes back-and-forth curvature estimates. Proposition
II.6.3 is an analog of Theorem I.12.2 and Proposition II.6.4 is an analog
of Theorem I.12.3.  The proofs are along the lines of the proofs of
Theorems I.12.2 and I.12.3, but are complicated by the possible presence
of surgeries.

The thick-thin decomposition for large-time slices is considered in
Section II.7. Using monotonicity arguments of Hamilton, it is
shown that as $t \rightarrow \infty$ the metric on the $w$-thick part $M^+(w,t)$
becomes closer and closer to having constant negative sectional
curvature. Using a hyperbolic rigidity argument of Hamilton, it is stated
that the hyperbolic pieces stabilize in the sense that
there is a finite collection $\{ (H_i, x_i) \}_{i=1}^k$ of
pointed finite-volume $3$-manifolds of constant sectional
curvature $- \: \frac14$ so that for large $t$, the metric
$\widehat{g}(t) \: = \: \frac{1}{t} \: g(t)$ on the $w$-thick part
$M^+(w,t)$ approaches the metric on the $w$-thick part
of $\bigcup_{i=1}^k H_i$. It is stated that the cuspidal tori (if any) of
the hyperbolic pieces  are incompressible in $M$. To show this
(following Hamilton), if there is a compressing $3$-disk then one
takes a minimal such $3$-disk, say of area $A(t)$, and shows from a
differential inequality for $A(\cdot)$ that for large $t$ the function
$A(t)$ is negative, which is a contradiction.

Theorem II.7.4, a statement in Riemannian geometry,
characterizes the thin part $M^-(w,t)$, for small $w$ and
large $t$, as a graph manifold.  The main hypothesis of the theorem
is that for each point $x$, there is a radius $\rho = \rho(x)$
so that the ball $B(x,\rho)$ has volume at most $w \rho^3$ and
sectional curvatures bounded below by $- \: \rho^{-2}$. In this
sense the manifold is locally volume collapsed with respect to a lower
sectional curvature bound.

Section II.8 contains an alternative proof of the incompressibility
of cuspidal tori, using the functional $\lambda_1(g) =
\lambda_1(-4 \triangle + R)$.
(At the beginning of Section \ref{II.8}, we give a simpler argument
using the functional $R_{min}(g) \: \vol(M,g)^{\frac23}$.)
More generally, the functional $\lambda_1(g)$ is
used to define a topological invariant that determines the nature
of the geometric decomposition. First, the manifold $M$ admits a
Riemannian metric $g$ with $\lambda_1(g) > 0$ if and only if it
admits a Riemannian metric with positive scalar curvature, which
in turn is equivalent to saying that $M$ is diffeomorphic to a
connected sum of $S^1 \times S^2$'s and round quotients of
$S^3$. If $M$ does not admit a Riemannian metric with
$\lambda_1 > 0$, let $\overline{\lambda}$ be the supremum of
$\lambda_1(g) \cdot \vol(M, g)^{\frac23}$ over all Riemannian metrics $g$
on $M$. If $\overline{\lambda} = 0$ then $M$ is a graph manifold.
If $\overline{\lambda} < 0$ then the geometric decomposition of $M$
contains a nonempty hyperbolic piece, with total volume
$\left( - \: \frac23 \: \overline{\lambda} \right)^{\frac32}$. The proofs
of these statements use the monotonicity of
$\lambda(g(t)) \cdot \vol(M, g(t))^{\frac23}$, when it is nonpositive,
under a smooth Ricci flow. The main work is to show
that in the case of a Ricci flow with surgery, one can choose $\delta(\cdot)$
so that $\lambda(g(t)) \cdot \vol(M, g(t))^{\frac23}$ is
arbitrarily close to being nondecreasing in $t$.

\section{II. Notation and terminology}
\label{notationterminology}

$B(x,t,r)$ denotes the open metric ball of radius $r$, with respect to the metric
at time $t$, centered at $x$.

$P(x,t,r, \Delta t)$ denotes a parabolic neighborhood, that is the set of
all points $(x^\prime, t^\prime)$ with $x^\prime \in B(x,t,r)$ and
$t^\prime \in [t, t+\Delta t]$ or $t^\prime \in [t+\Delta t, t]$, depending on
the sign of $\Delta t$.

\begin{definition}[smooth closeness and epsilon-necks]
  \label{def:kl-notation-terminology-smooth-closeness-epsilon-necks}
  \dcref{closenessdef}
  \group{Kleiner--Lott Surgery Foundations}
  \source{kleiner-lott}
 \label{closenessdef}
We say that a Riemannian manifold
$(M_1, g_1)$ has distance $\le \: \epsilon$ in the $C^N$-topology
to another Riemannian manifold $(M_2, g_2)$ if
there is a diffeomorphism $\phi \: : \: M_2 \rightarrow M_1$ so
that $\sum_{|I| \: \le \: N} \frac{1}{|I|!} \:
\parallel \nabla^I (\phi^* g_1 - g_2) \parallel_\infty \: \le
\: \epsilon$.
An open set $U$ in a Riemannian $3$-manifold $M$ is an
{\em $\eps$-neck}
if modulo rescaling, it has
distance less than $\eps$,
in the $C^{[1/\epsilon]+1}$-topology, to the product of the
round $2$-sphere of scalar curvature $1$ (and therefore
Gaussian curvature $\frac12$)
with an interval $I$ of length
greater than
$2 \eps^{-1}$.
If a point $x \in M$ and a neighborhood $U$ of $x$ are specified then we
will understand that ``distance''
refers to the pointed topology,
where the basepoint in $S^2 \times I$ projects to the center of $I$.

We make a similar definition of $\epsilon$-closeness in the
spacetime case, where $\nabla^I$ now includes time derivatives.
A subset of the form $U\times [a,b]\subset M\times [a,b]$,
where $U\subset M$ is open, sitting
in the spacetime of a Ricci flow is a {\em strong $\eps$-neck}
if after parabolic rescaling and time shifting, it has distance less
than $\eps$ to the
product Ricci flow defined on the time
interval $[-1,0]$ which,  at its final time,
is isometric to the  product of a  round $2$-sphere of scalar
curvature $1$
with an interval of length greater than $2\eps^{-1}$. (Evidently,
the time-$0$ slice of the product has $3$-dimensional scalar curvature
equal to $1$.)
\end{definition}

Our definition of an $\epsilon$-neck
differs in an insubstantial way from that on p. 1 of II.
In the definition of \cite{Perelman2},
a ball $B(x, t, \epsilon^{-1} r)$ is called
an $\epsilon$-neck if, after rescaling the metric with a factor
$r^{-2}$, it is $\epsilon$-close,
i.e. has distance less than $\epsilon$,
to {\em the corresponding subset of}
the standard neck
$S^2 \times I$...
(italicized words added by us).
(The issue is that a large metric ball in the cylinder $\R \times S^2$
does not have a smooth boundary.)
Clearly after a slight change
of the constants, an $\epsilon$-neck in our sense is contained in an
$\epsilon$-neck in the sense of \cite{Perelman2},
and vice versa. An important fact
is that the notion of $(x, t)$ being contained in an $\epsilon$-neck
is an open condition with respect to the pointed $C^{[1/\epsilon]+1}$-topology
on Ricci flow solutions.

With an $\epsilon$-approximation
$f \: : \: S^2 \times I \rightarrow U$
being understood, a cross-sectional sphere in $U$ will mean the
image of $S^2 \times \{ \lambda \}$ under $f$, for some
$\lambda \in (- \epsilon^{-1}, \epsilon^{-1})$.
Any curve $\gamma$ in $U$ that intersects both
$f(S^2 \times \{ \epsilon^{-1} \})$ and  $f(S^2 \times \{ -\epsilon^{-1} \})$ must
intersect each cross-sectional
sphere. If $\gamma$ is a minimizing  geodesic and $\epsilon$ is small
enough then $\gamma$ will intersect each cross-sectional sphere
exactly once.

There is a typo in the definition of a strong $\eps$-neck in \cite{Perelman2} :
the parabolic neighborhood should be $P(x,t,\eps^{-1}r,-r^2)$,
i.e. it should go backward in time rather than forward. We note that the
time interval involved in the
definition of strong $\eps$-neck, i.e. $1$ after rescaling, is different than the
rescaled time interval $\epsilon^{-1}$ in Theorem \ref{thmI.12.1}.

In the next definition, $\I$ is an open interval and  $B^3$ is an open ball.

\begin{definition}[epsilon-tubes horns and caps]
  \label{def:kl-notation-terminology-epsilon-tubes-horns-caps}
  \dcref{notation-terminology:definition-2}
  \group{Kleiner--Lott Surgery Foundations}
  \source{kleiner-lott}

A metric on
$S^2 \times \I$ such that each point is contained in some $\epsilon$-neck
is called an {\em $\epsilon$-tube}, or an {\em $\epsilon$-horn}, or a {\em double
$\epsilon$-horn}, if the scalar curvature stays bounded on both ends,
stays bounded on one end and tends to infinity on the other, or
tends to infinity on both ends, respectively.

A metric on $B^3$ or $\R P^3 - \overline{B^3}$, such that each point
outside some compact subset is contained in an $\epsilon$-neck, is
called an $\epsilon$-cap or a capped $\epsilon$-horn, if the
scalar curvature stays bounded or tends to infinity on the end,
respectively.
\end{definition}

An example of an $\epsilon$-tube is $S^2 \times (- \epsilon^{-1}, \epsilon^{-1})$ with the
product metric.
For a relevant example of an $\epsilon$-horn, consider the metric
\begin{equation}
g \: = \: dr^2 \: + \: \frac{1}{8\ln \frac{1}{r}} \: r^2 \: d\theta^2
\end{equation}
on $(0, R) \times S^2$, where $d\theta^2$ is the metric on $S^2$ with $R=1$.
From \cite{Angenent-Knopf}, the metric $g$ models a rotationally symmetric
neckpinch. Rescaling around $r_0$, we put
$s \: = \: \sqrt{8 \ln \frac{1}{r_0}} \: \left( \frac{r}{r_0} - 1 \right)$ and find
\begin{equation}
\frac{8 \ln \frac{1}{r_0} }{r_0^2} \: g \: = \:
ds^2 \: + \: \left( 1 +
{ \frac{1}{\sqrt{8 \ln \frac{1}{r_0} }}  }
\left( 2 +
\frac{1}{\ln \frac{1}{r_0}}
\right) s + O(s^2) \right) \: d\theta^2.
\end{equation}
For small $r_0$, if we take
$\epsilon \sim \left( \ln \frac{1}{r_0} \right)^{- \: \frac14}$ then the region with
$s \in \left( - \: \epsilon^{-1}, \epsilon^{-1} \right)$
will be $\epsilon$-biLipschitz close
to the standard cylinder. Note that as $r_0 \rightarrow 0$,
the constant $\epsilon$ improves; this is related to
Lemma \ref{lemmaII.4.3}.


An $\epsilon$-cap is the result of capping off an $\epsilon$-tube by a
$3$-ball or $\R P^3 - \overline{B^3}$
with an arbitrary metric.  A capped $\epsilon$-horn
is the result of capping off an $\epsilon$-horn by a
$3$-ball or $\R P^3 - \overline{B^3}$ with an arbitrary metric.

\begin{remark}[positivity for notation and terminology]
  \label{rem:kl-notation-terminology-positivity-notation-terminology}
  \dcref{remepsilon}
  \group{Kleiner--Lott Surgery Foundations}
  \source{kleiner-lott}

\label{remepsilon}
Throughout the rest of these notes, $\epsilon$ denotes a small
positive constant that is meant to be universal.  The precise value
of $\epsilon$ is unspecified.  If the statement of a lemma or
theorem invokes $\epsilon$ then the statement is meant to be true
uniformly with respect to the other variables, provided $\eps$ is
sufficiently small. When going through
the proofs one is allowed to make $\epsilon$ small enough
so that the arguments work, but one is only allowed to make a
finite number of such reductions.
\end{remark}

\begin{lemma}[neck geometry lemma]
  \label{lem:kl-notation-terminology-neck-geometry-lemma}
  \dcref{notation-terminology:lemma-4}
  \group{Kleiner--Lott Surgery Foundations}
  \source{kleiner-lott}

Let $U$ be
an $\epsilon$-neck in an $\epsilon$-tube (or horn) and let
$S$ be a cross-sectional
sphere in $U$. Then $S$ separates the two ends of the tube (or horn).
\end{lemma}
\begin{proof}
Let $W$ denote the tube (or horn).
As any point $m \in W$ lies in some $\epsilon$-neck,
there is a unique lowest eigenvalue of the Ricci operator
$\Ric  \in \End(T_m W)$ at $m$. Let $\xi_m \subset T_mW$ be the
corresponding eigenspace.  As $m$ varies, the
$\xi_m$'s form a smooth
line field $\xi$ on $W$, to which $S$ is transverse.

Suppose that $S$ does not separate the two ends of $W$.
Then $S$ represents a trivial element of
$\HH_2(S^2 \times I) \cong \pi_2(S^2 \times I)$ and
there is an embedded $3$-disk $D \subset W$ for which
$\partial D = S$. This contradicts the fact that
the line field $\xi$ is transverse to $S$ and extends over $D$.
\end{proof}


\section{II.1. Three-dimensional $\kappa$-solutions} \label{II.1}

This section is concerned with properties of three-dimensional oriented
$\kappa$-solutions. For brevity, in the rest of these notes we will generally omit
the phrases ``three-dimensional'' and ``oriented''.

If $(M, g(\cdot))$ is a $\kappa$-solution then its
topology is easy to describe.  By definition, $(M, g(t))$ has nonnegative
sectional curvature.  If it does not have strictly positive curvature
then the universal cover splits off a line (see Theorem \ref{splitting}),
from which
it follows (using Corollary \ref{2Dsoliton} and the $\kappa$-noncollapsing)
that $(M, g(\cdot))$ is a standard shrinking cylinder
$\R \times S^2$ or its $\Z_2$ quotient $\R \times_{\Z_2} S^2$.
If $(M, g(t))$ has strictly positive curvature and $M$ is compact
then it is diffeomorphic to a spherical space form
\cite{Hamiltonnnnn}.
If $(M, g(t))$ has strictly positive curvature and $M$ is noncompact
then it is diffeomorphic to $\R^3$ \cite{Cheeger-Gromoll}. The lemmas
in this section give more precise geometric information.
Recall that $M_\epsilon$ consists of the points in a $\kappa$-solution
which are not the center of an $\epsilon$-neck.

\begin{lemma}[existence for kappa solutions]
  \label{lem:kl-three-dimensional-solutions-existence-kappa-solutions}
  \dcref{noncompactsoliton}
  \group{Kleiner--Lott Surgery Foundations}
  \source{kleiner-lott}
 \label{noncompactsoliton}
If $(M, g(t))$ is a time slice of a
noncompact $\kappa$-solution and
$M_\epsilon \neq \emptyset$ then there is a compact
submanifold-with-boundary $X \subset M$ so that
$M_\epsilon \subset X$,  $X$ is diffeomorphic to $\overline{B^3}$ or
$\R P^3 - B^3$, and $M- \Int(X)$ is
diffeomorphic to $[0, \infty) \times S^2$.
\end{lemma}
\begin{proof}
If $(M,g(t))$ does not have positive sectional curvature and
$M_\epsilon \neq \emptyset$ then $M$ must be isometric to
$\R \times_{\Z_2} S^2$, in which case the lemma is easily seen to
be true with $X$ diffeomorphic to $\R P^3 - B^3$.
Suppose that $(M, g(t))$ has positive sectional curvature.
Choose $x \in M_\epsilon$. Let $\gamma \: : \: [0, \infty)
\rightarrow M$ be a ray with $\gamma(0) = x$. As $M_\epsilon$ is
compact, there is some $a > 0$ so that if $t > a$ then
$\gamma(t) \notin M_\epsilon$.
We can cover $(a, \infty)$ by open
intervals $V_j$ so that $\gamma \Big|_{V_j}$ is a geodesic segment
in an $\epsilon$-neck of
rescaled length approximately
$2 \epsilon^{-1}$.
Then we can find a cover of $(a, \infty)$ by
linearly ordered open intervals
$U_i$, refining the previous cover, so that \\
1. The rescaled length of $\gamma \Big|_{U_i}$ is approximately
$\frac{1}{10} \: \epsilon^{-1}$. \\
2. Choosing some $x_i \in U_i \cap U_{i+1}$, the
rescaled length (with rescaling at $x_i$) of
$\gamma \Big|_{U_i \cap U_{i+1}}$ is approximately
$\frac{1}{40} \: \epsilon^{-1}$ and
$\gamma \Big|_{U_i \cap U_{i+1}}$ lies in an $\epsilon$-neck $W_i$
centered at $x_i$.

Let $\phi_i$
be projection on the first factor in the assumed diffeomorphism
$W_i \cong  (- \epsilon^{-1}, \epsilon^{-1}) \times S^2$.
If $\epsilon$ is sufficiently small then
the composition $\phi_i \circ \gamma \big|_{U_i}  \: : \: U_i \rightarrow
(- \epsilon^{-1}, \epsilon^{-1})$ is a diffeomorphism onto its image.
Put $N_i = \phi_i^{-1}(\Image(\phi_i \circ \gamma \big|_{U_i})) \subset
W_i$ and
$p_i \: = \: (\phi_i \circ \gamma \big|_{U_i})^{-1} \circ \phi_i \: : \:
N_i \rightarrow U_i \subset (a, \infty)$. Then $p_i$ is $\epsilon$-close
to being a Riemannian submersion and
on overlaps $N_i \cap N_{i+1}$, the maps $p_i$ and $p_{i+1}$ are
$C^K$-close. Choosing an appropriate
partition of unity $\{b_i\}$ subordinate to the $U_i$'s, if $\epsilon$ is
small then the function
$f = \sum_i b_i p_i$ is a submersion from $\bigcup_i N_i$ to
$(a, \infty)$. The fiber is seen to be $S^2$. Given $t \in (a, \infty)$,
put $X = M - f^{-1}(t, \infty)$. Then $M - \Int(X) = f^{-1}([t, \infty))$
is diffeomorphic to $[0, \infty) \times S^2$.

Recall from Section \ref{pf11.7} that we have an exhaustion of
$M$ by certain convex compact subsets.  As $M$ is one-ended, the
subsets have connected boundary.
As in Section \ref{pf11.7}, if the boundary of such a subset intersects an
$\epsilon$-neck then the intersection will be a nearly
cross-sectional $2$-sphere in the $\epsilon$-neck. Hence with
an appropriate choice of $t$, the set $X$ will be isotopic to one
of our convex subsets and so diffeomorphic to a closed $3$-ball.
\end{proof}

\begin{lemma}[kappa solutions lemma]
  \label{lem:kl-three-dimensional-solutions-kappa-solutions-lemma}
  \dcref{onlyneck}
  \group{Kleiner--Lott Surgery Foundations}
  \source{kleiner-lott}
 \label{onlyneck}
If $(M, g(t))$ is a time slice of a $\kappa$-solution with
$M_\epsilon = \emptyset$ then the
Ricci flow is the evolving round cylinder $\R \times S^2$.
\end{lemma}
\par\medskip\noindent\textit{Source argument (not certified as complete).}\ \begingroup\itshape
By assumption, each point $(x,t)$ lies in an $\epsilon$-neck.
If $\epsilon$ is sufficiently small then piecing
the necks together, we conclude that $M$ must be diffeomorphic to
$S^1 \times S^2$ or $\R \times S^2$; see the proof of
Lemma \ref{noncompactsoliton} for a similar argument. Then the universal cover
$\widetilde{M}$ is
$\R \times S^2$. As it has nonnegative sectional curvature and
two ends, Toponogov's theorem implies that $(\widetilde{M},
\widetilde{g}(t))$ splits off
an $\R$-factor. Using the strong maximum principle, the Ricci flow
on $\widetilde{M}$
splits off an $\R$-factor; see Theorem \ref{splitting}.
Using Corollary \ref{2Dsoliton}, it follows that $(\widetilde{M},
\widetilde{g}(t))$ is the evolving round cylinder $\R \times S^2$.
From the $\kappa$-noncollapsing, the quotient $M$ cannot be
$S^1 \times S^2$.
\endgroup\par\medskip

A $\kappa$-solution has an asymptotic soliton
(Section \ref{I.11.2}) that is either compact or noncompact.
If the asymptotic soliton of a compact $\kappa$-solution
$(M, g(\cdot))$ is also compact then it must be a shrinking
quotient of the round $S^3$ \cite{Hamiltonnnnn}, so the same is true of
$M$.

\begin{lemma}[compactness for kappa solutions]
  \label{lem:kl-three-dimensional-solutions-compactness-kappa-solutions}
  \dcref{S3RP3}
  \group{Kleiner--Lott Surgery Foundations}
  \source{kleiner-lott}
 \label{S3RP3}
If a $\kappa$-solution $(M, g(\cdot))$
is compact and has a noncompact asymptotic soliton
then $M$ is diffeomorphic to $S^3$ or $\R P^3$.
\end{lemma}
\begin{proof}
  \uses{cor:kl-necklike-behavior-at-infinity-three-dimensional-solution-strong-version-existence-kappa-solutions,lem:kl-three-dimensional-solutions-existence-kappa-solutions,thm:kl-compactness-space-three-dimensional-solutions-compactness-kappa-solutions}
We use Corollary \ref{neckstructure} in Section \ref{strongversion}.
First, we claim that the time slices of the type-D $\kappa$-solutions of
Corollary \ref{neckstructure} have a
universal upper bound on $\max_M R \cdot \diam(M)^2$.
To see this, we can rescale at the point $x \in M_\epsilon$ by
$R(x)$, after which the diameter is bounded above by
$2 \sqrt{\alpha}$. We then use Theorem \ref{thmI.11.7}
to get an upper bound on the rescaled
scalar curvature, which proves the claim. Given an upper bound on
$\max_M R \cdot \diam(M)^2$, the asymptotic soliton
cannot be noncompact.

Thus we are in case C of Corollary \ref{neckstructure}.
Take a sequence $t_i \rightarrow - \infty$ and choose points $x_i, y_i
\in M_\epsilon(t_i)$ as in Corollary \ref{neckstructure}.C.
Rescale by
$R(x_i, t_i)$ and take a subsequence that converges to a pointed
Ricci flow solution
$(M_\infty, (x_\infty, t_\infty))$. The limit $M_\infty$ cannot be
compact, as otherwise we would have a uniform upper bound on
$R \: \cdot \: \diam^2$ for $(M, g(t_i))$, which would
contradict the existence of the noncompact asymptotic soliton.
Thus $M_\infty$ is a noncompact $\kappa$-solution.
We can find compact sets $X_i \subset M$
containing $B(x_i, \alpha R(x_i, t_i)^{- \: \frac12})$ so that
$\{X_i\}$ converges to a set $X_\infty \subset M_\infty$ as in
Lemma \ref{noncompactsoliton}.
Taking a further subsequence, we find similar compact
sets $Y_i \subset M$ containing
$B(y_i, \alpha R(y_i, t_i)^{- \: \frac12})$ so that $\{Y_i\}$ converges to
a set $Y_\infty \subset M_\infty$ as in Lemma
\ref{noncompactsoliton}. In particular, for large $i$, $X_i$ and $Y_i$ are
each diffeomorphic to either $\overline{B^3}$ or $\R P^3 - B^3$.
Considering a minimizing geodesic
segment $\overline{x_i y_i}$ as in the statement of
Corollary \ref{neckstructure}.C, we
can use an argument as in the proof of Lemma
\ref{noncompactsoliton} to construct a submersion from
$M - (X_i \cup Y_i)$ to an interval, with fiber $S^2$. Hence $M$
is diffeomorphic to the result of gluing $X_i$ and $Y_i$ along
a $2$-sphere.
As $M$ has finite fundamental group, no more
than one of $X_i$ and $Y_i$ can be diffeomorphic to
$\R P^3 \: - \: {B^3}$. Thus $M$ is diffeomorphic to
$S^3$ or $\R P^3$.
\end{proof}

From Lemma \ref{propI.11.9}, every ancient solution which is a
$\kappa$-solution {\em for some $\kappa$} is either a $\kappa_0$-solution
or a metric quotient of the round $S^3$.

\begin{lemma}[existence for kappa solutions]
  \label{lem:kl-three-dimensional-solutions-existence-kappa-solutions-4}
  \dcref{derivativeestimates}
  \group{Kleiner--Lott Surgery Foundations}
  \source{kleiner-lott}
 \label{derivativeestimates}
There is a universal constant
$\eta$ such that at each point of every ancient solution that is
a $\kappa$-solution for some $\kappa$, we have estimates
\begin{equation} \label{derestimates}
|\nabla R| < \eta R^{\frac32}, \: \: \: \: \: \: |R_t| < \eta R^2.
\end{equation}
\end{lemma}
\par\medskip\noindent\textit{Source argument (not certified as complete).}\ \begingroup\itshape
This is obviously true for metric quotients of the round $S^3$. For
$\kappa_0$-solutions it follows from the compactness result in
Theorem \ref{thmI.11.7}, after rescaling the scalar curvature at
the given point to be $1$.
\endgroup\par\medskip


It is sometimes useful to rewrite (\ref{derestimates}) as a pair of estimates
on the spacetime derivatives of the quantity $R^{-1}$ at points where $R\neq 0$:
\begin{equation}
\label{eqnderestimatesrinverse}
|\nabla (R^{-\frac12})| < \frac{\eta}{2}, \: \: \: \: \: \: |(R^{-1})_t| < \eta .
\end{equation}


\begin{lemma}[existence for kappa solutions]
  \label{lem:kl-three-dimensional-solutions-existence-kappa-solutions-5}
  \dcref{nbhd}
  \group{Kleiner--Lott Surgery Foundations}
  \source{kleiner-lott}
 \label{nbhd}
For every sufficiently small $\epsilon > 0$ one can
find $C_1 = C_1(\epsilon)$ and $C_2 = C_2(\epsilon)$
such that for each point $(x, t)$ in every
$\kappa$-solution there is a radius
$r \in [ R(x,t)^{-1/2}, C_1 R(x,t)^{-1/2}]$ and a neighborhood
$B$,
$\overline{B(x,t,r)} \subset B \subset B(x,t,2r)$,
which falls into one
of the four categories: \\
(a) $B$ is a strong $\epsilon$-neck
(more precisely, $B$ is the slice of a
strong $\epsilon$-neck at its maximal time, and an appropriate parabolic
neighborhood of $B$
satisfies the condition to be a strong $\epsilon$-neck),
or \\
(b) $B$ is an $\epsilon$-cap, or \\
(c) $B$ is a closed manifold, diffeomorphic to $S^3$ or $\R P^3$, or\\
(d) $B$ is a closed manifold of constant positive sectional curvature.

Furthermore:

\begin{itemize}
\item The scalar curvature in $B$ at time $t$ is between
$C_2^{-1} R(x,t)$ and $C_2 R(x,t)$.

\item The volume of $B$ in cases (a), (b) and (c) is
greater than $C_2^{-1} R(x,t)^{- \: \frac32}$.

\item In case (b), there is an $\eps$-neck $U\subset B$ with compact complement in $B$
(i.e. the end of $B$ is entirely contained in the $\eps$-neck) such that the
distance from $x$ to $U$ is at least $10000R(x,t)^{-1/2}$.

\item In case (c) the
sectional curvature in $B$ at time $t$ is greater than $C_2^{-1} R(x,t)$.
\end{itemize}
\end{lemma}

\begin{remark}[structural remark on three-dimensional -solutions]
  \label{rem:kl-three-dimensional-solutions-structural-remark-three-dimensional-solutions}
  \dcref{three-dimensional-solutions:remark-6}
  \group{Kleiner--Lott Surgery Foundations}
  \source{kleiner-lott}

The statement of the lemma is slightly stronger than the corresponding
statement in II.1.5, in that
we have $r \ge  R(x,t)^{-1/2}$ as opposed to $r > 0$.
\end{remark}

\begin{proof}
  \uses{lem:kl-three-dimensional-solutions-existence-kappa-solutions,lem:kl-three-dimensional-solutions-kappa-solutions-lemma}
We may assume that we are talking about a
$\kappa_0$-solution, as if $M$ is a metric quotient of a round sphere then
it falls into category (d) for any $r > \pi (R(x_i, t_i)/6)^{-1/2}$
(since then
$M \: = \: \overline{B(x_i, t_i, r)} \: = \: B(x_i, t_i, 2r)$).

Fix a small $\epsilon$ and suppose that the claim
is not true.  Then there is a sequence of $\kappa_0$-solutions $M_i$ that
together provide a counterexample.  That is, there is a sequence
$C_i \rightarrow \infty$ and a sequence of points $(x_i, t_i) \in
M_i \times (-\infty, 0]$ so that for any $r \in
[R(x_i, t_i)^{-1/2}, C_i R(x_i, t_i)^{-1/2} ]$ one cannot
find a $B$ between
$\overline{B(x_i, t_i, r)}$
and $B(x_i, t_i, 2r)$ falling into
one of the four categories
and satisfying the subsidiary conditions with parameter $C_2 = C_i$.
Rescale the metric by $R(x_i, t_i)$ and take
a convergent subsequence of
$(M_i, (x_i, t_i))$ to obtain a limit $\kappa_0$-solution
$(M_\infty, (x_\infty, t_\infty))$. Then for any $r > 1$,
one cannot
find a $B_\infty$ between
$\overline{B(x_\infty, t_\infty, r)}$
and
$B(x_\infty, t_\infty, 2r)$ falling into
one of the four categories and satisfying the subsidiary conditions
for any parameter $C_2$.

If $M_\infty$ is compact then
for any $r$ greater than the diameter of the
time-$t_\infty$ slice of $M_\infty$,
$\overline{B(x_\infty, t_\infty, r)} = M_\infty = B(x_\infty, t_\infty, 2r)$
falls
into category (c) or (d). For the subsidiary conditions,
$M_\infty$ clearly has a lower volume bound, a positive lower
scalar curvature bound and an upper scalar
curvature bound. As a compact
$\kappa_0$-solution has positive sectional curvature,
$M_\infty$ also has a lower sectional curvature bound.
This is a contradiction.

If $M_\infty$ is noncompact then Lemma
\ref{noncompactsoliton} (or more precisely its proof) and
Lemma \ref{onlyneck}
imply that for some $r>1$, there will be a $B$ between
$\overline{B(x_\infty, t_\infty, r)}$
and $B(x_\infty, t_\infty, 2r)$
falling into category (a) or (b).
In case (b), by choosing the parameter $r$ sufficiently large,
the existence of the $\eps$-neck $U$
with the desired properties follows from the proof of
Lemma \ref{noncompactsoliton}.
For the other subsidiary conditions,
$B$ clearly has a lower volume bound, a positive lower
scalar curvature bound and an upper scalar
curvature bound.
This contradiction completes the proof of the lemma.
\end{proof}

\section{II.2. Standard solutions} \label{II.2}



The next few sections are concerned with the properties of  special
Ricci flow solutions on $M=\R^3$.
We fix a smooth rotationally symmetric   metric $g_0$ which
is the result of gluing a hemispherical-type cap to a
half-infinite cylinder of scalar curvature $1$.
Among other properties, $g_0$ is complete and has
nonnegative  curvature operator.
We also assume that $g_0$ has scalar curvature
bounded below by $1$.

\begin{remark}[structural remark on standard solutions]
  \label{rem:kl-standard-solutions-structural-remark-standard-solutions}
  \dcref{standard-solutions:remark-1}
  \group{Kleiner--Lott Surgery Foundations}
  \source{kleiner-lott}

In Section \ref{surgery} we will further specialize the initial metric
$g_0$ of the standard solution, for technical convenience in
doing surgeries.
\end{remark}

\begin{definition}[standard solution definition]
  \label{def:kl-standard-solutions-standard-solution-definition}
  \dcref{defstandardsolution}
  \group{Kleiner--Lott Surgery Foundations}
  \source{kleiner-lott}

\label{defstandardsolution}
A Ricci flow $(\R^3,g(\cdot))$  defined on a time interval $[0,a)$
is a {\em standard solution} if
it has  initial condition $g_0$, the curvature $|\Rm|$ is bounded
on compact time intervals $[0,a']\subset [0,a)$, and it cannot
be extended to a Ricci flow with the same properties
on a strictly longer time interval.
\end{definition}

It will turn out that every standard solution
is defined on the time interval $[0,1)$.

To motivate the next few sections, let us mention that the surgery
procedure will amount to gluing in a truncated copy of
$(\R^3, g_0)$. The metric on this added region will then evolve as
part of the Ricci flow that takes up after the surgery is performed.
We will need to understand
the behavior of the Ricci flow after performing a surgery. Near
the added region, this will be modeled by a standard solution.
Hence one first needs to understand the Ricci flow of a standard
solution.

The main results of II.2 concerning the Ricci flow on a standard solution
(Sections \ref{claim2}-\ref{standardlycompact}) are used in
II.4.5 (Lemma \ref{II.4.5}) to show that,
roughly speaking, the part of the manifold added by surgery
acquires a large scalar curvature soon after the surgery time.  This
is used crucially in II.5 (Sections \ref{II.5.2section} and \ref{surgeryflow})
to adapt the noncollapsing argument of I.7 to Ricci flows
with surgery.

Our order of presentation of the material in
II.2 is somewhat different than that of \cite{Perelman2}.
In Sections \ref{claim2}-\ref{secclaim5}, we
cover Claims 2, 4 and 5 of II.2. These are what's needed in the
sequel.  The other results of II.2, Claims 1 and 3, are concerned with proving the
uniqueness of the standard solution. Although it may seem
intuitively obvious that there should be a unique and
rotationally-symmetric standard solution, the argument
is not routine since the manifold is noncompact.

In fact, the uniqueness
is not really needed for the sequel. (For example, the method of proof of
Lemma \ref{II.4.5} produces {\em a} standard solution in a
limiting argument and it is enough to know certain properties
of this standard solution.) Because of this we will talk
about {\em a} standard solution rather than {\em the}
standard solution.

Consequently, we present the material so that we do not logically
need the uniqueness of the standard solution.
Having uniqueness does not
shorten the subsequent arguments any.
Of course, one can ask independently whether the standard
solution is unique. In Section \ref{claimII.2.1} we
show that a standard solution is rotationally symmetric.
In Section \ref{claim3} we sketch the argument for uniqueness.
Papers concerning the uniqueness of the standard solution are
\cite{Chen-Zhu,Lu-Tian}.

We end this section by collecting some basic facts about standard solutions.

\begin{lemma}[existence for kappa solutions]
  \label{lem:kl-standard-solutions-existence-kappa-solutions}
  \dcref{lemstandard0}
  \group{Kleiner--Lott Surgery Foundations}
  \source{kleiner-lott}
  \uses{def:kl-standard-solutions-standard-solution-definition,thm:kl-canonical-neighborhood-theorem-compactness-canonical-neighborhoods}

\label{lemstandard0}
Let $(\R^3,g(\cdot))$ be a standard solution.  Then

(1) The curvature operator of $g$ is nonnegative.

(2) All derivatives of curvature are bounded for small time,
independent of the standard solution.

(3) The scalar curvature satisfies $\lim_{t \rightarrow a^-}
\sup_{x \in \R^3} R(x,t) \: = \: \infty$.

(4) $(\R^3,g(\cdot))$ is $\kappa$-noncollapsed at scales below $1$ on any time
interval contained in $[0,2]$, where
$\kappa$ depends only on the choice of the initial condition $g_0$.

(5) $(\R^3,g(\cdot))$ satisfies the conclusion of Theorem \ref{thmI.12.1},
in the sense that for any $\Delta t > 0$, there is an $r_0 > 0$ so that
for any point $(x_0, t_0)$ with $t_0 \ge \Delta t$ and
$Q = R(x_0, t_0) \ge r_0^{-2}$, the solution in
$\{(x,t) \: : \: \dist_{t_0}^2(x, x_0) < (\epsilon Q)^{-1}, t_0 - (\epsilon Q)^{-1}
\le t \le t_0\}$ is, after scaling by the factor $Q$, $\epsilon$-close to the
corresponding subset of a $\kappa$-solution.

Moreover, any Ricci flow which satisfies all of the conditions of Definition
\ref{defstandardsolution} except maximality of the time interval can be
extended to a standard solution. In particular, using short-time
existence \cite[Theorem 1.1]{Shi2 (1989)}, there is at least one standard solution.
\end{lemma}
\begin{proof}
  \uses{thm:kl-canonical-neighborhood-theorem-compactness-canonical-neighborhoods,thm:kl-no-local-collapsing-theorem-ii-existence-kappa-solutions}
(1) follows from \cite[Theorem 4.14]{Shi3 (1989)}.

(2) follows from Appendix \ref{applocalder}.

(3)
In view of (1), this is equivalent to saying that
$\lim_{t \rightarrow a^-}
\sup_{x \in \R^3} |\Rm|(x,t) \: = \: \infty$. The argument for
this last assertion
is in \cite[Chapter 6.7.2]{Chow-Knopf}. The proof in
\cite[Chapter 6.7.2]{Chow-Knopf}
is for the compact case but using the derivative estimates of
Appendix \ref{applocalder}, the same argument works in
the present case.


(4) See Theorem \ref{nolocalcollapse}.

(5) See Theorem \ref{thmI.12.1}.

The final assertion of the lemma follows from the method of
proof of (3).
\end{proof}

\section{Claim 2 of II.2.
The blow-up time for a standard solution is $\leq 1$
} \label{claim2}

\begin{lemma}[compactness for standard solutions]
  \label{lem:kl-blow-up-time-standard-solution-compactness-standard-solutions}
  \dcref{lemclaim2}
  \group{Kleiner--Lott Surgery Foundations}
  \source{kleiner-lott}
 (cf. Claim 2 of II.2) \label{lemclaim2}
Let $[0,T_S)$ be the maximal time interval such that
the curvature of all standard solutions is uniformly
bounded for every compact subinterval $[0,a]\subset [0,T_S)$.
Then on the time interval $[0,T_S)$, the family of standard solution converges
uniformly at (spatial) infinity  to the standard Ricci flow
on the round infinite cylinder $S^2\times \R$ of scalar curvature one.
In particular, $T_S$ is  at most $1$.
\end{lemma}
\begin{proof}
  \uses{lem:kl-standard-solutions-existence-kappa-solutions,thm:kl-maximum-principles-existence-curvature-operators}
Let $\{\M_i\}_{i=1}^\infty$ be a sequence of standard solutions, and
let $\{x_i\}_{i=1}^\infty$ be a sequence tending to infinity
in the time-zero slice $M$.

By (2) of Lemma \ref{lemstandard0}, the gradient estimates in Appendix \ref{applocalder},
 and Appendix \ref{subsequence}, every subsequence
of $\{\M_i, (x_i,0)\}_{i=1}^\infty$ has a subsequence which
converges in the pointed smooth topology on the time interval $[0,T_S)$.
Therefore, it suffices to show that if $\{\M_i, (x_i,0)\}_{i=1}^\infty$
converges to some pointed Ricci flow
$(\M_\infty, (x_\infty, 0))$ then
$\M_\infty$ is round cylindrical flow.


Since $g_i(0)=g_0$ for all $i$,  the sequence of pointed time-zero slices
$\{(M,x_i, g_i(0))\}_{i=1}^\infty$ converges in the pointed smooth topology
to the round cylinder, i.e. $(M_\infty,g_\infty(0))$ is a round cylinder
of scalar curvature $1$. Each time slice $(M_\infty,g_\infty(t))$
is biLipschitz equivalent to $(M_\infty,g_\infty(0))$. In particular,
it has two ends. As it also has nonnegative sectional curvature,
Toponogov's theorem implies that $(M_\infty,g_\infty(t))$ splits off
an $\R$-factor. Using the strong maximum principle,
the Ricci flow $\M_\infty$ splits off an $\R$-factor;
see Theorem \ref{splitting}.
Then using the uniqueness of the Ricci flow on the round $S^2$, it
follows that $\M_\infty$ is a standard shrinking cylinder, which proves the lemma.

In particular, $T_S\leq 1$.
\end{proof}

\section{Claim 4 of II.2.
The blow-up time of a standard solution is $1$}
\label{secclaim4}

\begin{lemma}[standard solutions lemma]
  \label{lem:kl-blow-up-time-standard-solution-standard-solutions-lemma}
  \dcref{lemclaim4}
  \group{Kleiner--Lott Surgery Foundations}
  \source{kleiner-lott}
  \uses{lem:kl-blow-up-time-standard-solution-compactness-standard-solutions}
 (cf. Claim 4 of II.2) \label{lemclaim4}
Let $T_S$ be as in  Lemma \ref{lemclaim2}. Then $T_S=1$.
In particular, every standard solution survives until time $1$.
\end{lemma}
\begin{proof}
  \uses{lem:kl-blow-up-time-standard-solution-compactness-standard-solutions,lem:kl-canonical-neighborhood-theorem-canonical-neighborhoods-lemma-8,lem:kl-standard-solutions-existence-kappa-solutions,lem:kl-three-dimensional-solutions-existence-kappa-solutions-4,prop:kl-asymptotic-scalar-curvature-asymptotic-volume-ratio-existence-kappa-solutions,thm:kl-canonical-neighborhood-theorem-compactness-canonical-neighborhoods,thm:kl-statement-pseudolocality-theorem-pseudolocality-theorem}
First, there is an
$\alpha > 0$
so that $T_S > \alpha$ \cite[Theorem 1.1]{Shi2 (1989)}.
In what follows we will apply Theorem \ref{thmI.12.1}. The
hypothesis of Theorem \ref{thmI.12.1} says that the flow
should exist on a time interval of duration at least one,
but by
rescaling we can apply Theorem \ref{thmI.12.1} just as well
with the alternative hypothesis that the
flow exists on a time interval of duration
at least $\alpha$.

Suppose that $T_S < 1$.
Then there is a sequence of standard solutions
$\{\M_i\}_{i=1}^\infty$,
times $t_i \rightarrow T_S$ and points $(x_i,t_i)\in \M_i$ so that
$\lim_{i \rightarrow \infty} R(x_i, t_i) = \infty$.

We first
argue that no subsequence of the points $x_i$ can go to infinity (with respect
to the time-zero slice). Suppose, after relabeling the subsequence, that
$\{x_i\}_{i=1}^\infty$ goes to infinity.
From Lemma \ref{lemclaim2},
for any fixed $t^\prime < T_S$ the pointed solutions
$(M,(x_i, 0), g_i(\cdot))$, defined for $t \in [0, t^\prime]$,
approach that of the shrinking cylinder on the same time interval.
Lemma \ref{lemstandard0} and the characterization of high-curvature regions from
Theorem \ref{thmI.12.1} implies a
uniform bound on high-curvature regions of the time derivative of $R$, of the form
(\ref{derestimates}). Then taking $t^\prime$
sufficiently close to $T_S$, we get a contradiction.
We conclude that outside of a compact
region the curvature stays uniformly bounded as $t \rightarrow T_S$;
compare with the proof of Lemma \ref{claim1}.
(Alternatively, one could apply Theorem \ref{thmI.10.1} to compact
approximants, as is done in II.2.)

Thus we may assume that
the sequence $\{x_i\}_{i=1}^\infty$ stays in a compact region of the
time-zero slice. By Theorem \ref{thmI.12.1},
there is a sequence $\epsilon_i \rightarrow 0$
so that after rescaling the pointed solution
$(\M, (x_i, t_i))$ by $R(x_i, t_i)$, the result is $\epsilon_i$-close
to the corresponding subset of an ancient solution. By
Proposition \ref{propI.11.4},
the ancient solutions have vanishing asymptotic volume ratio.
Hence for every $\beta > 0$, there is some $L < \infty$ so that
in the original unscaled solution, for large $i$ we have
$\vol \left( B(x_i, t_i, L \: R(x_i, t_i)^{- \: \frac12} ) \right)
\: \le \: \beta \left( L \: R(x_i, t_i)^{- \: \frac12} \right)^3$.
Applying the Bishop-Gromov inequality to the time-$t_i$ slices, we
conclude that for any $D >0$,
$\lim_{i \rightarrow \infty}
D^{-3} \vol(B(x_i, t_i, D)) \: = \: 0$.
However, this contradicts the
previously-shown fact that the solution extends smoothly to time $T_S < 1$
outside of a compact set.

Thus $T_S=1$.
\end{proof}

\begin{lemma}[standard solutions lemma]
  \label{lem:kl-blow-up-time-standard-solution-standard-solutions-lemma-3}
  \dcref{scalarblowup}
  \group{Kleiner--Lott Surgery Foundations}
  \source{kleiner-lott}
 \label{scalarblowup}
The infimal scalar curvature on the time-$t$ slice
tends to infinity as $t\ra 1^-$ uniformly for all standard solutions.
\end{lemma}
\begin{proof}
  \uses{lem:kl-blow-up-time-standard-solution-compactness-standard-solutions,lem:kl-canonical-neighborhood-theorem-canonical-neighborhoods-lemma-8,thm:kl-canonical-neighborhood-theorem-compactness-canonical-neighborhoods}
Suppose the lemma failed and
let $\{(\M_i,(x_i,t_i))\}_{i=1}^\infty$
be a sequence of pointed standard solutions,
with
$\{R(x_i, t_i)\}_{i=1}^\infty$
uniformly bounded and $\lim_{i \rightarrow \infty} t_i = 1$.

Suppose first that after passing to a subsequence,
the points $x_i$ go to infinity in the time-zero
slice. From Lemma \ref{lemclaim2}, for any $t^\prime \in [0,1)$ we have
$\lim_{i \rightarrow \infty} R^{-1}(x_i, t^\prime) \: = \: 1-t^\prime$.
Combining this with the derivative estimate
$\Big| \frac{\partial R^{-1}}{\partial t} \Big| \: \le \: \eta$ at high
curvature regions gives a contradiction; compare with the proof
of Lemma \ref{claim1}.
Thus the points $x_i$ stay in a compact region.
We can now use
the bounded-curvature-at-bounded-distance argument in Step 2 of
the proof of Theorem \ref{thmI.12.1} to extract
a convergent subsequence of $\{(\M_i,(x_i,0))\}_{i=1}^\infty$
with a limit Ricci flow solution
$(\M_\infty,(x_\infty,0))$ that exists on
the time interval $[0,1]$.
(In this case, the nonnegative curvature of the blowup
region $W$ comes from the
fact that a standard solution has nonnegative curvature.)
As in Step 3 of the proof of Theorem
\ref{thmI.12.1}, $\M_\infty$ will have
bounded curvature for $t \in [0,1]$.
Note that $\M_\infty$ is a standard solution.
This contradicts Lemma \ref{lemclaim2}.
\end{proof}

\section{Claim 5 of II.2. Canonical neighborhood property for
standard solutions}
\label{secclaim5}

Let $p$ be the center of the hemispherical region in the time-zero slice.

\begin{lemma}[existence for canonical neighborhoods]
  \label{lem:kl-canonical-neighborhood-property-standard-solutions-existence-canonical-neighborhoods}
  \dcref{claim5}
  \group{Kleiner--Lott Surgery Foundations}
  \source{kleiner-lott}
  \uses{lem:kl-three-dimensional-solutions-existence-kappa-solutions-4,lem:kl-three-dimensional-solutions-existence-kappa-solutions-5}
 (cf. Claim 5 of II.2) \label{claim5}
Given  $\epsilon > 0$ sufficiently small, there are constants
$\eta = \eta(\epsilon)$, $C_1 = C_1(\epsilon)$ and
$C_2 = C_2(\epsilon)$ so that
every standard solution $\M$ satisfies the conclusions of
Lemmas \ref{derivativeestimates} and \ref{nbhd},
except that the $\epsilon$-neck neighborhood
need not be strong. (Here the constants do not depend on the
standard solution.) More precisely,
any point $(x,t)$ is covered by one of the following cases :

1.The time $t$ lies in $(\frac34, 1)$ and
$(x,t)$ has an $\epsilon$-cap neighborhood
or a strong $\epsilon$-neck neighborhood as in
Lemma \ref{nbhd}.

2. $x \in B(p, 0, \epsilon^{-1})$, $t \in [0, \frac34]$ and
$(x,t)$ has an $\epsilon$-cap neighborhood as in Lemma \ref{nbhd}.

3. $x \notin B(p, 0, \epsilon^{-1})$, $t \in [0, \frac34]$ and
there is an $\epsilon$-neck $B(x,t,\epsilon^{-1} r)$
such that the solution in $P(x,t,\epsilon^{-1}r, -t)$ is, after
scaling with the factor $r^{-2}$, $\epsilon$-close to the appropriate
piece of the evolving round infinite cylinder.

Moreover, we have an estimate $R_{min}(t) \: \ge \: \const \: (1-t)^{-1}$, where
the constant does not depend on the standard solution.
\end{lemma}
\begin{proof}
  \uses{lem:kl-blow-up-time-standard-solution-compactness-standard-solutions,lem:kl-blow-up-time-standard-solution-standard-solutions-lemma,lem:kl-blow-up-time-standard-solution-standard-solutions-lemma-3,lem:kl-standard-solutions-existence-kappa-solutions,lem:kl-three-dimensional-solutions-existence-kappa-solutions-4,lem:kl-three-dimensional-solutions-existence-kappa-solutions-5,thm:kl-canonical-neighborhood-theorem-compactness-canonical-neighborhoods}
We first show that the conclusion of Lemma \ref{nbhd} is satisfied.

In view of Lemma \ref{scalarblowup}, there is a $\delta > 0$ so that
if $t \in (1-\delta, 1)$ then
we can apply Theorem \ref{thmI.12.1}
and Lemma \ref{nbhd} to a point $(x,t)$ to see that the conclusions
of Lemma \ref{nbhd} are satisfied in this case.
If $t \in [0, 1-\delta]$ and $x$ is sufficiently far from $p$
(i.e. $\dist_0(x,p) \ge D$ for an appropriate $D$) then Lemma \ref{lemclaim2} implies
that $(x,t)$ has a strong $\epsilon$-neck neighborhood or
there is an $\epsilon$-neck $B(x,t,\epsilon^{-1} r)$
such that the solution in $P(x,t,\epsilon^{-1}r, -t)$ is, after
scaling with the factor $r^{-2}$, $\epsilon$-close to the appropriate
piece of the evolving round infinite cylinder.

(To elaborate a bit on the last possibility,
the issue here is that
there is no backward extension of the solution to $t<0$.
Because of this, if $t >0$ is close to $0$ then the backward
neighborhood
$P(x,t,\epsilon^{-1}r, -t)$ will not exist for rescaled time one,
as required to have a strong $\epsilon$-neck neighborhood.
Since $\inf_{x \in M} R(x,0) = 1$, we know from
(\ref{Rlowerbound}) that
$R(x,t) \: \ge \: \frac{1}{1- \frac23 t}$.
Then if $t > \frac35$, the time from the initial
slice to $(x,t)$, after rescaled by the scalar curvature, is bounded below by
$t \: \frac{1}{1- \frac23 t} \: > \: 1$. In particular,
if $t \ge \frac34$ then $r^2 t$ is at least one and
we are ensured that
the backward neighborhood
$P(x,t,\epsilon^{-1}r, -t)$ does contain a strong
$\epsilon$-neck neighborhood.)

If $t \in [0, 1-\delta]$ and
$\dist_0(x,p) < D$ then,  provided that
$D$ and $\epsilon$ are chosen appropriately, we can say that
$(x,t)$ has an $\epsilon$-cap neighborhood.

We now show that the conclusion of
Lemma \ref{derivativeestimates} is satisfied.
If $t \in [1-\delta, 1)$ then the conclusion follows from
Theorem \ref{thmI.12.1} and Lemma \ref{derivativeestimates}.
If $\delta^\prime > 0$ is sufficiently
small and $t \in [0, \delta^\prime]$ then the conclusion follows from Appendix
\ref{applocalder}. If $t \in [\frac12 \: \delta^\prime,
1 \: - \: \frac12 \: \delta]$
then we have an upper scalar curvature bound from Lemma \ref{lemclaim4}.
From Hamilton-Ivey pinching (see Appendix \ref{phiappendix}), this implies
a double-sided sectional curvature bound.
The conclusion of Lemma \ref{derivativeestimates},
when $t \in [\frac12 \: \delta^\prime,
1 \: - \: \frac12 \: \delta]$, now follows from
the Shi estimates of Appendix \ref{applocalder}.

The last statement of the lemma follows from the estimate
$\Big| \frac{\partial R^{-1}}{\partial t} \Big| \: \le \: \const$, which holds
for $t$ near $1$
(see Lemmas \ref{derivativeestimates},
\ref{lemstandard0}(5) and \ref{scalarblowup})
and then can be extended to all $t \in [0,1)$ (see Lemma \ref{claim2}).
From Lemma \ref{scalarblowup}, $\lim_{t \rightarrow 1} R^{-1}(x,t)=0$ for
every $x$.  Thus $R^{-1}(x, t )\: \leq \: \const \: (1-t)$ for any $(x,t)$.
Equivalently,
\begin{equation}
R(x,t) \: \geq \: \const \: (1-t)^{-1}.
\end{equation}
\end{proof}

\section{Compactness of the space of standard solutions}
\label{standardlycompact}

\begin{lemma}[compactness for standard solutions]
  \label{lem:kl-compactness-space-standard-solutions-compactness-standard-solutions}
  \dcref{standardcompactness}
  \group{Kleiner--Lott Surgery Foundations}
  \source{kleiner-lott}
 \label{standardcompactness}
The family ${\mathcal ST}$ of pointed standard solutions $\{(\M,(p,0))\}$
is compact with respect to
pointed smooth convergence.
\end{lemma}
\begin{proof}
  \uses{lem:kl-blow-up-time-standard-solution-compactness-standard-solutions,lem:kl-blow-up-time-standard-solution-standard-solutions-lemma}
This follows immediately from  Appendix \ref{subsequence}
and the fact that the constant $T_S$ from Lemma \ref{lemclaim2}
is equal to $1$, by Lemma \ref{lemclaim4}.
\end{proof}


\section{Claim 1 of II.2. Rotational symmetry of standard solutions} \label{claimII.2.1}

Consider a standard solution $(M,g(\cdot))$.
Since the time-zero metric $g_0$ is rotationally symmetric, it is clear
by separation of variables that there is a rotationally symmetric
solution for some time interval $[0, T)$.  In this section we
show that every standard solution is rotationally
symmetric for each $t \in [0,1)$.
Of course this would follow from the uniqueness of the
standard solution; see \cite{Chen-Zhu,Lu-Tian}. But the direct argument given
here is the first step toward a uniqueness proof as in \cite{Lu-Tian}.

\begin{lemma}[standard solutions lemma]
  \label{lem:kl-rotational-symmetry-standard-solutions-standard-solutions-lemma}
  \dcref{rotational-symmetry-standard-solutions:lemma-1}
  \group{Kleiner--Lott Surgery Foundations}
  \source{kleiner-lott}
 (cf. Claim 1 of II.2)
Any Ricci flow solution in the space ${\mathcal ST}$ is
rotationally symmetric for all $t \in [0,1)$.
\end{lemma}
\begin{proof}
  \uses{lem:kl-blow-up-time-standard-solution-compactness-standard-solutions}
We first describe an evolution equation for vector fields which
turns out to send Killing vector fields to Killing vector fields.
Suppose that a vector field $u = \sum_m u^m \partial_m$ evolves by
\begin{equation} \label{Killingevolution}
u^m_t \: = \:
u^{m \: \: k}_{\: \: \: \: ;\: \: \:  k} + R^m_{\: \: \: i} u^i.
\end{equation}
Then
\begin{align} \label{bigmess1}
\partial_t (u^m_{\: \: \: ; i}) \: & = \:  u^m_{t \: ; i} \: + \:
(\partial_t \Gamma^m_{\: \: ki}) \: u^k \\
& = (u^{m \: \: k}_{\: \:
\: \: ;\: \: \:  k} + R^m_{\: \: \: k} u^k)_{; \: i} \: + \:
(\partial_t \Gamma^m_{\: \: ki}) \: u^k \notag \\
& = u^{m \: \: k}_{\: \: \: \: ;\: \: \:  ki} + R^m_{\: \: \: k;i} u^k +
R^m_{\: \: \: k} u^k_{\: \: ;i} \: + \:
(\partial_t \Gamma^m_{\: \: ki}) \: u^k \notag \\
& = u^{m \: \: k}_{\: \: \: \: ;\: \: \:  ik} -
R^m_{\: \: \: lki} \: u^{l \: \: \: k}_{\: \: ;} - R^k_{\: \: \: lki} \:
u^{m \: \: \: l}_{\: \: \: \: ;}
+ R^m_{\: \: \: k;i} u^k +
R^m_{\: \: \: k} u^k_{\: \: ;i} \: + \:
(\partial_t \Gamma^m_{\: \: ki}) \: u^k \notag \\
& = u^{m \: \: \: \: \: \: k}_{\: \: ;\:ki} -
R^m_{\: \: \: lki} \: u^{l \: \: \: k}_{\: \: ;} - R_{li} \:
u^{m \: \: \: l}_{\: \: \: \: ;}
+ R^m_{\: \: \: k;i} u^k +
R^m_{\: \: \: k} u^k_{\: \: ;i} \: + \:
(\partial_t \Gamma^m_{\: \: ki}) \: u^k \notag \\
& = u^{m \: \: \: \: \: \: k}_{\: \: ;\:ik} -
(R^m_{\: \: \: lki} u^l)_;^{\: \: k} -
R^m_{\: \: \: lki} \: u^{l \: \: \: k}_{\: \: ;} - R_{ki} \:
u^{m \: \: \: k}_{\: \: \: \: ;}
+ R^m_{\: \: \: k;i} u^k +
R^m_{\: \: \: k} u^k_{\: \: ;i} \: + \:
(\partial_t \Gamma^m_{\: \: ki}) \: u^k \notag \\
& = u^{m \: \: \: \: \: \: k}_{\: \: ;\:ik} -
R^{m \: \: \: \: \: \: \: k}_{\: \: \: lki;} u^l -
2 R^m_{\: \: \: lki} \: u^{l \: \: \: k}_{\: \: ;} - R_{ik} \:
u^{m \: \: \: k}_{\: \: \: \: ;}
+ R^m_{\: \: \: k;i} u^k +
R^m_{\: \: \: k} u^k_{\: \: ;i} \: + \:
(\partial_t \Gamma^m_{\: \: ki}) \: u^k. \notag
\end{align}
Contracting the second Bianchi identity gives
\begin{equation} \label{bigmess2}
R_{mlki;}^{\: \: \: \: \: \: \: \: \: \: \: k} = R_{il;m} - R_{im;l}.
\end{equation}
Also,
\begin{align} \label{bigmess3}
\partial_t \Gamma^m_{\: \: ki} & = \partial_t ( g^{ml} \Gamma_{lki}) =
2 R^{ml} \Gamma_{lki} - g^{ml} (R_{lk,i} + R_{li,k} - R_{ik,l}) \\
& = - R^m_{\: \: \: k;i} - R^m_{\: \: \: i;k} +
R^{\: \: \: \: \: m}_{ik;}.  \notag
\end{align}
Substituting (\ref{bigmess2}) and (\ref{bigmess3}) in (\ref{bigmess1})
gives
\begin{equation}
\partial_t (u^m_{\: ; i}) \: = \:
u^{m \: \: \: \: \: \: k}_{\: \: ;\:ik} -
2 R^m_{\: \: \: lki} \: u^{l \: \: \: k}_{\: \: ;} - R_{ik} \:
u^{m \: \: \: k}_{\: \: \: \: ;} +
R^m_{\: \: \: k} u^k_{\: \: ;i}.
\end{equation}
Then
\begin{align}
\partial_t (u_{j ; i}) \: & = \:
\partial_t (g_{jm} u^m_{\: \: \: ; i}) \: = \:
-2 R_{jm} u^m_{\: \: \: ; i} + g_{jm} \: \partial_t (u^m_{\: \: \: ; i}) \\
& = \:
u^{\: \: \: \: \: \: \: \: \: k}_{j  ;\:ik} -
2 R_{jlki} \: u^{l \: \: \: k}_{\: \: ;} - R_{ik} \:
u^{\: \: \: \: k}_{j  \: ;} -
R_{jk} u^k_{\: \: ;i} \notag \\
& = \:
u^{\: \: \: \: \: \: \: \: \: k}_{j  ;\:ik} +
2 R_{ikjl} \: u^{l \: \: \: k}_{\: \: ;} - R_{ik} \:
u^{\: \: \: \: k}_{j  \: ;} -
R_{kj} u^k_{\: \: ;i}. \notag
\end{align}
Equivalently, writing $v_{ij} = u_{j ; i}$ gives
\begin{equation}
\partial_t v_{ij} \:  = \:
v^{\: \: \: \: \: \: \: k}_{ij  ;k} +
2 R^{\: \: k \: \: l}_{i \: \: j} \: v_{kl} - R_{i}^{\: \: k}
v_{kj} -
R^k_{\: \: j} v_{ik}. \notag
\end{equation}
Then putting $L_{ij} = v_{ij} + v_{ji}$ gives
\begin{equation}
\partial_t L_{ij} \:  = \:
L_{ij;k}^{\: \: \: \: \: \: \: k} +
2 R^{\: \: k \: \: l}_{i \: \: j} \: L_{kl} - R_{i}^{\: \: k}
L_{kj} -
R^k_{\: \: j} L_{ik}.  \notag
\end{equation}

For any $\lambda \in \R$, we have
\begin{equation}
\partial_t (e^{2\lambda t} \: L_{ij} L^{ij}) \:  = \:
2 \lambda \: (e^{2\lambda t} \: L_{ij} L^{ij}) \: + \:
(e^{2\lambda t} L_{ij} L^{ij})_{;k}^{\: \: \: \: k} -
2 \: e^{2\lambda t} \: L_{ij;k} \: L^{ij;k} + Q(\Rm, e^{\lambda t}L),
\end{equation}
where $Q(\Rm, L)$ is an algebraic expression that is
linear in the curvature tensor $\Rm$ and
quadratic in $L$.
Putting $M_{ij} = e^{\lambda t}L_{ij}$ gives
\begin{equation} \label{M2}
\partial_t (M_{ij} M^{ij}) \:  = \:
2 \lambda \: M_{ij} M^{ij} \: + \:
(M_{ij} M^{ij})_{;k}^{\: \: \: \: k} -
2 \: M_{ij;k} \: M^{ij;k} + Q(\Rm, M).
\end{equation}

Suppose that we have a Ricci flow solution $g(t)$, $t \in [0, T]$,
with $g(0) \: = \: g_0$. Let $u(0)$ be a rotational Killing vector field
for $g_0$. Let $u_\infty(0)$ be its restriction to (any) $S^2$, which we will
think of as the $2$-sphere at spatial infinity.
Solve (\ref{Killingevolution}) for $t \in [0, T]$ with $u(t)$ bounded at
spatial infinity for each $t$;
due to the asymptotics coming from Lemma \ref{lemclaim2}
(which is independent of the
rotational symmetry question), there is no problem
in doing so. Arguing as in the proof of Lemma \ref{lemclaim2},
one can show that
for any $t \in [0, T]$, at spatial infinity $u(t)$ converges to
$u_\infty(0)$.
Construct $M_{ij}(t)$ from $u(t)$. As $u(0)$ is a Killing vector field,
$M_{ij}(0) = 0$. For any $t \in [0, T]$,
at spatial infinity the tensor
$M_{ij}(t)$ converges smoothly to zero.
Suppose that $\lambda$ is sufficiently negative,
relative to the $L^\infty$-norm of the sectional curvature on the
time interval $[0, T]$.
We can apply the maximum principle to (\ref{M2}) to conclude that
$M_{ij}(t) = 0$ for all $t \in [0, T]$. Thus $u(t)$ is a Killing
vector field for all $t \in [0, T]$.

To finish the argument,
as Ben Chow pointed out, any Killing vector field $u$
satisfies
\begin{equation}
u^{m \: \: k}_{\: \: \: \: ;\: \: \:  k} + R^m_{\: \: \: i} u^i
\: = \: 0.
\end{equation}
To see this, we use the Killing field equation to write
\begin{align}
0 \: & = \: u_{m;k}^{\: \: \: \: \: \: \: \: k} \: + \:
u_{k;m}^{\: \: \: \: \: \: \: \: k} \: = \:
u_{m;k}^{\: \: \: \: \: \: \: \: k} \: + \:
u_{k;m}^{\: \: \: \: \: \: \: \: k} \: - \: u_{k; \: \: m}^{\: \: \: \: k}
\: = \:  u_{m; \: \: \: k}^{\: \: \: \: \: k} \: + \:
u_{\: \: ;mk}^{k} \: - \: u_{\: \: ;km}^{k} \\
& = \:
u_{m; \: \: k}^{\: \: \: \: \: \: k} \: - \: R^k_{\: \: imk} \: u^i \: = \:
u_{m; \: \: k}^{\: \: \: \: \: k} \: + \: R_{mi} \: u^i. \notag
\end{align}

Then from (\ref{Killingevolution}), $u_t^m \: = \: 0$ and the Killing vector
fields are not
changing at all.  This implies that $g(t)$ is rotationally
symmetric for all $t \in [0,T]$.
\end{proof}

\section{Claim 3 of II.2. Uniqueness of the standard solution} \label{claim3}

In this section, which is not needed for the sequel,
we outline an argument for the
uniqueness of the standard solution.
We do this for the convenience of the reader.
Papers on the uniqueness issue are
\cite{Chen-Zhu,Lu-Tian}.
Our argument is somewhat different than that of
\cite[Proof of Claim 3 of Section 2]{Perelman2},
which seems to have some unjustified statements.

In general,
suppose that we have
two Ricci flow solutions $\M \equiv (M, g(\cdot))$ and
$\widehat{\M} \equiv (M, \widehat{g}(\cdot))$ with bounded curvature
on each compact time interval and the same initial condition.
We want to show that they coincide.
As the set of times for which $g(t) = \widehat{g}(t)$
is closed, it suffices to show that it is relatively open.  Thus
it is enough to show that $g$ and $\widehat{g}$ agree on $[0, T)$
for some small $T$.

We will carry out the Deturck trick in this noncompact setting,
using a time-dependent background metric as in
\cite[Section 2]{Anderson-Chow}. The idea is to
define a $1$-parameter family of metrics $\{h(t)\}_{t \in [0, T)}$ by
$h(t) \: = \: \phi^{-1}(t)^* g(t)$, where
$\{\phi(t)\}_{t \in [0, T)}$ is a $1$-parameter family of
diffeomorphisms of $M$ whose generator is the negative of the time-dependent
vector field
\begin{equation} \label{Weqn}
W^i(t) \: = \: h^{jk} \: \left( \Gamma(h)^i_{\: jk} \: - \:
\Gamma(\widehat{g})^i_{\: jk} \right),
\end{equation}
with $\phi_0 \: = \: \Id$.
More geometrically, as in \cite[Section 6]{Hamilton}, we consider
the solution of the harmonic heat flow equation
$\frac{\partial F}{\partial t} \: = \: \triangle F$ for maps
$F : M \rightarrow M$ between the
manifolds $(M, g(t))$ and $(M, \widehat{g}(t))$, with $F(0) = \Id$.

We now specialize to the case when
$(M, g(\cdot))$ and
$(M, \widehat{g}(\cdot))$ come from standard solutions.
The technical issue, which we do not address here, is to show
that a solution to the harmonic heat flow
will exist for some time interval $[0, T)$ with uniformly
bounded derivatives; see \cite{Chen-Zhu}.
One is allowed to use the asymptotics of
Section \ref{claim2} here and
from Section \ref{claimII.2.1},
one can also assume that all of the metrics are
rotationally invariant. In the rest of this section we assume
the existence of such a solution $F$.
By further reducing the time interval if
necessary, we may assume that $F(t)$ is a diffeomorphism of $M$
for each $t \in [0, T)$. Then $h(t) \: = \: F^{-1}(t)^* g(t)$.
Clearly $h(0) \: = \: g(0) \: = \:
\widehat{g}(0)$.

By Section \ref{claim2}, $g$ and $\widehat{g}$ have the
same spatial asymptotics, namely that of the shrinking cylinder.
We claim that this is also true for $h$. That is, we claim that
$(\M, h(\cdot))$ converges
smoothly to the shrinking cylinder solution on $[0, T)$.
It suffices to show that $F$ converges smoothly to the identity
on $[0, T)$.
Suppose not.  Let $\{x_i\}_{i=1}^\infty$ be a sequence of
points in
the time-zero slice so that no subsequence of
the pointed spacetime maps $(F, (x_i, 0))$ converges to the identity.
Using the derivative
bounds, we can extract a subsequence that converges
to some $\widetilde{F} \: : \: [0, T) \times \R \times S^2 \rightarrow
\R \times S^2$ in the pointed smooth topology.
However, $\widetilde{F}$ will satisfy the harmonic heat flow
equation from the shrinking cylinder $\R \times S^2$ to itself, with
$\widetilde{F}(0)$ being the identity, and will have bounded derivatives.
The uniqueness of $\widetilde{F}$ follows by standard methods.
Hence $\widetilde{F}(t)$ is the identity for all $t \in [0, T)$,
which is a contradiction.

By construction, the family of metrics $\{h(t)\}_{t \in [0, T)}$ satisfies
the equation
\begin{equation} \label{heqn}
\frac{dh_{ij}}{dt} \: = \: - \: 2 \: R_{ij}(h) \: + \:
\nabla(h)_i W_j \: + \: \nabla(h)_j W_i.
\end{equation}
In local coordinates,
the right-hand side of (\ref{heqn}) is a polynomial in
$h_{ij}$, $h^{ij}$, $h_{ij,k}$ and $h_{ij,kl}$. The leading term in
(\ref{heqn}) is
\begin{equation} \label{hder}
\frac{dh_{ij}}{dt} \: = \: h^{kl} \: \partial_k \partial_l h_{ij}
\: + \: \ldots.
\end{equation}
A particular solution of (\ref{heqn}) is $h(t) \: = \:
\widehat{g}(t)$, since if we had $g \: = \: \widehat{g}$ then
we would have $W = 0$ and $\phi_t = \Id$.

Put $w(t) \: = \: h(t) - \widehat{g}(t)$. We claim that
$w$ satisfies an equation of the form
\begin{equation} \label{wder}
\frac{dw}{dt} \: = \: - \: \nabla(\widehat{g})^* \nabla(\widehat{g}) w
\: + \: Pw \: + \: Qw,
\end{equation}
where $P$ is a first-order operator and $Q$ is a zeroth-order
operator. To obtain the leading derivative terms in (\ref{wder}),
using (\ref{hder}) we write
\begin{align}
\frac{dw_{ij}}{dt} \: & = \: h^{kl} \: \partial_k \partial_l h_{ij} \: - \:
\widehat{g}^{kl} \: \partial_k \partial_l \widehat{g}_{ij} \: + \: \ldots \\
& = \:
\widehat{g}^{kl} \: \partial_k \partial_l w_{ij} \: + \:
\left( h^{kl} \: - \: \widehat{g}^{kl} \right) \:
\partial_k \partial_l h_{ij} \: + \: \ldots \notag \\
& = \:
\widehat{g}^{kl} \: \partial_k \partial_l w_{ij} \: - \:
\widehat{g}^{ka} \: w_{ab} \: h^{bl} \:
\partial_k \partial_l h_{ij} \: + \: \ldots \notag \\
& = \:
\widehat{g}^{kl} \: \partial_k \partial_l w_{ij} \: - \:
\widehat{g}^{ka} \: h^{bl} \:
\partial_k \partial_l h_{ij} \: w_{ab} \: + \: \ldots \notag
\end{align}
A similar procedure can be carried out for the lower order terms,
leading to (\ref{wder}). By construction the operators $P$ and $Q$
have smooth coefficients which, when expressed in terms of orthonormal frames,
will be bounded on $M$. In fact, as
$h$ and $\widehat{g}$
have the same spatial asymptotics, it follows from
\cite[Proposition 4]{Anderson-Chow} that
the operator on the right-hand side of (\ref{wder}) converges at
spatial infinity to
the Lichnerowicz Laplacian $\triangle_L(\widehat{g})$.

By assumption, $w(0) \: = \: 0$. We now claim that $w(t) \: = \: 0$
for all $t \in [0, T)$.
Let $K \subset M$ be a codimension-zero compact submanifold-with-boundary.
For any $\lambda \in \R$, we have
\begin{align}
& e^{2\lambda t} \: \frac{d}{dt} \: \left( \frac{1}{2} \: e^{-2\lambda t} \:
\int_K |w(t)|^2 \: \dvol_{\widehat{g}(t)} \right) \:  = \\
& \int_K \left[ \left( -\lambda - \frac{R}{2} \right) |w|^2 \: + \:
\langle w, \frac{dw}{dt} \rangle \right] \: \dvol_{\widehat{g}(t)} = \notag \\
&
\int_K \left[ \left( -\lambda - \frac{R}{2} \right) |w|^2 \: - \:
|\nabla(\widehat{g}) w|^2 \: +
w_{ab} P^{abcdi} \nabla(\widehat{g})_i w_{cd} \: + \:
\langle w,  Q w \rangle \right] \: \dvol_{\widehat{g}(t)}
\pm \notag \\
& \int_{\partial K} \langle w, \nabla_n w \rangle \:
\dvol_{\partial \widehat{g}(t)} = \notag \\
& \int_K \left[ \left( -\lambda - \frac{R}{2} \right) |w|^2 \: - \:
|\nabla(\widehat{g})_i w^{cd} - \frac{1}{2} P^{abcdi} w_{ab}|^2 \:
+ \frac{1}{4} |P^{abcdi} w_{ab}|^2 \: + \:
\langle w,  Q w \rangle \right] \: \dvol_{\widehat{g}(t)} \pm \notag \\
& \int_{\partial K} \langle w, \nabla_n w \rangle
\: \dvol_{\partial \widehat{g}(t)}. \notag
\end{align}
Choose
\begin{equation}
\lambda \: > \: \sup_{v \neq 0} \frac{
\frac{1}{4} |P^{abcdi} v_{ab}|^2 \: + \:
\langle v,  Q v \rangle
}{\langle v, v \rangle}.
\end{equation}
On any subinterval $[0, T^\prime] \subset [0, T)$, since
$w$ converges to zero at infinity
and $(M, \widehat{g}(t))$
is standard at infinity, by choosing $K$ appropriately we can make
$\int_{\partial K} \langle w, \nabla_n w \rangle
\: \dvol_{\partial \widehat{g}(t)}$ small.
It follows that there is an exhaustion
$\{K_i\}_{i=1}^\infty$ of $M$ so that
\begin{equation}
e^{2\lambda t} \: \frac{d}{dt} \: \left( \frac{1}{2} \: e^{-2\lambda t} \:
\int_{K_i} |w(t)|^2 \: \dvol_{\widehat{g}(t)} \right) \: \le \:
\frac{1}{i}
\end{equation}
for $t \in [0, T^\prime]$.
Then
\begin{equation}
\int_{K_i} |w(t)|^2 \: \dvol_{\widehat{g}(t)} \: \le \:
\frac{e^{2\lambda t}-1}{\lambda i}
\end{equation}
for all $t \in [0, T^\prime]$.
Taking $i \rightarrow \infty$ gives $w(t) = 0$.

Thus $h \: = \: \widehat{g}$. From (\ref{Weqn}), $W = 0$ and so
$h \: = \: g$. This shows that if $\M, \widehat{\M} \in
{\mathcal ST}$ then $\M = \widehat{\M}$.

\section{II.3. Structure at the first singularity time} \label{II.3}

This section is concerned with the structure of the Ricci flow solution at the
first singular time, in the case when the solution does go singular.

Let $M$ be a connected closed oriented $3$-manifold.  Let $g(\cdot)$ be a Ricci flow on
$M$ defined on a maximal time interval $[0,T)$ with $T < \infty$. One knows
that $\lim_{t \rightarrow T^-} \max_{x \in M} |\Rm|(x,t) \: = \: \infty$.

From Theorem \ref{nolocalcollapse} and Theorem \ref{thmI.12.1}, given
$\epsilon > 0$ there are numbers $r = r(\epsilon) > 0$ and
$\kappa = \kappa(\epsilon) > 0$ so that for any
point $(x,t)$ with $Q = R(x,t)  \ge r^{-2}$, the solution in
$P(x,t, (\epsilon Q)^{- \: \frac12}, (\epsilon Q)^{- 1})$ is (after rescaling
by the factor $Q$) $\epsilon$-close
to the corresponding subset of a $\kappa$-solution. By Lemma
\ref{derivativeestimates}, the estimate (\ref{derestimates})
holds at $(x,t)$,
provided $\epsilon$ is sufficiently small. In addition,
there is a neighborhood $B$ of $(x,t)$
as described in Lemma  \ref{nbhd}.  In particular, $B$ is a strong $\epsilon$-neck,
an $\epsilon$-cap or a closed manifold with positive sectional curvature.

If $M$ has positive sectional curvature at some time $t$ then it
is diffeomorphic
to a finite quotient of the round $S^3$ and shrinks to a point at time
$T$ \cite{Hamiltonnnnn}.
The topology of $M$ satisfies the conclusion of the geometrization
conjecture and $M$ goes extinct in a finite time. Therefore
for the remainder of this section we will assume that the sectional curvature
does not become everywhere positive.

We now look at the behavior of the Ricci flow as one approaches the
singular time $T$.

\begin{definition}[bounded-curvature regular set]
  \label{def:kl-structure-at-first-singularity-time-bounded-curvature-regular-set}
  \dcref{structure-at-first-singularity-time:definition-1}
  \group{Kleiner--Lott Singularities and Surgery}
  \source{kleiner-lott}

Define a subset $\Omega$ of $M$ by
\begin{equation}
\Omega \: = \: \{x \in M \: : \: \sup_{t \in [0,T)} |\Rm|(x,t) \: < \: \infty \}.
\end{equation}
\end{definition}

Suppose that $x \in M - \Omega$, so there is a sequence of times $\{t_i\}$ in $[0,T)$ with
$\lim_{i \rightarrow \infty} t_i =T$  and
$\lim_{i \rightarrow \infty} |\Rm|(x,t_i) = \infty$. As $\min_M R(\cdot, t)$ in
nondecreasing in $t$, the largest sectional curvature at $(x, t_i)$
goes to infinity as $i \rightarrow \infty$. Then by the $\Phi$-almost nonnegative sectional
curvature result of Appendix \ref{phiappendix}, $\lim_{i \rightarrow \infty} R(x, t_i) = \infty$.
From the time-derivative estimate of (\ref{derestimates}), $\lim_{t \rightarrow T^-}
R(x,t) = \infty$. Thus $x \in M-\Omega$ if and only if $\lim_{t \rightarrow T^-}
R(x,t) = \infty$.

\begin{lemma}[structure at the first singularity time lemma]
  \label{lem:kl-structure-at-first-singularity-time-structure-at-first-singularity-time-lemma}
  \dcref{omegaopen}
  \group{Kleiner--Lott Singularities and Surgery}
  \source{kleiner-lott}
 \label{omegaopen}
$\Omega$ is open in $M$.
\end{lemma}
\begin{proof}
Given $x \in \Omega$, using the time-derivative estimate in
(\ref{eqnderestimatesrinverse})
gives a bound of the form $|R(x,t)| \le C$ for $t \in [0,T)$.
Then using the spatial-derivative estimate in
(\ref{eqnderestimatesrinverse}) gives
a number $\widehat{r} > 0$ so that
so that $|R(\cdot,t)| \le 2C$ on $B(x,t,\widehat{r})$, for each
$t \in [0, T)$.
The $\Phi$-almost nonnegative sectional
curvature implies a bound of the form $|\Rm(\cdot,t)| \le C^\prime$ on
$B(x,t,\widehat{r})$, for each
$t \in [0, T)$.
Then the length-distortion estimate of Section \ref{secI.8.3}
implies that we can pick a neighborhood $N$ of $x$ so that
$|\Rm| \le C^\prime$ on $N \times [0,T)$.
Thus $N \subset \Omega$.
\end{proof}


\begin{lemma}[compactness for structure at the first singularity time]
  \label{lem:kl-structure-at-first-singularity-time-compactness-structure-at-first-singularity-time}
  \dcref{omeganoncompact}
  \group{Kleiner--Lott Singularities and Surgery}
  \source{kleiner-lott}
 \label{omeganoncompact}
Any connected component $C$ of $\Omega$ is noncompact.
\end{lemma}
\begin{proof}
Since $M$ is connected, if $C$ were
compact then it would  be all of $M$.  This contradicts the assumption that
there is a singularity at time $T$.
\end{proof}

We remark that {\it a priori}, the structure of $M - \Omega$ can be quite
complicated.  For example, it is not ruled out that an accumulating
collection of
$2$-spheres in $M$ can simultaneously shrink to points.
That is, $M - \Omega$ could have
a subset of the form $(\{ 0 \} \cup \{\frac{1}{i}\}_{i=1}^\infty )
\times S^2 \subset (-1, 1) \times S^2$, the picture being that
$\Omega$ contains a sequence of smaller and smaller adjacent double horns.
One could even imagine a
Cantor set's worth of $2$-spheres simultaneously shrinking, although conceivably
there may be additional arguments to rule out both of these cases.

\begin{lemma}[rigidity for structure at the first singularity time]
  \label{lem:kl-structure-at-first-singularity-time-rigidity-structure-at-first-singularity-time}
  \dcref{standard}
  \group{Kleiner--Lott Singularities and Surgery}
  \source{kleiner-lott}
 \label{standard}
If $\Omega = \emptyset$ then $M$ is diffeomorphic
to $S^3$, $\R P^3$, $S^1 \times S^2$ or $\R P^3 \# \R P^3$.
\end{lemma}
\begin{proof}
  \uses{lem:kl-three-dimensional-solutions-existence-kappa-solutions,lem:kl-three-dimensional-solutions-existence-kappa-solutions-5}
The time-derivative estimate in (\ref{eqnderestimatesrinverse}) implies that
for $t$ slightly less than $T$, we have
$R(x,t) \ge r^{-2}$ for all
$x \in M$. Thus at that time, every $x \in M$ has a neighborhood that is in an
$\epsilon$-neck or an $\epsilon$-cap, as described in Lemma \ref{nbhd}.
(Recall that we have already excluded the positively-curved case of the lemma.)

As in the proof of Lemma \ref{noncompactsoliton}, by splicing together the projection
maps associated with neck regions, one obtains an open subset  $U\subset M$
and a $2$-sphere
fibration $U\ra N$ where the fibers are nearly totally geodesic, and the complement
of $U$  is contained in a union of $\eps$-caps.   It follows that $U$
is connected. If
 there are any $\epsilon$-caps then there must be exactly two of them $U_1,\;U_2$,
and they may be chosen to intersect $U$ in connected  open sets $V_i=U_i\cap U$
which are isotopic to product regions in both $U$ and in the $U_i$'s.
 The caps being diffeomorphic to $\overline{B^3}$ or
$\R P^3 - B^3$, it follows that $M$ is diffeomorphic to $S^3$, $\R P^3$ or $\R P^3
\# \R P^3$ if $U\neq M$; otherwise  $M$ is diffeomorphic to an $S^2$ bundle over a circle,
and the orientability assumption implies that this bundle is diffeomorphic
to $S^1\times S^2$.
\end{proof}

In the rest of this section we assume that $\Omega \neq \emptyset$.
From the local derivative estimates of Appendix \ref{applocalder}, there is a smooth
Riemannian metric $\overline{g} \: =\: \lim_{t \rightarrow T^-} g(t) \Big|_\Omega$
on $\Omega$. Let $\overline{R}$ denote its scalar curvature.
Thus
the scalar curvature function extends to a continuous function on
the subset $(M\times [0,T))\cup (\Omega\times\{T\})\subset M\times [0,T]$.

\begin{lemma}[volume lemma]
  \label{lem:kl-structure-at-first-singularity-time-volume-lemma}
  \dcref{volumeomega}
  \group{Kleiner--Lott Singularities and Surgery}
  \source{kleiner-lott}
 \label{volumeomega}
$(\Omega, \overline{g})$ has finite volume.
\end{lemma}
\begin{proof}
From the lower scalar curvature bound of (\ref{Rlowerbound}) and the formula
$\frac{d}{dt} \vol(M, g(t)) \: = \: - \: \int_M R \: \dvol_M$,
we obtain an estimate of the form
$\vol(M, g(t)) \: \le \: \const + \const t^{\frac32}$, for
$t < T$.
The lemma follows.
\end{proof}

\begin{lemma}[existence for structure at the first singularity time]
  \label{lem:kl-structure-at-first-singularity-time-existence-structure-at-first-singularity-time}
  \dcref{lemrinverseextends}
  \group{Kleiner--Lott Singularities and Surgery}
  \source{kleiner-lott}

\label{lemrinverseextends}
There is a open neighborhood $V$ of $(M-\Omega)\times\{T\}$
in $M\times [0,T]$ such that $R^{-1}$  extends to a
continuous function on $V$ which vanishes on $(M-\Omega)\times\{T\}$.
\end{lemma}
\begin{proof}
  \uses{lem:kl-structure-at-first-singularity-time-structure-at-first-singularity-time-lemma}
As observed above Lemma \ref{omegaopen}, $x\in M-\Omega$ if and only if
$\lim_{t\ra T^-}\;R^{-1}(x,t)=0$. The lemma follows by applying
(\ref{eqnderestimatesrinverse}) to suitable spacetime paths.
\end{proof}


\begin{definition}[curvature sublevel of the regular set]
  \label{def:kl-structure-at-first-singularity-time-curvature-sublevel-regular-set}
  \dcref{structure-at-first-singularity-time:definition-7}
  \group{Kleiner--Lott Singularities and Surgery}
  \source{kleiner-lott}

For
$\rho < \frac{r}{2}$,
put $\Omega_\rho \: = \: \{x \in \Omega \: : \:
\overline{R}(x) \le \rho^{-2} \}$.
\end{definition}

\begin{lemma}[characterization for structure at the first singularity time]
  \label{lem:kl-structure-at-first-singularity-time-characterization-structure-at-first-singularity-time}
  \dcref{lemoverlinerisproper}
  \group{Kleiner--Lott Singularities and Surgery}
  \source{kleiner-lott}

\label{lemoverlinerisproper}
The function $\overline{R}:\Omega\ra \R$ is proper; equivalently,
if $\{x_i\}\subset \Omega$ is a sequence which leaves every compact
subset of $\Omega$, then $\lim_{i \rightarrow \infty} \overline{R}(x_i) = \infty$.
In particular,
$\Omega_\rho$ is a compact subset of $M$ for every $\rho<r$.
\end{lemma}
\begin{proof}
Suppose $\{x_i\}\subset \Omega$ is a sequence such that
$\{\overline{R}(x_i)\}$ is bounded.  After passing to a subsequence,
we may assume that $\{x_i\}$ converges to some point $x_\infty\in M$.
But $R^{-1}$ is well-defined and continuous on $V$, and vanishes on
$(M-\Omega)\times\{T\}$, so we must have $x_\infty\in \Omega$.
\end{proof}

We now consider the connected components of $\Omega$ according to whether
they intersect $\Omega_\rho$ or not. First, let $C$ be a connected component of
$\Omega$ that does not intersect $\Omega_\rho$.
Given $x \in C$, there is a neighborhood $B_x$ of $x$ which is $\epsilon$-close
to a region as described in Lemma \ref{nbhd}.
From Lemma \ref{omeganoncompact}, the neighborhood $B_x$ cannot be of
type (c) or (d) in the terminology of Lemma \ref{nbhd}.

We now introduce some terminology.

If a manifold $Z$ is diffeomorphic to $\R^3$ or $\R P^3-\ol{B^3}$
then any embedded $2$-sphere
$\Si\subset Z$ separates $Z$ into two connected subsets, one
of which has compact closure and the other contains the end
of $Z$. We refer to the first component as the {\em compact
side} and the other component as the {\em noncompact side}.

An open subset $R$ of a Riemannian manifold is a {\em good cylinder} if:

\begin{itemize}
\item It is
$\eps$-close, modulo rescaling, to a segment of a round
cylinder of scalar curvature $1$.

\item The diameter of $R$ is approximately
$100$ times its cross-section.

\item Every point in $R$, lies
in an $\eps$-neck in the ambient Riemannian manifold.
\end{itemize}
From Lemma \ref{nbhd},
every $\eps$-cap neighborhood $B_x$
contains a good cylinder lying in the $\eps$-neck at the end
of $B_x$.


\begin{lemma}[neck geometry lemma]
  \label{lem:kl-structure-at-first-singularity-time-neck-geometry-lemma}
  \dcref{doublehorn}
  \group{Kleiner--Lott Singularities and Surgery}
  \source{kleiner-lott}
  \uses{lem:kl-three-dimensional-solutions-existence-kappa-solutions-5}
 \label{doublehorn}
Suppose that for all $x \in C$,
the neighborhood $B_x$ can be taken to be a strong
$\epsilon$-neck as in case (a) of Lemma \ref{nbhd}.
Then $C$ is a double $\epsilon$-horn.
\end{lemma}
\begin{proof}
  \uses{lem:kl-structure-at-first-singularity-time-existence-structure-at-first-singularity-time,lem:kl-three-dimensional-solutions-existence-kappa-solutions}
Each point $x$ has an $\epsilon$-neck neighborhood.
We can glue these $\epsilon$-necks together to form a
submersion from $C$ to a $1$-manifold, with fiber $S^2$;
cf. the proof of Lemma \ref{noncompactsoliton}. (We can
do the gluing by successively adding on good
cylinders, where the intersections of successive cylinders
have diameter approximately 10 times the diameter of the cross-sections.)
In view of Lemma
\ref{lemrinverseextends}, it follows in this case that
$C$ is a double $\epsilon$-horn.
\end{proof}


\bigskip
\begin{lemma}[existence for cap geometry]
  \label{lem:kl-structure-at-first-singularity-time-existence-cap-geometry}
  \dcref{structure-at-first-singularity-time:lemma-10}
  \group{Kleiner--Lott Singularities and Surgery}
  \source{kleiner-lott}
  \uses{lem:kl-three-dimensional-solutions-existence-kappa-solutions-5}

Suppose that there is some $x \in C$ whose neighborhood
$B_x$ is an $\epsilon$-cap as in case (a) of Lemma \ref{nbhd}.
Then $C$ is a capped $\epsilon$-horn.
\end{lemma}
\begin{proof}
  \uses{lem:kl-structure-at-first-singularity-time-compactness-structure-at-first-singularity-time,lem:kl-structure-at-first-singularity-time-volume-lemma,lem:kl-three-dimensional-solutions-existence-kappa-solutions-5}

Put $p_1=x$.  Let $R$ be a good cylinder in the $\eps$-neck
at the end of $B_{p_1}$.  Now glue on successive good cylinders
to $R$,
as in the proof of the preceding lemma,
going away from $p_1$. \\ \\
Case 1 :
Suppose this gluing process can be continued
indefinitely.  Then taking the union of $B_{p_1}$ with
the good cylinders, we obtain an open subset $W$ of $C$ which
is diffeomorphic to $\R^3$ or $\R P^3-\ol{B^3}$.
We claim that $W$ is a closed subset of $\Om$.
If not then there is a sequence $\{x_k\}_{k=1}^\infty
\subset W$ converging to some
$x_\infty\in \Om-W$. This implies that $\{\ol{R}(x_k)\}_{k=1}^\infty$
remains bounded. In view of the overlap condition between successive
good cylinders, a subsequence of
$\{x_k\}_{k=1}^\infty$ lies in an infinite
number of mutually disjoint good cylinders, whose volumes have a positive lower
bound (because of the upper scalar curvature bound at the points $x_k$).
This contradicts Lemma \ref{volumeomega}.

Thus $W$ is open and closed in $\Om$. Hence $W=C$ and
we are done. \\ \\
Case 2 :
Now suppose that the gluing process cannot be continued beyond some
good cylinder $R_1$.  Then there must be a point $p_2\in R_1$
such that $B_{p_2}$ is an $\eps$-cap. Also, note that
the union $W_1$ of $B_{p_1}$ with the good cylinders is diffeomorphic
to $\R^3$ or $\R P^3-\overline{B^3}$, and that $R_1$ has compact
complement in $W_1$.   Let $\Si\subset R_1$ be a cross-sectional $2$-sphere
passing through $p_2$.

We first claim that if $V$ is the compact side of $\Si$ in $W_1$,
then $V$ coincides with the
compact side $V'$ of $\Si$ in $B_{p_2}$.  To see this, note that $V$ and $V'$
are both connected open sets disjoint from $\Si$, with topological
frontiers $\partial V=\partial V'=\Si$.  Then
$V - V^\prime \: = \: V \cap (C - (V^\prime \cup \Sigma))$
and we obtain two
open decompositions
\begin{equation}
\label{eqndecompv}
V=(V\cap V')\sqcup (V-V'),\quad V'=(V\cap V')\sqcup (V'-V).
\end{equation}
If $V\cap V'=\emptyset$, then $\ol{V}\cup\ol{V'}$ is a union of
two compact manifolds with the same boundary $\Si$, and disjoint interiors.
Hence it is an open
and closed subset of the connected component $C$, which contradicts
Lemma \ref{omeganoncompact}.  Thus $V\cap V'$ is nonempty.
By (\ref{eqndecompv}) and the connectedness of $V$ and $V^\prime$,
we get $V\subset V'$ and $V'\subset V$,
so $V=V'$ as claimed.

Next, we claim that if $R_2\subset B_{p_2}$ is a good cylinder
with compact complement in $B_{p_2}$, then $R_2$ is disjoint from $W_1$.
To see this,  note that  $R_2$ is disjoint from $\Si$ because
$p_2 \in \Si$ and the diameter
of $\Si$ is close to $\pi (\ol{R}(p_2)/6)^{-1/2}$, whereas by Lemma \ref{nbhd}
there is an $\eps$-neck $U\subset B_{p_2}$ with compact complement in
$B_{p_2}$, at distance at least $9000 \ol{R}(p_2)^{-1/2}$ from $p_2$.
Thus $R_2$ must lie in the noncompact side of $\Si$ in $B_{p_2}$, and
hence is disjoint from $V$.
As the good cylinder
$R_1 \ni p_2$ lies within $B(p_2, 1000 \ol{R}(p_2)^{-1/2}) \subset \Omega$,
it follows that $R_2$ is also disjoint from $R_1$,
so $R_2$ is disjoint from $W_1=V\cup R_1\subset B_{p_2}$.

We continue adding good cylinders to $R_2$ as long as we can.
If we come
to another cap point $p_3$ then we jump to its cap
$B_{p_3}$ and continue the process.
When so doing, we encounter
successive cap points $p_1, p_2, \ldots$ with associated caps
$B_{p_1}\subset B_{p_2} \subset \ldots$ and disjoint good cylinders
$R_1, R_2, \ldots$.
Since the ratio
$\frac{\sup_{B_{p_k}}\,\ol{R}}{\inf_{B_{p_k}}\,\ol{R}}$ has an
{\it a priori}
bound by Lemma \ref{nbhd},
in view of the disjoint good cylinders in $B_{p_k}$
we get $\vol(B_{p_k})\geq \const  k\,\ol{R}(p_1)^{-3/2}$.
Then Lemma \ref{volumeomega} gives
an upper bound on $k$.  Hence we encounter a finite
number of cap points.
Arguing as in Case 1,
we conclude that $C$ is a capped $\eps$-horn.
\end{proof}


We note that there could be an infinite number of connected components of
$\Omega$ that do not intersect $\Omega_\rho$.

Now suppose that $C$ is a connected component of $\Omega$ that intersects
$\Omega_\rho$. As $C$ is noncompact, there must be some point $x \in C$ that is
not in $\Omega_\rho$. Again, any such $x$ has a neighborhood $B$ as in Lemma
\ref{nbhd}.
If one of the boundary components of $B$ intersects
$\Omega_{2\rho}$ then we terminate the process in that direction.
For the directions of the boundary components of $B$ that do not
intersect $\Omega_{2\rho}$,
we perform the above algorithm of looking for an adjacent
$\epsilon$-neck, etc.
The only difference from before
is that in at
least one direction any such
sequence of overlapping $\epsilon$-necks will be finite, as it must
eventually intersect
$\Omega_{2\rho}$.
(In the other direction it may terminate
in $\Omega_{2\rho}$, in an $\epsilon$-cap, or not terminate at all.)
Once a cross-sectional $2$-sphere intersects $\Omega_{2\rho}$, if
$\epsilon$ is small then the entire $2$-sphere lies in $\Omega_\rho$.
Thus any connected component of $C - (C \cap \Omega_\rho)$
is contained in an $\epsilon$-tube with both boundary components in $\Omega_\rho$,
an $\epsilon$-cap with boundary in $\Omega_\rho$ or an $\epsilon$-horn with boundary in
$\Omega_\rho$. We note that $\Omega_\rho$ need not have a nice boundary.

There is an {\it a priori} $\rho$-dependent lower bound for the volume of any such
connected component of $C - (C \cap \Omega_\rho)$, in view of the fact that
it contains $\epsilon$-necks that adjoin $\Omega_\rho$.
From Lemma \ref{volumeomega},
there is a finite number of connected components of $\Omega$ that
intersect $\Omega_\rho$. Any such connected component has a finite
number of ends, each being an $\epsilon$-horn. Note that the
$\epsilon$-horns can be made disjoint, each with a quantitative lower
volume bound.

The surgery procedure, which will be described in detail in Section
\ref{II.4.4}, is performed as follows.  First, one throws away all
connected components of $\Omega$ that do not intersect $\Omega_\rho$.
For each connected component $\Omega_j$ of $\Omega$ that intersects
$\Omega_\rho$ and for each $\epsilon$-horn of $\Omega_j$, take
a cross-sectional sphere that lies far in the $\epsilon$-horn.
Let $X$ be what's left after cutting the $\epsilon$-horns
at these $2$-spheres and removing the tips. The
(possibly-disconnected) postsurgery manifold $M^\prime$ is the
result of capping off $\partial X$ by $3$-balls.

We now discuss how to reconstruct the original manifold
$M$ from $M^\prime$.

\begin{lemma}[structure at the first singularity time lemma]
  \label{lem:kl-structure-at-first-singularity-time-structure-at-first-singularity-time-lemma-11}
  \dcref{reconstruct}
  \group{Kleiner--Lott Singularities and Surgery}
  \source{kleiner-lott}
 \label{reconstruct}
$M$ is the result of taking connected sums of components of $M^\prime$
and possibly taking additional connected sums with a finite number of
$S^1 \times S^2$'s and $\R P^3$'s.
\end{lemma}
\begin{proof}
  \uses{lem:kl-three-dimensional-solutions-existence-kappa-solutions-5}

At a time shortly before $T$,
each point of $M - X$ has a neighborhood as in Lemma
\ref{nbhd}. The components of $M-X$ are $\epsilon$-tubes and
$\epsilon$-caps.  Writing $M^\prime = X \cup \bigcup B^3$ and
$M = X \cup (M-X)$, one builds $M$ from
$M^\prime$ as follows.
If the boundary of an $\epsilon$-tube of $M-X$ lies in
two disjoint components of $X$ then it gives rise to
a connected sum of two components of $M^\prime$.  If the
boundary of an $\epsilon$-tube lies in a single connected
component of $X$ then it gives rise to the
connected sum of the corresponding component of $M^\prime$
with a new copy of $S^1 \times S^2$.
If an $\epsilon$-cap
in $M-X$
is a $3$-ball it does not have any effect
on $M^\prime$. If
an $\epsilon$-cap is $\R P^3 - B^3$ then it gives rise to the
connected sum of the corresponding component of $M^\prime$
with a new copy of $\R P^3$. The lemma follows.
\end{proof}

\begin{remark}[estimates for diameter]
  \label{rem:kl-structure-at-first-singularity-time-estimates-diameter}
  \dcref{structure-at-first-singularity-time:remark-12}
  \group{Kleiner--Lott Singularities and Surgery}
  \source{kleiner-lott}

We do not assume that
the diameter of $(M, g(t))$ stays bounded as $t \rightarrow T$;
it is an open question whether this is the case.
\end{remark}

\section{Ricci flow with surgery:
the general setting}
\label{Ricci flow with surgery}

In this section we introduce some notation and terminology in order to
treat Ricci flows with surgery.

The principal purpose of sections II.4 and II.5 is to show that
one can prescribe the surgery procedure in such a way that Ricci flow
with surgery is well-defined for all time.  This involves showing
that

$\bullet$ One can give a sufficiently precise description of
the formation of singularities so that one can envisage defining a geometric
surgery.
In the case of the formation of the first singularity,
such a description was given in Section \ref{II.3}.

$\bullet$ The sequence of surgery times cannot accumulate.

\noindent
The argument in Section \ref{II.3} strongly uses both the $\kappa$-noncollapsing
result
of Theorem \ref{nolocalcollapse} and the characterization
of the geometry in a spacetime region around a point $(x_0, t_0)$ with
large scalar curvature, as given in Theorem \ref{thmI.12.1}.
The proofs of both of these results use the
smoothness of the solution at times before $t_0$. If surgeries occur
before $t_0$ then one must have strong control on the scales at which the
surgeries occur, in order to extend the arguments of Theorems
\ref{nolocalcollapse} and \ref{thmI.12.1}.
This forces one to consider time-dependent scales.

Section II.4 introduces Ricci flow with surgery, in varying degrees of
generality. Our treatment of this material follows Perelman's.
We have added some terminology to help formalize the surgery process.
There is some arbitrariness in this formalization, but the version given
below seems adequate.

For later  use, we now summarize the relevant notation
that we introduce.  More precise  definitions will be given below.
We will avoid using new notation as much as possible.

$\bullet$ $\M$ is a Ricci flow with surgery.

$\bullet$ $\M_t$ is the time-$t$ slice of $\M$.

$\bullet$ $\M_{\reg}$ is the set of regular points of $\M$.

$\bullet$ If $T$ is a singular time then $M_T^-$ is the limit of time slices
$\M_t$ as $t \rightarrow T^-$ (called $\Omega$ in II.4.1) and
$M_T^+$ is the outgoing time slice (for example,  the result of performing
surgery on $\Omega$). If $T$ is a nonsingular time then
$\M_T^- = \M_T^+ = \M_T$.

The basic notion of a Ricci flow with surgery is simply a sequence of
Ricci flows which ``fit together'' in the sense that the final
(possibly singular) time slice
of each flow is isometric, modulo surgery, to the initial time slice
of the next one.
\begin{definition}[ricci flow with surgery definition]
  \label{def:kl-ricci-flow-surgery-general-setting-ricci-flow-surgery-definition}
  \dcref{flowwithsurgery}
  \group{Kleiner--Lott Singularities and Surgery}
  \source{kleiner-lott}
   \label{flowwithsurgery}
A {\em Ricci flow with surgery} is given by

$\bullet$   A collection of Ricci flows
$\{(M_k\times[t_k^-,t_k^+),g_k(\cdot))\}_{1\leq k \leq N}$,
where $N\leq \infty$, $M_k$ is a compact (possibly empty) manifold,
$t_k^+=t_{k+1}^-$ for  all
$1\leq k<N$, and the flow $g_k$
goes singular at $t_k^+$
for each $k < N$. We allow $t_N^+$ to be $\infty$.

$\bullet$ A collection of limits
$\{(\Om_k,\bar g_k)\}_{1\leq k\leq N}$, in the sense of Section \ref{II.3},
at the respective final times $t_k^+$ that are
singular if $k < N$. (Recall that
$\Omega_k$ is an open subset of
$M_k$.)

$\bullet$ A collection of isometric embeddings
$\{\psi_k:X_k^+\ra X_{k+1}^-\}_{1\leq k<N}$
where $X_k^+\subset \Om_k$ and $X_{k+1}^-\subset M_{k+1}$,
$1 \le k< N$, are compact $3$-dimensional submanifolds with
boundary.   The $X_k^\pm$'s are the subsets which survive the
transition from one flow to the next, and the $\psi_k$'s give
the identifications between them.

We will say that $t$ is a {\em singular time} if $t=t_k^+=t_{k+1}^-$
for some $1\leq k<N$, or $t=t_N^+$ and the metric goes singular
at time $t_N^+$.
\end{definition}

A Ricci flow with surgery does not necessarily have to have any
real surgeries, i.e. it could be a smooth nonsingular flow.
Our definition allows Ricci flows with surgery that are more
general than those appearing in the argument for geometrization,
where the transitions/surgeries have a very special form.  Before
turning to these more special flows in Section \ref{II.4.4},
we first discuss some
basic features of Ricci flow with surgery.

It will be convenient to associate a (non-manifold) spacetime $\M$ to
the Ricci flow with surgery.  This is constructed by taking the disjoint union
of  the smooth manifolds-with-boundary
\begin{equation}
\left(M_k\times [t_k^-,t_k^+)\right)\cup \left(\Om_k\times\{t_k^+\}\right)
\subset M_k\times [t_k^-,t_k^+]
\end{equation}
 for $1\leq k\leq N$
and making identifications using the $\psi_k$'s as gluing
maps.   We denote the quotient space by $\M$ and the quotient map
by $\pi$.  We will sometimes also use $\M$ to refer to the
whole Ricci flow with surgery structure,
rather than just the associated spacetime.
The {\em time-$t$ slice $\M_t$ of $\M$} is the image
of the time-$t$ slices of the constituent Ricci
flows under the quotient map.

If $t=t_k^+$ is a singular time then we put
$\M_t^-=\pi(\Omega_k\times\{t_k^+\})$;  if in addition $t\neq t_N^+$ then
we put $\M_t^+=\pi(M_{k+1}\times\{t_{k+1}^-\})$. If $t$ is not a singular time
then we put  $\M_t^+=\M_t^-=\M_t$. We refer to
$\M_t^+$ and $\M_t^-$ as the forward and backward time
slices, respectively.

Let us summarize the structure of $\M$ near a singular time $t = t_k^+ = t_{k+1}^-$.
The backward time slice $\M_t^-$ is a copy of $\Omega_k$. The forward
time slice $\M_t^+$ is a copy of $M_{k+1}$. The time slice
$\M_t$ is the result of gluing $\Omega_k$ and $M_{k+1}$ using $\psi_k$.
Thus it is the disjoint union of $\Omega_k - X_k^+$, $M_{k+1} - X_{k+1}^-$ and
$X_k^+ \cong X_{k+1}^-$.  If $s > 0$ is small then in going from $M_{t-s}$ to
$M_{t+s}$, the topological change is that we remove
$M_k - X_k^+$ from $M_k$ and add $M_{k+1} - X_{k+1}^-$.

We let $\M_{(t,t')} = \cup_{\bar t\in(t,t')} \M_{\bar t}$ denote
the  {\em time slab} between $t$ and $t'$, i.e.
the union of the time slices between $t$ and $t'$.
The closed time
slab $\M_{[t,t']}$ is defined to be the closure of $\M_{(t,t')}$
in $\M$, so $\M_{[t,t']}=\M_t^+\cup\M_{(t,t')}\cup \M_{t'}^-$.
We  (ab)use  the notation
$(x,t)$ to denote a point $x\in\M$ lying in the time $t$ slice $\M_t$,
even though
$\M$ may no longer be a product.

The spacetime $\M$
has three types of points:  \\
1. The  $4$-manifold
points, which include all points at
nonsingular times in
$(t^-_1, t^+_N)$ and all
points in $\pi(\interior(X_k^+)\times \{t_k^+\})$
(or $\pi(\interior(X_{k}^-)\times \{t_{k}^-\})$) for $1 \le k < N$,
\\
2. The boundary points of $\M$, which are the images in $\M$ of
$M_1\times \{t_1^-\}$,
$\Omega_N \times \{t^+_N\}$,
$(\Om_k - X_k^+)\times \{t_k^+\}$ for
$1\leq k<N$, and $(M_k - X_k^-)\times \{t_k^-\}$ for $1<k\leq N$, and \\
3. The ``splitting'' points, which
 are the images in $\M$ of $\D X_k^+\times\{t_k^+\}$
for $1\leq k<N$.

Here the classification of points is according
to the smooth structure, not  the topology. In fact, $\M$ is  a topological
manifold-with-boundary.
We say that $(x,t)$ is {\em regular} if it is either a $4$-manifold
point, or it lies in the initial time slice $\M_{t_1^-}$ or final time slice
$\M_{t_N^+}$.
Let  $\M_{\reg}$ denote the set of regular points. It has a natural
smooth structure since the gluing maps $\psi_k$, being isometries
between smooth Riemannian manifolds, are smooth maps.


Note that the Ricci
flows on the $M_k$'s define a Riemannian metric $g$ on the ``horizontal''
subbundle of the tangent bundle of $\M_{\reg}$.  It follows from the
definition of the Ricci flow that $g$ is actually smooth on $\M_{\reg}$.


We metrize each time slice $\M_t$, and the forward and backward time
slices $\M_t^\pm$, by infimizing the path length of piecewise smooth paths.
We allow our distance functions to be infinite,
since the infimum will be infinite when points lie in different components.
If $(x,t)\in \M_t$ and $r>0$ then we let $B(x,t,r)$ denote the
corresponding metric ball. Similarly,  $B^\pm(x,t,r)$ denotes
the ball in $\M_t^\pm$ centered at $(x,t)\in\M_t^\pm$.
A ball $B(x,t,r)\subset\M_t$ is {\em proper} if the distance
function $d_{(x,t)}:B(x,t,r)\ra [0,r)$ is a proper function;
a proper ball ``avoids singularities'', except possibly at its
frontier.
Proper balls $B^\pm(x,t,r)\subset\M_t^\pm$ are defined likewise.

An {\em admissible curve} in $\M$ is a path $\ga:[c,d]\ra \M$,
with $\ga(t)\in \M_t$ for all $t\in [c,d]$, such that
for each $k$, the part of $\ga$ landing in $\M_{[t_k^-t_k^+]}$
lifts to a smooth map into
$M_k\times[t_k^-,t_k^+)\;\cup \;\Om_k\times\{t_k^+\}$.  We will
use $\dot{\ga}$ to denote the ``horizontal'' part of the velocity
of an admissible curve $\ga$.
If $t < t_0$,
a point $(x,t)\in\M$ is
{\em accessible from $(x_0,t_0)\in\M$} if there is an admissible curve
running from $(x,t)$ to $(x_0,t_0)$.  An
admissible curve $\ga:[c,d]\ra\M$ is {\em static} if
its lifts to the product spaces have constant first component. That is,
the points in the image of a  static curve are ``the same'',
modulo the passage of time and identifications taking place
at surgery times.
 A {\em barely
admissible curve} is an admissible curve $\ga:[c,d]\ra\M$
such that the  image is not contained
in $\M_c^+ \cup \M_{\reg}\cup \M_d^-$.
If $\ga:[c,d]\ra\M$ is barely admissible then
there is a surgery time $t=t_k^+=t_{k+1}^-\in(c,d)$ such that
$\ga(t)$ lies in
\begin{equation}
\pi(\D X_k^+\times\{t_k^+\})
=\pi(\D X_{k+1}^-\times\{t_{k+1}^-\}).
\end{equation}

If $(x,t)\in\M_t^+$, $r>0$, and $\De t>0$ then we define the
forward {\em parabolic region }
$P(x,t,r,\De t)$ to be the union of
(the images of) the static
admissible curves $\ga:[t,t^\prime] \ra \M$ starting in $B^+(x,t,r)$, where
$t^\prime \le t+\De t$. That is, we take the union of all the maximal
extensions of all static
curves, up to time $t + \De t$, starting from the
initial time slice $B^+(x,t,r)$.
When $\De t<0$, the parabolic region $P(x,t,r,\De t)$ is defined
similarly using static admissible curves ending in $B^-(x,t,r)$.

If $Y\subset \M_t$, and $t\in[c,d]$ then we say that $Y$
is {\em unscathed in $[c,d]$} if every point $(x,t)\in Y$
lies on  a static curve defined on the time
interval $[c,d]$.  If, for instance, $d=t$ then this will force
$Y\subset \M_t^-$.
The term ``unscathed'' is intended to capture the idea that
the set is unaffected by singularities and surgery.
(Sometimes Perelman uses the phrase
``the solution is defined in $P(x,t,r,\De t)$'' as synonymous with
``the solution is unscathed in $P(x,t,r,\De t)$'', for example in
the definition of canonical neighborhood in II.4.1.)
We  may use the
notation  $Y\times [c,d]$ for the set of points lying on static
curves $\ga:[c,d]\ra\M$ which pass through $Y$, when
$Y$ is unscathed on $[c,d]$.
Note that if $Y$ is open and
unscathed on $[c,d]$ then we can think of the Ricci flow on $Y\times [c,d]$
as an ordinary (i.e. surgery-free) Ricci flow.

The definitions of $\eps$-neck, $\eps$-cap, $\eps$-tube and
(capped/double) $\eps$-horn
from Section \ref{notationterminology}
do not require modification for
a Ricci flow with surgery, since they are just special types of Riemannian
manifolds; they will turn up as subsets of forward or backward
time slices of a Ricci flow with surgery. A {\em strong $\eps$-neck}
is a subset  of the form $U\times [c,d]\subset \M$, where $U\subset \M_d^-$
is an open set  that is unscathed on the interval $[c,d]$,
which is a strong $\eps$-neck in the sense of
Section \ref{notationterminology}.


\section{II.4.1. {\it A priori} assumptions}

This section introduces the notion of canonical neighborhood.

The following definition captures the geometric structure that
emerges by combining Theorem \ref{thmI.12.1} and its extension
to Ricci flows with surgery (Section \ref{II.5})
with the geometric description of $\kappa$-solutions.   The idea
is that blowups either yield $\kappa$-solutions, whose
structure is well understood from Section \ref{II.1}, or there are surgeries nearby
in the recent  past, in which case the local geometry resembles
that of the standard solution.   Both alternatives produce canonical
neighborhoods.

\begin{definition}[canonical neighborhood assumption]
  \label{def:kl-priori-assumptions-canonical-neighborhood-assumption}
  \dcref{canonicalnbhddef}
  \group{Kleiner--Lott Singularities and Surgery}
  \source{kleiner-lott}
  \uses{lem:kl-canonical-neighborhood-property-standard-solutions-existence-canonical-neighborhoods,lem:kl-three-dimensional-solutions-existence-kappa-solutions-4,lem:kl-three-dimensional-solutions-existence-kappa-solutions-5}
(Canonical neighborhoods, cf. Definition in II.4.1)
\label{canonicalnbhddef}
Let $\eps>0$ be
small enough so that Lemmas \ref{nbhd} and \ref{claim5} hold.
Let $C_1$ be the maximum of $30 \epsilon^{-1}$ and the $C_1(\eps)$'s of Lemmas
\ref{nbhd} and \ref{claim5}. Let $C_2$ be the maximum of the
$C_2(\epsilon)$'s of Lemmas
\ref{nbhd} and \ref{claim5}.
Let  $r:[a,b]\ra (0,\infty)$ be a positive nonincreasing function.
A Ricci flow with surgery $\M$ defined on the time interval $[a,b]$ satisfies
the {\em $r$-canonical neighborhood assumption} if every
$(x,t)\in \M_t^\pm$ with scalar curvature $R(x,t)\geq r(t)^{-2}$
has a canonical neighborhood in the corresponding (forward/backward)
time slice, as in Lemma \ref{nbhd}.  More precisely, there is
an  $\hat r\in (R(x,t)^{-\frac12},C_1R(x,t)^{-\frac12})$
and an open
set $U\subset \M_t^\pm$ with
$\ol{B^\pm(x,t,\hat r)}\subset U\subset B^\pm(x,t,2\hat r)$ that
falls into one of the
following categories :

(a) $U\times [t-\De t,t]\subset\M$  is a strong $\eps$-neck
for some $\De t>0$.  (Note that after parabolic rescaling the scalar
curvature at $(x,t)$ becomes $1$, so the scale factor must be
$\approx R(x,t)$, which implies that $\De t\approx R(x,t)^{-1}$.)

(b) $U$ is an $\eps$-cap which, after rescaling, is $\epsilon$-close to the
corresponding piece of a $\kappa_0$-solution or a time slice of a standard
solution (cf. Section \ref{II.2}).

(c) $U$ is a closed manifold diffeomorphic to $S^3$ or $RP^3$.


(d) $U$ is  $\epsilon$-close to
a closed manifold of constant positive sectional
curvature.

\noindent
Moreover, the scalar curvature in $U$ lies between $C_2^{-1}R(x,t)$
and $C_2R(x,t)$.
In cases (a), (b), and  (c), the volume of $U$ is greater than
$C_2^{-1}R(x,t)^{-\frac32}$. In case (c), the infimal
sectional curvature of $U$ is greater than $C_2^{-1}R(x,t)$.

Finally, we require that
\begin{equation}
\label{II.(1.3)estimates}
|\nabla R(x, t)|<\eta R(x, t)^{\frac32}, \: \: \:  \: \: \: \:
\left|\frac{\D R}{\D t} (x, t)\right|<\eta R(x, t)^2,
\end{equation}
where $\eta$ is the constant from (\ref{derestimates}).  Here the time
dervative $\frac{\D R}{\D t} (x, t)$ should be interpreted
as a one-sided derivative when the point $(x,t)$ is added or removed
during surgery at time $t$.
\end{definition}

\begin{remark}[continuation for canonical neighborhoods]
  \label{rem:kl-priori-assumptions-continuation-canonical-neighborhoods}
  \dcref{priori-assumptions:remark-2}
  \group{Kleiner--Lott Singularities and Surgery}
  \source{kleiner-lott}
  \uses{lem:kl-performing-surgery-continuing-flows-existence-canonical-neighborhoods,rem:kl-performing-surgery-continuing-flows-structural-remark-canonical-neighborhoods}

Note that the smaller of the
two balls in $\ol{B^\pm(x,t,\hat r)}\subset U\subset B^\pm(x,t,2\hat r)$ is closed,
in order to
make it easier to check the openness of the canonical neighborhood condition.
The requirement that $C_1$ be at least $30 \epsilon^{-1}$ will be used
in the proof of Lemma \ref{extendit}; see Remark \ref{30}.
\end{remark}

\begin{remark}[continuation for kappa solutions]
  \label{rem:kl-priori-assumptions-continuation-kappa-solutions}
  \dcref{priori-assumptions:remark-3}
  \group{Kleiner--Lott Singularities and Surgery}
  \source{kleiner-lott}

For convenience, in case (b) we have added the
extra condition that $U$ is $\epsilon$-close to the
corresponding piece of a $\kappa_0$-solution or
a time slice of a standard solution.
One does not need this
extra condition, but it is consistent to add it.
We remark that when surgery is performed according to the recipe
of Section \ref{II.4.4}, if a point $(p,t)$ lies in $\M_t^+   -  \M_t^-$ (i.e.
it  is ``added'' by surgery) then
it will sit in an $\eps$-cap, because $\M_t^+$ will resemble
a standard solution from Section \ref{II.2} near $(p,t)$.   Points lying
somewhat further out on the capped neck will belong to a strong
$\eps$-neck which extends backward in time prior to the surgery.
 \end{remark}

The next condition, which will ultimately be guaranteed by the
Hamilton-Ivey curvature pinching result and careful surgery,
is also essential in blowup arguments \`a la Section
\ref{I.12.1}.
\begin{definition}[curvature pinching assumption]
  \label{def:kl-priori-assumptions-curvature-pinching-assumption}
  \dcref{phipinchingdef}
  \group{Kleiner--Lott Singularities and Surgery}
  \source{kleiner-lott}
($\Phi$-pinching)
\label{phipinchingdef}
Let  $\Phi \in C^\infty(\R)$ be a positive nondecreasing
function such that for positive $s$,
$\frac{\Phi(s)}{s}$ is a decreasing function which tends to zero
as $s \rightarrow \infty$.  The Ricci flow with surgery $\M$
satisfies the {\em $\Phi$-pinching assumption} if for all
$(x,t) \in \M$, one has $\Rm(x,t) \geq -\Phi(R(x,t))$.
\end{definition}

We remark that the notion of $\Phi$-pinching here is somewhat
different from Perelman's $\phi$-pinching.  The purpose of
this definition is to distill out the properties of the
Hamilton-Ivey pinching condition which are needed in the rest of the
proof.

\begin{definition}[surgery a priori assumptions]
  \label{def:kl-priori-assumptions-surgery-priori-assumptions}
  \dcref{apriori}
  \group{Kleiner--Lott Singularities and Surgery}
  \source{kleiner-lott}
  \uses{def:kl-priori-assumptions-canonical-neighborhood-assumption,def:kl-priori-assumptions-curvature-pinching-assumption}
 \label{apriori}
A Ricci flow with surgery satisfies the {\em {\it a priori} assumptions}
if it satisfies the $\Phi$-pinching and $r$-canonical
neighborhood assumptions
on the time interval of the flow. Note that
the {\it a priori} assumptions depend on $\epsilon$, the function $r(t)$ of
Definition \ref{canonicalnbhddef}
and the function $\Phi$ of Definition \ref{phipinchingdef}.
\end{definition}



\section{II.4.2. Curvature bounds from the {\it a priori} assumptions}

In this section we state some technical lemmas about Ricci flows with surgery
that satisfy the {\em {\it a priori}} assumptions of the previous section.

The first one is the surgery
analog of  Lemma \ref{claim1}.

\begin{lemma}[estimates for curvature]
  \label{lem:kl-curvature-bounds-from-priori-assumptions-estimates-curvature}
  \dcref{claim1II.4.2}
  \group{Kleiner--Lott Singularities and Surgery}
  \source{kleiner-lott}
  \uses{def:kl-priori-assumptions-canonical-neighborhood-assumption}
(cf. Claim 1 of II.4.2)
\label{claim1II.4.2}
Given $(x_0,t_0)\in\M$ put $Q = |R(x_0,t_0)|+r(t_0)^{-2}$.
Then $R(x,t)\leq 8Q$ for all
$(x,t) \in P(x_0,t_0,\frac12\eta^{-1}Q^{-\frac12},
-\frac18\eta^{-1}Q^{-1})$, where $\eta$ is the constant from
(\ref{II.(1.3)estimates}).
\end{lemma}
\begin{proof}
  \uses{def:kl-priori-assumptions-canonical-neighborhood-assumption,lem:kl-canonical-neighborhood-theorem-canonical-neighborhoods-lemma-8}
The lemma follows from the estimates (\ref{II.(1.3)estimates}).
One integrates these derivative bounds along a subinterval of
a path that goes
in $B(x_0,t_0,\frac12\eta^{-1}Q^{-\frac12})$ and then
backward in time along a static path. See the proof of
Lemma \ref{claim1}.
We also use the fact that if $t^\prime \le t_0$ and
$R(x^\prime, t^\prime) \ge Q$ then the inequalities
(\ref{II.(1.3)estimates}) are valid at $(x^\prime, t^\prime)$,
since $r(\cdot)$ is nonincreasing.
\end{proof}

The next lemma expresses the main consequence of
Claim 2 of II.4.2.

\begin{lemma}[canonical neighborhoods lemma]
  \label{lem:kl-curvature-bounds-from-priori-assumptions-canonical-neighborhoods-lemma}
  \dcref{claim2II.4.2}
  \group{Kleiner--Lott Singularities and Surgery}
  \source{kleiner-lott}
(cf. Claim 2 of II.4.2)
\label{claim2II.4.2}
If $\epsilon$ is small enough then the following holds.
Suppose that $\M$ is a Ricci flow with surgery that satisfies the
$\Phi$-pinching assumption.
Then for any $A<\infty$ and $\widehat{r} > 0$ there exist
$\xi=\xi(A) > 0$ and $K=K(A,\widehat{r}) < \infty$
with the following property.
Suppose that $\M$ also satisfies the $r$-canonical neighborhood assumption
for some function $r(\cdot)$.
Then for any time $t_0$, if
$(x_0,t_0)$ is a point so that $Q = R(x_0,t_0) > 0$ satisfies
$\frac{\Phi(Q)}{Q} < \xi$
and $(x, t_0)$ is a point so that
$\dist_{t_0}(x_0, x) \le A Q^{- \: \frac12}$ then $R(x, t_0) \le KQ$,
where $K = K(A, r(t_0))$.
\end{lemma}
\begin{proof}
  \uses{def:kl-priori-assumptions-canonical-neighborhood-assumption,thm:kl-canonical-neighborhood-theorem-compactness-canonical-neighborhoods}
The proof is similar to the  proof of Step 2 of Theorem
\ref{thmI.12.1}.
(The canonical neighborhood assumption
replaces Step 1 of the proof of Theorem \ref{thmI.12.1}.)
Assuming that the lemma fails, one obtains a
piece of a nonflat metric cone as a blowup limit.
Using the canonical neighborhood assumption, one concludes
that the corresponding points in $\M$
have a neighborhood of type (a), i.e. a strong
$\epsilon$-neck, since
the neighborhoods of type (b), (c) and (d) of Definition
\ref{canonicalnbhddef} are not close to a piece of metric
cone. A strong $\epsilon$-neck, has the time interval needed to apply
the strong maximum
principle as in Step 2 of the proof of Theorem \ref{thmI.12.1}, in
order to get a contradiction.
\end{proof}

\section{II.4.3. $\delta$-necks in $\epsilon$-horns}

In this section we show that an $\epsilon$-horn has a
self-improving property as one goes down the horn.
For any $\delta > 0$, if the scalar curvature at a point
is sufficiently large then the point actually lies in a
$\delta$-neck.


In the statement of the
next lemma we will write $\Omega$
synonymously with the $\M_T^-$ of Section
\ref{Ricci flow with surgery}.

\begin{lemma}[existence for Ricci flow with surgery]
  \label{lem:kl-necks-horns-existence-ricci-flow-surgery}
  \dcref{lemmaII.4.3}
  \group{Kleiner--Lott Singularities and Surgery}
  \source{kleiner-lott}
  \uses{def:kl-priori-assumptions-surgery-priori-assumptions}
(cf. Lemma II.4.3) \label{lemmaII.4.3}
\label{hexists}

Given the pinching function $\Phi$, a number
$\widehat{T} \in (0, \infty)$, a positive nonincreasing
function $r \: : \: [0, \widehat{T}] \rightarrow \R$ and a number $\delta \in
\left( 0, \frac12 \right)$,
there is a
nonincreasing
function $h\::\:[0,\widehat{T}]\ra \R$ with
$0 < h(T) < \delta^2 r(T)$
so that the following property is satisfied.
Let
$\M$ be a Ricci flow with surgery defined on $[0, T)$, with $T < \widehat{T}$,
which satisfies the {\it a priori}
assumptions (Definition \ref{apriori}) and which goes singular
at time $T$. Let $(\Omega, \overline{g})$ denote the time-$T$ limit, in the sense of
Section \ref{II.3}.
Put $\rho = \de\; r(T)$ and
\begin{equation}
\Om_\rho=\{(x,t)\in \Omega \mid \overline{R}(x,T)\leq \rho^{-2}\}.
\end{equation}
\noindent
Suppose that $(x,T)$
lies in an $\eps$-horn $\h\subset \Omega$ whose boundary is contained in
$\Om_\rho$.
Suppose also that  $\overline{R}(x,T)\geq h^{-2}(T)$.
Then the parabolic
region $P(x,T,\de^{-1}\overline{R}(x,T)^{-\frac12},-\overline{R}(x,T)^{-1})$ is
contained in a strong $\de$-neck. (As usual, $\epsilon$ is
a fixed constant that is small enough so that the result holds
uniformly  with respect to the other variables.)
\end{lemma}
\begin{proof}
  \uses{def:kl-priori-assumptions-canonical-neighborhood-assumption,lem:kl-blow-up-time-standard-solution-compactness-standard-solutions,lem:kl-curvature-bounds-from-priori-assumptions-canonical-neighborhoods-lemma,lem:kl-structure-at-first-singularity-time-characterization-structure-at-first-singularity-time}
Fix $\delta \in (0,1)$. Suppose that the claim is not true.
Then there is a sequence of
Ricci flows with surgery $\M^\al$  and points $(x^\alpha, T^\alpha) \in
\M^\al$ with $T^\alpha < \widehat{T}$ such that \\
1. $\M^\al$ satisfies
the $\Phi$-pinching
and $r$-canonical neighborhood assumptions, \\
2. $\M^\al$ goes singular at time $T^\al$,\\
3. $(x^\al,T^\al)$ belongs to an $\eps$-horn $\h^\al\subset \Omega^\alpha$
whose boundary is contained in $\Om_\rho^\al$, and \\
4. $\overline{R}(x^\al,T^\al) \rightarrow \infty$, but \\
5. For each $\alpha$,
 $P(x^\al,T^\al,\de^{-1}\overline{R}(x^\al,T^\al)^{-\frac12},
 -\overline{R}(x^\al,T^\al)^{-1})$
is not contained in a strong $\delta$-neck.

Recall that when $\eps$ is small enough, any cross-sectional $2$-sphere
sitting in an $\eps$-neck $V\subset \h^\al$
separates the ends of $\h^\al$;
see Section \ref{notationterminology}.  We may
find a properly embedded minimizing geodesic $\ga^\al\subset \h^\al$
which joins the two
ends of $\h^\al$.   As $\ga^\al$ must intersect a cross-sectional
$2$-sphere containing $(x^\al,T^\al)$, it must pass within distance
$\leq 10\overline{R}(x^\al,T^\al)^{-\frac12}$ of $(x^\al,T^\al)$,
when $\eps$ is small.
Let $y^\al$ be the endpoint of $\gamma^\al$ contained in
$\Omega^\al_\rho$ and let $\widehat{y}^\al$ be the first point,
moving along $\gamma^\al$ from the noncompact end of $\h^\al$
toward $y^\al$,
where $\overline{R}(\widehat{y}^\al, T^\al) = \rho^{-2}$. As the gradient bound
$|\nabla \overline{R}^{- \: \frac12}| \le \frac12 \eta$ is valid along
$\gamma^\al$ starting from $\widehat{y}^\al$ and going out the noncompact end
(since such points on $\ga^\al$ have scalar curvature greater than
$r(T^\al)^{-2}$), we have $\dist_{T^\al}(x^\al, y^\al) \ge
\dist_{T^\al}(x^\al, \widehat{y}^\al) \ge \frac{2}{\eta}
\left(\rho - \overline{R}(x^\al, T^\al)^{- \: \frac12} \right)$.
Let $L^\al$ denote the time-$T^\al$ distance from $x^\al$ to the other
end of $\h^\al$. Since $\overline{R}$ goes to infinity as one exits the end,
Lemma \ref{claim2II.4.2} implies that $\lim_{\al \rightarrow \infty}
\overline{R}(x^\al, T^\al)^{\frac12} L^\al \: = \: \infty$.
From the existence of $\ga^\al$, whose length in either direction from
$x^\al$ is large compared to $\overline{R}(x^\al, T^\al)^{-\frac12}$,
it is clear that for large $\al$,
the canonical neighborhood of $(x^\al, T^\al)$ must be
of type (a) or (b) in the terminology of Definition \ref{canonicalnbhddef}.
By Lemmas \ref{lemoverlinerisproper} and \ref{claim2II.4.2},
we also know that for any fixed $\sigma < \infty$, for large
$\al$ the ball $B(x^\al, T^\al, \sigma \overline{R}(x^\al, T^\al)^{-\frac12})$
has compact closure in the time-$T^\al$ slice of $\M^\al$.

By Lemma \ref{claim2II.4.2}, after rescaling the metric on the
time-$T^\alpha$ slice by $\overline{R}(x^\al,T^\al)$ we have uniform curvature
bounds on distance balls.
We also have a uniform lower bound on the injectivity radius at
$(x^\al,T^\al)$ of the rescaled solution, in view of its
canonical neighborhood.  Hence after passing to
a subsequence, we may take a pointed
smooth complete limit $(M^\infty, x^\infty, g_\infty)$
of the time-$T^\al$ slices, where the derivative bounds needed to
take a smooth limit come from the canonical neighborhood assumption.
By the $\Phi$-pinching
assumption, $M^\infty$ will have nonnegative curvature.

After passing to a subsequence, we can also assume that
the $\ga^\al$'s converge to
a minimizing geodesic $\gamma$ in $M^\infty$ that passes within distance $10$
from $x^\infty$.
The
rescaled length of $\gamma^\al$ from $x^\al$ to $y^\al$ is
bounded below by $\frac{2}{\eta} \left( \overline{R}(x^\al, T^\al)^{\frac12} \rho - 1 \right)$,
which tends to infinity as $\al \rightarrow \infty$.
We have shown that the rescaled length of $\gamma^\al$ from
$x^\al$ to the other end of $\h^\al$ also tends to infinity as
$\al \rightarrow \infty$.
It follows that $\gamma$ is bi-infinite.
Thus by Toponogov's
theorem, $M^\infty$ splits off an $\R$-factor. Then for large $\al$,
the canonical neighborhood of $(x^\al,T^\al)$
must be an $\epsilon$-neck, and $M^\infty = \R \times S^2$
for some positively curved metric on $S^2$.  In particular,
$M^\infty$ has scalar curvature uniformly bounded above.

Any point
$\widehat{x} \in M^\infty$ is a limit of points
$(\widehat{x}^\alpha, T^\alpha) \in \M^\alpha$. As
$R_\infty(\widehat{x}) > 0$ and $R(x^\al,T^\al) \rightarrow \infty$, it
follows that $R(\widehat{x}^\alpha, T^\alpha) \rightarrow \infty$. Then
for large $\alpha$, $(\widehat{x}^\al,T^\al)$ is in a canonical
neighborhood which, in view of the $\R$-factor in $M^\infty$, must
be a strong $\epsilon$-neck.  From the upper bound on the scalar
curvature of $M^\infty$, along with the time interval involved in
the definition of a strong $\epsilon$-neck, it follows that
we can parabolically rescale the pointed flows
$(\M^\al,x^\al,T^\al)$ by $R(x^\al,T^\al)$, shift
time and
extract a smooth pointed limiting Ricci flow
$(\M^\infty,x^\infty,0)$ which is defined on a time
interval $(\xi,0]$, for some $\xi<0$.

In view of the strong $\epsilon$-necks around the
points $(\widehat{x}^\alpha, T^\alpha)$, if we take
$\xi$ close to zero
then we are ensured that the Ricci flow
$(\M^\infty,x^\infty,0)$ has positive
scalar curvature $R_\infty$. Given $(\widehat{x}, t) \in \M^\infty$,
as $R_\infty(\widehat{x}, t) > 0$ and $R(x^\al,T^\al) \rightarrow \infty$,
the $\Phi$-pinching implies that
the time-$t$ slice $\M^\infty_t$ has nonnegative curvature at
$\widehat{x}$. Thus $\M^\infty$ has nonnegative curvature.
The time-$0$ slice $\M^\infty_0$
splits off an $\R$-factor,
which means that the same will be true of all time slices;
cf. the proof of Lemma \ref{lemclaim2}. Hence
$\M^\infty$ is a product Ricci flow.

Let $\xi$ be the minimal negative number so that after parabolically rescaling
the pointed flows
$(\M^\al,x^\al,T^\al)$ by $R(x^\al,T^\al)$, we can extract a limit
Ricci flow $(\M^\infty, (x^\infty,0), g_\infty(\cdot))$ which is the product of $\R$ with
a positively curved Ricci flow on $S^2$, and is defined on the
time interval $(\xi, 0]$.  We claim that $\xi =  - \infty$. Suppose
not, i.e. $\xi > - \infty$. Given $(\widehat{x}, t) \in \M^\infty$,
as $R_\infty(\widehat{x}, t) > 0$ and
$R(x^\al,T^\al) \rightarrow \infty$, it follows that
$(x, t)$ is a limit of points $(\widehat{x}^\al, t^\al) \in \M^\alpha$
that lie in canonical
neighborhoods.  In view of the $\R$-factor in $\M^\infty$,
for large $\al$ these
canonical neighborhoods must be strong $\epsilon$-necks.
This implies in particular that
$\left( R_\infty^{-2} \: \frac{\partial R_\infty}{\partial t}\right)
(\widehat{x},t) > 0$, so
$\frac{\partial R_\infty}{\partial t}(\widehat{x},t) > 0$.
Then there
is a uniform upper bound $Q$ for the scalar curvature on $\M^\infty$.
Extending  backward from a time-$(\xi + \frac{1}{100Q})$ slice,
we can construct a limit Ricci flow that exists on some time interval
$(\xi^\prime, 0]$ with $\xi^\prime < \xi$. As before, using the
strong $\epsilon$-neck condition and the $\Phi$-pinching, if
$\xi^\prime$ is sufficiently close to $\xi$ then we are ensured
that the Ricci flow on $(\xi^\prime, 0]$ is the product of $\R$
with a positively curved Ricci flow on $S^2$. This is a contradiction.

Thus we obtain an ancient solution $\M^\infty$
with the property that each point
$(x, t)$ lies in a strong $\epsilon$-neck. Removing the $\R$-factor gives an
ancient solution on $S^2$. In view of the fact that each time slice is
$\epsilon$-close to the round $S^2$, up to rescaling, it follows that
the ancient solution on $S^2$ must be the standard shrinking
solution (see Sections \ref{I.11.3} and \ref{altpf11.3}).
Then $\M^\infty$ is the standard shrinking solution
on $\R \times S^2$. Hence for an infinite number of
$\alpha$,  $P(x^\al,T^\al,\de^{-1}R(x^\al,T^\al)^{-\frac12},
-R(x^\al,T^\al)^{-1})$ is in fact in a strong $\delta$-neck, which is a
contradiction.
\end{proof}

 \begin{remark}[monotonicity for neck geometry]
  \label{rem:kl-necks-horns-monotonicity-neck-geometry}
  \dcref{necks-horns:remark-2}
  \group{Kleiner--Lott Singularities and Surgery}
  \source{kleiner-lott}
  \uses{lem:kl-necks-horns-existence-ricci-flow-surgery}

If a given $h$ makes Lemma \ref{hexists} work for a given function $r$ then
one can check that logically, $h$ also works for any $r^\prime$
with $r^\prime \ge r$. Because of this, we may assume that $h$ only
depends on $\min r=r(T)$ and is monotonically nondecreasing as a
function of $r(T)$. Similarly, if a given $h$ makes Lemma \ref{hexists}
work for a given value of $\delta$ then $h$ also works for any
$\delta^\prime$ with $\delta^\prime  \ge\delta$.  Thus we may assume
that $h$ is monotonically nondecreasing as a function of $\delta$.
\end{remark}

\section{Surgery and the pinching condition} \label{surgery}

This section describes how one can take a $\de$-neck
satisfying the time-$t$ Hamilton-Ivey pinching condition,
and perform surgery so as to obtain a new manifold which also satisfies
the time-$t$ pinching condition, and which is $\de'$-close to the
standard solution modulo rescaling.
Here $\de'$ is a nonexplicit function of $\de$ but satisifes the important
property that  $\de'(\de)\ra 0$ as $\de\ra 0$.



The main geometric idea which handles the delicate part of the surgery
procedure is contained in the following lemma.  It says that one can ``round off''
the boundary of an approximate round half-cylinder so as to simultaneously
increase the scalar curvature  and  the minimum of sectional curvature at each
point.

As the statement of the following lemma involves the curvature
operator, we state our conventions. If $M$ has constant sectional
curvature $k$ then the curvature operator acts on $2$-forms as
multiplication by $2k$.  This is consistent with the usual
Ricci flow literature, e.g. \cite{Chow-Knopf}.

Recall that $\epsilon$ is our global parameter, which is taken
sufficiently small.

\begin{lemma}[Ricci flow with surgery lemma]
  \label{lem:kl-surgery-pinching-condition-ricci-flow-surgery-lemma}
  \dcref{lemroundoff}
  \group{Kleiner--Lott Singularities and Surgery}
  \source{kleiner-lott}

\label{lemroundoff}
Let $g_{cyl}$ denote the round cylindrical metric  of scalar
curvature $1$ on $\R \times S^2$. Let $z$ denote the coordinate
in the $\R$-direction.
Given $A>0$, suppose that $f \: : \: (-A,0]\ra \R$ is a
smooth function such that

$\bullet$ $f^{(k)}(0)=0$ for all $k\geq 0$.

$\bullet$  On $(-A,0)$,
\begin{equation}
f(z) <0\,,\quad f^\prime(z) >0\,,\quad f^{\prime \prime}(z)<0.
\end{equation}

$\bullet$ $\|f\|_{C^2}<\eps$.

$\bullet$  For every $z \in (-A,0)$,
\begin{equation}
\label{eqnderivcomparision}
\max\left(|f(z)|,\left|f^\prime(z) \right|\right)
\leq \eps\:\left| f^{\prime \prime}(z) \right|.
\end{equation}

Then if  $h_0$ is a smooth metric on $(-A,0] \times S^2$
with $\|h_0-g_{cyl}\|_{C^2}<\eps$  and we set
 $h_1= e^{2f(z)}h_0$, it follows that for all $p \in (-A,0) \times S^2$ we have
$R_{h_1}(p) \: > \: R_{h_0}(p) \: - \: f^{\prime\prime}(z(p))$.
Also, if $\la_1(p)$ denotes the lowest eigenvalue of the curvature operator at $p$ then $\la_1^{h_1}(p)>\la_1^{h_0}(p) \: - \: f^{\prime\prime}(z(p))$.
\end{lemma}
\begin{proof}
We will use the variational characterization of $\lambda_1(p)$ :
\begin{equation} \label{Rvar}
\lambda_1(p) \: = \:
\inf_{\omega \neq 0} \frac{\omega_{ij} \: R^{ij}_{\: \: \: kl} \: \omega^{kl}}{\omega_{ij} \: \omega^{ij}}
\end{equation}
where $\omega \in \Lambda^2(T_pM)$.
We will also use the following formulas about curvature quantities for
conformally related metrics in dimension $3$ :
\begin{equation} \label{confscalar}
R_{h_1} \: = \: e^{-2f} \: \left( R_{h_0} \:  - \: 4 \: \triangle f \: - \: 2 \: |\nabla f|^2 \right)
\end{equation}
and
\begin{equation} \label{confcurv}
R^{ij}_{\: \: \: kl}(h_1) \:  = \: e^{-2f} \left( R^{ij}_{\: \: \: kl}(h_0) \: - \: \widetilde{f}^i_{\: \: k} \:
\delta^j_{\: \: l}
 + \: \widetilde{f}^i_{\: \: l} \:
\delta^j_{\: \: k}
 + \: \widetilde{f}^j_{\: \: k} \:
\delta^i_{\: \: l}
 - \: \widetilde{f}^j_{\: \: l} \:
\delta^i_{\: \: k} \: - \: |\nabla f|^2
\left( \delta^i_{\: \: k} \: \delta^j_{\: \: l} \: - \:
\delta^i_{\: \: l} \: \delta^j_{\: \: k} \right)
\right),
\end{equation}
where $\widetilde{f}_{ij} \: = \: f_{;ij} \: - \: f_{;i} \: f_{;j}$. That is,
$\widetilde{f} \: = \: \Hess(f) \: - \: df \: \otimes \: df$. The right-hand sides of these
expressions are
computed using the metric $h_0$.

To motivate the proof, let us first consider the linearization of these expressions around
$h_0$. Keeping only the linear terms in $f$ gives
to leading order,
\begin{equation}
R_{h_1} \: \sim \: R_{h_0} \:  {-2f} \:  R_{h_0} \:  - \: 4 \: \triangle f
\end{equation}
and
\begin{equation}
R^{ij}_{\: \: \: kl}(h_1) \:  \sim \: R^{ij}_{\: \: \: kl}(h_0) \: - \: {2f} \:  R^{ij}_{\: \: \: kl}(h_0)
\: - \: {f}^{ \: i}_{; \: \: k} \:  \delta^j_{\: \: l}
\: + \: {f}^{ \: i}_{; \: \: l} \:  \delta^j_{\: \: k}
\: + \: {f}^{ \: j}_{; \: \: k} \:  \delta^i_{\: \: l}
\: - \: {f}^{ \: j}_{; \: \: l} \:  \delta^i_{\: \: k}.
\end{equation}
From the assumptions, $R_{h_0} \sim 1$ and $f < 0$ on $(-A,0) \times S^2$.
As $h_0$ is close to $g_{cyl}$, we have $\triangle f \sim f^{\prime \prime}(z)$,
so ${-2f} \:  R_{h_0} \:  - \: 4 \: \triangle f \: \ge \: - \: f^{\prime \prime}(z)$.
Similarly, in the case of $g_{cyl}$ a minimizer
$\omega$ in (\ref{Rvar}) is of the  form $\omega \: = \: X \wedge \partial_z$,
where $X$ is a unit vector in the $S^2$-direction.  As $h_0$ is close to
$g_{cyl}$, a minimizing $\omega$ for $h_0$ will be close to something of
the form $X \wedge \partial_z$. Then
\begin{equation}
\lambda_1^{h_1} \sim \lambda_1^{h_0} \: - \: 2f \: \lambda_1^{h_0}\: - \: 2 \:
f^{\prime \prime}(z) \: \ge \: \lambda_1^{h_0} \: + \: 2 f(z) \: |\lambda_1^{h_0}|\: - \: 2 \:
f^{\prime \prime}(z).
\end{equation}
As $h_0$ is close to $g_{cyl}$, $\lambda_1^{h_0}$ is close to
$\lambda_1^{g_{cyl}} = 0$. Then we can
use (\ref{eqnderivcomparision}) to say that $2 f(z) \: |\lambda_1^{h_0}|\: - \: 2 \:
f^{\prime \prime}(z) \: \ge \: - \: f^{\prime \prime}(z)$.

The remaining issue is to show that the increase in $R$ and $\lambda_1$
coming from the linear approximation is still approximately valid in the
nonlinear case, provided that $\epsilon$ is sufficiently small. For this,
we have to show that the increase from the linear approximation dominates
the error terms that we have neglected.

To deal with the scalar curvature first, from (\ref{confscalar}) we have
\begin{align}
R_{h_1} \: & = \: e^{-2f} \left( R_{h_0} \: - \: 4 \: \triangle_{g_{cyl}} f \right) \: + \:
4 e^{-2f}  \: (\triangle_{g_{cyl}} f \: - \: \triangle f) \:- \: 2 \: e^{-2f} \: |\nabla f|^2 \\
& \ge \: R_{h_0} \: - \: 4 \: f^{\prime \prime}(z)  \: + \:
4 e^{-2f}  \: (\triangle_{g_{cyl}} f \: - \: \triangle f) \:- \: 2 \: e^{-2f} \: |\nabla f|^2. \notag
\end{align}
Next, there is an estimate of the form
\begin{align}
|\triangle_{g_{cyl}} f \: - \: \triangle f| \: & \le \: \const \parallel h_0 - g_{cyl}
\parallel_{C^2} \: \left( |f(z)| \: + \: |f^\prime(z)| \: + \:
|f^{\prime \prime}(z)| \right) \\
& \le \:
\const \: \epsilon \: \left( |f(z)| \: + \: |f^\prime(z)| \: + \:
|f^{\prime \prime}(z)| \right). \notag
\end{align}
As $e^{-2f} \: \le \: e^{2 \epsilon}$, if $\epsilon$ is small then
\begin{equation}
\left| e^{-2f}  \: (\triangle_{g_{cyl}} f \: - \: \triangle f) \right| \: \le \: \const \:
\epsilon \: \left( |f(z)| \: + \: |f^\prime(z)| \: + \:
|f^{\prime \prime}(z)| \right)
\end{equation}
Similarly,
\begin{equation}
e^{-2f} \: |\nabla f|^2 \: \le \: \const  \:
|f^\prime(z)|^2 \: \le \: \const \: \epsilon \: |f^\prime(z)|.
\end{equation}
When combined with (\ref{eqnderivcomparision}), if $\epsilon$ is taken
sufficiently small
then
\begin{equation}
- \: 4 \: f^{\prime \prime}(z)  \: + \:
4 e^{-2f}  \: (\triangle_{g_{cyl}} f \: - \: \triangle f) \:- \: 2 \: e^{-2f} \: |\nabla f|^2
\: \ge \: - \: f^{\prime \prime}(z).
\end{equation}
This shows the desired estimate for $R_{h_1}(p)$.

To estimate $\lambda_1^{h_1}$ we use
(\ref{Rvar}) and (\ref{confcurv}) to write
\begin{equation} \label{min1}
\lambda_1^{h_1}(p) \: = \: e^{-2f(z)} \: \left(
\inf_{\omega \neq 0} \frac{\omega_{ij} \: R^{ij}_{\: \: \: kl}(h_0) \: \omega^{kl}
\: - \: 4 \: \omega_{ij} \widetilde{f}^i_{\: \: k} \omega^{kj}}{\omega_{ij} \: \omega^{ij}} \: - \: 2 \: |\nabla f|^2(z) \right).
\end{equation}
Comparing with
\begin{equation}
\lambda_1^{h_0}(p) \: = \:
\inf_{\omega \neq 0} \frac{\omega_{ij} \: R^{ij}_{\: \: \: kl}(h_0) \: \omega^{kl}}{
\omega_{ij} \: \omega^{ij}}
\end{equation}
gives
\begin{equation}
\lambda_1^{h_0}(p) \: \le \: e^{2f(z)} \: \lambda_1^{h_1}(p) \: + \:
\frac{4 \: \omega_{ij} \widetilde{f}^i_{\: \: k} \omega^{kj}}{\omega_{ij} \: \omega^{ij}} \: + \: 2 \: |\nabla f|^2(z),
\end{equation}
where $\omega$ is a minimizer in (\ref{min1}), or
\begin{equation}
\lambda_1^{h_1}(p) \: \ge \: e^{- 2f(z)} \: \lambda_1^{h_0}(p) \: - \:
4 \: e^{- 2f(z)} \:
\frac{\omega_{ij} \widetilde{f}^i_{\: \: k} \omega^{kj}}{\omega_{ij} \: \omega^{ij}}
\: - \: 2 \: e^{- 2f(z)} \: |\nabla f|^2(z).
\end{equation}
Using the variational formula (\ref{Rvar}), one can show that
$|\lambda_1^{h_0}(p)| \: \le \: \const \epsilon$.
From eigenvalue perturbation theory
\cite[Chapter 12]{Reed-Simon},
$\omega$ will be of the
form $X \wedge \partial_z \: + \: O(\epsilon)$ for some unit vector $X$ tangential
to $S^2$. Then we get an estimate
\begin{equation}
\lambda_1^{h_1}(p) \: \ge \: \lambda_1^{h_0}(p) \: - \:
2 f^{\prime \prime}(z) \: - \:
\const \: \epsilon \: \left( |f(z)| \: + \: |f^\prime(z)| \: + \:
|f^{\prime \prime}(z)| \right).
\end{equation}
From (\ref{eqnderivcomparision}), if $\epsilon$ is taken sufficiently small
then $\lambda_1^{h_1}(p) \: - \: \lambda_1^{h_0}(p) \: \ge \:  - \:
f^{\prime \prime}(z)$.
\end{proof}

Recall that the initial condition $\s_0$ for the standard solution  is
an $O(3)$-symmetric metric
$g_0$
 on $\R^3$ with nonnegative curvature operator, whose
end is isometric to a round half-cylinder of scalar curvature $1$.
To facilitate the surgery procedure, we will assume that some metric
ball around
the $O(3)$-fixed point has constant positive curvature. Outside of
this ball we use radial coordinates
$(z, \theta) \in (-B,\infty) \times S^2$, with
$g_{0} \: = \: e^{2F(z)} g_{cyl}$.
Here $g_{cyl}$ is the round cylindrical metric of scalar curvature one
and $F \in  C^\infty(-B,\infty)$.

\begin{lemma}[positivity for Ricci flow with surgery]
  \label{lem:kl-surgery-pinching-condition-positivity-ricci-flow-surgery}
  \dcref{capoff}
  \group{Kleiner--Lott Singularities and Surgery}
  \source{kleiner-lott}
  \uses{lem:kl-surgery-pinching-condition-ricci-flow-surgery-lemma}
 \label{capoff}
Given $A > 0$, we can choose $B > A$ and
$F  \in  C^\infty(-B,\infty)$ so that \\
1. $F\equiv 0$ on $[0,\infty) \times S^2$. \\
2. The restriction of $F$ to $(-A,0] \times S^2$
satisfies the hypotheses of Lemma \ref{lemroundoff}.\\
3. The metric $e^{2F(z)} g_{cyl}$ on $(-B,\infty) \times S^2$
has nonnegative sectional curvature and extends
smoothly to a metric on $\R^3$
by adding a ball of constant positive curvature at
$\{-B\} \times S^2$.
\end{lemma}
\begin{proof}
  \uses{lem:kl-surgery-pinching-condition-ricci-flow-surgery-lemma}
For a metric of the form $e^{2F(z)} g_{cyl}$, one computes that
the sectional curvatures are
$ - \: e^{- \: 2F} \: F^{\prime \prime}$ and $e^{- \: 2F}
\left( \frac12 \: - \: (F^\prime)^2 \right)$.
In particular, the conditions for positive sectional curvature
are $F^{\prime \prime} < 0$ and $|F^\prime| <
\frac{1}{\sqrt{2}}$.

The $3$-sphere of constant sectional curvature $k^2$,
with two points removed, has a metric given by
\begin{equation}
F_k(z) \: = \:  \log \left( \frac{\sqrt{2}}{k} \right) \: + \:
\frac{1}{\sqrt{2}} \: z \: - \: \log \left( 1 + e^{\sqrt{2} z} \right).
\end{equation}
(Shifting $z$ gives other metrics of constant curvature $k^2$.
We have normalized so that the $z=0$ slice is the slice of maximal area.)
Note that the derivative
\begin{equation}
D(z) \: = \: \frac{1}{\sqrt{2}} \: - \:
\sqrt{2} \: \frac{e^{\sqrt{2} z}}{1+e^{\sqrt{2} z}}
\end{equation}
is independent of $k$.

Given $A > 0$, we take $F$ to be $0$ on $[0, \infty)$ and of the form
$c_1 \: e^{c_2/z}$ on $(-A, 0]$. We can take the constant
$c_1 > 0$ sufficiently small and the constant $c_2 < \infty$ sufficiently large
so that the hypotheses of Lemma \ref{lemroundoff} are satisfied. It
remains to smoothly cap off
$([-A, \infty) \times S^2, e^{2F(z)} g_{cyl})$ with something of positive
sectional curvature.

With our given choice of $F \big|_{(-A, 0]}$, we have
$F^\prime(-A) \in (0, \epsilon)$. As
$\lim_{z \rightarrow -\infty}
D(z) \: = \: \frac{1}{\sqrt{2}}$,
we can choose $B > A$
so that $D(-B) > F^\prime(-A)$. As $F^{\prime \prime}(-A) < 0$ and
$D^{\prime}(-B) < 0$, we can extend $F^\prime$ to a smooth function
$\widetilde{D} \: : \: (-B, \infty) \rightarrow \left( 0,
\frac{1}{\sqrt{2}} \right)$
which has $\widetilde{D}^\prime < 0$ and which
coincides with $D$ on a small interval $(-B, -B + \delta)$.
Putting
\begin{equation}
F(z) \: = \: F(0) \: + \: \int_0^z \widetilde{D}(w) \: dw,
\end{equation}
we obtain
$F \in C^\infty(B, \infty)$ which coincides with $F_k$ on
$(-B, -B + \delta)$, for some $k > 0$. Then we can glue on a round metric
ball of constant curvature $k^2$ to $\{-B\} \times S^2$, in order to
obtain the desired metric.
\end{proof}

In the statement of the next lemma we continue with the metric
constructed in Lemma \ref{capoff}.
\begin{lemma}[existence for standard solutions]
  \label{lem:kl-surgery-pinching-condition-existence-standard-solutions}
  \dcref{preservingphipinching}
  \group{Kleiner--Lott Singularities and Surgery}
  \source{kleiner-lott}
  \uses{def:kl-almost-nonnegative-curvature-positivity-curvature-operators}

\label{preservingphipinching}
There exists $\de'=\de'(\de)$ with $\lim_{\de\ra 0}\de'(\de)=0$ and a constant
$\de_0>0$
 such that the following holds.  Suppose that $\de<\de_0$, $x\in \{0\} \times S^2$
and $h_0$ is a Riemannian metric
on $(-A,\frac{1}{\de}) \times S^2$ with $R(x) > 0$ such that:

\medskip
$\bullet$ $h_0$ satisfies the time-$t$ Hamilton-Ivey pinching condition of
Definition \ref{HamIvey}.

$\bullet$  $R(x)h_0$ is
$\de$-close to $g_{cyl}$ in the $C^{[\frac{1}{\de}]+1}$-topology.

\bigskip
Then there is a smooth metric
$h$ on $\R^3 = D^3 \: \cup \: \left( (-B,\frac{1}{\de}) \times S^2 \right)$
such that

\medskip
$\bullet$ $h$ satisfies the time-$t$ pinching condition.

$\bullet$ The restriction of $h$ to $[0,\frac{1}{\de}) \times S^2$ is $h_0$.

$\bullet$  The restriction of $R(x)h$ to $(-B,-A) \times S^2$ is
$g_{0}$, the initial metric of a standard solution.


$\bullet$  The restriction of $R(x)h$ to $D^3$ has constant curvature $k^2$.

$\bullet$ $R(x)h$ is $\de'$-close to $e^{2F}g_{cyl}$ in the $C^{[\frac{1}{\de'}]+1}$-topology on
$\left( -B, \frac{1}{\delta} \right) \times S^2$.

\end{lemma}
\par\medskip\noindent\textit{Source argument (not certified as complete).}\ \begingroup\itshape
Put
\begin{equation}
U_1 \: = \:  (-B,-\frac{A}{2}) \times S^2,\quad U_2 \: = \:  (-A,\frac{1}{\de}) \times S^2
\end{equation}
and let $\{\al_1,\al_2\}$ be a $C^{\infty}$ partition of unity subordinate to the open
cover $\{U_1,U_2\}$ of $(-B,\frac{1}{\de}) \times S^2$.  We set
\begin{equation}
h \: = \:  \al_1 \: R(x)^{-1} \: g_{0} \: + \: \al_2 \: e^{2F} \: h_0
\end{equation}
on $\left( -B, \frac{1}{\delta} \right) \times S^2$ and cap it off with a
$3$-ball of constant curvature $k$, as in Lemma \ref{capoff}.


Given $\de'$, we claim that if $\de$ is sufficiently small then the conclusion of
the lemma holds.
The only part of the lemma that is not obvious is the pinching
condition.  Note that on $\left( -\frac{A}{2}, \frac{1}{\delta} \right) \times S^2$ the metric
$h$ agrees with $e^{2F}h_0$ and hence, when $\de$ is sufficiently
small, the pinching
condition will hold on $\left( -\frac{A}{2}, \frac{1}{\delta} \right) \times S^2$
by Lemmas \ref{lemroundoff} and \ref{lempinchingmonotonic}.
On the other hand, when $\de$ is sufficiently small, the
restrictions of the metrics
$g_{0} \: =\: e^{2F}g_{cyl}$
and $R(x)e^{2F}h_0$ to $(-A,- \frac{A}{2}) \times S^2$
will be very close and will have strictly positive
curvature. (The positive curvature for $e^{2F}h_0$ also follows from Lemma
\ref{lemroundoff}; if $\delta$ is small enough then $\lambda_1^{h_0}$ will be close
to zero, while $- f^{\prime \prime}(z)$ is strictly positive for
$z \in (-A,- \frac{A}{2})$.) Thus $h$ will have positive curvature on $(-B,- \frac{A}{2}) \times S^2$
and the pinching condition will hold there.
\endgroup\par\medskip

We have now fixed the initial condition
$g_{0}$
for a
standard solution, along with the procedure to meld
$g_{0}$
to
an approximate cylinder.

\section{II.4.4. Performing surgery and continuing flows}
\label{II.4.4}
This section discusses the surgery procedure  and shows how to prolong a Ricci
flow with surgery, provided that the {\it a priori} assumptions hold.

\begin{definition}[Ricci flow with cutoff]
  \label{def:kl-performing-surgery-continuing-flows-ricci-flow-cutoff}
  \dcref{rdeltacutoffflowdef}
  \group{Kleiner--Lott Singularities and Surgery}
  \source{kleiner-lott}
  \uses{def:kl-priori-assumptions-surgery-priori-assumptions,lem:kl-necks-horns-existence-ricci-flow-surgery,lem:kl-surgery-pinching-condition-existence-standard-solutions}
(Ricci flow with cutoff)
\label{rdeltacutoffflowdef}
Suppose that $a\geq 0$ and let $\M$ be a Ricci flow with surgery
defined on $[a,b]$ that satisfies the {\it a priori}
assumptions of Definition \ref{apriori}.
Let $\de:[a,b]\ra (0,\delta_0)$ be a nonincreasing function, where
$\delta_0$ is the parameter of Lemma \ref{preservingphipinching}.
Then $\M$ is a
{\em  Ricci flow with $(r, \de)$-cutoff} if
at each singular time $t$, the forward time slice
$\M_{t}^+$ is obtained from the backward time slice $\Omega =
\M_{t}^-$
by applying the following procedure:

A.  Discard each component of $\Omega$ that does not
intersect
\begin{equation}
\Om_{\rho} = \{(x,t)\in \Omega \mid R(x,t)\leq\rho^{-2}\},
\end{equation}
where $\rho= \de(t)r(t)$.

B. In each $\eps$-horn
$\h_{ij}$ of each of the remaining components
$\Om_i$, find a point $(x_{ij},t)$
such that $R(x_{ij},t)=h^{-2}$, where
$h=h(t)$
is as in Lemma \ref{hexists}.

C. Find a strong $\de$-neck $U_{ij}\times [t-h^2,t]$
containing
$P(x_{ij},t,\de^{-1}R(x_{ij},t)^{-\frac12},-R(x_{ij},t)^{-1})$;  this
is guaranteed to exist by Lemma \ref{hexists}.

D.  For each $ij$, let $S_{ij}\subset U_{ij}$ be
a cross-sectional $2$-sphere containing $(x_{ij},t)$.
Cut $\bigcup_i\,\Om_i$
along the $S_{ij}$'s and throw away the tips of the horns
$\h_{ij}$,
to obtain a compact manifold-with-boundary $X$ having a
spherical boundary component for each $ij$.

E. Glue caps onto $X$, using Lemma \ref{preservingphipinching},
to obtain the closed manifold $\M_{t}^+$.
\end{definition}

For concreteness, we take the parameter
$A$ of Lemma \ref{preservingphipinching} to be $10$.
The neighborhood of a boundary
component of $X$ is parametrized as $[-A, \delta^{-1}) \times S^2$,
with $S_{ij} \: = \: \{-A \} \times S^2$.
The metric on $[0, \delta^{-1}) \times S^2$ is unaltered by the surgery procedure.
The corresponding region in the new manifold $\M_{t}^+$, minus
a metric ball of constant curvature, is
parametrized by $(-B, \delta^{-1}) \times S^2$.
Put $S_{ij}^\prime \: = \: \{0\} \times S^2 \subset \M_{t}^-$.
We will consider the part added by surgery on $\h_{ij}$
to be the $3$-disk in $\M_{t}^+$ bounded by $S_{ij}^\prime$.
In terms of Definition \ref{flowwithsurgery}, if $t = t_k^+$ then
the subset $X_k^+$ of $\Omega_k = \M_{t}^-$ has
boundary $\bigcup_{ij} S_{ij}^\prime$. The added part
$\M_{t}^+ - X_{k+1}^-$ is a union of $3$-balls.

\begin{remark}[existence for canonical neighborhoods]
  \label{rem:kl-performing-surgery-continuing-flows-existence-canonical-neighborhoods}
  \dcref{performing-surgery-continuing-flows:remark-2}
  \group{Kleiner--Lott Singularities and Surgery}
  \source{kleiner-lott}

Our definition of surgery differs slightly from that in
\cite{Perelman2}. The paper \cite{Perelman2} has two
extra steps involving throwing away certain components
of the postsurgery manifold.
We omit these steps in order to simplify the definition of
surgery, but there is no real loss either way.

First, in the setup of \cite[Section 4.4]{Perelman2},
any component of $\M_{t}^+$ that is $\epsilon$-close
to a metric quotient of the round $S^3$ is thrown away. The
motivation of \cite{Perelman2} was to not have to include
these in the list of canonical neighborhoods. Such
components
are topologically standard.  We do include such manifolds
in the list of canonical neighborhoods and do not
throw them away in the surgery procedure.

Second,
when considering the long-time behavior of Ricci flow
in \cite[Section 7]{Perelman2},
any component of $\M_{t}^+$ which admits a
metric of nonnegative scalar curvature is thrown away.
The motivation for this extra step is that any such
component admits a metric that is either flat or has
finite extinction time.  In either case one concludes that
the component is a graph manifold and, for the purposes
of the geometrization conjecture, is standard.
(Recall the definition of graph manifolds from
Appendix \ref{geom}.)
Again, we do not throw away such components.
\end{remark}


Note that the definition of Ricci flow with $(r, \de)$-cutoff also depends
on the function $r(t)$ through the {\it a priori} assumption. We now state
how the topology of the time slice changes when going backward through the
singular time  $t$.  Recall that for $t'<t$ close to $t$, the time slices
$\M_{t'}$ are all diffeomorphic; we refer to this diffeomorphism type
as the {\em presurgery manifold}, and the forward time slice
$\M_t^+$ as the {\em postsurgery manifold}.

\begin{lemma}[rigidity for Ricci flow with surgery]
  \label{lem:kl-performing-surgery-continuing-flows-rigidity-ricci-flow-surgery}
  \dcref{reconstruct2}
  \group{Kleiner--Lott Singularities and Surgery}
  \source{kleiner-lott}
 \label{reconstruct2}
The presurgery manifold may be obtained
from the postsurgery manifold by applying the following
operations finitely many times:

\begin{itemize}
\item Replacing two connected components with their connected sum.
\item Taking the connected sum of a connected component
with  $S^1\times S^2$ or $\R P^3$.
\item Taking the disjoint union with an additional
$S^1 \times S^2$ or an isometric quotient of the round $S^3$.
\end{itemize}
\end{lemma}
\par\medskip\noindent\textit{Source argument (not certified as complete).}\ \begingroup\itshape
The proof is basically the same as that of Lemma \ref{reconstruct}. The only
difference is that we must take
into account the compact components of $\Omega$ that do not intersect
$\Omega_\rho$; these are thrown away in Step A.
(Such components did not occur in Lemma \ref{reconstruct} because
in Lemma \ref{reconstruct} we were dealing with the first surgery for
the Ricci flow on the initial connected manifold;
see Lemma \ref{omeganoncompact}, which is valid for the first surgery time.)
Any such component is
diffeomorphic to
$S^1 \times S^2$, $\R P^3 \# \R P^3$ or
a quotient of the round $S^3$, in view of the canonical
neighborhood assumption; see the proof of Lemma \ref{standard}.
\endgroup\par\medskip

\begin{remark}[structural remark on Ricci flow with surgery]
  \label{rem:kl-performing-surgery-continuing-flows-structural-remark-ricci-flow-surgery}
  \dcref{volumelossremark}
  \group{Kleiner--Lott Singularities and Surgery}
  \source{kleiner-lott}

\label{volumelossremark}
When $\de>0$ is sufficiently small, we will have
$\vol(\M_t^+)<\vol(\M_t^-)-h(t)^3$
for each surgery time $t\in (a,b)$.
This is because each component that is discarded in step D contains
at least
``half'' of the $\de$-neck $U_{ij}$, which has volume at least $\const \de^{-1}
h(t)^3$,
while the cap added has volume at most $\const h(t)^3$.
\end{remark}

\begin{remark}[existence for Ricci flow with surgery]
  \label{rem:kl-performing-surgery-continuing-flows-existence-ricci-flow-surgery}
  \dcref{performing-surgery-continuing-flows:remark-5}
  \group{Kleiner--Lott Singularities and Surgery}
  \source{kleiner-lott}

For a Ricci flow with surgery whose original manifold is nonaspherical and
irreducible,
one wants to know that the Ricci flow goes extinct within a finite time
\cite{Colding-Minicozzi,Colding-Minicozzi2,Perelman3}.
Consider the effect of a first surgery, say at time $t$. Among the connected
components of the postsurgery
manifold $\M_{t}^+$, one will
be diffeomorphic to the presurgery manifold
 and the others will be $3$-spheres.
Let ${\mathcal N}_{t}^+$ be a component  of $\M_{t}^+$ that is diffeomorphic to
the presurgery manifold.
By the nature of the surgery procedure, there is a
 function
$\xi$ defined on a small interval $(t-\alpha, t)$ so that
$\lim_{t' \rightarrow t} \xi(t') = 1$ and for $t' \in (t - \alpha, t)$,
there is a homotopy-equivalence from $(\M_{t'}, g(t'))$ to ${\mathcal N}_{t}^+$ that
expands distances by at most $\xi(t')$. Following the
subsequent evolution of
${\mathcal N}_{t}^+$,
there is a similar statement for
the later singular times. This fact is needed in
\cite{Colding-Minicozzi,Colding-Minicozzi2,Perelman3}
in order to control the decay of
a certain area functional as one goes through a surgery.
\end{remark}

We discuss how to continue Ricci flows after surgery.
We recall that in Definition \ref{flowwithsurgery}  of a Ricci flow with surgery defined
on an interval $[a,c]$, the final time slice $\M_c$ consists of a single
manifold $\M_c^- = \Omega$ that may or may not be singular.

\begin{lemma}[existence for canonical neighborhoods]
  \label{lem:kl-performing-surgery-continuing-flows-existence-canonical-neighborhoods}
  \dcref{extendit}
  \group{Kleiner--Lott Singularities and Surgery}
  \source{kleiner-lott}
  \uses{def:kl-almost-nonnegative-curvature-positivity-curvature-operators}
(Prolongation of Ricci flows with cutoff)
\label{extendit}
Take the function $\Phi$  to be the time-dependent pinching
function  associated to
Definition \ref{HamIvey}
in Appendix \ref{phiappendix}.
Suppose that $r$ and $\de$ are nonincreasing positive functions
defined on $[a,b]$. Let $\M$ be a Ricci flow with
$(r, \de)$-cutoff defined on an interval $[a,c]\subset [a,b]$.
Provided $\sup\de$ is sufficiently small, either
\begin{enumerate}
\item $\M$ can be prolonged to a Ricci flow with
$(r, \de)$-cutoff defined on $[a,b]$, or
\item There is an extension of $\M$ to a Ricci flow with
surgery
defined  on an interval $[a,T]$ with $T \in (c, b]$, where \\
a. The restriction of  the flow to any subinterval $[a, T^\prime]$,
$T^\prime < T$,
is a Ricci flow with $(r, \de)$-cutoff, but\\
b. The $r$-canonical neighborhood assumption fails at some
point $(x,T)\in \M_T^-$.
\end{enumerate}
In particular, the only obstacle to prolongation of Ricci
flows with $(r,\de)$-cutoff is the potential breakdown of the
$r$-canonical neighborhood assumption.
\end{lemma}
\begin{proof}
  \uses{def:kl-almost-nonnegative-curvature-positivity-curvature-operators,def:kl-performing-surgery-continuing-flows-ricci-flow-cutoff,def:kl-priori-assumptions-canonical-neighborhood-assumption,lem:kl-canonical-neighborhood-property-standard-solutions-existence-canonical-neighborhoods,lem:kl-surgery-pinching-condition-existence-standard-solutions}
Consider the time slice of $\M_c = \M_c^-$ at time $c$.
If it is
singular then we perform steps A-E of Definition \ref{rdeltacutoffflowdef}
to produce $\M_c^+$; otherwise we set $\M_c^+=\M_c^-$.
Since the surgery is done using Lemma \ref{preservingphipinching},
provided $\de>0$ is sufficiently small,
the forward time slice $\M_c^+$ will satisfy the $\Phi$-pinching
assumption.

We claim that the $r$-canonical neighborhood
assumption holds in $\M_c^+$.
More precisely, if a point
$(x,c) \in \M_c^+$ lies within a distance of $10 \epsilon^{-1} h$ from
the added part $\M_c^+ - \M_c^-$ then it lies in an
$\epsilon$-cap, while if $(x,c)$ lies at distance greater than
$10 \epsilon^{-1} h$ from $\M_c^+ - \M_c^-$ and has
scalar curvature greater than $r(c)^{-2}$ then it lies in a
canonical neighborhood that was present in the presurgery manifold
$\M_c^-$.
(We are assuming that $\epsilon < \frac{1}{100}$.)
In view of  Lemma \ref{claim5},
the only point to observe is that points at distance roughly
$10 \epsilon^{-1} h$ lie in $\epsilon$-necks, as they are unaltered
by the surgery and they were in $\delta$-necks before the surgery.
This gives the $\epsilon$-neck needed to define an
$\epsilon$-cap.

We now prolong $\M$  by Ricci flow with initial
condition $\M_c^+$. If the flow extends smoothly
up to time $b$ then we are done because either the canonical neighborhood
assumption holds up to time $b$ yielding (a), or it fails at some
time in the interval $(c,b]$, and we have (b).
Otherwise, there is some time $t_{\sing} \le b$ at which it goes singular.
We add the singular limit $\Omega$ at time $t_{\sing}$
to obtain a Ricci flow with surgery defined on $[a, t_{\sing}]$.
From Lemma \ref{preservingphipinching},
$\M$ satisfies
the Hamilton-Ivey pinching condition of
Definition \ref{HamIvey} on $[a,t_{\sing}]$.
As the function $r$ is nonincreasing in $t$, it follows from
Definition \ref{canonicalnbhddef} that the set of times
$t\in [c,t_{\sing}]$ for which the $r$-canonical neighborhood assumption
holds is relatively open to the right (i.e. if the $r$-canonical
neighborhood assumption holds at time $t \in [c, t_{\sing})$ then
it also holds within some interval
$[t, t^\prime)$). Thus the set of times
$t\in [c,t_{\sing}]$ for which the $r$-canonical neighborhood assumption
holds is either an interval $[c, T)$, with $T \le t_{\sing}$, or
$[c, t_{\sing}]$.
If the set of such times $t$ is
$(c, T)$ for some $T \le t_{\sing}$ then the lemma holds.
Otherwise, the $r$-canonical neighborhood assumption
holds at $t_{\sing}$.  In this case we repeat the construction
with $c$ replaced by $t_{\sing}$, and iterate if necessary. Either we
will reach time $b$ after a finite number  of iterations, or we
will reach a time $T$ satisfying (2),
or we will hit an infinite number of singular times
before time $b$. However, the last possibility cannot
occur.  A singular time corresponds to a component going extinct or
to a surgery. The number of components going extinct before time $b$
can be bounded in terms of the number of surgeries before time $b$,
so it suffices to show that the latter is finite.
Each surgery removes a volume of at least $h^3$, but
the lower bound on the scalar curvature during the flow,
coming from the maximum principle,
gives a finite upper bound on the total volume growth during the complement
of the singular times.
\end{proof}

\begin{remark}[structural remark on canonical neighborhoods]
  \label{rem:kl-performing-surgery-continuing-flows-structural-remark-canonical-neighborhoods}
  \dcref{30}
  \group{Kleiner--Lott Singularities and Surgery}
  \source{kleiner-lott}
  \uses{def:kl-priori-assumptions-canonical-neighborhood-assumption}
 \label{30}
The condition $C_1 \ge 30 \epsilon^{-1}$ in
Definition \ref{canonicalnbhddef}
was in order to ensure that the $\epsilon$-cap coming from a
surgery satisfies the requirements to be a canonical neighborhood.
\end{remark}

\section{II.4.5. Evolution of a surgery cap}
\label{secII.4.5}

Let $\M$ be a Ricci flow with $(r,\de)$-cutoff.
The next result says that provided $\de$ is small,
after a surgery at scale $h$
there is a ball $B$  of radius $Ah\gg h$ centered in the
surgery cap whose evolution is close to that
of a standard solution for an elapsed time close to $h^2$,
unless another surgery occurs  during which
the entire ball is thrown away.   Note that the elapsed time $h^2$
corresponds, modulo parabolic rescaling, to the duration
of the standard solution.

\begin{lemma}[existence for standard solutions]
  \label{lem:kl-evolution-surgery-cap-existence-standard-solutions}
  \dcref{II.4.5}
  \group{Kleiner--Lott Singularities and Surgery}
  \source{kleiner-lott}
  \uses{lem:kl-necks-horns-existence-ricci-flow-surgery}
(cf. Lemma II.4.5)
\label{II.4.5}\\
For any $A<\infty$, $\th\in (0,1)$ and $\hat r>0$,
one can find $\hat\de=\hat\de(A,\th,\hat r)>0$
with the following property.  Suppose that we
have a Ricci flow with $(r, \de)$-cutoff
defined on a time interval $[a,b]$ with  $\min r=r(b)\geq \hat r$.
Suppose that there is a surgery time
$T_0\in (a,b)$, with $\de(T_0)\leq \hat\de$. Consider a given
surgery at the surgery time and let
$(p,T_0)\in\M_{T_0}^+$
be the center of the surgery cap.
Let $\hat h= h(\de(T_0),\eps,r(T_0),\Phi)$ be the surgery
scale given by Lemma \ref{hexists} and put
$T_1=\min(b,T_0+\th \hat h^2)$.
Then one of the two following possibilities occurs : \\
(1) The solution is unscathed on $P(p,T_0,A\hat h,T_1-T_0)$.
The pointed solution there (with respect to the basepoint
$(p, T_0)$) is, modulo parabolic
rescaling, $A^{-1}$-close to the pointed flow on
$U_0\times [0,(T_1-T_0)\hat h^{-2}]$, where
$U_0$ is an open subset of the initial time slice $\s_0$
of a standard solution
$\s$ and the basepoint is
the center $c$ of the cap in $\s_0$. \\
(2) Assertion (1) holds with $T_1$ replaced by some
$t^+ \in [T_0, T_1)$, where $t^+$ is a surgery time. Moreover,
the entire
ball $B(p,T_0,A\hat h)$ becomes extinct at time $t^+$,
i.e.
$P(p,T_0,A\hat h,t_+-T_0)\cap \M_{t_+}\subset \M_{t_+}^- - \M_{t_+}^+$.
\end{lemma}
\begin{proof}
  \uses{def:kl-priori-assumptions-canonical-neighborhood-assumption,lem:kl-blow-up-time-standard-solution-standard-solutions-lemma,lem:kl-evolution-surgery-cap-existence-standard-solutions-2,lem:kl-length-distortion-estimates-distance-distortion-lemma,lem:kl-necks-horns-existence-ricci-flow-surgery,lem:kl-standard-solutions-existence-kappa-solutions,lem:kl-surgery-pinching-condition-existence-standard-solutions}
We give a proof with the same ingredients as the proof in
\cite{Perelman2}, but which
is slightly rearranged.
We first show the following result, which is almost the same as
Lemma \ref{II.4.5}.
\begin{lemma}[existence for standard solutions]
  \label{lem:kl-evolution-surgery-cap-existence-standard-solutions-2}
  \dcref{intermediate}
  \group{Kleiner--Lott Singularities and Surgery}
  \source{kleiner-lott}
  \uses{lem:kl-necks-horns-existence-ricci-flow-surgery}
 \label{intermediate}
For any $A<\infty$, $\th\in (0,1)$ and $\hat r>0$,
one can find $\hat\de=\hat\de(A,\th,\hat r)>0$
with the following property.  Suppose that we
have a Ricci flow with $(r, \de)$-cutoff
defined on a time interval $[a,b]$ with  $\min r=r(b)\geq \hat r$.
Suppose that there is a surgery time
$T_0\in (a,b)$, with $\de(T_0)\leq \hat\de$.  Consider a given
surgery at the surgery time and let
$(p,T_0)\in\M_{T_0}^+$
be the center of the surgery cap.
Let $\hat h= h(\de(T_0),\eps,r(T_0),\Phi)$ be the surgery
scale given by Lemma \ref{hexists} and put
$T_1=\min(b,T_0+\th \hat h^2)$.
Suppose that the solution is unscathed on $P(p,T_0,A\hat h,T_1-T_0)$.
Then the pointed solution there (with respect to the basepoint
$(p, T_0)$) is, modulo parabolic
rescaling, $A^{-1}$-close to the pointed flow on
$U_0\times [0,(T_1-T_0)\hat h^{-2}]$, where
$U_0$ is an open subset of the initial time slice $\s_0$
of a standard solution
$\s$ and the basepoint is
the center $c$ of the cap in $\s_0$.
\end{lemma}
\begin{proof}
  \uses{def:kl-priori-assumptions-canonical-neighborhood-assumption,lem:kl-blow-up-time-standard-solution-standard-solutions-lemma,lem:kl-length-distortion-estimates-distance-distortion-lemma,lem:kl-standard-solutions-existence-kappa-solutions,lem:kl-surgery-pinching-condition-existence-standard-solutions}
Fix $\theta$ and $\hat{r}$.
Suppose that the lemma is not true. Then for some $A > 0$,
there is a sequence $\{\M^\al, (p^\al,T_0^\al) \}_{\al = 1}^\infty$ of
pointed Ricci flows with $(r^\al, \delta^\al)$-cutoff
that together provide a
counterexample. In particular,\\
1. $\lim_{\al \rightarrow \infty}
\delta^\al (T_0^\al) \: = \: 0$. \\
2. $\M^\al$ is unscathed on
$P(p^\al, T_0^\al, A {\hat h}^\al, T_1^\al - T_0^\al)$.\\
3. If
$(\widehat{\M}^\al, (\hat{p}^\al,0))$ is the pointed Ricci
flow arising from $(\M^\al, (p^\al,T_0^\al))$ by
a time shift of $T_0^\al$ and a
parabolic rescaling by ${\hat h}^\al$ then
$P(\hat{p}^\al, 0, A, (T_1^\al - T_0^\al)(\hat{h}^\al)^{-2})$
is not $A^{-1}$-close to a pointed
subset of a standard solution.

Put $T_2 \: = \: \liminf_{\al \rightarrow \infty} \: (T_1^\al - T_0^\al)(\hat{h}^\al)^{-2}$. (We do
not exclude that $T_2 = 0$.) Clearly $T_2 \le \: \theta$.
After passing to a subsequence, we can assume that
$T_2 \: = \: \lim_{\al \rightarrow \infty} \: (T_1^\al - T_0^\al)(\hat{h}^\al)^{-2}$.
Let $T_3$ be the supremum of the set of times
$\tau \in [0, T_2]$ with the property that we can apply
Appendix \ref{subsequence}, if we want,
to take a convergent subsequence of the pointed solutions
$(\widehat{\M}^\al, (\hat{p}^\al,0))$ on the time interval $[0, \tau]$ to get
a limit solution with bounded curvature.
(In applying Appendix \ref{subsequence}, we use the case $l >0$ of
Appendix \ref{applocalder} to get bounds on the
curvature derivatives near time $0$. In particular, if
$T_2 > 0$ then $T_3 > 0$.)
From the nature of the
surgery gluing in Lemma \ref{preservingphipinching},
since $\lim_{\al \rightarrow \infty}
\delta^\al (T_0^\al) \: = \: 0$
we know that we can at least
take a limit of the pointed solutions
$(\widehat{\M}^\al, (\hat{p}^\al,0))$ on the time interval $[0,0]$,
so $T_3$ is well-defined.

\begin{lemma}[Ricci flow with surgery lemma]
  \label{lem:kl-evolution-surgery-cap-ricci-flow-surgery-lemma}
  \dcref{evolution-surgery-cap:sublemma-3}
  \group{Kleiner--Lott Singularities and Surgery}
  \source{kleiner-lott}

$T_3 \: = \: T_2$.
\end{lemma}
\begin{proof}
  \uses{def:kl-priori-assumptions-canonical-neighborhood-assumption,lem:kl-blow-up-time-standard-solution-standard-solutions-lemma,lem:kl-length-distortion-estimates-distance-distortion-lemma,lem:kl-standard-solutions-existence-kappa-solutions}
Suppose not. Consider the interval $[0, T_3)$ (where we define
$[0, 0)$ to be $\{0\}$).
Given $\sigma \in (0, T_2 - T_3]$, for any subsequence
of $\{\M^\al, (\hat{p}^\al,0)\}_{\al=1}^\infty$ (which we relabel as
$\{\M^\al, (\hat{p}^\al,0)\}_{\al=1}^\infty$)
either \\
1. There is some $\lambda > 0$ and an
infinite number of $\al$
for which the set $B(\hat{p}^{\al}, 0, \lambda)$ becomes
scathed on $[0, T_3 + \sigma]$, or \\
2. For each $\lambda > 0$ the set
$P(\hat{p}^{\al}, 0, \lambda, T_3 + \sigma)$
is unscathed for large $\al$, but for each
$\Lambda > 0$ there is some
$\lambda_\Lambda > 0$ such that
$\limsup_{\al \rightarrow \infty}
\sup_{P(\hat{p}^{\al}, 0, \lambda_\Lambda, T_3 + \sigma)} |\Rm| \: \ge \: \Lambda$.

By Appendix \ref{subsequence}, after
passing to a subsequence, there is a complete limit solution
$(\widehat{\M}^\infty, (\hat{p}^\infty,0))$ defined on the time
interval $[0, T_3)$ with bounded curvature on compact time intervals.
Relabel the subsequence by $\al$. By Lemma \ref{lemstandard0},
$(\widehat{\M}^\infty, (\hat{p}^\infty,0))$ must be the same as the
restriction of some standard solution to $[0, T_3)$.
From Lemma \ref{lemclaim4}, the  curvature of $\widehat{\M}^\infty$ is uniformly
bounded on $[0, T_3)$; therefore by the canonical neighborhood assumption and equation
(\ref{II.(1.3)estimates}),
we can choose $\sigma \in (0, T_2 - T_3]$ and $\Lambda^\prime > 0$
so that for any $\lambda > 0$,  we have
$\limsup_{\al \rightarrow \infty}
\sup_{P(\hat{p}^{\al}, 0, \lambda, T_3 + \sigma)} |\Rm| \: \le \: \Lambda^\prime$.
However, $\lim_{\al \rightarrow \infty}
\delta^\al (T_0^\al) \: = \: 0$ and surgeries only occur near the
centers of $\delta$-necks. From the curvature bound on the time
interval $[0, T_3 + \sigma]$ and the length distortion
estimates of Lemma \ref{distancedist}, for a given $\lambda$ the
balls $B(\hat{p}^{\al}, 0, \lambda)$ will stay within
a uniformly bounded distance from $\hat{p}^{\al}$ on the time interval
$[0, T_3 + \sigma]$. Hence they
cannot be scathed on $[0, T_3 + \sigma]$ for an infinite number of $\al$,
as the collar length of the $\delta$-neck around the supposed surgery locus
would be large enough to prohibit the cap point $\hat{p}^{\al}$
from being within a bounded
distance from the surgery locus.
This, along with the fact that $\limsup_{\al \rightarrow \infty}
\sup_{P(\hat{p}^{\al}, 0, \lambda, T_3 + \sigma)} |\Rm| \: \le \: \Lambda^\prime$
for all $\lambda > 0$,
gives a contradiction.
\end{proof}

Returning to the original sequence
$\{\M^\al, (p^\al,T_0^\al) \}_{\al = 1}^\infty$ and its rescaling
$\{ \widehat{\M}^\al, (\hat{p}^\al,0) \}_{\al = 1}^\infty$, we can now
take a subsequence that converges on the time interval $[0, T_2)$,
again necessarily to a standard solution.  Then there will be
an infinite subsequence
$\{ \widehat{\M}^{\al_\beta}, (\hat{p}^{\al_\beta},0) \}_{\beta = 1}^\infty$
of
$\{ \widehat{\M}^{\al}, (\hat{p}^{\al},0) \}_{\al = 1}^\infty$,
with
$\lim_{\beta \rightarrow \infty}
(T_1^{\al_\beta} - T_0^{\al_\beta})(\hat{h}^{\al_\beta})^{-2} \: = \: T_2$,
so that $P(\hat{p}^{\al_\beta}, 0, A, (T_1^{\al_\beta} - T_0^{\al_\beta})
(\hat{h}^{\al_\beta})^{-2})$
is $A^{-1}$-close to a pointed
subset of a standard solution (by the canonical
neighborhood assumption, equation
(\ref{II.(1.3)estimates}) and Appendix \ref{applocalder}).
This is a contradiction.
\end{proof}

We now finish the proof of Lemma \ref{II.4.5}. If the solution is
unscathed on $P(p,T_0,A\hat h,T_1-T_0)$ then we can apply
Lemma \ref{intermediate} to see
that we are in case (1) of the conclusion of Lemma \ref{II.4.5}.
Suppose, on the other hand,
that the solution is scathed on $P(p,T_0,A\hat h,T_1-T_0)$. Let
$t^+$ be the largest $t$ so that the solution is unscathed on
$P(p,T_0,A\hat h,t-T_0)$.
We can apply
Lemma \ref{intermediate} to
see that conclusion (1) of Lemma \ref{II.4.5}
holds with $T_1$ replaced by $t^+$.
As surgery is always performed near the middle of a $\delta$-neck,
if $\hat \delta << A^{-1}$ then the final time slice in the parabolic
neighborhood $P(p,T_0,A\hat h,t^+-T_0)$ cannot intersect a
$2$-sphere where a surgery is going to be performed.
The only other possibility is that
the entire ball $B(p,T_0,A\hat h)$ becomes extinct at time $t^+$.
\end{proof}

\section{II.4.6. Curves that penetrate the surgery region}
\label{II.4.6}

Let $\M$ be a Ricci flow with $(r,\de)$-cutoff.
The next result, Corollary \ref{corollaryII.4.6}, says that if $\de$ is sufficiently small
then an admissible curve $\ga$ which comes close
to a surgery cap at a surgery time will have a large value
of $\int_\gamma \left( R(\gamma(t))+|\dot{\ga}(t)|^2 \right) \: dt$.
Note that the latter quantity is not quite the same as
${\mathcal L}(\gamma)$, and is invariant under parabolic rescaling.

Corollary \ref{corollaryII.4.6} is used in the extension of
Theorem \ref{nolocalcollapse} to Ricci flows with
surgery.   The idea is that if $\de$ is small
and $L(x,t)$ isn't too large then any $\L$-minimizing sequence of
admissible curves joining the basepoint $(x_0,t_0)$ to $(x,t)$
must avoid surgery regions, and will therefore
accumulate on a minimizing $\L$-geodesic.

\begin{corollary}[Ricci flow with surgery corollary]
  \label{cor:kl-curves-that-penetrate-surgery-region-ricci-flow-surgery-corollary}
  \dcref{corollaryII.4.6}
  \group{Kleiner--Lott Singularities and Surgery}
  \source{kleiner-lott}
  \uses{lem:kl-evolution-surgery-cap-existence-standard-solutions,lem:kl-necks-horns-existence-ricci-flow-surgery}
 (cf. Corollary II.4.6)
\label{corollaryII.4.6}
For any $l<\infty$ and $\hat r>0$, we can find $A=A(l,\hat r)<\infty$ and
$\th=\th(l,\hat r)$ with the following property.
Suppose that we are in the situation of Lemma \ref{II.4.5}, with
$\delta(T_0) < \hat{\delta}(A, \th, \hat r)$.
As usual,
$\hat h$ will be the surgery scale coming from Lemma \ref{hexists}.
Let $\ga:[T_0,T_\ga]\ra \M$
be an admissible curve, with
$T_\ga\in (T_0,T_1]$.
Suppose that
$\ga(T_0)\in B(p,T_0,\frac{A\hat h}{2})$,
$\ga([T_0,T_\ga))\subset P(p,T_0,A\hat h, T_\ga-T_0)$, and
either

a. $T_\ga=T_1=T_0+\th(\hat h)^2$,

or

b. $\ga(T_\ga)\in\D B(p,T_0,A\hat h)\times [T_0,T_\ga]$.

Then
\begin{equation}
\int_{T_0}^{T_\ga} \left( R(\gamma(t),t)+|\dot{\ga}(t)|^2 \right) \: dt
\: > \: l.
\end{equation}
\end{corollary}
\begin{proof}
  \uses{lem:kl-canonical-neighborhood-property-standard-solutions-existence-canonical-neighborhoods,lem:kl-evolution-surgery-cap-existence-standard-solutions}
For the moment, fix $A<\infty$ and $\th\in (0,1)$.
Choose $\hat\de=\hat\de(A,\th,\hat r)$
so as to satisfy
Lemma \ref{II.4.5}.   Let $\M$, $(p,T_0)$, etc.,
be as in the hypotheses of Lemma \ref{II.4.5}. Let $\ga:[T_0,T_\ga]\ra\M$
be a curve as in the hypotheses of the Corollary.
From Lemma \ref{II.4.5}, we know that there is a standard solution
$\s$ such that  the parabolic region
$P(p,T_0,A\hat h,T_\ga-T_0) \subset \M$, with basepoint $(p,T_0)$, is
(after parabolic rescaling by $\hat h^{-2}$) $A^{-1}$-close
to a  pointed flow $U_0\times [0,\hat T_\ga] \subset \s$, the latter
having basepoint $(c,0)$. Here
$U_0\subset \s_0$ and $\hat T_\ga= (T_\ga-T_0)\hat h^{-2}$.
Then the image of $\ga$, under the diffeomorphism implicit in the definition
of $A^{-1}$-closeness, gives rise to a smooth curve
$\ga_0:[0,\hat T_\ga]\ra U_0\times [0,\hat T_\ga]$ so that
(if $A$ is sufficiently large) :

\begin{gather}
\ga_0(0)\in B(c,0,\frac35 A),\\
\int_{T_0}^{T_\ga}|\dot{\ga}|^2dt
\geq \frac12 \int_0^{\hat T_\ga}|\dot{\ga}_0|^2dt,\\
\int_{T_0}^{T_\ga}R({\ga}(t),t)dt
\geq \frac12 \int_0^{\hat T_\ga}R({\ga}_0(t),t)dt,
\end{gather}
and

(a) $\hat T_\ga=\th$,

or

(b) $\ga_0(\hat T_\ga)\not\in P(c,0,\frac45 A,\hat T_\ga)$.

\noindent
In case (a) we have, by Lemma \ref{claim5},
\begin{equation}
\int_{T_0}^{T_\ga} R(\ga(t), t)dt\geq
\frac12 \int_0^\th R(\ga_0(t),t)dt
\geq \frac12\int_0^\th \const(1-t)^{-1}dt=\const \log(1-\th).
\end{equation}
If we choose $\th$ sufficiently close to $1$ then in this case,
we can ensure that
\begin{equation}
\int_{T_0}^{T_\ga} \left( R(\gamma(t), t)+|\dot{\ga}(t)|^2 \right) \: dt
\: \ge \:
\int_{T_0}^{T_\ga} R(\gamma(t), t) \: dt \: > \: l.
\end{equation}

In case (b), we may use the fact that the Ricci curvature
of the standard solution is everywhere
nonnegative,
and hence the metric tensor is nonincreasing with time.  So if
$\pi:\s=\s_0\times [0,1)\ra\s_\th$ is projection
to the time-$\th$ slice and we put $\eta= \pi\circ\ga_0$
then
\begin{align}
\int_{T_0}^{T_\ga}|\dot{\ga}(t)|^2dt \:
& \geq \: \frac12 \int_0^{\hat T_\ga}|\dot{\ga_0}(t)|^2dt \:
\geq \: \frac12 \int_0^{\hat T_\ga}|\dot{\eta}(t)|^2dt
\: \geq \:
\frac{1}{2 \hat T_\ga} \left(d(\eta(0),\eta(\hat T_\ga))\right)^2 \\
& \ge \:
\frac{1}{2} \left(d(\eta(0),\eta(\hat T_\ga))\right)^2. \notag
\end{align}
With our given value of $\th$, in view of (b), if we take $A$ large enough
then we can ensure that $\frac12 \left(d(\eta(0),\eta(\hat T_\ga))\right)^2
\: > \: l$.  This proves the lemma.
\end{proof}


\section{II.4.7. A technical estimate}

The next result is a technical result that will not be used in the sequel.

\begin{corollary}[existence for a technical estimate]
  \label{cor:kl-technical-estimate-existence-technical-estimate}
  \dcref{Rbound}
  \group{Kleiner--Lott Singularities and Surgery}
  \source{kleiner-lott}
  \uses{lem:kl-evolution-surgery-cap-existence-standard-solutions}
(cf. Corollary II.4.7)
For any $Q<\infty$ and $\hat r>0$, there is a $\th=\th(Q,\hat r)\in (0,1)$
with the following property.  Suppose that we are in the situation of
Lemma \ref{II.4.5}, with
$\delta(T_0) < \hat{\delta}(A, \th, \hat r)$
and $A>\eps^{-1}$.  If $\ga:[T_0,T_x]\ra \M$ is a static curve
starting in $B(p,T_0,A\hat h)$, and
\begin{equation}
\label{Rbound}
Q^{-1}R(\ga(t))\leq R(\ga(T_x))\leq Q(T_x-T_0)^{-1}
\end{equation}
for all $t\in [T_0,T_x]$,  then $T_x\leq T_0+\th\hat h^2$.
\end{corollary}
\begin{remark}[estimates for standard solutions]
  \label{rem:kl-technical-estimate-estimates-standard-solutions}
  \dcref{technical-estimate:remark-2}
  \group{Kleiner--Lott Singularities and Surgery}
  \source{kleiner-lott}
  \uses{cor:kl-technical-estimate-existence-technical-estimate}

The hypothesis (\ref{Rbound}) in the corollary means
that in the scale of the scalar curvature $R(\gamma(T_x))$ at the
endpoint $\gamma(T_x)$, the scalar curvature on $\gamma$ is bounded
and the elapsed time of $\gamma$ is bounded. The conclusion says that
given these bounds, the elapsed time is strictly less than that of
the corresponding rescaled standard solution.
\end{remark}
\begin{proof}
  \uses{cor:kl-technical-estimate-existence-technical-estimate,lem:kl-canonical-neighborhood-property-standard-solutions-existence-canonical-neighborhoods,lem:kl-evolution-surgery-cap-existence-standard-solutions}
If $T_x>T_0+\th\hat h^2$ then by Lemma \ref{claim5} and \ref{II.4.5},
\begin{equation}
R(\ga(T_0+\th\hat h^2))\ge \const(1-\th)^{-1}\hat h^{-2}.
\end{equation}
Thus by (\ref{Rbound}) we get
\begin{equation}
Q^{-1}\const (1-\th)^{-1}\hat h^{-2}\leq R(\ga(T_x))\leq Q(T_x-T_0)^{-1},
\end{equation}
or
\begin{equation}
T_x-T_0\leq \const Q^2(1-\th)\hat h^2,
\end{equation}
If we choose $\th$ close enough to $1$ then $\const Q^2(1-\th)\hat h^2$ is less than
$\th\hat h^2$, which gives a contradiction.
\end{proof}


\section{II.5. Statement of the
the existence theorem for
 Ricci flow with surgery} \label{II.5}

 Our presentation of this material follows Perelman's, except for
some shuffling of the material.  We will be using some terminology
introduced in Section \ref{Ricci flow with surgery}, as well
as results about the $L$-function and noncollapsing from
Sections \ref{lfunctionandsurgerysection} and \ref{II.5.2section}.

\begin{definition}[normalized definition]
  \label{def:kl-statement-existence-theorem-ricci-flow-surgery-normalized-definition}
  \dcref{statement-existence-theorem-ricci-flow-surgery:definition-1}
  \group{Kleiner--Lott Surgery Existence}
  \source{kleiner-lott}

A compact Riemannian $3$-manifold is {\em normalized} if
$|\Rm|\leq 1$ everywhere, and the volume of every unit
ball is at least half the volume of the Euclidean unit
ball.
\end{definition}

We will use the fact that a smooth normalized
Ricci flow, with bounded curvature on compact time intervals,
satisfies the Hamilton-Ivey pinching condition of
Definition \ref{HamIvey}.

The main result of the surgery procedure is Proposition \ref{surgeryflowexists}
(cf. II.5.1), which
implies that one can choose positive nonincreasing
functions $r:\R_+\ra (0,\infty)$,
$\de:\R_+\ra (0,\infty)$ such that the Ricci
flow with $(r,\de)$-surgery
flow starting with any normalized initial condition will be defined
for all time.

The actual statement is structured to facilitate
a proof by induction:


\begin{proposition}[monotonicity for kappa solutions]
  \label{prop:kl-statement-existence-theorem-ricci-flow-surgery-monotonicity-kappa-solutions}
  \dcref{surgeryflowexists}
  \group{Kleiner--Lott Surgery Existence}
  \source{kleiner-lott}
(cf. Proposition II.5.1)
\label{surgeryflowexists}
There exist decreasing sequences $0<r_j<\eps^2$,
$\kappa_j>0$, $0<\bar\de_j<\eps^2$ for $1\leq j<\infty$,
such that for any normalized initial data and any nonincreasing function
$\de:[0,\infty)\ra(0,\infty)$ such that $\de<\bar\de_j$
on $[2^{j-1}\eps,2^j\eps]$, the Ricci flow with $(r,\de)$-cutoff
is defined for all time and is $\kappa$-noncollapsed
at scales below $\eps$.
\end{proposition}
\noindent
Here, and in the rest of this section, $r$ and $\kappa$
will always denote functions defined on an interval
$[0,T]\subseteq [0,\infty)$ with the property that
$r(t)=r_j$ and $\kappa(t)=\kappa_j$
for all $t\in [0,T]\cap [2^{j-1}\eps,2^j\eps)$.
By ``$\kappa$-noncollapsed at scales below $\epsilon$'', we mean
that for each $\rho < \epsilon$
and all $(x, t) \in \M$
with $t \ge \rho^2$,
whenever $P(x, t, \rho, - \rho^2)$ is unscathed and
$|\Rm| \: \le \: \rho^{-2}$ on $P(x, t, \rho, - \rho^2)$,
then we also have
$\vol(B(x,t,\rho)) \: \ge \: \kappa(t) \rho^3$.

Recall that
$\eps$ is a ``global'' parameter which is assumed to be
small, i.e. all statements involving $\eps$ (explicitly or
otherwise) are true provided $\eps$ is sufficiently small.
Proposition \ref{surgeryflowexists} does not impose any
serious new constraints on $\epsilon$.
For example, instead of using the time intervals
$\{[2^{j-1}\eps,2^j\eps]\}_{j=1}^\infty$,
we could have taken any collection of adjoining time intervals
starting at a small positive time. Also, we just need some
fixed upper bound on $r_j$ and $\overline{\delta}_j$.
We will follow \cite{Perelman2} and write these somewhat
arbitrary constants in terms of the single global parameter $\epsilon$.
Note also that having normalized initial data sets a length
scale for the Ricci flow.

The phrase ``the Ricci flow with $(r,\de)$-cutoff
is defined for all time'' allows for the possibility that the entire
manifold goes extinct, i.e. that after some time we are talking
about the flow on the empty set.

In the rest of this section we give a sketch of the proof.  The details
are in the subsequent sections.

Given positive nonincreasing functions $r$ and $\de$,
if one has a normalized initial condition $(M,g(0))$ then there
will be a maximal time interval on which the Ricci flow with
$(r,\de)$-cutoff is defined.   This interval can be finite
only if it is of the form $[0,T)$ for some $T<\infty$, and the
Ricci flow with $(r,\de)$-cutoff on $[0,T)$ extends to a Ricci flow with surgery
on $[0,T]$ for which the $r$-canonical neighborhood
assumption fails at time $T$; see Lemma \ref{extendit}. The main point
here is that the $r$-canonical neighborhood assumption
allows one to run the flow forward up to the singular time, and then perform
surgery, while volume considerations rule out an accumulation of
surgery times.   Thus the crux of the proof is showing that the
functions $r$
and $\de$ can be chosen so that the $r$-canonical neighborhood assumption
will continue to hold, and the Ricci flow with surgery
satisfies a noncollapsing
condition.

The strategy is to  argue
by induction on $i$ that $r_i$, $\bar\de_i$, and $\kappa_i$
can be chosen (and $\bar\de_{i-1}$ can be adjusted)
so that the statement of the proposition holds on the
the finite time interval $[0,2^i\eps]$.
In the induction step, one establishes the canonical
neighborhood assumption using an argument by contradiction similar
to the proof of Theorem \ref{thmI.12.1}. (We recommend that the reader review this
before proceeding).   The main difference between the proof of
Theorem \ref{thmI.12.1} and that of Proposition \ref{surgeryflowexists} is that the
non-collapsing assumption, the key ingredient
that allows one to implement the blowup
argument, is no longer available as a direct consequence of Theorem \ref{nolocalcollapse},
due to the presence of surgeries.

We now discuss the augmentations to the non-collapsing argument
of Theorem \ref{nolocalcollapse} necessitated by surgery; this is treated in detail in sections
\ref{lfunctionandsurgerysection} and \ref{II.5.2section}.
We first recall Theorem \ref{nolocalcollapse} and its proof:
  if a parabolic ball $P(x_0,t_0,r_0,-r_0^2)$
in Ricci flow (without surgery) is
sufficiently collapsed then one
uses the $L$-function
with basepoint $(x_0,t_0)$, and the $\L$-exponential map based
at $(x_0,t_0)$, to get a contradiction.   One considers the reduced
volume of a suitably chosen time slice $\M_t$. There is a
positive lower bound on the reduced volume coming
from the selection of a point where the reduced distance
is at most $\frac{3}{2}$, which in turn comes from an application of the
maximum principle to the $L$-function. On the other hand, there
is an upper bound on the reduced volume, which the collapsing
forces to be small, thereby giving the contradiction. The upper bound comes
from the monotonicity of the weighted Jacobian of the $\L$-exponential
map.   In fact,  this upper bound works without
significant modification  in the presence of surgery,
provided one considers only the reduced volume contributed
by those points in the time $t$ slice
which may be joined to $(x_0,t_0)$ by minimizing $\L$-geodesics
lying in the regular part of spacetime (see Lemma \ref{localreducedvolume}).

To salvage the lower bound on the reduced volume, the basic idea is that
by making the surgery parameter $\de$ small, one can force the
$\L$-length of any curve passing close to the surgery locus to
be large (Lemma \ref{lplusbig}).   This implies that if $(x,t)$ is a point where $L$
isn't too large, then there will  necessarily be an
$\L$-geodesic from $(x_0,t_0)$ to $(x,t)$.  To construct
the minimizer, one takes  a sequence of admissible curves
from $(x,t)$ to $(x_0,t_0)$ with
$\L$-length tending to the infimum, and argues that they must stay away from the
surgeries; hence
they remain in a compact part of
spacetime, and subconverge to a minimizer.
Therefore the calculations
from Sections \ref{I.7}-\ref{I.7.3}
will be valid near such a point $(x,t)$. The
maximum principle can then be applied as before
to show that the minimum of the reduced
length is $\leq \frac{3}{2}$ on each time slice (see Lemma \ref{stillleq3/2}).

To be more precise, if one makes the surgery parameter $\de(t^\prime)$ small
for a surgery at a given time $t^\prime$ then one can force the $\L$-length of
any curve passing close to the time-$t^\prime$ surgery locus to be large, provided that the
endtime $t_0$ of the curve is not too large compared to $t^\prime$.
(If $t_0$ is much larger than $t^\prime$ then the curve may spend a long time
in regions of negative scalar curvature after time $t^\prime$. The ensuing negative
effect on $\L$ could overcome the
positive effect of the small surgery parameter.)
In the proof of Theorem
\ref{nolocalcollapse}, in order to show noncollapsing at time $t_0$, one went all
the way back to a time slice near the initial time and found a point there where
$l$ was at most $\frac32$. There would be a problem in using this method
for Ricci flows with surgery - we would have to constantly redefine $\delta(t^\prime)$
to handle the case of larger and larger $t_0$. The resolution is to not go back to
a time slice
near the initial time slice. Instead, in order to
show $\kappa$-noncollapsing in the time slice
$[2^i \eps, 2^{i+1} \eps]$, we will want to get a lower bound on the
reduced volume for
a time $t$-slice with $t$ lying in the preceding time interval
$[2^{i-1}\eps,2^i\eps]$.   As we inductively have control over the geometry
in the time slice $[2^{i-1}\eps,2^i\eps]$, the argument works equally well.

Finally, as mentioned, after obtaining the {\it a priori} $\kappa$-noncollapsing
estimate on the interval $[2^{i}\eps,2^{i+1}\eps]$, one proves that the
$r$-canonical neighborhood assumption holds at time
$T \in [2^{i}\eps,2^{i+1}\eps]$. One difference here is
that because of possible nearby surgeries, there are two ways
to obtain the canonical neighborhood : either from closeness to
a $\kappa$-solution, as in the proof of
Theorem \ref{thmI.12.1}, or from closeness to a standard
solution.



\section{The $L$-function  of I.7 and Ricci flows with surgery}
\label{lfunctionandsurgerysection}


In this section we examine several points which arise when one adapts
the noncollapsing argument of Theorem \ref{nolocalcollapse} to Ricci flows with surgery.
This material is implicit background for
Lemma \ref{II.5.2} and Proposition \ref{propII.6.3}.  We will use notation
and terminology introduced in Section \ref{Ricci flow with surgery}.

Let $\M$ be a Ricci flow with surgery,
and fix a point $(x_0,t_0)\in\M$.   One may
define the $\L$-length of an  admissible curve $\ga$ from
$(x_0,t_0)$ to some $(x,t)$, for $t<t_0$, using the formula
\begin{equation}
\L(\ga)=\int_t^{t_0}\sqrt{t_0-\bar t}\,\left(R+|\dot{\ga}|^2\right)d\bar t,
\end{equation}
where $\dot{\ga}$ denotes the spatial part of the velocity of $\ga$.
One defines the $L$-function on $\M_{(-\infty, t_0)}$ by setting
$L(x,t)$ to be the infimal $\L$-length of the admissible curves
from $(x_0,t_0)$ to $(x,t)$ if such an admissible curve exists, and
infinity otherwise. We note that if $(x,t)$ is in a surgery time slice
$\M_t^-$ and is actually removed by the surgery then there will not be an
admissible curve from $(x_0, t_0)$ to $(x,t)$.

If $\ga$ is an admissible curve lying in $\M_{\reg}$ then
the first variation formula applies.
Hence an admissible curve in $\M_{\reg}$ from $(x_0, t_0)$ to $(x,t)$
whose ${\mathcal L}$-length
equals $L(x,t)$ will satisfy the ${\mathcal L}$-geodesic
equation.
If $\ga$ is a stable $\L$-geodesic in $\M_{\reg}$ then
the proof of the monotonicity along $\ga$ of
the weighted Jacobian
$\tau^{-\frac{3}{2}}\exp(-l(\tau))J(\tau)$ remains valid.
Similarly, if $U\subset \M_{(-\infty,t_0)}$ is an open set such that
every $(x, t)\in U$ is accessible from $(x_0,t_0)$
by a minimizing $\L$-geodesic (i.e. an $\L$-geodesic of $\L$-length
$L(x,t)$) contained in $\M_{\reg}$,
then the arguments of Section \ref{I.(7.15)} imply that the
differential inequality
\begin{equation} \label{newbarrier}
\bar L_\tau+\De \bar L\leq 6
\end{equation}
holds in $U$, in the barrier sense, where $\tau = t_0 - t$,
$\bar L \: = \: 2 \sqrt{\tau} \: L$ and
$l = \frac{\bar L}{4\tau}$.

\begin{lemma}[existence for Ricci flow with surgery]
  \label{lem:kl-reduced-l-geometry-existence-ricci-flow-surgery}
  \dcref{minimizersexist}
  \group{Kleiner--Lott Surgery Existence}
  \source{kleiner-lott}
[Existence of $\L$-minimizers]
\label{minimizersexist}
Let $\M$ be a Ricci flow with surgery
defined on $ [a,b]$. Suppose that $(x_0,t_0)\in \M$ lies in the backward
time slice $\M_{t_0}^-$.

(1) For each $(x,t)\in \M_{[a,t_0)}$ with $L(x,t)<\infty$, there exists
an $\L$-minimizing admissible path $\ga:[t,t_0]\ra\M$ from $(x,t)$ to $(x_0,t_0)$
which satisfies the
$\L$-geodesic equation at every time $\overline{t}\in (t,t_0)$ for which
$\ga(\overline{t})\in\M_{\reg}$.


(2) $L$ is lower semicontinuous on
$\M_{[a,t_0)}$
and continuous
on
$\M_{\reg}\cap \M_{[a,t_0)}$. (Note that $\M_a^+ \subset \M_{\reg}$.)

(3) Every sequence $(x_j,t_j)\in \M_{[a,t_0)}$ with $\limsup_j L(x_j,t_j)<\infty$
has  a convergent subsequence.
\end{lemma}
\par\medskip\noindent\textit{Source argument (not certified as complete).}\ \begingroup\itshape
(1)   Let
$\{\ga_j:[t,t_0]\ra\M\}_{j=1}^\infty$ be a sequence of
admissible curves from $(x,t)$ to $(x_0,t_0)$ such that
$\lim_{j\ra\infty}\L(\ga_j)=L(x,t)<\infty$. By restricting the sequence,
we may assume that
$\sup_j \L(\ga_j)<2L(x,t)$.
We claim that there is
a subsequence of the $\ga_j$'s  that

(a) converges uniformly to some $\ga_\infty:[t,t_0]\ra\M$,

and

(b) converges weakly to $\ga_\infty$ in $W^{1,2}$ on any  subinterval
$[t',t'']\subset [t,t_0)$ such that $[t',t'']$ is free of singular times.

\noindent
To see this, note that on any time interval $[c,d]\subset [t,t_0)$
which is free of singular times, one may apply the Schwarz inequality
to the $\L$-length, along with the fact that the metrics on the time
slices $\M_t$, $t \in [c,d]$, are uniformly biLipschitz to each other,
to conclude that the $\ga_j$'s are uniformly
H\"older-continuous on $[c,d]$.  We know
that $\ga_j(t')$ lies in $\M_{t'}^-\cap \M_{t'}^+$ for each surgery
time $t'\in (t,t_0)$, and so one can use similar reasoning to get
H\"older control on a short time interval of the form $[t'',t']$.
Using a change of variable as in  (\ref{changetos}), one obtains
uniform H\"older control near $t_0$ after reparametrizing with $s$.
It follows that the $\ga_j$'s are equicontinuous and map into a
compact part of spacetime,  so
Arzela-Ascoli applies;  therefore, by passing  to a subsequence we
may assume  that (a) holds.

To show (b),  we apply weak compactness to the sequence
\begin{equation}
\{\ga_j\restr_{[t',t'']}\};
\end{equation}
\noindent
this is justified by the fact that  the paths $\ga_j\restr_{[t',t'']}$'s remain in a part
of $\M$ with  bounded geometry.  Thus we may assume that our sequence $\{\ga_j\}$
converges uniformly on $[t,t_0]$ and weakly on every subinterval $[t',t'']$ as in (b).
By weak lower semicontinuity of $\L$-length, it follows that the $W^{1,2}$-path
$\ga_\infty$ has $\L$-length $\leq L(x,t)$.  Since any $W^{1,2}$ path  may be approximated
in $W^{1,2}$ by  admissible curves with the same endpoints, it follows that $\ga_\infty$
 minimizes $\L$-length among $W^{1,2}$ paths, and therefore
it restricts to a smooth solution of the $\L$-geodesic equation on each time
interval $[t',t'']\subset [t,t_0]$ such that $(t',t'')$ is free of singular times.
Hence $\ga_\infty$ is an $\L$-minimizing admissible curve.

(2) Pick $(x,t)\in\M_{[a,t_0)}$.  To verify lower semicontinuity at
$(x,t)$ we suppose the sequence $\{(x_j,t_j)\}\subset \M_{[a,t_0)}$
converges to $(x,t)$
and $\liminf_{j\ra\infty}L(x_j,t_j)<\infty$.   By (1) there is a sequence
$\{\ga_j\}$ of $\L$-minimizing admissible curves, where
$\ga_j$ runs from $(x_j,t_j)$ to $(x_0,t_0)$.  By the reasoning above,
a subsequence of $\{\ga_j\}$ converges uniformly and weakly in $W^{1,2}$ to a $W^{1,2}$
curve $\ga_\infty:[t,t_0]\ra\M$ going from $(x,t)$ to $(x_0,t_0)$, with
\begin{equation}
\L(\ga_\infty)\leq \liminf_{j\ra\infty}\L(\ga_j).
\end{equation}
Therefore $L(x,t)\leq \liminf_{j\ra\infty} L(x_j,t_j)$, and we have established
semicontinuity.  If
$(x,t)\in \M_{\reg}$,
the opposite inequality obviously
holds, so in this case $(x,t)$ is a point of continuity.

(3) Because $\{L(x_j,t_j)\}$ is uniformly bounded,  any sequence $\{\ga_j\}$ of $\L$-minimizing
paths with $\ga_j(t_j)=(x_j,t_j)$ will be equicontinuous, and hence by Arzela-Ascoli
a subsequence converges uniformly.  Therefore a subsequence of $\{(x_j,t_j)\}$ converges.
\endgroup\par\medskip

The fact that (\ref{newbarrier}) can hold locally allows one
to appeal -- under appropriate conditions -- to the maximum
principle as in Section \ref{I.(7.15)} to prove
that $\min l\leq \frac{3}{2}$ on
every time slice.
Recall that $l \: = \: \frac{L}{2\sqrt{\tau}} \: = \:
\frac{\overline{L}}{4\tau}$.


\begin{lemma}[existence for reduced length]
  \label{lem:kl-reduced-l-geometry-existence-reduced-length}
  \dcref{stillleq3/2}
  \group{Kleiner--Lott Surgery Existence}
  \source{kleiner-lott}

\label{stillleq3/2}
Suppose that $\M$ is a Ricci flow with surgery defined on $ [a,b]$.
Take
$t_0\in (a,b]$
and $(x_0,t_0)\in \M_{t_0}^-$.  Suppose that
for every $t \in [a, t_0)$, every admissible curve
$[t,t_0]\ra\M$ ending at $(x_0,t_0)$ which does not lie
in
$\M_{\reg}\cup  \M_{t_0}^-$
has reduced length strictly greater than $\frac{3}{2}$. Then
there is a point $(x,a)\in \M_a^+$ where $l(x,a)\leq \frac{3}{2}$.
\end{lemma}
\begin{remark}[structural remark on reduced length]
  \label{rem:kl-reduced-l-geometry-structural-remark-reduced-length}
  \dcref{reduced-l-geometry:remark-4}
  \group{Kleiner--Lott Surgery Existence}
  \source{kleiner-lott}

In the lemma we consider the Ricci flow with surgery to begin at time $a$.
Hence $\M_{\reg}\cup  \M_{t_0}^- \: = \: \M_a^+ \cup \M_{\reg}\cup  \M_{t_0}^-$
and so the hypothesis of the lemma is a statement about the reduced
lengths of barely admissible curves, in the sense of Section
\ref{Ricci flow with surgery}.
\end{remark}


\begin{proof}
  \uses{lem:kl-reduced-l-geometry-existence-ricci-flow-surgery,lem:kl-reduced-l-geometry-positivity-reduced-l-geometry}
As in the case when there are no surgeries,
the proof relies on the maximum principle and a continuity argument.

Let $\be:\M_{[a,t_0)}\ra \R\cup \{\infty\}$ be the function
\begin{equation}
\be \: = \: \bar L-6\tau=4\tau\left(l-\frac{3}{2}\right),
\end{equation}
where as usual, $\tau(x,t) \: = \: t_0 - t$.
Note that for each $\tau\in (0,t_0-a]$, the function
$\be$ attains a minimum $\be_{\min}(\tau)<\infty$ on the
slice $\M_{t_0 - \tau}$,  because by (2) of Lemma \ref{minimizersexist},
it is continuous on the compact manifold
$\M_{t_0 - \tau}^+$ (as seen by changing the parameter $a$ of
Lemma \ref{minimizersexist} to $t_0 - \tau$),
and $\be\equiv\infty$ on
$\M_{t_0 - \tau}\;-\;\M_{t_0 - \tau}^+$.
Thus it suffices to show that $\be_{\min}(t_0-a)\leq 0$.

From Lemma \ref{neg}, $\beta_{\min}(\tau)<0$  for $\tau>0$ small.
Let $\tau_1\in (0,t_0-a]$ be the supremum of the
$\bar \tau\in(0,t_0-a]$ such that $\be_{\min}< 0$ on the interval $(0,\bar\tau)$.

We claim that

(a) $\be_{\min}$ is continuous on $(0,\tau_1)$

and

(b) The upper right $\tau$-derivative of $\be_{\min}$ is nonpositive
on $(0,\tau_1)$.

To see (a), pick $\tau\in (0,\tau_1)$, suppose
that $\{\tau_j\}\subset (0,\tau_1)$ is a sequence
converging to $\tau$ and choose $(x_j,t_0 - \tau_j)\in \M_{t_0 - \tau_j}^+$
such that $\be(x_j,t_0 - \tau_j)=\be_{\min}(\tau_j)< 0$.
By Lemma \ref{minimizersexist}
part (3), the sequence $\{(x_j,t_0 - \tau_j)\}$ subconverges to some
$(x,t_0 - \tau)\in\M_{t_0 - \tau}^+$
for which $\be(x,t_0 - \tau)\leq \liminf_{j\ra\infty}\be(x_j,t_0 - \tau_j)$.
Thus
$\be_{\min}$ is lower semicontinuous at $\tau$.  On the other hand,
since $\be_{\min}(\tau)< 0$, the minimum of $\be$ on $\M_{t_0 - \tau}$
will be attained at a point
$(x,t_0 - \tau)\in \M_{t_0 - \tau}^+$ lying
in the interior of $\M_{t_0 - \tau}^-\cap \M_{t_0 - \tau}^+$,
as $\be>0$ elsewhere on $\M_{t_0 - \tau}$
(by Lemma \ref{minimizersexist} and the
hypothesis on admissible curves).  Therefore $\be$ is continuous at
$(x,{t_0 - \tau})$,
which implies that $\be_{\min}$ is upper semicontinuous at $\tau$.
This gives (a).

Part (b) of the claim follows from the fact that if $\tau\in (0,\tau_1)$ and
the minimum of $\be$ on $\M_{t_0 - \tau}$ is attained at
$(x,t_0 - \tau)$ then $l(x,\tau)< \frac32$, so there is a neighborhood $U$
of $(x,{t_0 - \tau})$ such that the inequality
\begin{equation}
\frac{\D \be}{\D \tau}+\De \be \leq 0
\end{equation}
holds in the barrier sense on $U$ (by Lemma \ref{minimizersexist} and the hypothesis
on admissible curves).
Hence the upper right derivative $\frac{d}{ds} \Big|_{s=0} \beta(x, \tau+s)$ is
nonpositive, so the upper right $\tau$-derivative of
$\beta_{\min}(\tau)$ is also nonpositive.


The claim implies that $\be_{\min}$ is  nonincreasing on $(0,\tau_1)$,
and so
$\limsup_{\tau\ra \tau_1^-}\be(\tau)<0$.
By parts (2) and (3) of Lemma \ref{minimizersexist},
we have $\be_{\min}(\tau_1)<0$, and the minimum is attained at some
$(x,{t_0 - \tau_1})\in
\M_{\reg}$.
(Recall that $\M_a \subset \M_{\reg}$.)
This implies that $\tau_1=t_0-a$, for otherwise $\be_{\min}(\tau)$
would be strictly negative for $\tau\geq\tau_1 $ close to $\tau_1$,
contradicting the definition of $\tau_1$.
\end{proof}

The notion of local collapsing can be adapted to Ricci flows with surgery,
as follows.
\begin{definition}[kappa noncollapsing]
  \label{def:kl-reduced-l-geometry-kappa-noncollapsing}
  \dcref{surgerycollapsedef}
  \group{Kleiner--Lott Surgery Existence}
  \source{kleiner-lott}

\label{surgerycollapsedef}
Let $\M$ be a Ricci flow with surgery defined on $[a,b]$. Suppose that
$(x_0,t_0)\in\M$ and $r>0$ are such that $t_0-r^2\geq a$,
$B(x_0,t_0,r)\subset\M_{t_0}^-$ is a proper ball
and the parabolic ball $P(x_0,t_0,r,-r^2)$ is unscathed.
Then {\em $\M$ is $\kappa$-collapsed at $(x_0,t_0)$ at scale
$r$} if $|\Rm|\leq r^{-2}$ on $P(x_0,t_0,r,-r^2)$ and
$\vol(B(x_0,t_0,r))<\kappa r^3$; otherwise it is $\kappa$-noncollapsed.
\end{definition}

We make use of the following variant of the noncollapsing
argument from Section \ref{I.7.3}.

\begin{lemma}[existence for reduced volume]
  \label{lem:kl-reduced-l-geometry-existence-reduced-volume}
  \dcref{localreducedvolume}
  \group{Kleiner--Lott Surgery Existence}
  \source{kleiner-lott}
(Local version of reduced volume comparison)
\label{localreducedvolume}
There is a function $\kappa':\R_+\ra \R_+$, satisfying
$\lim_{\kappa\ra 0}\kappa'(\kappa)=0$, with the following property.
Let $\M$ be a Ricci flow with surgery defined on $[a,b]$. Suppose
that we are given $t_0\in (a,b]$,
$(x_0,t_0)\in \M_{t_0}\cap \M_{\reg}$, $t\in [a,t_0)$ and
$r \in (0,\sqrt{t_0-t})$.  Let $Y$ be the set of points $(x,t)\in \M_t$
that are accessible from $(x_0,t_0)$ by means of minimizing $\L$-geodesics
which remain in $\M_{\reg}$.  Assume
in addition that
$\M$ is $\kappa$-collapsed at $(x_0,t_0)$ at scale $r$, i.e.
$P(x_0,t_0,r,-r^2) \cap \M_{[t_0 - r^2, t_0)} \subset\M_{\reg}$,
$|\Rm|\leq r^{-2}$ on
$P(x_0,t_0,r,-r^2)$, and $\vol(B(x_0,t_0,r))<\kappa r^3$.
Then  the reduced volume of $Y$ is at most
$\kappa'(\kappa)$.
\end{lemma}
\begin{proof}
  \uses{thm:kl-no-local-collapsing-theorem-ii-existence-kappa-solutions}
Let $\hat Y\subset T_{x_0}\M_{t_0}$ be the set of vectors
$v\in T_{x_0}\M_{t_0}$
such that there is a minimizing $\L$-geodesic $\ga:[t,t_0]\ra \M_{\reg}$
 running from $(x_0,t_0)$ to some point in $Y$, with
\begin{equation}
\lim_{\bar t\ra t_0}\sqrt{t_0-\bar t}\;\dot{\ga}(\bar t)=-v.
\end{equation}

\noindent
The calculations from Sections \ref{Lremarks}-\ref{monoV}
apply to $\L$-geodesics sitting in $\M_{\reg}$.
In particular, the monotonicity of the weighted Jacobian
$\tau^{-\frac{n}{2}}\exp(-l(\tau))J(\tau)$ holds.  Now one repeats
 the proof of Theorem \ref{nolocalcollapse}, working with the set $\hat Y$
instead of  the set of initial velocities of {\em all} minimizing
$\L$-geodesics.
\end{proof}


\section{Establishing noncollapsing in the presence of surgery}
\label{II.5.2section}

The key result of this section, Lemma \ref{II.5.2},
gives conditions under which
one can deduce noncollapsing on a time interval $I_2$, given
a noncollapsing bound on a preceding interval $I_1$ and
lower bounds on $r$ on $I_1\cup I_2$.

\begin{definition}[noncollapsing definition]
  \label{def:kl-establishing-noncollapsing-presence-surgery-noncollapsing-definition}
  \dcref{establishing-noncollapsing-presence-surgery:definition-1}
  \group{Kleiner--Lott Surgery Existence}
  \source{kleiner-lott}

The {\em ${\mathcal L}_+$-length} of an admissible curve $\gamma$ is
\begin{equation}
{\mathcal L}_+(\gamma, \tau) \: = \:
\int_{t_0 - \tau}^{t_0} \sqrt{t_0 - t} \: \left(
R_+(\gamma(t), t) \: + \: |\dot{\gamma}(t)|^2 \right) \: dt,
\end{equation}
where $R_+(x,t) = \max(R(x,t), 0)$.
\end{definition}

\begin{lemma}[existence for noncollapsing]
  \label{lem:kl-establishing-noncollapsing-presence-surgery-existence-noncollapsing}
  \dcref{lplusbig}
  \group{Kleiner--Lott Surgery Existence}
  \source{kleiner-lott}
(Forcing $\L_+$ to be large, cf. Lemma II.5.3)
\label{lplusbig}

For all $\La< \infty$, $\bar r>0$ and $\hat r>0$,
there is  a constant  $F_0=F_0(\La,\bar r,\hat r)$
with the following property.  Suppose that

$\bullet$  $\M$ is a Ricci flow  with $(r,\de)$-cutoff
defined on an interval containing $[t,t_0]$, where
 $r([t,t_0])\subset [\hat r,\eps]$,

$\bullet$ $r_0\geq \bar r$,
 $B(x_0,t_0,r_0)$ is a proper ball which is unscathed on $[t_0-r_0^2,t_0]$, and
$|\Rm|\leq r_0^{-2}$ on $P(x_0,t_0,r_0,-r_0^2)$,

$\bullet$ $\ga:[t,t_0]\ra\M$ is an admissible curve
ending at $(x_0, t_0)$
whose image
is not contained in $\M_{\reg}\cup\M_{t_0}$, and

$\bullet$ $\de < F_0(\La,\bar r,\hat r)$
on $[t,t_0]$.

\noindent
Then $\L_+(\ga)>\La$.
\end{lemma}
\begin{proof}
  \uses{cor:kl-curves-that-penetrate-surgery-region-ricci-flow-surgery-corollary,lem:kl-evolution-surgery-cap-existence-standard-solutions}
The idea is that the hypotheses on  $\ga$  imply that it must touch the
part of the manifold added during surgery at some time $\bar t\in[t,t_0]$.
Then either $\ga$ has to move very fast at times close to $\bar t$ or
$t_0$, or it
will stay in the surgery region while it develops large
scalar curvature. In the first case $\L_+(\ga)$ will be large because
of the $|\dot{\ga}|^2$ term in the formula for $\L_+$,
and in the second case it will be large because of
the $R(\ga)$ term.

First, we can assume that $F_0$ is small enough so that
$F_0 \: < \: \sqrt{\frac{\bar r}{100 \epsilon}}$.
Then since
\begin{equation}
\label{h<<r}
\max_{[t,t_0]}h(t)\leq \left(\max_{[t,t_0]}\de\right)^2
\left(\max_{[t,t_0]}r(t)\right)\leq F_0^2\eps,
\end{equation}
we have $\max_{[t,t_0]}h(t)
<\frac{\bar r}{100}\leq  \frac{r_0}{100}$.

Put
$\De t=10^{-10} \bar r^4\La^{-2}$.
It always suffices to prove
the lemma for a larger value of $\La$,
so without loss of generality we can assume that
$\De t\leq \bar r^2\leq r_0^2$. Set
\begin{equation}
\label{aandthetaconditions}
A=A((\De t)^{-\frac12}\La,\hat r),\quad
\th=\th((\De t)^{-\frac12}\La,\hat r)
\end{equation}
where $A(\cdot,\cdot)$ and $\th(\cdot,\cdot)$
are  the functions from Corollary \ref{corollaryII.4.6}.
That is, we will eventually be applying Corollary \ref{corollaryII.4.6}
with $l \: = \: (\De t)^{-\frac12}\La$.
We impose the additional constraint on $F_0$ that
\begin{equation}
\label{deltasmall}
 F_0\leq \hat\de(A+2,\th,\hat r)
\end{equation}
on the interval $[t, t_0]$,
where $\hat\de $ is the function from Lemma \ref{II.4.5}.

As $\gamma$ is admissible but is not contained in $\M_{\reg}\cup\M_{t_0}$,
it must pass through the boundary of a surgery cap at some time
in the interval $[t, t_0)$ or it must start in the interior of a
surgery cap at time $t$.
By dropping an initial segment of $\ga$
if necessary, we may assume that $\ga(t)$ lies in a
surgery cap.

Let $x$ denote the tip of the surgery cap.
Note that
\begin{equation}
\label{p'sdisjoint}
P(x_0,t_0,r_0,-r_0^2)\cap P(x,t,Ah(t),\th h^2(t))=\emptyset
\end{equation}
since by Lemma \ref{II.4.5} the scalar curvature on
$P(x,t,Ah(t),\th h^2(t))$ is at least
$\frac{h^{-2}}{2}>\frac{10^4}{2}r_0^{-2}$, while $|\Rm|\leq r_0^{-2}$
on $P(x_0,t_0,r_0,-r_0^2)$.  Therefore when going backward in time
from $(x_0,t_0)$, $\ga$ must leave the parabolic region
$P(x_0,t_0,r_0,-r_0^2)$
before it arrives at $(x,t)$.
If it exits at a time $\tilde t > t_0-\De t$ then
applying the Schwarz inequality we get
\begin{equation}
\int_{\tilde{t}}^{t_0} \sqrt{t_0 - s} \:
|\dot{\gamma}(s)|^2 \: ds \: \ge \:
\left( \int_{\tilde{t}}^{t_0}
|\dot{\gamma}(s)| \: ds \right)^2
\left( \int_{\tilde{t}}^{t_0} (t_0 - s)^{-1/2}
\: ds \:\right)^{-1}
\ge \: \frac{1}{100} \: r_0^2 (\Delta t)^{-1/2}>\La,
\end{equation}
where the factor of $\frac{1}{100}$ comes from
the length distortion estimate
of Section \ref{secI.8.3},
using the fact that $|\Rm| \le r_0^{-2}$ on $P(x_0, t_0, r_0, -r_0^2)$.
So we can restrict to the case when $\ga$ exits
$P(x_0,t_0,r_0,-\De t)$ through the
initial time slice  at time $t_0-\De t$. In particular,
by (\ref{p'sdisjoint}),
$\gamma$ must exit the parabolic region $P(x,t,Ah(t),\th h^2(t))$
by time $t_0-\De t$.

By Lemma \ref{II.4.5}, the parabolic region $P(x,t,Ah,\theta h^2)$ is
either unscathed, or it coincides (as a set) with the parabolic
region $P(x,t,Ah,s)$ for some $s\in (0,\th h^2)$ and the entire
final time slice $P(x,t,Ah,s)\cap \M_{t+s}$
of $P(x,t,Ah,s)$ is thrown away by a surgery
at time $t+s$.

One possibility is that  $\gamma$ exits $P(x,t,Ah,\theta h^2)$
through the final time slice. If this is the case then $P(x,t,Ah,\theta h^2)$
must be unscathed (as otherwise the final face is removed by surgery at
time $t+s < t+\theta h^2$ and $\gamma$ would have nowhere to go after this time), so
$\gamma$ lies in $P(x,t,Ah,\theta h^2)$ for the entire time interval
$[t, t+\theta h^2]$.

The other possibility is that $\ga$ leaves $P(x,t,Ah,\theta h^2)$ before the
final time slice of $P(x,t,Ah, \theta h^2)$, in which case it exits
the ball $B(x,t,Ah)$ by time $t + \theta h^2$.

Corollary \ref{corollaryII.4.6} applies to either of these
two possibilities. Putting
\begin{equation}
T_\ga= \sup\{\bar t\in [t,t+\theta h^2]\mid \ga([t,\bar t])\subset
P(x,t,Ah,\theta h^2)\}
\end{equation}
and using the fact that $T_\ga \le t_0 - \De t$,
we have
\begin{align}
& \int_{t}^{t_0} \sqrt{t_0 - s} \left( R_+(\gamma(s), s) \: + \:
|\dot{\gamma}(s)|^2 \right) \: ds \: \ge \:
\int_{t}^{T_\gamma}
\sqrt{t_0 - s} \left( R_+(\gamma(s), s) \: + \:
|\dot{\gamma}(s)|^2 \right) \: ds \: \ge \\
& (\Delta t)^{1/2} \int_{t}^{T_\gamma}
\left( R_+(\gamma(s), s) \: + \:
|\dot{\gamma}(s)|^2 \right) \: ds \: \ge \:
(\Delta t)^{1/2} \: l \: = \: \La, \notag
\end{align}
where the last inequality comes from Corollary \ref{corollaryII.4.6}
and the choice
of $A,\th$, and $\de$ in (\ref{aandthetaconditions}) and (\ref{deltasmall}).
This completes the proof.
\end{proof}

\begin{lemma}[vanishing for noncollapsing]
  \label{lem:kl-establishing-noncollapsing-presence-surgery-vanishing-noncollapsing}
  \dcref{Rev}
  \group{Kleiner--Lott Surgery Existence}
  \source{kleiner-lott}
 \label{Rev}
If $\M$ is a Ricci flow with surgery, with normalized initial
condition at time zero, then for all $t \ge 0$, $R(x,t) \: \ge \:
- \: \frac{3}{2} \: \frac{1}{t+\frac14}$.
\end{lemma}
\begin{proof}
From the initial conditions, $R_{min}(0) \: \ge \: -6$.
If the Ricci flow is smooth then
(\ref{Rlowerbound}) implies
that $R_{min}(t) \: \ge \: - \: \frac{3}{2} \: \frac{1}{t+\frac14}$.
If there is a surgery
at time $t_0$ then $R_{min}$ on $\M_{t_0}^+$ equals
$R_{min}$ on $\M_{t_0}^-$, as surgery is done in
regions of high scalar curvature.  The lemma follows by applying
(\ref{Rlowerbound}) on the time intervals between
the singular times.
\end{proof}

In the statement of the next lemma, one has successive
time intervals $[a,b)$ and $[b,c)$. As a mnemonic we
use the subscript $-$ for quantities attached to the
earlier interval $[a,b)$, and $+$ for those associated with $[b,c)$.
We will also assume that the global parameter $\epsilon$ is small
enough that the $\Phi$-pinching condition
implies that whenever $|\Rm(x,t)|\geq \eps^{-2}$, then
$R(x,t)>\frac{|\Rm(x,t)|}{100}$.
(We  remind the reader of the role of the parameter
$\epsilon$; see Remark \ref{remepsilon}.)

\begin{lemma}[existence for kappa solutions]
  \label{lem:kl-establishing-noncollapsing-presence-surgery-existence-kappa-solutions}
  \dcref{II.5.2}
  \group{Kleiner--Lott Surgery Existence}
  \source{kleiner-lott}
(Noncollapsing estimate) (cf. Lemma II.5.2)
\label{II.5.2}

Suppose $\epsilon \ge r_- \ge r_+ >0$, $\kappa_->0$, $E_->0$ and $E<\infty$.
Then there are constants $\overline{\delta} =
\overline{\delta}(r_-,r_+,\kappa_-,E_-,E)$
and $\kappa_+=\kappa_+(r_-,\kappa_-,E_-,E)$ with
the following property.  Suppose that

$\bullet$  $a<b<c,\quad b-a\geq E_-,\quad c-a\leq E$,

$\bullet$  $\M$ is a Ricci flow with $(r,\de)$-cutoff
with normalized initial condition
defined on a time interval containing $[a,c)$,

$\bullet$ $r\geq r_-$ on $[a,b)$ and $r\geq r_+$ on $[b,c)$,

$\bullet$ $r \le \epsilon$ ,

$\bullet$ $\M$ is $\kappa_-$-noncollapsed at scales below
$\eps$ on $[a,b)$ and

$\bullet$ $\de\leq \bar\de$ on $[a,c)$,


Then $\M$ is $\kappa_+$-noncollapsed at scales below $\eps$
on $[b,c)$.
\end{lemma}
\begin{remark}[estimates for kappa solutions]
  \label{rem:kl-establishing-noncollapsing-presence-surgery-estimates-kappa-solutions}
  \dcref{establishing-noncollapsing-presence-surgery:remark-5}
  \group{Kleiner--Lott Surgery Existence}
  \source{kleiner-lott}

The important point to notice here is that $\bar\de$ is allowed
to depend on the lower bound $r_+$ on $[b,c)$, but the
noncollapsing constant $\kappa_+$ does not depend on $r_+$.
\end{remark}
\begin{proof}
  \uses{def:kl-priori-assumptions-canonical-neighborhood-assumption,lem:kl-canonical-neighborhood-property-standard-solutions-existence-canonical-neighborhoods,lem:kl-curvature-bounds-from-priori-assumptions-canonical-neighborhoods-lemma,lem:kl-curvature-bounds-from-priori-assumptions-estimates-curvature,lem:kl-establishing-noncollapsing-presence-surgery-existence-kappa-solutions,lem:kl-establishing-noncollapsing-presence-surgery-existence-noncollapsing,lem:kl-establishing-noncollapsing-presence-surgery-vanishing-noncollapsing,lem:kl-evolution-surgery-cap-existence-standard-solutions,lem:kl-performing-surgery-continuing-flows-existence-canonical-neighborhoods,lem:kl-reduced-l-geometry-existence-reduced-length,lem:kl-reduced-l-geometry-existence-reduced-volume,lem:kl-reduced-l-geometry-existence-ricci-flow-surgery,lem:kl-standard-solutions-existence-kappa-solutions,lem:kl-structure-at-first-singularity-time-characterization-structure-at-first-singularity-time,prop:kl-statement-existence-theorem-ricci-flow-surgery-monotonicity-kappa-solutions,thm:kl-canonical-neighborhood-theorem-compactness-canonical-neighborhoods,thm:kl-no-local-collapsing-theorem-ii-existence-kappa-solutions}
In the proof, we can assume that $\frac{r_+}{100} \le \sqrt{E_-/3}$.
If this were not the case then we could prove the lemma with
$r_+$ replaced by $100 \sqrt{E_-/3}$. Then the lemma would
also hold for the original value of $r_+$.

First, from Lemma \ref{Rev}, $R \ge -6$ on
$\M_{[a,c)}$.

Suppose that $r_0\in (0,\eps)$,  $(x_0,t_0)\in\M_{[b,c)}$,
  $B(x_0,t_0,r_0)$
is a proper ball  unscathed on the interval $[t_0-r_0^2,t_0]$, and
$|\Rm|\leq r_0^{-2}$ on $P(x_0,t_0,r_0,-r_0^2)$.


We first assume that
$r_0\leq \sqrt{E_-/3}$
and $r_0\geq \frac{r_+}{100}$.


We will consider $\L$-length, $\L_+$-length, etc in
$\M_{[a,t_0)}$ with basepoint at $(x_0,t_0)$. Suppose that
$\widehat{t} \in [a,t_0)$.
Then for any admissible curve
$\ga:[\widehat{t},t_0]\ra \M_{[a,t_0]}$
ending at $(x_0,t_0)$, we have

\begin{gather}
\label{llplus}
\L(\ga)\leq \L_+(\ga)\leq \L(\ga)+\int_a^c6\sqrt{c-t}\: dt
\leq \L(\ga)+4E^{\frac32} \\
\mbox{and}\quad l(x,t)\geq\frac{L_+-4E^{\frac32}}{2E^{\frac12}}. \notag
\end{gather}

Assume that $\bar\de\leq F_0(4E^{\frac12}+4E^{\frac32},\frac{r_+}{100},r_+)$
where $F_0$ is the function from Lemma \ref{lplusbig}.
Then by (\ref{llplus}) and Lemma \ref{lplusbig},
we conclude that any admissible
curve
$[\widehat{t},t_0]\ra \M_{[a,t_0]}$
ending at $(x_0,t_0)$
which does not lie in $\M_{\reg}\cup \M_{t_0}$
has reduced length bounded below by $2=\frac32+\frac12$.  By Lemma
\ref{stillleq3/2} there is an admissible curve $\ga:[a,t_0]\ra\M$
ending at $(x_0,t_0)$ such that
\begin{equation}
\L(\ga)=L(\ga(a))=2\sqrt{t_0-a}\;l(\ga(a))\leq 3\sqrt{t_0-a},
\end{equation}
so by (\ref{llplus}) it follows that
\begin{equation}
\label{lplusupper}
\L_+(\ga)\leq 3\sqrt{t_0-a}+4E^{\frac32}\leq 3\sqrt{E}+4E^{\frac32}.
\end{equation}


Set
\begin{equation}
t_1= a+\frac{b-a}{3},\quad t_2= a+\frac{2(b-a)}{3}
\end{equation}
and
\begin{equation}
\rho= \left(3\sqrt{E}+4E^{\frac32}\right)
\left(\frac13 E_-\right)^{-\frac32}.
\end{equation}
By construction, $t_2 \le t_0 - r_0^2$.
Note that there is a $\bar t\in [t_1,t_2]$ such that
$R(\ga(\bar t))\leq \rho$.
Otherwise we would get
\begin{equation}
\L_+(\ga) > \int_{t_1}^{t_2}\sqrt{t_0-t}\;R_+(\ga(t)) \: dt
\geq \sqrt{\frac13 E_-}\int_{t_1}^{t_2}\rho \: dt
\geq \sqrt{\frac13 E_-}\left(\frac13 E_-\right)\rho=3\sqrt{E}+4E^{\frac32},
\end{equation}
contradicting (\ref{lplusupper}).

Put $\overline{x} = \gamma(\overline{t})$.
By Lemma \ref{claim1II.4.2}, there is an estimate of the form
\begin{equation}
\label{Rbound2}
R\leq \const s^{-2}
\end{equation}
on the parabolic ball
$\hat P = P(\overline{x}, \bar t,s,-s^2)$ with
$s^{-2} \: = \: \const \: (\rho \: + \: r_-^{-2})$.
Appealing to Hamilton-Ivey curvature pinching as
usual, we get that $|\Rm|\leq \const s^{-2}$ in $\hat P$.
If $\overline{t} - s^2 < a$ then
we shrink $s$ (as little as possible)
to ensure that $\hat P\subset\M_{[a,b)}$.
Provided that  $\bar \de$ is less than a small constant
$c_1=c_1(r_-,E_-,E)$,
we can guarantee that $\hat P$ is unscathed, by forcing the
curvature in a surgery cap to exceed our bound (\ref{Rbound2})
on $R$.  Put
$U= \M_{\bar t-\frac12 s^2}\cap \hat P$.
If $s < \epsilon$ then
the $\kappa_-$-noncollapsing assumption on $[a,b)$  gives a lower bound
on $\vol(B(\bar x,\bar t,s))s^{-3}$. If $s \ge \epsilon$ then
the $\kappa_-$-noncollapsing assumption gives a lower bound
on $\vol(B(\bar x,\bar t,\epsilon/2)) (\epsilon/2)^{-3}$.
In either, case, we get a lower bound on $\vol(B(\bar x,\bar t,s))$ and
hence a lower bound
$\vol(U)\geq v=v(r_-,\kappa_-,E_-,E)$.
Now every point
in $U$ can be joined to $(x_0,t_0)$ by  a curve of
$\L_+$-length at most $\La_+=\La_+(r_-,E_-,E)$,
by concatenating an admissible curve
$[\bar t-\frac12 s^2,\bar t]\ra \M$ (of controlled
$\L_+$-length) with $\ga\restr_{[\bar t,t_0]}$.
Shrinking $\bar\de$ again, we can apply
Lemmas \ref{minimizersexist} and \ref{lplusbig} with
(\ref{llplus}) to ensure that every point in
$U$ can be joined to $(x_0,t_0)$ by a minimizing
$\L$-geodesic lying in $\M_{\reg}\cup \M_{t_0}$.
Lemma \ref{localreducedvolume}
then implies that
\begin{equation}
\label{kappa1}
\vol(B(x_0,t_0,r_0))r_0^{-3}\geq \kappa_1=\kappa_1(r_-,\kappa_-,E_-,E).
\end{equation}

(We briefly recall the argument. We have a parabolic ball around
$(\bar x, \bar t)$, of small but controlled size, on which we have
uniform curvature bounds.  The lower volume bound coming from
the $\kappa_-$-noncollapsing assumption
on $[a,b)$ means that we have bounded geometry on the parabolic ball.
As we have a fixed upper bound on $l(\bar x, \bar t)$, we can
estimate from below the reduced volume of the accessible points
$Y \subset \M_{\bar t-\frac12 s^2}^+$. Then we obtain a lower
bound on $\vol(B(x_0,t_0,r_0))r_0^{-3}$ as in Theorem
\ref{nolocalcollapse}.)

This completes the proof of the lemma when
$r_0\leq \sqrt{E_-/3}$
and $r_0\geq \frac{r_+}{100}$.

Now suppose that $r_0>\sqrt{E_-/3}$.  Applying our
noncollapsing estimate (\ref{kappa1}) to the ball of radius
$\sqrt{E_-/3}$ gives
\begin{equation}
\vol(B(x_0,t_0,r_0))r_0^{-3}\geq \left(\vol(B(x_0,t_0,\sqrt{E_-/3}) (E_-/3)^{-\frac32}
\right)\frac{(E_-/3)^{\frac32}}{r_0^3}
\geq \kappa_1\frac{(E_-/3)^{\frac32}}{\eps^3}=\kappa_2,
\end{equation}
where $\kappa_2=\kappa_2(r_-,\kappa_-,E_-,E)$.

The next sublemma deals with the case when $r_0< \frac{r_+}{100}$.
\begin{lemma}[kappa solutions lemma]
  \label{lem:kl-establishing-noncollapsing-presence-surgery-kappa-solutions-lemma}
  \dcref{r0/100}
  \group{Kleiner--Lott Surgery Existence}
  \source{kleiner-lott}

\label{r0/100}
If $r_0< \frac{r_+}{100}$ then $\vol(B(x_0,t_0,r_0))r_0^{-3}\geq
\kappa_3=\kappa_3(r_-,\kappa_-,E_-,E)$.
\end{lemma}
\begin{proof}
  \uses{def:kl-priori-assumptions-canonical-neighborhood-assumption,lem:kl-establishing-noncollapsing-presence-surgery-existence-kappa-solutions,lem:kl-standard-solutions-existence-kappa-solutions}
Let $s$ be the maximum of the numbers $\bar s\in [r_0,\frac{r_+}{100}]$
such that $B(x_0,t_0,\bar s)$ is unscathed on $[t_0-\bar s^2,t_0]$,
and $|\Rm|\leq \bar s^{-2}$ on $P(x_0,t_0,\bar s,-\bar s^{-2})$.
Then  either

(a) Some point $(x,t)$ on the frontier of $P(x_0,t_0,s,-s^2)$
lies in a surgery cap

or

(b) Some point $(x,t)$ in the closure of $P(x_0,t_0,s,-s^2)$
has $|\Rm|=s^{-2}$


or

(c) $s=\frac{r_+}{100}$.

In case (a) the scalar curvature
at $(x,t)$ will satisfy $R(x,t)\in (\frac{h^{-2}}{2},10h^{-2})$,
since $(x,t)$ lies in the cap at the surgery time.  Since
$|\Rm(x,t)|\leq s^{-2}$ we conclude that $s \le \const h(t)$.
If $\bar\de$ is small then the pointed time slice
$(\M_t,(x,t))$ will be close, modulo scaling by $h(t)$, to
the initial condition of the standard solution with
the basepoint somewhere in the cap.  Using the
fact that the time slices of $P(x_0,t_0,s,-s^2)$ have
comparable metrics, along with the fact that
$r_0 \le s  \le \const h(t)$, we get a lower bound
$\vol(B(x_0,t_0,r_0))r_0^{-3}\geq \mbox{const}$.

In case (b), we have a static curve $\ga:[t,t_0]\ra\M$
such that $\ga(t)=(x,t)$, $\ga(t_0)\in \ol{B(x_0,t_0,s)}$,
and $|\Rm(x,t)|=s^{-2}\geq 10^4r_+^{-2}$.
Hence by $\Phi$-pinching,
$R(x,t)\geq \frac{1}{100} s^{-2} \geq 100 r_+^{-2}$
(cf. the remark just before the statement of
Lemma \ref{II.5.2}).
With reference to the constant $\eta$ of
(\ref{II.(1.3)estimates}),
put $\sigma \: = \: 10^{-6} \: \min \left( 1,
\frac{1}{\eta} \right)$.
Let
$\al:[0,\widehat{\rho}]\ra \ol{B(x_0,t_0,s)}\subset \M_{t_0}^-$
be a minimizing geodesic from $(x_0,t_0)$ to
$\ga(t_0)\in \ol{B(x_0,t_0,s)}$. We can find a point
$z$ along $\alpha$ with $\dist_{t_0}(z, \gamma(t_0)) \le \sigma s$ and
$\dist_{t_0}(z, x_0) \le (1 - \sigma) s$.
Let $\bar\ga:[t,t_0]\ra\M$
be the static curve ending at $z$
and put $(\bar x,t)= \bar\ga(t)$.  In brief,
we get $(\bar x,t)$ by ``pulling $(x,t)$ slightly
inward from the boundary''.

From the distance distortion estimate of Section \ref{secI.8.3},
we can say that $\dist_{t} (\bar x, x) \le 10^6 \sigma s$.
Then applying (\ref{II.(1.3)estimates}) along a minimizing
time-$t$ curve from $\bar x$ to $x$, we conclude that
\begin{equation}
|R^{- \: \frac12}(\bar x, t) \: - \: R^{- \: \frac12}(x, t)| \: \le \:
\frac12 \: \eta \: \dist_{t} (\bar x, x) \: \le \: \frac12 \: s,
\end{equation}
so $R^{- \: \frac12}(\bar x, t) \: \le \:
R^{- \: \frac12}(x, t) \: + \: \frac12 \: s \: \le 20s$ and hence
$R(\bar x, t) \: \ge \: \frac{1}{400} \: s^{-2} \ge \: 25 \: r_+^{-2}$.
In particular, $(\bar x,t)$ has a canonical neighborhood.
We also know that $R(\bar x, t) \le 6 |\Rm (\bar x, t)| \le 6 s^{-2}$.
It follows from the definition of canonical neighborhoods
(see Definition \ref{canonicalnbhddef}) that
there is some universal constant so that
$\vol(B(\bar x,t,10^{-9}s) ) \ge \const s^3$.
(We recall that from Lemma \ref{lemstandard0},
there is a $\kappa>0$ such that a standard solution
is $\kappa$-noncollapsed as scales $<1$.)
The distance distortion
estimate ensures that $B(\bar x,t,10^{-9}s) \subset
B(z,t_0,10^{-6}s) \subset B(x_0,t_0,s)$. Then the standard volume
distortion estimate implies that
$\vol(B(x_0,t_0,s)) \: \ge \: \const
\vol(B(\bar x,t,10^{-9}s) ) \ge \const s^3$,
again for some universal constant.
Finally we use Bishop-Gromov volume comparison to get
$\vol(B(x_0,t_0,r_0))r_0^{-3}\geq \const \vol(B(x_0,t_0,s))s^{-3}$.

In case (c) we apply (\ref{kappa1}), replacing the
$r_0$ parameter there by $s$, and Bishop-Gromov
volume comparison as in case (b).
\end{proof}


\section{Construction of the Ricci flow with surgery}  \label{surgeryflow}

The proof is by induction on $i$.
To start the induction process,
we observe that
the initial normalization
$|\Rm|\leq 1$ at $t=0$
implies that a smooth solution exists for some definite time
\cite[Corollary 7.7]{Chow-Knopf}.
The curvature bound on this time interval, along with the
volume assumption on the initial time balls,
implies that the solution is $\kappa$-noncollapsed below  scale $1$
and satisfies the  $\rho$-canonical neighborhood assumption
vacuously for small $\rho>0$.

Now assume inductively that $r_j$, $\kappa_j$, and $\bar\de_j$
have been selected for $1\leq j\leq i$, thereby defining
the functions $r$, $\kappa$, and $\bar\delta$
on $[0, 2^i\eps]$, such that for
any nonincreasing function $\de$ on $[0,2^i\eps]$ satisfying
$0<\de(t)\leq \bar\de(t)$, if one has normalized
initial data then the Ricci flow with $(r,\de)$-cutoff
is defined on $[0,2^i\eps]$ and is $\kappa(t)$-noncollapsed at
scales $<\eps$.

We will determine $\kappa_{i+1}$ using Lemma \ref{II.5.2}.
So suppose it is not possible to choose $r_{i+1}$
and $\bar\de_{i+1}$ (and, if necessary, make $\bar\de_i$ smaller)
so that if we put
$\kappa_{i+1}= \kappa_+(r_i,\kappa_i,2^{i-1}\eps,
3(2^{i-1})\eps)$ (where $\kappa_+$ denotes the function from
Lemma \ref{II.5.2})
then the inductive statement above holds with $i$ replaced by $i+1$.
Then given sequences $r^\al\ra 0$ and $\bar\de^\al\ra 0$, for each
$\al$ there
must be a counterexample, say $(M^{\al},g^{\al}(0))$,
to the statement with $r_{i+1} = r^{\al}$ and
$\bar\de_i = \bar\de_{i+1} = \bar\de^{\al}$. We assume that
\begin{equation}
  \bar\de^\al < \hat\de(\al,1-\frac{1}{\al},r^\al)
\end{equation}
where $\hat\de$ is the quantity from Lemma \ref{II.4.5}; this
will guarantee that for any $A<\infty$ and $\th\in (0,1)$
we may apply Lemma \ref{II.4.5} with parameters $A$ and $\th$
for sufficiently large $\al$.  We also assume that
\begin{equation}
\label{F_1condition}
\bar\de^{\al}<\bar\de(r_i,r^\al,\kappa_i,2^{i-1}\eps,3(2^{i-1})\eps)
\end{equation}
where $\bar\de(r_i,r^\al,\kappa_i,2^{i-1}\eps,3(2^{i-1})\eps)$ is from Lemma \ref{II.5.2}.
 By Lemma \ref{extendit}  each initial condition $(M^{\al},g^{\al}(0))$
will prolong to a Ricci flow
with surgery   $\M^{\al}$ defined on a time interval $[0,T^{\al}]$
with
$T^{\al}\in (2^i\eps,\infty]$,
which restricts to a Ricci flow with $(r,\de)$-cutoff
on any proper subinterval  $[0,\tau]$ of $[0,T^{\al}]$,
 but for which the
$r^\al$-canonical neighborhood assumption fails at some point
$(\bar x^{\al},T^{\al})$ lying in the backward time
slice
$\M^{\al-}_{T^{\al}}$.
(It is implicit in this statement that $R(\bar x^{\al},T^{\al}) \ge
\frac{1}{(r^\al)^2}$.)
Since $(M^{\al},g^{\al}(0))$ violates the
theorem, we must have $T^\al\in (2^i\eps,2^{i+1}\eps]$. By (\ref{F_1condition})
and Lemma \ref{II.5.2}, it follows that  $\M^{\al}$ is
$\kappa_{i+1}$-noncollapsed at scales below $\eps$
on the interval $[2^i \epsilon,T^\al)$, where
$\kappa_{i+1}= \kappa_+(r_i,\kappa_i,2^{i-1}\eps,
3(2^{i-1})\eps)$ and $\kappa_+$ denotes the function from
Lemma \ref{II.5.2}.

Let $(\widehat\M^\al,(\bar x^\al,0))$ be the pointed Ricci flow
with surgery obtained from $(\M^\al,(\bar x^\al,T^\al))$
by shifting time by $T^\al$ and parabolically rescaling by
$R(\bar x^\al,T^\al)$.
We also remove the part of
$(\widehat\M^\al,(\bar x^\al,0))$ after time zero and
we take the time-zero slice $\widehat\M^\al_0$ to be
diffeomorphic to $\M^{\al -}_{T^\al}$.
In brief, the rest of the proof goes as follows.  If surgeries
occur further and further away from $(\bar x^\al,0)$
in spacetime as $\al\ra\infty$, then the reasoning of Theorem
\ref{thmI.12.1} applies and we obtain a $\kappa$-solution as a limit.
This would contradict the fact that $(\bar x^\al,T^\al)$ does
not have a canonical neighborhood.  Thus there must be
surgeries in a parabolic ball of a fixed size centered
at $(\bar x^\al,0)$, for arbitrarily large $\al$. Then one
argues using
Lemma \ref{II.4.5} that the solution will be close to
the (suitably rescaled and time-shifted) standard solution,
which again leads to a canonical neighborhood and a
contradiction.

We now return to the proof.
Recall that a metric ball $B$ is
proper if the distance function from the center is a proper
function on $B$. If $T$ is a surgery time for a
Ricci flow with surgery then
a metric ball in $\M_T^-$ need not be proper.

Note that by continuity,
every  point in $\widehat\M^\al$ whose scalar
curvature is strictly greater than that of $(\bar x^\al,0)$ has a
neighborhood as in Definition \ref{canonicalnbhddef}, except
that the error estimate
is $2 \epsilon$ instead of $\eps$.

\begin{lemma}[Ricci flow with surgery lemma]
  \label{lem:kl-construction-ricci-flow-surgery-ricci-flow-surgery-lemma}
  \dcref{construction-ricci-flow-surgery:sublemma-1}
  \group{Kleiner--Lott Surgery Existence}
  \source{kleiner-lott}

For all $\la<\infty$,
the ball $B(\bar x^\al,0,\la)\subset\widehat\M^\al_0$
is proper for sufficiently
large $\al$.
\end{lemma}
\begin{proof}
  \uses{lem:kl-curvature-bounds-from-priori-assumptions-canonical-neighborhoods-lemma,lem:kl-structure-at-first-singularity-time-characterization-structure-at-first-singularity-time}
As in Lemma \ref{claim2II.4.2}, for each $\rho < \infty$
the scalar curvature on
$B(\bar x^\al, 0, \rho) \subset \widehat\M^{\al}_0$
is uniformly bounded in terms of $\al$.
(In carrrying out the proof of Lemma \ref{claim2II.4.2},
we now use the aforementioned property of having
canonical neighborhoods of quality $2\epsilon$.)
Thus $B(\bar x^\al, 0, \rho)$ has compact closure in
$\widehat\M^{\al}_0$
(see Lemma \ref{lemoverlinerisproper}), from which the sublemma follows.
\end{proof}

Let $T_1 \in [-\infty,0]$ be the infimum of the set of numbers
$\tau'\in (-\infty,0]$ such that for all $\la<\infty$,
the ball $B(\bar x^\al,0,\la)\subset\widehat\M^\al_0$
is proper,
and unscathed on $[\tau',0]$ for sufficiently
large $\al$.

\begin{lemma}[convergence for Ricci flow with surgery]
  \label{lem:kl-construction-ricci-flow-surgery-convergence-ricci-flow-surgery}
  \dcref{bdrysmoothness}
  \group{Kleiner--Lott Surgery Existence}
  \source{kleiner-lott}
 \label{bdrysmoothness}
After passing to a subsequence
if necessary, the pointed flows  $(\widehat\M^\al,(\bar x^\al,0))$
converge on the time interval $(T_1, 0]$ to a Ricci flow (without
surgery) $(\M^\infty,(\bar x^\infty,0))$ with
a smooth complete nonnegatively-curved Riemannian metric
on each time slice, and
scalar curvature globally bounded above by
some number $Q<\infty$.
(We interpret $(0, 0]$ to mean $\{0\}$ rather than the
empty set.)
\end{lemma}
\begin{proof}
  \uses{lem:kl-curvature-bounds-from-priori-assumptions-canonical-neighborhoods-lemma,lem:kl-curvature-bounds-from-priori-assumptions-estimates-curvature,lem:kl-evolution-surgery-cap-existence-standard-solutions,prop:kl-statement-existence-theorem-ricci-flow-surgery-monotonicity-kappa-solutions,thm:kl-canonical-neighborhood-theorem-compactness-canonical-neighborhoods}
Suppose first that $T_1 < 0$.
Then the arguments of Theorem \ref{thmI.12.1} apply in the
time interval $(T_1, 0]$, to give the Ricci flow (without
surgery) $(\M^\infty,(\bar x^\infty,0))$.
Since $r^{\al} \rightarrow 0$, Hamilton-Ivey pinching
implies that $\M^\infty$ will have nonnegative curvature.
The fact that the
canonical neighborhood assumption,
with $\epsilon$ replaced by $2 \epsilon$,
holds for each $\widehat\M^\al$ allows us
to deduce
that the scalar curvature of $\M^\infty$ is globally bounded above by
some number $Q<\infty$; compare with Section \ref{pf11.7}
and Step 4 of the proof of Theorem \ref{thmI.12.1}.

Now suppose that $T_1 = 0$.
The argument is similar to Steps 2 and 3 of the proof of Theorem
\ref{thmI.12.1}.
As in Step 2,
or more precisely as in  Lemma \ref{claim2II.4.2},
for each $\rho < \infty$
the scalar curvature on
$B(\bar x^\al, 0, \rho) \subset \widehat\M^{\al}_0$
is uniformly bounded in terms of $\al$.
Given $\rho < \infty$ and
$(x^\al,0) \in B(\bar x^\al, 0, \rho) \subset \widehat\M^\al$,
if the parabolic region of
Lemma \ref{claim1II.4.2} (centered around
$(x^\al, 0) \in \widehat\M^\al$) is unscathed then we can
apply
the $2 \epsilon$-canonical neighborhood assumption
on $\widehat\M^\al$, Lemma \ref{claim1II.4.2} and Appendix \ref{applocalder}
to derive bounds on the curvature derivatives at
$(x^\al,0) \in \widehat\M^\al_0$
that depend on $\rho$ but are independent of $\al$.
If the parabolic region is scathed then we can apply Lemma
\ref{II.4.5}, along with our scalar curvature bound at
$(x^\al,0)$,
to again obtain uniform bounds on the curvature derivatives at
$(x^\al,0)$.
Hence there is a subsequence of
the pointed Riemannian manifolds
$\{(\widehat\M^\al_0, \bar x^\al)\}_{\al = 1}^\infty$
that converges to a smooth complete pointed
Riemannian manifold $(\M^\infty_0, \bar x^\infty)$.
As in the previous case, it will have
bounded nonnegative sectional curvature.

This proves the lemma.  Alternatively, in the case $T_1 = 0$
one can argue directly that if the parabolic region of
Lemma \ref{claim1II.4.2} is scathed then
$(\overline{x}^\al, T^\al)$ has a canonical neighborhood;
see the rest of the proof of Proposition \ref{surgeryflowexists}.
\end{proof}

If we can show that
$T_1 \: = \: -\infty$ then $(\M^\infty,(\bar x^\infty,0))$
will be a $\kappa$-solution, which will contradict the assumption
that $(\bar x^\al,T^\al)$ does not admit a canonical neighborhood.
Suppose that $T_1 \: > \: -\infty$.
We know that for all
$\tau' \in (T_1, 0]$
and $\la<\infty$, the scalar curvature in $P(\bar x^\al,0,\la,\tau')$
is bounded by $Q+1$ when $\al$ is sufficiently large.  By
Lemma \ref{claim1II.4.2}, there exists $\sigma< T_1$ such that for all
$\la<\infty$, if (for large $\al$)
the solution $\widehat\M^\al$ is unscathed on
$P(\bar x^\al,0,\la,t^\al)$ for some $t^\al > \sigma$ then
\begin{equation}
\label{2q+1}
R(x,t)< 8(Q+2)\quad\mbox{for all}\quad (x,t)\in P(\bar x^\al,0,\la,t^\al).
\end{equation}
(In applying Lemma \ref{claim1II.4.2}, we use the fact that in the unscaled
variables,
$r(T^\al)^{-2} \le R(\bar x^\al,T^\al)$ by assumption, along with
the fact that $r(\cdot)$ is a nonincreasing function.)


By the definition of $T_1$, and after passing
to a subsequence if necessary, there exist $\la<\infty$ and a sequence
$\ga^\al:[\si^\al,0]\ra\widehat\M^\al$  of static curves so that  \\
1. $\ga^\al(0)\in B(\bar x^\al,0,\la)$ and \\
2. The point
$\ga^\al(\si^\al)$ is inserted during surgery at time $\si^\al>\si$.

For each $\al$, we may assume that $\si^\al$ is the largest number having
this property.
Put
$\xi^\al= \si^\al+\left( h^\al(\si^\al)\right)^2$.
(In the notation of \cite{Perelman2},
$\left( h^\al(\si^\al)\right)^2$ would be written
as $R(\bar x, \bar t) \: h^2(T_0)$. Note that we have no {\it a priori} control
on $h^\al(\si^\al)$.) Then $\xi^\al$
is the blowup time of the rescaled and shifted
standard solution that Lemma \ref{II.4.5}
compares with $(\widehat\M^\al,\ga^\al(\si^\al))$.  We claim
that $\liminf_{\al\ra\infty}\xi^\al>0$. Otherwise, Lemma
\ref{II.4.5} would imply that after passing to a subsequence,
there are regions of
$\widehat\M^\al$, starting from time $\si^\al$, that are better and better
approximated by rescaled and shifted standard solutions whose blowup times
$\xi^\al$ have a limit that is nonpositive,
thereby contradicting (\ref{2q+1}). Lemma \ref{II.4.5}, along with
the fact that $R({\bar x}^\al, 0) = 1$, also gives a uniform upper bound
on $\xi^\al$.

Now Lemma \ref{II.4.5} implies that for large $\al$,
the restriction of
$\widehat\M^\al$ to the time interval $[\si^\al,0]$ is well
approximated by the restriction to $[\si^\al,0]$
of a rescaled and shifted standard
solution. Then Lemma \ref{claim5} implies that $(\bar x^\al,T^\al)$
has a canonical neighborhood.
The canonical neighborhood may be either a strong $\epsilon$-neck
or an $\epsilon$-cap. (Note
a strong $\epsilon$-neck may arise when an
$\epsilon$-neck around
$(\bar x^\al,T^\al)$ extends smoothly backward in time to form
a strong $\epsilon$-neck that incorporates part of the Ricci flow
solution that existed before the surgery time $\si^\al$.)

This is a contradiction.
\end{proof}

\section{II.6. Double sided curvature bound in the thick part}
\label{II.6}

Having shown that for a suitable choice of the functions $r$ and $\de$,
the Ricci flow with $(r,\de)$-cutoff exists for all time and for every
normalized initial condition, one wants to understand
its implications.
The main results in II.6 are noncollapsing and curvature estimates
which form the basis of the analysis of the large-time behavior given
in II.7.

\begin{lemma}[compactness for thick--thin geometry]
  \label{lem:kl-double-sided-curvature-bound-thick-part-compactness-thick-thin-geometry}
  \dcref{pscempty}
  \group{Kleiner--Lott Long-Time Curvature}
  \source{kleiner-lott}
 \label{pscempty}
If $\M$ is a Ricci flow with $(r,\de)$-cutoff on a compact manifold and
$g(0)$ has positive scalar curvature then the solution goes extinct
after a finite time, i.e. $\M_T \: = \: \emptyset$ for some $T > 0$.
\end{lemma}
\begin{proof}
  \uses{lem:kl-establishing-noncollapsing-presence-surgery-vanishing-noncollapsing}
We apply (\ref{Rlowerbound}). This
formula is initially derived for smooth flows but because surgeries are
performed in regions of high scalar curvature, it is also valid for a Ricci
flow with surgery; cf. the proof of Lemma \ref{Rev}.
It follows that the flow goes extinct by time
$\frac{3}{2R_{min}(0)}$.
\end{proof}

\begin{lemma}[compactness for Ricci flow with surgery]
  \label{lem:kl-double-sided-curvature-bound-thick-part-compactness-ricci-flow-surgery}
  \dcref{finitetime}
  \group{Kleiner--Lott Long-Time Curvature}
  \source{kleiner-lott}
 \label{finitetime}
If $\M$ is a Ricci flow with surgery that goes extinct after a finite time,
then the initial (compact connected orientable) $3$-manifold is diffeomorphic to a
connected sum of $S^1 \times S^2$'s and quotients of the round
$S^3$.
\end{lemma}
\begin{proof}
  \uses{lem:kl-performing-surgery-continuing-flows-rigidity-ricci-flow-surgery}
This follows from Lemma  \ref{reconstruct2}.
\end{proof}

According to \cite{Colding-Minicozzi,Colding-Minicozzi2} and
\cite{Perelman3}, if
none of the prime factors in the Kneser-Milnor decomposition of
the initial manifold are aspherical
then the Ricci flow with surgery again goes extinct
after a finite time.
Along with Lemma \ref{finitetime}, this proves the
Poincar\'e Conjecture.

Passing to Ricci flow solutions that may not go extinct after a finite time,
the main result of II.6 is the following :

\begin{corollary}[existence for Ricci flow with surgery]
  \label{cor:kl-double-sided-curvature-bound-thick-part-existence-ricci-flow-surgery}
  \dcref{II.6.8}
  \group{Kleiner--Lott Long-Time Curvature}
  \source{kleiner-lott}
 \label{II.6.8}
(cf. Corollary II.6.8) For any $w > 0$ one can find
$\tau = \tau(w) > 0$, $K = K(w) < \infty$,
$\overline{r} = \overline{r}(w) > 0$ and $\theta = \theta(w) > 0$ with
the following property. Suppose we have a solution to the Ricci flow
with $(r, \delta)$-cutoff on the time interval $[0, t_0]$, with
normalized initial data.
Let $h_{\max}(t_0)$ be the maximal surgery radius on $[t_0/2, t_0]$.
(If there are no surgeries on $[t_0/2, t_0]$ then $h_{\max}(t_0) = 0$.)
Let $r_0$ satisfy \\
1. $\theta^{-1}(w) h_{\max}(t_0) \: \le \: r_0 \: \le \:
\overline{r} \sqrt{t_0}$.\\
2. The ball $B(x_0, t_0, r_0)$ has sectional curvatures at least
$- \: r_0^{-2}$ at each point.\\
3. $\vol(B(x_0, t_0, r_0)) \: \ge \: w r_0^3$.\\

Then the solution is unscathed in $P(x_0, t_0, r_0/4, - \tau r_0^2)$
and satisfies $R \: < \: Kr_0^{-2}$ there.
\end{corollary}

Corollary \ref{II.6.8} is an analog of Corollary \ref{corI.12.4}, but there
are some differences.  One minor difference is that Corollary \ref{corI.12.4} is
stated as the contrapositive of Corollary \ref{II.6.8}. Namely,
Corollary \ref{corI.12.4} assumes that $- \: r_0^{-2}$ is achieved as a sectional
curvature in $B(x_0, t_0, r_0)$, and its conclusion is that
$\vol(B(x_0, t_0, r_0)) \: \le \: w r_0^3$. The relation
with Corollary \ref{II.6.8} is the following. Suppose that
assumptions 1 and 2 of Corollary \ref{II.6.8} hold. If
$- \: r_0^{-2}$ is achieved somewhere as a sectional
curvature in $B(x_0, t_0, r_0)$ then Hamilton-Ivey pinching implies
that the scalar curvature is very large at that point,
which contradicts the conclusion
of Corollary \ref{II.6.8}. Hence assumption 3 of Corollary \ref{II.6.8}
must not be satisfied.

A more substantial difference is that the smoothness of the flow
in Corollary \ref{corI.12.4}
is guaranteed by the setup, whereas in Corollary \ref{II.6.8} we must
prove that the solution is unscathed in $P(x_0, t_0, r_0/4, - \tau r_0^2)$.

The role of the parameter $\overline{r}$ in Corollary \ref{II.6.8} is
essentially to guarantee that we can use Hamilton-Ivey pinching
effectively.


\section{II.6.5. Earlier scalar curvature bounds on smaller balls from
lower curvature bounds and a later volume bound} \label{II.6.5}

For terminology, with reference to Section \ref{Ricci flow with surgery},
by a time-dependent family $\bigcup_{t \in [c,d]} B(x, t, r)$
of metric balls we mean first that there is a static
curve $\gamma : [c,d] \rightarrow \M$, whose intersection with each
$\M_t$ will be denoted by $x$, and second that there is
a subset ${\mathcal U}$ of $\M$ so that
\begin{enumerate}
\item If $t \notin [c,d]$ then ${\mathcal U} \cap \M_t = \emptyset$.
\item If $t \in [c,d]$ is not
a singularity time then ${\mathcal U} \cap \M_t$ is the $r$-ball around $x$ in
$\M_t$.
\item If $t \in [c,d]$ is a surgery time $t_k^+$ then
${\mathcal U} \cap \M_t$ is
the image in $\M_t$ of an $r$-ball around $x$ in
$\Omega_k$, that lies entirely in $X_k^+ \subset \Omega_k$.
\end{enumerate}
At the expense of being redundant, if these conditions are
satisfied then we will say that
we have an {\em unscathed} time-dependent family of metric balls,
to emphasize that the metric balls do not touch the surgery regions.
The restriction of the Ricci-flow-with-surgery to ${\mathcal U}$
is a smooth Riemannian metric $g$ on the ``horizontal'' subbundle
of $T{\mathcal U}$.

We first state a consequence of Corollary \ref{I.11.6end}.

\begin{lemma}[estimates for scalar curvature]
  \label{lem:kl-earlier-scalar-curvature-bounds-smaller-balls-from-lower-curvature-bound-estimates-scalar-curvature}
  \dcref{I.11.6(a)}
  \group{Kleiner--Lott Long-Time Curvature}
  \source{kleiner-lott}
 \label{I.11.6(a)}
Given $w > 0$, there exist $\tau_0 = \tau_0(w) > 0$ and
$K_0^\prime = K_0^\prime(w) < \infty$
with the following property. Suppose that we have
a Ricci flow with $(r,\de)$-cutoff such that \\
1. $\bigcup_{t \in
[-\tau r_0^2, 0]} B(x_0, t, r_0)$ is an unscathed time-dependent family of
metric balls, where $\tau \: \le \: \tau_0$. \\
2. The sectional curvatures are bounded below by $- \: r_0^{-2}$
on the
above family of balls. \\
3. $\vol(B(x_0, 0, r_0)) \: \ge \: w r_0^3$. \\

Then
\begin{enumerate}
\item
$R \: \le \: K_0^\prime \: \tau^{-1} \: r_0^{-2}$
on $\bigcup_{ t \in [- \tau r_0^2/2, 0]} B(x_0, t, r_0/2)$, and
\item $\vol(B(x_0, -\tau r_0^2, r_0/2))$ is at least $\frac{1}{10}$
of the volume of the Euclidean ball of the same radius.
\end{enumerate}
\end{lemma}

Here we have changed the conclusion of Corollary \ref{I.11.6end} to obtain an
upper curvature bound on
$\bigcup_{ t \in [- \tau r_0^2/2, 0]} B(x_0, t, r_0/2)$
instead of $\bigcup_{ t \in [- \tau r_0^2/2, 0]} B(x_0, t, r_0/4)$,
but this clearly follows from the arguments of the proof of Corollary
\ref{I.11.6end}. We have also added a lower volume bound to the conclusion,
which follows from the proof of Corollary \ref{corI.11.6}(b), provided
that $\tau_0$ is sufficiently small.

An analog of Corollary \ref{II.6.8} is the following
Lemma \ref{II.6.5(a)}, which is stated as Lemma II.6.5(a)
in \cite{Perelman2}. The lemma is used there to prove Corollary \ref{II.6.8}.
Our proof of Corollary \ref{II.6.8} will use Lemma \ref{I.11.6(a)} but will
not use Lemma \ref{II.6.5(a)}.
We include the proof of Lemma \ref{II.6.5(a)} for completeness, even
though it will not be used in the sequel.

\begin{lemma}[estimates for scalar curvature]
  \label{lem:kl-earlier-scalar-curvature-bounds-smaller-balls-from-lower-curvature-bound-estimates-scalar-curvature-4}
  \dcref{II.6.5(a)}
  \group{Kleiner--Lott Long-Time Curvature}
  \source{kleiner-lott}
 \label{II.6.5(a)}
(cf. Lemma II.6.5(a))
Given $w > 0$, there exist $\tau_0 = \tau_0(w) > 0$ and $K_0 =
K_0(w) < \infty$ with the following property. Suppose that we have
a Ricci flow with $(r,\de)$-cutoff such that \\
1. The parabolic neighborhood
$P(x_0, 0, r_0, -\tau r_0^2)$ is unscathed,
where $\tau \: \le \: \tau_0$. \\
2. The sectional curvatures are bounded below by $- \: r_0^{-2}$ on
$P(x_0, 0, r_0, -\tau r_0^2)$. \\
3. $\vol(B(x_0, 0, r_0)) \: \ge \: wr_0^3$. \\

Then $R \: \le \: K_0 \: \tau^{-1} \: r_0^{-2}$ on
$P(x_0, 0, r_0/4, -\tau r_0^2/2)$.
\end{lemma}
\begin{proof}
  \uses{lem:kl-earlier-scalar-curvature-bounds-smaller-balls-from-lower-curvature-bound-estimates-scalar-curvature,lem:kl-length-distortion-estimates-distance-distortion-lemma}
If $\tau_0$ is sufficiently small, then for
$t \in [- \tau r_0^2, 0]$ and
$(x,t) \in B(x_0, t, 9r_0/10)$, the lower
curvature bound
$\Rm \: \ge \: - \: r_0^{-2}$ on $P(x_0, 0, r_0, -
\tau r_0^2)$ implies that $(x,0) \in B(x_0, 0, r_0)$ (more precisely,
that $(x,t)$ lies on a static curve with one endpoint in $B(x_0,0,r_0)$,
or equivalently, that $(x,t)\in P(x_0,0,r_0,-\tau r_0^2)$).
Thus $\bigcup_{ t \in
[-\tau r_0^2, 0]} B(x_0, t, 9r_0/10)
\subset P(x_0, 0, r_0, - \tau r_0^2)$ and so
$\Rm \: \ge \:-r_0^{-2}\:\ge\: - (9r_0/10)^{-2}$
 on
$\bigcup_{t \in
[-\tau (9r_0/10)^2, 0]} B(x_0, t, 9r_0/10)$.

Applying
Lemma
\ref{I.11.6(a)} with $r_0$ replaced by
$9r_0/10$, and slightly redefining $w$, gives that
$R \: \le \: K_0^\prime \: \tau^{-1} \: (9r_0/10)^{-2}$
on $\bigcup_{t \in [- \frac34 \tau (9r_0/10)^2, 0]} B(x_0, t, 9r_0/20)$.
Then the length distortion estimate of Lemma \ref{distancedist} implies that for
sufficiently small $\tau_0$, if
$(x,0) \in B(x_0, 0, r_0/4)$
then
$(x,t) \in B(x_0, t, 9r_0/20)$
for
$t \in [-\tau r_0^2/2, 0]$. That is,
\begin{equation}
P(x_0, 0, r_0/4, -\tau r_0^2/2) \subset
\bigcup_{t \in [- \tau r_0^2/2, 0]} B(x_0, t, 9r_0/20).
\end{equation}
In applying the length distortion estimate we use the fact that
the change in distance is estimated by
$\Delta d \: \le \:
\const \: \sqrt{K_0^\prime \: \tau^{-1} \: (9r_0/10)^{-2}} \: \cdot \:
\tau r_0^2/2$
which, for small $\tau_0$, is a small fraction
of $r_0$.

Thus we have shown that
$R \: \le \: K_0^\prime \: \tau^{-1} \: (9r_0/10)^{-2}$
on $P(x_0, 0, r_0/4, -\tau r_0^2/2)$.
This proves the lemma.
\end{proof}

The formulation of \cite[Lemma II.6.5]{Perelman2} specializes Lemma
\ref{II.6.5(a)} to the case
$w \: = \: 1 - \epsilon$. It includes the statement
\cite[Lemma II.6.5(b)]{Perelman2} saying that $\vol(B(x_0, -\tau
r_0^2, r_0/4))$ is at
least $\frac{1}{10}$ of the volume of the Euclidean ball of the same
radius.  This follows from the proof of Corollary \ref{corI.11.6}(b), provided
that $\tau_0$ is sufficiently small.

There is an evident analogy between Lemma \ref{II.6.5(a)} and
Corollary \ref{II.6.8}. However, there is the important difference that
Corollary \ref{II.6.8} (along with Corollary \ref{corI.12.4}) only
assumes a lower sectional curvature bound at the final time slice.

\section{II.6.6. Locating small balls whose subballs have almost
Euclidean volume}

The result of this section is a technical lemma about volumes
of subballs.

\begin{lemma}[existence for sectional curvature]
  \label{lem:kl-locating-small-balls-whose-subballs-almost-euclidean-volume-existence-sectional-curvature}
  \dcref{Alexlemma2}
  \group{Kleiner--Lott Long-Time Curvature}
  \source{kleiner-lott}
 \label{Alexlemma2} (cf. Lemma II.6.6)
For any $\widehat{\epsilon}, w>0$
there exists $\theta_0 = \theta_0(\widehat{\epsilon},w)$
such that if $B(x,1)$ is a metric ball of volume at least $w$, compactly
contained in a manifold without boundary with sectional
curvatures at least $-1$, then there exists a subball
$B(y,\theta_0) \subset B(x,1)$ such that every subball
$B(z,r) \subset B(y,\theta_0)$ of any radius has volume at
least $(1-\widehat{\epsilon})$ times the volume of the Euclidean ball
of the same radius.
\end{lemma}
The proof is similar to that of Lemma \ref{Alexlemma}.
Suppose that the claim is not true.
Then there is a sequence of Riemannian
manifolds $\{M_i\}_{i=1}^\infty$ and balls $B(x_i, 1)
\subset M_i$ with compact closure
so that $\Rm \Big|_{B(x_i, 1)} \: \ge \: - \: 1$ and
$\vol(B(x_i, 1)) \: \ge \: w$, along with
a sequence $r^\prime_i \rightarrow 0$ so
that each subball
$B(x_i^\prime, r_i^\prime) \subset B(x_i, 1)$ has a subball
$B(x_i^{\prime \prime}, r^{\prime \prime}_i) \subset
B(x_i^\prime, r_i^\prime)$
with $\vol(B(x_i^{\prime \prime}, r_i^{\prime \prime})) \: < \:
(1 - \widehat{\epsilon}) \: \omega_3
(r^{\prime \prime})^n$. After taking a subsequence, we can assume that
$\lim_{i \rightarrow \infty} (B(x_i, 1), x_i) \: = \: (X, x_\infty)$ in the
pointed Gromov-Hausdorff topology, where $(X, x_\infty)$ is a pointed
Alexandrov space with curvature bounded below by $-1$. From
\cite[Theorem 10.8]{Burago-Gromov-Perelman}, the Riemannian volume
forms $\dvol_{M_i}$ converge weakly to the three-dimensional Hausdorff
measure $\mu$ of $X$. If $x_\infty^\prime$ is a regular point of
$X$ then there is some $\delta > 0$ so that
$B(x_\infty^\prime, \delta)$ has compact closure in $X$ and
for all $r < \delta$,
$\mu(B(x_\infty^\prime,
r)) \: \ge \: (1-\frac{\widehat{\epsilon}}{10}) \: \omega_3 r^3$. Fixing
such an $r$ for the moment,
for large $i$ there are balls
$B(x_i^\prime, r) \subset B(x_i, 1)$
with
$\vol(B(x_i^\prime, r)) \: \ge \: (1-\frac{\widehat{\epsilon}}{5}) \: \omega_3
r^3$.
Recalling the sequence $\{r^\prime_i\}$,
by hypothesis there is a subball
$B(x_i^{\prime \prime},
r_i^{\prime \prime}) \subset B(x_i^\prime, r_i^\prime)$ with
$\vol(B(x_i^{\prime \prime}, r_i^{\prime \prime})) \: < \:
(1 - \widehat{\epsilon}) \: \omega_3
(r_i^{\prime \prime})^3$.
Clearly
$B(x_i^\prime,
r) \: \subset \: B(x_i^{\prime \prime}, r + r_i^\prime)$.
From the Bishop-Gromov inequality,
\begin{equation}
\frac{\vol(B(x_i^{\prime \prime}, r + r_i^\prime))}{\vol(B(x_i^{\prime \prime},
r_i^{\prime \prime}))} \: \le \:
\frac{\int_0^{r + r_i^\prime} \sinh^2(s) \:
ds}{\int_0^{r_i^{\prime \prime}} \sinh^2(s) \: ds}.
\end{equation}
Then
\begin{equation}
\vol(B(x_i^\prime,
r)) \: \le \:
\vol(B(x_i^{\prime \prime}, r + r_i^\prime)) \: \le \:
(1-{\widehat{\epsilon}}) \: \omega_3 \:
\frac{(r_i^{\prime \prime})^3}{\int_0^{r_i^{\prime \prime}} \sinh^2(s) \: ds}
\: \int_0^{r + r_i^\prime} \sinh^2(s) \:
ds.
\end{equation}
For large $i$ we obtain
\begin{equation}
\vol(B(x_i^\prime,
r)) \: \le \:
(1-\frac{\widehat{\epsilon}}{2}) \: \omega_3 \: \cdot \: 3
\: \int_0^{r} \sinh^2(s) \:
ds.
\end{equation}
Then if we choose $r$ to be
sufficiently small, we contradict the fact that
$\vol(B(x_i^\prime, r)) \: \ge \: (1-\frac{\widehat{\epsilon}}{5}) \: \omega_3
r^3$ for all $i$.

\begin{remark}[existence for volume]
  \label{rem:kl-locating-small-balls-whose-subballs-almost-euclidean-volume-existence-volume}
  \dcref{locating-small-balls-whose-subballs-almost-euclidean-volume:remark-2}
  \group{Kleiner--Lott Long-Time Curvature}
  \source{kleiner-lott}
  \uses{lem:kl-locating-small-balls-whose-subballs-almost-euclidean-volume-existence-sectional-curvature}

By similar reasoning, for every $L>1$ one may find
$\th_1=\th_1(\widehat{\epsilon},L)$ such that
under the hypotheses of Lemma \ref{Alexlemma2}, there is a subball
$B(y,\th_1)\subset B(x,1)$ which is $L$-biLipschitz to the Euclidean
unit ball.
\end{remark}

\section{II.6.8. Proof of the double sided curvature bound in the
thick part, modulo two propositions}

In this section we explain how Corollary \ref{II.6.8}
follows from Lemma \ref{Alexlemma2} and two other
propositions, which will be proved in subsequent sections.
We first state the other propositions,
which are Propositions \ref{propII.6.3} and \ref{propII.6.4}.

\begin{proposition}[existence for canonical neighborhoods]
  \label{prop:kl-proof-double-sided-curvature-bound-thick-part-modulo-two-propositions-existence-canonical-neighborhoods}
  \dcref{propII.6.3}
  \group{Kleiner--Lott Long-Time Curvature}
  \source{kleiner-lott}
  \uses{def:kl-priori-assumptions-canonical-neighborhood-assumption}
 \label{propII.6.3}
(cf. Proposition II.6.3)
For any
$A < \infty$
one can find positive constants
$\kappa(A)$, $K_1(A)$, $K_2(A)$, $\overline{r}(A)$,
such that for any $t_0 < \infty$
there exists $\overline{\delta}_{A}(t_0)
> 0$, decreasing in $t_0$,
with the following property.  Suppose that
we have a Ricci flow with $(r, \delta)$-cutoff on a time interval
$[0, T]$, where $\delta(t) < \overline{\delta}_A(t)$ on $[t_0/2, t_0]$, with
normalized initial data.  Assume that \\
1. The solution is unscathed on a parabolic ball
$P(x_0, t_0, r_0, -r_0^2)$, with $2r_0^2 < t_0$. \\
2. $|\Rm| \le \frac{1}{3r_0^2}$ on $P(x_0, t_0, r_0, -r_0^2)$. \\
3. $\vol(B(x_0, t_0, r_0)) \: \ge \: A^{-1} r_0^3$.\\

Then\\
(a) The solution is $\kappa$-noncollapsed on scales less than $r_0$
in $B(x_0, t_0, Ar_0)$.\\
(b) Every point $x \in B(x_0, t_0, Ar_0)$ with $R(x, t_0) \ge K_1 r_0^{-2}$
has a canonical neighborhood in the sense of
Definition \ref{canonicalnbhddef}.\\
(c) If $r_0 \: \le \: \overline{r} \sqrt{t_0}$ then $R \: \le \: K_2
r_0^{-2}$ in $B(x_0, t_0, Ar_0)$.
\end{proposition}

Proposition \ref{propII.6.3}(a) is an analog of Theorem
\ref{8.2thm}.

(The reason for the ``$3$'' in the hypothesis
$|\Rm| \le \frac{1}{3r_0^2}$ comes from Remark \ref{8.2fix}.)
Proposition \ref{propII.6.3}(c) is an analog of Theorem
\ref{thmI.12.2}, but the hypotheses are slightly different.
In Proposition \ref{propII.6.3} one assumes a lower bound on the volume of the
time-$t_0$ ball $B(x_0, t_0, r_0)$, while in Theorem \ref{thmI.12.2}
one assumes assumes a lower bound on the volume of
the time-$(t_0 - r_0^2)$ ball $B(x_0, t_0 - r_0^2, r_0)$. In
view of the curvature assumption on $P(x_0, t_0, r_0, - r_0^2)$,
the hypotheses are
essentially
equivalent.


Conclusions (a), (b) and (c) of Proposition \ref{propII.6.3}
are similar to the conclusions of Theorem \ref{8.2thm},
Lemma \ref{12.2lemma} and Theorem \ref{thmI.12.2}, respectively.
Conclusions (a) and (b) of Proposition \ref{propII.6.3}
are also related
to what was proved in Proposition \ref{surgeryflowexists} to construct the
Ricci flow with surgery. The difference
is that the noncollapsing and canonical neighborhood results of Proposition
\ref{surgeryflowexists}
are statements at or below the scale $r(t)$, whereas Proposition \ref{propII.6.3}
is a statement about much larger scales, comparable to $\sqrt{t_0}$.
We note that the parameter $\overline{\delta}_A$ in
Proposition \ref{propII.6.3} is independent of the function $\delta$
used to define the Ricci flow with $(r, \delta)$-cutoff.


In the proof of the next proposition we will apply Lemma \ref{Alexlemma2} with
$\widehat{\epsilon}$ equal to the global parameter $\epsilon$,
so we will write $\theta(w)$ instead of $\theta(\epsilon, w)$.


\begin{proposition}[existence for thick--thin geometry]
  \label{prop:kl-proof-double-sided-curvature-bound-thick-part-modulo-two-propositions-existence-thick-thin-geometry}
  \dcref{propII.6.4}
  \group{Kleiner--Lott Long-Time Curvature}
  \source{kleiner-lott}
 \label{propII.6.4}
(cf. Proposition II.6.4)
There exist $\tau, \overline{r}, C_1 > 0$ and $K < \infty$ with the
following property. Suppose that we have a Ricci flow with $(r, \delta)$-cutoff
on the time interval $[0, t_0]$, with normalized initial data.
  Let $r_0$ satisfy
$2C_1 h_{\max}(t_0) \: \le \: r_0 \: \le \: \overline{r} \sqrt{t_0}$,
where $h_{\max}(t_0)$ is the maximal cutoff radius for surgeries in
$[t_0/2, t_0]$. (If there are no surgeries on $[t_0/2, t_0]$ then
$h_{\max}(t_0) = 0$.)

Assume \\
1. The ball $B(x_0, t_0, r_0)$ has sectional curvatures at least $- r_0^{-2}$
at each point. \\
2. The volume of any subball $B(x, t_0, r) \subset B(x_0, t_0, r_0)$ with
any radius $r > 0$ is at least $(1-\epsilon)$ times the volume of the
Euclidean ball of the same radius.\\

Then the solution is unscathed on
$P(x_0, t_0, r_0/4, - \tau r_0^2)$
and satisfies $R \: < \: K r_0^{-2}$ there.
\end{proposition}

Proposition \ref{propII.6.4} is an analog of Theorem \ref{thmI.12.3}. However,
there is the important difference that in Proposition \ref{propII.6.4}
we have to prove that no surgeries occur within
$P(x_0, t_0, r_0/4, - \tau r_0^2)$.


Assuming the validity of Propositions \ref{propII.6.3} and \ref{propII.6.4},
suppose that the hypotheses of Corollary \ref{II.6.8} are satisfied.
We will allow ourselves to shrink the parameter $\overline{r}$ in order
to apply Hamilton-Ivey pinching when needed.
Put $r_0^\prime \: = \: \theta_0(w) \: r_0$, where $\theta_0(w)$ is
from Lemma \ref{Alexlemma2}.
By Lemma \ref{Alexlemma2}, there is a subball $B(x_0^\prime, t_0, r_0^\prime) \subset
B(x_0, t_0, r_0)$ such that every subball of
$B(x_0^\prime,t_0, r_0^\prime)$ has volume at least
$(1-\epsilon)$ times the volume of the Euclidean ball of the same
radius.  As the sectional curvatures are bounded below by $- r_0^{-2}$ on
$B(x_0, t_0, r_0)$, they are bounded below by $- (r_0^\prime)^{-2}$ on
$B(x^\prime_0, t_0, r_0^\prime)$. By an appropriate choice of the
parameters $\theta(w)$ and $\overline{r}$ of Corollary \ref{II.6.8},
in particular taking $\theta(w) \: \le \: \frac{\theta_0(w)}{2C_1}$, we can
ensure that Proposition \ref{propII.6.4}
applies to $B(x^\prime_0, t_0, r_0^\prime)$.
Then the solution is unscathed on
$P(x_0^\prime, t_0, r_0^\prime/4, - \tau (r_0^\prime)^2)$ and satisfies
$|\Rm| \: \le \: K \: (r_0^\prime)^{-2}$ there, where the lower
bound on $\Rm$ comes from Hamilton-Ivey pinching.
With $\tau$ being the parameter of Proposition \ref{propII.6.4} and
putting
$r_0^{\prime \prime} \: = \: \min(K^{-1/2}, \tau^{1/2}, \frac14) \: r_0^\prime$,
for all $t_0^{\prime \prime} \in
[t_0 - (r_0^{\prime \prime})^2, t_0]$ the
solution is unscathed on $P(x_0^\prime, t_0^{\prime \prime},
r_0^{\prime \prime}, -(r_0^{\prime \prime})^2)$ and satisfies
$|\Rm| \: \le \: (r_0^{\prime \prime})^{-2}$ there.
From the curvature bound $\Rm \: \ge \: - \: (r_0^\prime)^{-2}$ on
$P(x_0^\prime, t_0, r_0^\prime/4, - \tau (r_0^\prime)^2)$
(coming from pinching)
and the fact that $B(x^\prime_0, t_0, r_0^{\prime \prime})$
has almost Euclidean volume, we obtain a bound
$\vol(B(x^\prime_0, t_0^{\prime \prime},
r_0^{\prime \prime})) \: \ge \: \const \:
(r_0^{\prime \prime})^{3}$.
Applying Proposition \ref{propII.6.3} with
$A = \frac{100r_0}{r_0^{\prime \prime}}$ gives
$R \: \le \: K_2 (r_0^{\prime \prime})^{-2}$ on
$B(x_0, t_0^{\prime \prime}, 10 r_0) \subset
B(x_0^\prime, t_0^{\prime \prime}, 100 r_0)$,
for all $t_0^{\prime \prime} \in
[t_0 - (r_0^{\prime \prime})^2, t_0]$.
Writing this as
$R \: \le \: \const r_0^{-2}$,
if we further restrict $\theta(w)$ to be sufficiently
small then we can ensure that $R \: \le \: \const \: \theta^2(w) \: h^{-2}
\: \le \: .01 \: h^{-2}$. As surgeries only occur
at spacetime points $(x,t)$ where
$R(x,t) \sim h(t)^{-2}$,
there are no surgeries on
$\bigcup_{t_0^{\prime \prime} \in
[t_0 - (r_0^{\prime \prime})^2, t_0]} B(x_0, t_0^{\prime \prime}, 10 r_0)$.
Using length distortion estimates,
we can find a parabolic neighborhood
$P(x_0, t_0, r_0/4, - \tau r_0^2) \subset
\bigcup_{t_0^{\prime \prime} \in
[t_0 - (r_0^{\prime \prime})^2, t_0]} B(x_0, t_0^{\prime \prime}, 10 r_0)$
for some fixed $\tau$.
This proves Corollary \ref{II.6.8}.





\section{II.6.3. Canonical neighborhoods and later curvature bounds on bigger balls from
curvature and volume bounds}
We now prove Proposition \ref{propII.6.3}.
We first recall its statement.
\begin{proposition}[existence for canonical neighborhoods]
  \label{prop:kl-canonical-neighborhoods-later-curvature-bounds-bigger-balls-from-curvatu-existence-canonical-neighborhoods}
  \dcref{canonical-neighborhoods-later-curvature-bounds-bigger-balls-from-curvatu:proposition-1}
  \group{Kleiner--Lott Long-Time Curvature}
  \source{kleiner-lott}
  \uses{def:kl-priori-assumptions-canonical-neighborhood-assumption}

(cf. Proposition II.6.3)
For any $A > 0$ one can find positive constants $\kappa(A)$, $K_1(A)$, $K_2(A)$,
$\overline{r}(A)$,
such that for any $t_0 < \infty$
there exists $\overline{\delta}_{A}(t_0)
> 0$, decreasing in $t_0$, with the following property.  Suppose that
we have a Ricci flow with $(r, \delta)$-cutoff on a time interval
$[0, T]$, where $\delta(t) < \overline{\delta}_A(t)$ on $[t_0/2, t_0]$, with
normalized initial data.  Assume that \\
1. The solution is unscathed on a parabolic neighborhood
$P(x_0, t_0, r_0, -r_0^2)$, with $2r_0^2 < t_0$. \\
2.
$|\Rm| \le \frac13 r_0^{-2}$
on $P(x_0, t_0, r_0, -r_0^2)$. \\
3. $\vol(B(x_0, t_0, r_0)) \: \ge \: A^{-1} r_0^3$.\\

Then\\
(a) The solution is $\kappa$-noncollapsed on scales less than $r_0$
in $B(x_0, t_0, Ar_0)$.\\
(b) Every point $x \in B(x_0, t_0, Ar_0)$ with $R(x, t_0) \ge K_1 r_0^{-2}$
has a canonical neighborhood in the sense of
Definition \ref{canonicalnbhddef}.\\
(c) If $r_0 \: \le \: \overline{r} \sqrt{t_0}$ then $R \: \le \: K_2
r_0^{-2}$ in $B(x_0, t_0, Ar_0)$.
\end{proposition}
\par\medskip\noindent\textit{Source argument (not certified as complete).}\ \begingroup\itshape
The proof of part (a) is analogous to the proof of
Theorem \ref{8.2thm}. The proof of part (b) is analogous to
the proof of Lemma \ref{12.2lemma}. The proof of part (c) is analogous
to the proof of Theorem \ref{thmI.12.2}.
We will be brief on the parts of the proof of Proposition \ref{propII.6.3}
that are along the same lines as was done before, and will
concentrate on the differences.


For part (a),
we first remark
 that the $\kappa$-noncollapsing that we want does not follow from
the noncollapsing estimate used in the proof of Proposition
\ref{surgeryflowexists}, which would give a
time-dependent $\kappa$.
So suppose that $(x,t_0) \in B(x_0, t_0, Ar_0)$, $\rho < r_0$,
the parabolic neighborhood $P(x, t_0, \rho, - \rho^2)$ is
unscathed and $|\Rm| \le \rho^{-2}$ there. We want to get
a lower bound on $\rho^{-3} \: \vol(B(x, t_0, \rho))$.
We first reduce the case $\rho \: < \: \frac{r(t_0)}{100}$ to the
case $\rho \:\geq\: \frac{r(t_0)}{100}$.

Suppose that $\rho <
\frac{r(t_0)}{100}$ and let $s$ be the largest
number so that the parabolic neighborhood
$P(x, t_0, s, - s^2)$ is unscathed and $|\Rm| \le s^{-2}$ there.
Clearly $s \ge \rho$.
If $s \: < \: \frac{r(t_0)}{100}$ then we obtain a lower
bound on $\rho^{-3} \: \vol(B(x, t_0, \rho))$
as in Sublemma \ref{r0/100};
a canonical
neighborhood of type (d) with small volume cannot occur,
in view of condition 3 of Proposition \ref{propII.6.3}.

If $s \: \ge \: \frac{r(t_0)}{100}$ and $\frac{r(t_0)}{100} < r_0$
then once we have proved part (a) of the
proposition at scale $\frac{r(t_0)}{100}$, the Bishop-Gromov
inequality will give a lower bound on $\rho^{-3} \: \vol(B(x, t_0, \rho))$
and prove part (a) of the proposition at scale $\rho$.

Suppose that
$s \: \ge \: \frac{r(t_0)}{100} \ge r_0$. Let $D$ be the largest
radius so that $|\Rm| \le r_0^{-2}$ on $B(x,t_0, D)$.
Clearly $D \ge s \ge \rho$.
If $D \ge
(A+1) r_0$ then assumption 3. of the proposition implies that
$\vol(B(x,t_0,(A+1)r_0) \ge A^{-1} r_0^{-3}$. Since $\rho < r_0$,
the Bishop-Gromov inequality implies a lower bound on
$\rho^{-3} \: \vol(B(x, t_0, \rho))$.

Finally, if $D < (A+1)r_0$
then there is some $x_1$ with $d_{t_0}(x,x_1) = D$ so that
$|\Rm(x_1,t_0)| = r_0^{-2} \ge 10000 (r(t_0))^{-2}$. Slightly moving
$x_1$ inward toward $x$, there is a point
$x_2 \in B(x_1, t_0, D-cr(t_0))$ with
$|\Rm(x_2,t_0)| \ge 5000 (r(t_0))^{-2}$, for a universal constant $c$;
cf. case (b) of the proof of Sublemma \ref{r0/100}.
If $\overline{r}$ is small then Hamilton-Ivey pinching implies that
$(x_2, t_0)$ is in a canonical neighborhood, which gives an estimate
\begin{equation}
\vol(B(x, t_0, D)) \ge \vol(B(x_2, t_0, cr(t_0))) \ge
\const (r(t_0))^3 \ge 10^6 \const r_0^3.
\end{equation}
Since $\rho \le D$, the Bishop-Gromov inequality now implies a lower bound on
$\rho^{-3} \: \vol(B(x, t_0, \rho))$, depending on $A$.

This shows that it suffices to consider scales $\rho$ that are at
least $r(t_0)/100$.
To continue,
we recall the idea of the proof of Theorem \ref{8.2thm}.
With the notation of Theorem \ref{8.2thm},
after rescaling so that $r_0 = t_0 = 1$,
we had a point $x \in B(x_0, 1, A)$
around which we wanted to prove noncollapsing.  Defining $l$ using curves
starting at $(x, 1)$, we wanted to find a point
$(y, 1/2) \in
B(x_0, 1/2, 1/2)$ so that
$l(y, 1/2)$ was bounded above by a universal constant.  Given such a point,
we concatenated a minimizing ${\mathcal L}$-geodesic (from $(x, 1)$
to $(y, 1/2)$)
with curves emanating backward in time from
$(y, 1/2)$. Then the bounded geometry near $(y, 1/2)$
allowed us
to estimate from below the reduced volume at a time slightly less than
$1/2$.

We knew that there was some point $y \in M$ so that
$l(y, 1/2) \: \le \: \frac{3}{2}$, but the issue in
Theorem \ref{8.2thm} was to find
a point $(y, 1/2) \in
B(x_0, 1/2, 1/2)$ with $l(y, 1/2)$ bounded above
by a universal constant. The idea was to take the proof that
some point $y \in M$ has
$l(y, 1/2) \: \le \: \frac{3}{2}$ and localize it near $x_0$.
The proof of Theorem \ref{8.2thm} used the function
$h(y, t) = \phi(d(y, t) - A (2t-1)) (\overline{L}(y, 1-t) + 7)$.
Here $\phi$ was a certain nondecreasing function that is one on $(- \infty,
1/20)$ and infinite on $[1/10, \infty)$, and
$\overline{L}(q, \tau) \: = \: 2 \sqrt{\tau} \: L(q, \tau)$.
Clearly $\min h(\cdot, 1) \: \le \: 7$ and $\min h(\cdot, 1/2)$ is
achieved in $B(x_0, 1/2, 1/10)$. The equation
$\square h \: \ge \: - (6+ C(A)) h$ implied that
$\frac{d}{dt} \min h \: \ge \: -(6 + C(A)) \min h$, and so
$(\min h)(t) \: \le \:
7 \: e^{(6+C(A))(1-t)}$.

In the present case,
if one knew that the possible contribution of a barely admissible curve to
$h(y, t)$ was greater than $7 \: e^{(6+C(A))(1-t)} + \epsilon$
then one could still apply the maximum principle to find a point $(y, 1/2)$
with $h(y, 1/2) \: \le \: 7 \: e^{(6+C(A))/2}$.
For this, it
suffices to know that the possible contribution of a barely admissible curve
to $\overline{L}(q, \tau)$ can be bounded below by a sufficiently large
number.  However, Lemma \ref{lplusbig}  only says that we can make the
contribution of a barely admissible curve to $L$ large
(using the lower scalar curvature bound to pass from
${\mathcal L}_+$ to ${\mathcal L}$).
Because of the factor $2 \sqrt{\tau}$ in the definition of
$\overline{L}(q, \tau)$, we cannot necessarily say that its contribution
to $\overline{L}(q, \tau)$ is large. To salvage the argument,
the idea is to redefine $h$ and
redo the proof of Theorem \ref{8.2thm} in order to get an extra factor
of $\sqrt{\tau}$ in $\min h$.

(The use of Lemma \ref{lplusbig} is similar
to what was done in the proof of Proposition \ref{surgeryflowexists}.
However, there is a difference in scales.  In Proposition
\ref{surgeryflowexists} one was working at a microscopic scale
in order to construct the Ricci flow with surgery. The function
$\overline{\delta}(t)$ in Proposition \ref{surgeryflowexists} was relevant
to this scale. In the present
case we are working at the macroscopic scale $r_0 \sim \sqrt{t_0}$ in
order to analyze the long-time behavior of the Ricci flow with
surgery.  The function $\delta_A(t)$ of Proposition
\ref{propII.6.3} is relevant to this scale.  Thus we will end up further
reducing the surgery function $\overline{\delta}(t)$ of Proposition \ref{surgeryflowexists}
in order to be able to apply Proposition \ref{propII.6.3}.)


By assumption,
$|\Rm| \: \le \: 1$ at $t = 0$.
From Lemma \ref{Rev}, $R \: \ge \: - \: \frac32 \: \frac{1}{t+1/4}$.
Then for $t \in [t_0 - r_0^2/2, t_0]$,
\begin{equation}
R \: r_0^2 \: \ge \: - \: \frac32 \: \frac{r_0^2}{t_0 - r_0^2/2}
\: \ge \: - \: \frac32 \: \frac{r_0^2}{2r_0^2 - r_0^2/2} \: = \: -1.
\end{equation}
After rescaling so that $r_0 = 1$, the time interval
$[t_0 - r_0^2, t_0]$ is shifted to $[0, 1]$.
Then for $t \in [\frac12, 1]$, we certainly have
$R \: \ge \: -3$.

From this, if $0 < \tau \le \frac12$ then
$\overline{L}(y, \tau) \: \ge \: - 6 \: \sqrt{\tau} \:
\int_0^\tau \sqrt{v} \: dv \: = \: - 4 \tau^2$, so
$\hat{L}(y, \tau) \: \equiv \: \overline{L}(y, \tau) \: + \: 2 \sqrt{\tau}
\: > \: 0$.

Putting
\begin{equation}
h(y, \tau) \: = \: \phi(d_t(x_0, y) - A (2t-1)) \: \hat{L}(y, \tau)
\end{equation}
and using the fact that $\frac{d}{dt} \sqrt{\tau} \: = \:
- \frac{d}{d\tau} \sqrt{\tau} \: = \: - \: \frac{1}{2 \sqrt{\tau}}$,
the computations of the proof of Theorem \ref{8.2thm} give
\begin{align}
\square h \: & \ge \: - \: \left( \overline{L} + 2 \sqrt{\tau} \right)
\: C(A) \: \phi \:
\: - \: 6 \: \phi \: - \: \frac{1}{\sqrt{\tau}} \: \phi \\
& = \: - \: C(A) h \: - \: \left( 6 \: + \: \frac{1}{\sqrt{\tau}}
\right) \: \phi.
\notag
\end{align}
Then if  $h_0(\tau) \: = \: \min h(\cdot, \tau)$, we have
\begin{align}
\frac{d}{d\tau} \left( \log \left( \frac{h_0(\tau)}{\sqrt{\tau}} \right)
\right) \: & = \: h_0^{-1} \: \frac{dh_0}{d\tau} \: - \: \frac{1}{2\tau}
\: \le \: C(A) \: + \: \left( 6 \: + \: \frac{1}{\sqrt{\tau}} \right)
\: \frac{\phi}{h_0} \: - \: \frac{1}{2\tau} \\
& = \: C(A) \: + \: \left( 6 \: + \: \frac{1}{\sqrt{\tau}} \right)
\: \frac{1}{\overline{L} + 2 \sqrt{\tau}} \: - \: \frac{1}{2\tau} \notag \\
& = \: C(A) \: + \:
\frac{6 \sqrt{\tau} \: + \: 1}{\sqrt{\tau} \overline{L} + 2 {\tau}} \:
- \: \frac{1}{2\tau}. \notag
\end{align}
As $\overline{L} \: \ge \: - \: 4 \tau^2$,
\begin{equation}
\frac{d}{d\tau} \left( \log \left( \frac{h_0(\tau)}{\sqrt{\tau}} \right)
\right) \:  \le \:
C(A) \: + \:
\frac{6 \sqrt{\tau} \: + \: 1}{2 {\tau} - 4 \tau^2 \sqrt{\tau}} \:
- \: \frac{1}{2\tau} \: \le \: C(A) \: + \: \frac{50}{\sqrt{\tau}}.
\end{equation}

As $\tau \rightarrow 0$, the Euclidean space computation gives
$\overline{L}(q, \tau) \sim |q|^2$, so
$\lim_{\tau \rightarrow 0} \frac{h_0(\tau)}{\sqrt{\tau}} \: = \: 2$.
Then
\begin{equation}
{h_0(\tau)} \: \le \:
2 \: \sqrt{\tau} \: \exp(C(A) \tau \: + \: 100 \sqrt{\tau}).
\end{equation}
This estimate has the desired extra factor of $\sqrt{\tau}$.

It now suffices to show that for a barely admissible curve $\gamma$
that hits a surgery region at time $1-\tau$,
\begin{equation}
\int_0^\tau \sqrt{v} \: \left( R(\gamma(1-v), v) +
|\dot{\gamma}(v)|^2
\right) \: dv \: \ge \:
\exp(C(A) \tau \: + \: 100 \sqrt{\tau}) \: + \:
\epsilon,
\end{equation}
where $0 < \tau \le \frac12$. Choosing $\overline{\delta}_A(t_0)$ small
enough, this follows from Lemma \ref{lplusbig}
along with the lower scalar curvature bound.
Then we can apply the maximum principle and follow the proof of
Theorem \ref{8.2thm}. In our case the bounded geometry near $(y, 1/2) \in
B(x_0, 1/2, 1/2)$ comes from assumptions 2 and 3 of the
Proposition. The function $\overline{\delta}_A$ is now determined.
After reintroducing the scale $r_0$,
this proves part (a) of the proposition.

The proof of part (b)
is similar to the proofs of
Lemma \ref{12.2lemma}
and Proposition \ref{surgeryflowexists}.
Suppose that for some $A > 0$ the claim is not true. Then there
is a sequence of Ricci flows $\M^\al$ which together provide a
counterexample.  In particular, some point $(x^\al, t^\al) \in
B(x_0^\al, t_0^\al, A r_0^\al)$ has $R(x^\al, t^\al) \ge K_1^\al
(r_0^\al)^{-2}$ but does not have a canonical neighborhood, where
$K_1^\al \rightarrow \infty$ as $\al \rightarrow \infty$.
Because of the canonical neighborhood
assumption, we must have $K_1^\al \: (r_0^{\al})^{-2} \: \le \:
r(t_0^{\al})^{-2}$.
Then $2 K_1^{\al} \: \le \: K_1^{\al} \: t_0^{\al} (r_0^{\al})^{-2}
\: \le \: t_0^{\al} \: r(t_0^{\al})^{-2}$. Since $K_1^{\al}
\rightarrow \infty$ and
the function $t \rightarrow t r(t)^{-2}$ is bounded on any finite
$t$-interval, it follows that $t_0^\al \rightarrow \infty$.
Applying point selection to each
$\M^\al$ and removing the superscripts, there are points
$\overline{x} \in B(x_0, \overline{t}, 2Ar_0)$ with
$\overline{t} \in [t_0 - r_0^2/2, t_0]$ such that
$\overline{Q} \equiv R(\overline{x}, \overline{t}) \: \ge \:
K_1 r_0^{-2}$ and $(\overline{x}, \overline{t})$ does not have a
canonical neighborhood, but each point $(x, t) \in \overline{P}$
with $R(x, t) \: \ge \: 4 \overline{Q}$ does have a canonical
neighborhood, where $\overline{P} \: = \:
\{(x, t) \: : \: d_t(x_0, x) \: \le \: d_{\overline{t}}(x_0,
\overline{x}) + K_1^{1/2} \overline{Q}^{-1/2}, t \in
[\overline{t} - \frac14 K_1 \overline{Q}^{-1}, \overline{t}]\}$.
From (a), we have noncollapsing in $\overline{P}$. Rescaling by
$\overline{Q}^{-1}$, we have bounded curvature at bounded distances
from $\overline{x}$; see Lemma \ref{claim2II.4.2}.
Then we can extract a pointed limit $X_\infty$, which we think of as
a time zero slice, that
will have nonnegative sectional curvature.
(The required pinching for the last statement comes from the assumption that
$2r_0^2 < t_0$, along with the fact that $K_1^\al \rightarrow \infty$.)
The fact that
points $(x, t) \in \overline{P}$
with $R(x, t) \: \ge \: 4 \overline{Q}$ have a canonical
neighborhood implies that regions of large scalar curvature in $X_\infty$
have canonical neighborhoods, from which one can
deduce as in Section \ref{pf11.7} that
the sectional curvatures of $X_\infty$ are globally bounded above
by some $Q_0 > 0$.
Then for each $\overline{A}$,
Lemmas \ref{distancedist} and \ref{claim1II.4.2}
imply that for large $\al$, the parabolic neighborhood
$P(\overline{x}, \overline{t}, \overline{A} \: \overline{Q}^{-1/2},
- \epsilon \eta^{-1} Q_0^{-1} \overline{Q}^{-1})$ is
contained in $\overline{P}$.
(Here $\epsilon$ is a small parameter, which we absorb in the
global parameter $\epsilon$.)
In applying Lemma \ref{distancedist} we use the curvature
bound near $x_0$ coming from the hypothesis of the proposition along with
the curvature bound near $\overline{x}$ just derived; cf. the proof
of Lemma \ref{12.2lemma}.
In addition, we claim that
$P(\overline{x}, \overline{t}, \overline{A} \: \overline{Q}^{-1/2},
- \epsilon \eta^{-1} Q_0^{-1} \overline{Q}^{-1})$ is
unscathed. This is proved as in Section \ref{surgeryflow}.
Recall that the idea is to show that a surgery in
$P(\overline{x}, \overline{t}, \overline{A} \: \overline{Q}^{-1/2},
- \epsilon \eta^{-1} Q_0^{-1} \overline{Q}^{-1})$ implies that
$(\overline{x}, \overline{t})$ lies in a canonical neighborhood,
which contradicts our assumption. In the argument
we use the fact that $t_0^\al \rightarrow \infty$
implies $\delta(t_0^\al) \rightarrow 0$ in order to rule out
surgeries; this is the replacement for the condition
$\overline{\delta}^{\alpha} \rightarrow 0$ that was used in
Section \ref{surgeryflow}.


We extend $X_\infty$ to the maximal backward-time limit and
obtain an ancient $\kappa$-solution, which contradicts the assumption
that the points $(\overline{x}, \overline{t})$ did not have
canonical neighborhoods. This proves part (b) of the proposition.

To prove part (c),
we can rescale $t_0$ to $1$ and then apply
Lemma \ref{claim2II.4.2}; see the end of the proof of Theorem
\ref{thmI.12.2}.
The $\Phi$-pinching that we use comes from the Hamilton-Ivey
estimate of (\ref{estimate}).
We recall that in the proof of Theorem \ref{thmI.12.2}
we need to get nonnegative curvature in the region $W$ near the
blowup point;
this comes from the fact that $\overline{r}^\al \rightarrow 0$ in the
contradiction argument, along with the Hamilton-Ivey pinching.
\endgroup\par\medskip

This proves the proposition.
In what follows, we will want to apply it freely for arbitrary
$A$, provided that $t_0$ is large enough. To do so,
we reduce the function $\overline{\delta}$ used to define the
Ricci flow with $(r, \delta)$-cutoff, if necessary, in order to ensure
that $\overline{\delta}(t) \: \le \: \overline{\delta}_{2t}(2t)$.
Here $\overline{\delta}_{2t}(2t)$ is the quantity
$\overline{\delta}_{A}(2t)$ from Proposition \ref{propII.6.3}
evaluated at $A = 2t$.



\section{II.6.4. Earlier scalar curvature bounds on smaller balls from
lower curvature bounds and volume bounds, in the presence of possible
surgeries}

In this section we prove Proposition \ref{propII.6.4}.
We first recall its statement.
\begin{proposition}[existence for scalar curvature]
  \label{prop:kl-earlier-scalar-curvature-bounds-smaller-balls-from-lower-curvature-bound-existence-scalar-curvature}
  \dcref{earlier-scalar-curvature-bounds-smaller-balls-from-lower-curvature-bound:proposition-5}
  \group{Kleiner--Lott Long-Time Curvature}
  \source{kleiner-lott}

There exist $\tau, \overline{r}, C_1 > 0$ and $K < \infty$
with the
following property. Suppose that we have a Ricci flow with $(r, \delta)$-cutoff
on the time interval $[0, t_0]$, with normalized initial data.  Let
$r_0$ satisfy
$2C_1 h_{\max}(t_0) \: \le \: r_0 \: \le \: \overline{r} \sqrt{t_0}$,
where $h_{\max}(t_0)$ is the maximal cutoff radius for surgeries in
$[t_0/2, t_0]$. (If there are no surgeries on $[t_0/2, t_0]$ then we put
$h_{\max}(t_0) = 0$.)
Assume \\
1. The ball $B(x_0, t_0, r_0)$ has sectional curvatures at least $- r_0^{-2}$
at each point. \\
2. The volume of any subball $B(x, t_0, r) \subset B(x_0, t_0, r_0)$ with
any radius $r > 0$ is at least $(1-\epsilon)$ times the volume of the
Euclidean ball of the same radius.\\

Then the solution is unscathed on
$P(x_0, t_0, r_0/4, - \tau r_0^2)$ and satisfies $R < K r_0^{-2}$
there.
\end{proposition}

Proposition \ref{propII.6.4}
is an analog of Theorem \ref{thmI.12.3}.  However, the proof of Proposition
\ref{propII.6.4} is more complicated, due to the need to deal with possible
surgeries.

\begin{proof}
The constants $C_1$, $K$ and $\tau$ are fixed numbers, but the
requirements on them will
be specified during the proof. The number $\overline{r}$ will emerge
from the proof, via a contradiction argument.

The first step is to prove an analog of the proposition in which
the parabolic neighborhood of the conclusion is replaced by
a time-dependent family of metric balls.

\begin{lemma}[existence for scalar curvature]
  \label{lem:kl-earlier-scalar-curvature-bounds-smaller-balls-from-lower-curvature-bound-existence-scalar-curvature-6}
  \dcref{extra}
  \group{Kleiner--Lott Long-Time Curvature}
  \source{kleiner-lott}
 \label{extra}
There exists $\tau^\prime > 0$
with the
following property. Suppose that we have a Ricci flow with $(r, \delta)$-cutoff
on the time interval $[0, t_0]$, with normalized initial data.  Let
$r_0$ satisfy
$2C_1 h_{\max}(t_0) \: \le \: r_0 \: \le \: \overline{r} \sqrt{t_0}$,
where $h_{\max}(t_0)$ is the maximal cutoff radius for surgeries in
$[t_0/2, t_0]$. (If there are no surgeries on $[t_0/2, t_0]$ then we put
$h_{\max}(t_0) = 0$.)
Assume
\begin{enumerate}
\item The ball $B(x_0, t_0, r_0)$ has sectional curvatures at least $- r_0^{-2}$
at each point.
\item The volume of any subball $B(x, t_0, r) \subset B(x_0, t_0, r_0)$ with
any radius $r > 0$ is at least $(1-\epsilon)$ times the volume of the
Euclidean ball of the same radius.
\end{enumerate}

Then there is an unscathed time-dependent family
of metric balls
$\bigcup_{t \in [t_0-\tau^\prime r_0^2, t_0]} B(x_0, t, r_0/2)$
on which the supremum of $R$ is less than $K r_0^{-2}$.
\end{lemma}
\begin{proof}
  \uses{lem:kl-curvature-bounds-from-priori-assumptions-estimates-curvature,lem:kl-earlier-scalar-curvature-bounds-smaller-balls-from-lower-curvature-bound-estimates-scalar-curvature,lem:kl-earlier-scalar-curvature-bounds-smaller-balls-from-lower-curvature-bound-estimates-scalar-curvature-4,lem:kl-length-distortion-estimates-distance-distortion-lemma,lem:kl-locating-small-balls-whose-subballs-almost-euclidean-volume-existence-sectional-curvature,prop:kl-proof-double-sided-curvature-bound-thick-part-modulo-two-propositions-existence-canonical-neighborhoods}
Again, $\tau^\prime$ is a fixed number but the
requirements on it will
be specified during the proof.
The idea of the proof of the lemma
is to put oneself in a setting in
which one can apply Lemma \ref{II.6.5(a)}.

To argue by contradiction,
suppose that we have a sequence of
Ricci flows with $(r, \delta)$-cutoff $\M^\alpha$ having
balls $B(x_0^\al, t_0^\al, r_0^\al)$ that
satisfy the assumptions of the lemma, with $\overline{r}^\alpha
\rightarrow 0$, but do not satisfy the conclusion.

\begin{lemma}[scalar curvature lemma]
  \label{lem:kl-earlier-scalar-curvature-bounds-smaller-balls-from-lower-curvature-bound-scalar-curvature-lemma}
  \dcref{earlier-scalar-curvature-bounds-smaller-balls-from-lower-curvature-bound:sublemma-7}
  \group{Kleiner--Lott Long-Time Curvature}
  \source{kleiner-lott}

$r_0^\al > r(t_0^\al)$ for all but a finite number of $\al$.
\end{lemma}
\begin{proof}
  \uses{lem:kl-curvature-bounds-from-priori-assumptions-estimates-curvature,lem:kl-earlier-scalar-curvature-bounds-smaller-balls-from-lower-curvature-bound-existence-scalar-curvature-6}
If not then $r_0^\al \le r(t_0^\al)$ for
infinitely many
$\al$. After passing to a subsequence, we can assume
that $r_0^\al \le r(t_0^\al)$ for all $\al$. We  claim that
$R \: \le \: (r_0^\al)^{-2}$ on $B(x_0^\al, t_0^\al, \frac{3r_0^\al}{4})$.
If not then $R(x^\al, t_0^\al) \: > \: (r_0^\al)^{-2}$ for some
$x^\al \in B(x_0^\al, t_0^\al, \frac{3r_0^\al}{4})$, and so
$R(x^\al, t_0^\al) \: > \: r(t_0^\al)^{-2}$.
This implies that
$(x^\al, t_0^\al)$ is in a canonical neighborhood, which
contradicts the almost-Euclidean-volume assumption on subballs of
$B(x_0^\al, t_0^\al, r_0^\al)$.

Thus $R \: \le \: (r_0^\al)^{-2}$ on
$B(x_0^\al, t_0^\al, \frac{3r_0^\al}{4})$.
Lemma \ref{claim1II.4.2}
implies that
$R \: \le \: 16 \: (r_0^\al)^{-2}$ on
$P(x^\al_0, t^\al_0,  \frac{3r_0^\al}{4},
- \: \frac{1}{16} \eta^{-1} (r_0^\al)^2)$.
Furthermore, if \\
(*.1) $C_1 \: \ge \: 100$ \\
then $R \: \le \: \frac{1}{2 h^2}$ on $P(x_0^\al, t_0^\al, \frac{3r_0^\al}{4},
- \frac{1}{16} \eta^{-1} (r_0^\al)^2)$.
As surgeries only occur when $R \: \ge \: h^{-2}$,
there cannot be any surgeries in the region.

If $t \in [t_0^\al - \frac{1}{16} \eta^{-1} (r_0^\al)^2, t_0^\al]$ then
\begin{equation}
\frac{(r_0^\al)^2}{t} \: \le \:
\frac{(r_0^\al)^2}{t_0^\al - \frac{1}{16} \eta^{-1} (r_0^\al)^2} \: = \:
\frac{(r_0^\al)^2/t_0^\al}{1 - \frac{1}{16} \eta^{-1} (r_0^\al)^2/t_0^\al},
\end{equation}
The fact that $\lim_{\alpha \rightarrow \infty} \overline{r}^\al = 0$
implies that the Hamilton-Ivey pinching estimate improves with $\al$.
In particular, for large $\al$, we have $|\Rm| \le
16 \: (r_0^\al)^{-2}$ on
$P(x^\al_0, t^\al_0,  \frac{3r_0^\al}{4},
- \: \frac{1}{16} \eta^{-1} (r_0)^\al)^2)$.
By the distance distortion estimates of Section \ref{secI.8.3},
if $\tau \le \min \left( \frac{1}{16} \eta^{-1},
\frac{1}{32} \log(\frac32) \right)$ then
\begin{equation}
\bigcup_{t \in [t_0^\al -\tau (r^\al_0)^2, t^\al_0]}
B(x^\al_0, t, r^\al_0/2) \subset P(x^\al_0, t^\al_0,  \frac{3r_0^\al}{4},
- \: \frac{1}{16} \eta^{-1} (r_0)^\al)^2).
\end{equation}
Hence if we have \\
(*.2) $K \: \ge \: 200$ and \\
(*.3) $\tau \le \min \left( \frac{1}{16} \eta^{-1},
\frac{1}{32} \log(\frac32) \right)$\\
then the Ricci flows $\M^\al$ and the balls
$B(x_0^\al, t_0^\al, r_0^\al)$ do satisfy the conclusion of
Lemma \ref{extra}, contrary to assumption. This proves the
sublemma.
\end{proof}

Continuing with the proof of Lemma \ref{extra}, we can assume that
for all $\al$, we have $r_0^\al \: > \: r(t_0^\al)$.
For the flow $\M^\alpha$,
let $t_0^\alpha$ be the first time so that for some radius $r_0^\al$,
the ball $B(x_0^\al, t_0^\al, r_0^\al)$ satisfies the hypotheses of the
lemma, but
$\bigcup_{t \in [t_0^\al -\tau^\prime (r^\al_0)^2, t^\al_0]}
B(x^\al_0, t, r^\al_0/2)$ fails to be unscathed or the
supremum of $R$ on that region is at least
$K (r_0^\al)^{-2}$.
Let $r_0^\al$ be the smallest such radius;
such a radius exists since
$r_0^\alpha \: > \: r(t_0^\alpha)$.

Let
$\widehat{\tau}\geq 0$
be the supremum of the numbers $\widetilde{\tau}$
with the property that for large $\alpha$,
$\bigcup_{t \in [t_0^\al -\widetilde{\tau} (r^\al_0)^2, t^\al_0]}
B(x^\al_0, t, r^\al_0)$ is unscathed and
$R \ge - (r_0^\al)^{-2}$ there.

\begin{lemma}[estimates for scalar curvature]
  \label{lem:kl-earlier-scalar-curvature-bounds-smaller-balls-from-lower-curvature-bound-estimates-scalar-curvature-8}
  \dcref{earlier-scalar-curvature-bounds-smaller-balls-from-lower-curvature-bound:sublemma-8}
  \group{Kleiner--Lott Long-Time Curvature}
  \source{kleiner-lott}
  \uses{lem:kl-earlier-scalar-curvature-bounds-smaller-balls-from-lower-curvature-bound-estimates-scalar-curvature}

$\widehat{\tau}$ is bounded below by the parameter
$\tau_0$ of Lemma \ref{I.11.6(a)}, where we take $w = 1-\epsilon$ in
Lemma \ref{I.11.6(a)}.
\end{lemma}
\begin{proof}
  \uses{lem:kl-earlier-scalar-curvature-bounds-smaller-balls-from-lower-curvature-bound-estimates-scalar-curvature,lem:kl-length-distortion-estimates-distance-distortion-lemma,lem:kl-locating-small-balls-whose-subballs-almost-euclidean-volume-existence-sectional-curvature,prop:kl-proof-double-sided-curvature-bound-thick-part-modulo-two-propositions-existence-canonical-neighborhoods}
Suppose that $\widehat{\tau} < \tau_0$.
Put $\widehat{t}^\al \: = \: t_0^\al \: - \: (1 - \epsilon^\prime) \:
\widehat{\tau} \: (r_0^\al)^2$, where $\epsilon^\prime$ will eventually be
taken to be a small positive number. Applying Lemma \ref{I.11.6(a)} to the
solution on $\bigcup_{t \in [t_0^\al - \: (1 - \epsilon^\prime) \:
\widehat{\tau} \: (r_0^\al)^2, t_0^\al]} B(x_0^\al, t, r_0^\al)$,
the volume
of $B(x_0^\al, \widehat{t}^\al, r_0^\al/2)$ is at least $\frac{1}{10}$ of the
volume of the Euclidean ball of the same radius.
From Lemma \ref{Alexlemma2}, there is a subball
$B(x_1^\al, \widehat{t}^\al, r^\al)
\subset
B(x_0^\al, \widehat{t}^\al, r_0^\al/2)$ of radius $r^\al \: = \:
\theta_0(1/10) \: r_0^\al/2$
with the property that all of its subballs have volume at least
$(1-\epsilon)$ times the volume of the Euclidean ball of the same radius.
The sectional curvature on $B(x_1^\al, \widehat{t}^\al, r^\al)$
is bounded below by
$- \: (r^\al)^{-2}$.
Since $B(x_1^\al, \widehat{t}^\al, r^\al)$ has an
earlier time or smaller radius
than $B(x_0^\al, t_0^\al, r_0^\al)$, it
follows that $\bigcup_{t \in [\widehat{t}^\al - \tau^\prime (r^\al)^2,
\widehat{t}^\al]} B(x_1^\al, t, r^\al/2)$
is unscathed and has $R \: < \: K (r^\al)^{-2}$ there.
For $t \in [\widehat{t}^\al - \tau^\prime (r^\al)^2,
\widehat{t}^\al]$, we have
\begin{equation}
\frac{(r^\al)^2}{t} \: \le \:
\frac{(r^\al)^2}{t_0^\al - (1-\epsilon) \widehat{\tau} (r_0^\al)^2 -
\tau^\prime (r^\al)^2} =
\frac{(r^\al)^2}{(r_0^\al)^2}
\frac{(r_0^\al)^2}{t_0^\al}
\frac{1}{1 - (1-\epsilon) \widehat{\tau} \frac{(r_0^\al)^2}{t_0^\al}
- \tau^\prime \frac{(r^\al)^2}{(r_0^\al)^2} \frac{(r_0^\al)^2}{t_0^\al}
}.
\end{equation}
The fact that $\lim_{\al \rightarrow \infty} \overline{r}^\al = 0$
implies that the Hamilton-Ivey pinching improves with $\al$. In
particular, for large $\al$, $|\Rm| \: < \: K (r^\al)^{-2}$ on
$\bigcup_{t \in [\widehat{t}^\al - \tau^\prime (r^\al)^2,
\widehat{t}^\al]} B(x_1^\al, t, r^\al/2)$. Putting
$\widetilde{r}_0^\al \: = \: K^{-1/2} \: r^\al$,  if \\
(*.4) $K^{-1} \: \le \: \frac{1}{10} \: \tau^\prime$ \\
then $|\Rm| \: \le \: (\widetilde{r}_0^\al)^{-2}$ on
$\bigcup_{t \in [\widetilde{t}^\al - (\widetilde{r}_0^\al)^2,
\widetilde{t}^\al]} B(x_1^\al, t, \widetilde{r}_0^\al)$ for
$\widetilde{t}^\al \in [\widehat{t}^\al - \frac12 \tau^\prime (r^\al)^2,
\widehat{t}^\al]$.
Taking $A \: = \: 100 r_0^\al/\widetilde{r}_0^\al$, Proposition
\ref{propII.6.3}(c)
now implies that
for large $\al$, we have
$R \: \le \: K_2(A) \: (\widetilde{r}^\al_0)^{-2}$ on
$B(x_1^\al, \widetilde{t}^\al, 100 r_0^\al)$.
Provided that \\
(*.5) $K_2(A) \: K
\: (\theta_0(1/10))^{-2} \: \le \: \frac{1}{1000} \: C_1^{2}$ \\
we will have $K_2(A) \: (\widetilde{r}^\al_0)^{-2} \: < \: \frac12 \: h^{-2}$
and so such balls will avoid the surgery regions.
Then the length distortion estimates
of Lemma \ref{distancedist} imply that there is some explicit
$c \in (0, \tau^\prime/2)$ so that
$\bigcup_{t \in [\widehat{t}^\al - c (r_0^\al)^2, \widehat{t}^\al]}
B(x_0^\al, t, r_0^\al) \subset
\bigcup_{t \in [\widehat{t}^\al - c (r_0^\al)^2, \widehat{t}^\al]}
B(x_1^\al, t, 100 r_0^\al)$,
and hence
$R \: \le \: K_2(A) \: (\widetilde{r}_0^\al)^{-2}$ on
$\bigcup_{t \in [\widehat{t}^\al - c (r_0^\al)^2, \widehat{t}^\al]}
B(x_0^\al, t, r_0^\al)$.
Hamilton-Ivey pinching now
implies that for large $\al$,
$\Rm \: \ge \: - \: (r_0^\al)^{-2}$ on
$\bigcup_{t \in [\widehat{t}^\al - c (r_0^\al)^2, \widehat{t}^\al]}
B(x_0^\al, t, r_0^\al)$.
As $c$ can be taken independent
of the small number $\epsilon^\prime$, taking $\epsilon^\prime \rightarrow 0$
we contradict the maximality of
$\widehat{\tau}$.
\end{proof}

Continuing with the proof of the lemma,
we can now apply Lemma \ref{I.11.6(a)} to obtain
$R \: \le \: K_0^\prime \: \tau_0^{-1}
\: (r_0^\al)^{-2}$ on
$\bigcup_{t \in [t_0^\al - \: \tau_0 (r_0^\al)^2/2, t_0^\al]}
B(x_0^\al, t, r_0^\al/2)$. This will
give a contradiction provided that \\
(*.6) $K_0^\prime \: \tau_0^{-1} \: < \: K/2$, \\
(*.7) $\tau^\prime \: < \: \tau_0/2$ and \\
(*.8) $K_0^\prime \: \tau_0^{-1} \: \le \: \frac{1}{1000} C_1^2$, \\
where the last condition ensures that
$\bigcup_{t \in [t_0^\al-\tau^\prime (r_0^\al)^2, t_0^\al]}
B(x_0^\al, t, r_0^\al/2)$
does not hit the surgery regions.

We first choose $\tau^\prime$ to satisfy (*.7).
We then choose $K$ to satisfy (*.4) and (*.6).
Finally, we choose $C_1$ to
satisfy (*.5) and (*.8).
This proves the lemma.
\end{proof}

We now finish the proof of Proposition \ref{propII.6.4}.
By the distance distortion estimates of Section \ref{secI.8.3}, if\\
(*.9) $\tau \le \min \left( \tau^\prime, \frac{\log(2)}{2K} \right)$\\
then $P(x_0, t_0, r_0/4, -\tau r_0^2) \subset
\bigcup_{t_0 - \tau^\prime r_0^2} B(x_0, t, r_0/2)$.
In addition to the conditions on the parameters coming from
Lemma \ref{extra}, we can assume that $C_1$ satisfies (*.1),
$K$ satisfies (*.2), and
$\tau$ satisfies (*.3) and (*.9). The proposition follows.
\end{proof}

\begin{remark}[existence for scalar curvature]
  \label{rem:kl-earlier-scalar-curvature-bounds-smaller-balls-from-lower-curvature-bound-existence-scalar-curvature}
  \dcref{earlier-scalar-curvature-bounds-smaller-balls-from-lower-curvature-bound:remark-9}
  \group{Kleiner--Lott Long-Time Curvature}
  \source{kleiner-lott}
  \uses{prop:kl-proof-double-sided-curvature-bound-thick-part-modulo-two-propositions-existence-thick-thin-geometry}

The proof of Proposition \ref{propII.6.4} outlined in
\cite[Pf. of II.6.4]{Perelman2}
uses parabolic balls throughout.  Richard Bamler pointed out
to us that there is an apparent problem with this approach,
due to the issue of length distortion.
He also indicated that the problem can be circumvented
by using time-dependent families of metric balls.
\end{remark}

\begin{remark}[structural remark on scalar curvature]
  \label{rem:kl-earlier-scalar-curvature-bounds-smaller-balls-from-lower-curvature-bound-structural-remark-scalar-curvature}
  \dcref{hmax}
  \group{Kleiner--Lott Long-Time Curvature}
  \source{kleiner-lott}
  \uses{cor:kl-double-sided-curvature-bound-thick-part-existence-ricci-flow-surgery}
 \label{hmax}
In subsequent sections we will want to know that for any
$w > 0$, with the notation
of Corollary \ref{II.6.8}, we have
$\theta^{-1}(w) \: h_{\max}(t_0) \: \le \: r(t_0)$ if $t_0$ is
sufficiently large (as a function of $w$).
We can always achieve this by lowering the function $\overline{\delta}(\cdot)$
used to define the Ricci flow with $(r, \delta)$-cutoff so that
$\lim_{t_0 \ra \infty} \frac{h_{\max}(t_0)}{r(t_0)} \: = \: 0$.
We will assume hereafter that this is the case.
\end{remark}

\section{II.7.1. Noncollapsed pointed limits are hyperbolic} \label{II.7.1}

In this section we start the analysis of the long-time decomposition
into hyperbolic and graph manifold pieces. In the section,
$\M$ will denote a Ricci flow with $(r,\delta)$-cutoff whose initial time slice
$(\M_0, g(0))$ is compact and has normalized metric.

From Lemma \ref{pscempty}, if
$g(0)$ has positive scalar curvature then the solution goes extinct
in a finite time.  From Lemma \ref{finitetime}, these manifolds are understood
topologically. If $g(0)$ has nonnegative scalar curvature then
either it acquires positive scalar curvature or it is flat, so again
the topological type is understood.
Hereafter we assume that the flow
does not become extinct, and
that $R_{\min}<0$ for all $t$.

\begin{lemma}[monotonicity for noncollapsing]
  \label{lem:kl-noncollapsed-pointed-limits-hyperbolic-monotonicity-noncollapsing}
  \dcref{noncollapsed-pointed-limits-hyperbolic:lemma-1}
  \group{Kleiner--Lott Long-Time Topology}
  \source{kleiner-lott}

$V(t) \: \left( t + \frac14 \right)^{- \: \frac32}$ is nonincreasing in $t$.
\end{lemma}
\begin{proof}
  \uses{lem:kl-establishing-noncollapsing-presence-surgery-vanishing-noncollapsing}
Suppose first that the flow is nonsingular. In the case the lemma follows
from Lemma \ref{Rev} and the equation
\begin{equation} \label{dVdt}
\frac{dV}{dt} \: = \: - \: \int_M R \: dV \: \le \: - \: R_{min} \: V.
\end{equation}

If there are surgeries then it only has the effect of causing further
decrease in $V$.
\end{proof}

\begin{definition}[noncollapsing definition]
  \label{def:kl-noncollapsed-pointed-limits-hyperbolic-noncollapsing-definition}
  \dcref{VR}
  \group{Kleiner--Lott Long-Time Topology}
  \source{kleiner-lott}
 \label{VR}
Put $\overline{V} \: = \: \lim_{t \rightarrow \infty} V(t) \: \left( t + \frac14 \right)^{- \: \frac32}$
and $\hat{R}(t) \: = \: R_{min}(t) \: V(t)^{\frac23}$.
\end{definition}

\begin{lemma}[noncollapsing lemma]
  \label{lem:kl-noncollapsed-pointed-limits-hyperbolic-noncollapsing-lemma}
  \dcref{Rhatev}
  \group{Kleiner--Lott Long-Time Topology}
  \source{kleiner-lott}
 \label{Rhatev}
On any time interval which is free of singular times,
and on which   $R_{min}(t) \: \le \: 0$ for all $t$
(which we are assuming), we have
\begin{equation} \label{Rhat}
\frac{d\hat{R}}{dt} \: \ge \: \frac23 \: \hat{R} \: V^{-1} \:
\int_M (R_{min} - R) \: dV.
\end{equation}
\end{lemma}
\begin{proof}
From (\ref{3dR}), $\frac{dR_{min}}{dt} \: \ge \: \frac23 \: R_{min}^2$.
Then
\begin{equation}
\frac{d\hat{R}}{dt} \: =\: \frac{dR_{min}}{dt} \: V^{\frac23} \: + \:
\frac23 \: R_{min} \: V^{- \: \frac13} \: \frac{dV}{dt} \: \ge \:
\frac23 \: R_{min}^2 \: V^{\frac23} \: - \:
\frac23 \: R_{min} \: V^{- \: \frac13} \: \int_M R \: dV,
\end{equation}
from which the lemma follows.
\end{proof}

\begin{corollary}[monotonicity for noncollapsing]
  \label{cor:kl-noncollapsed-pointed-limits-hyperbolic-monotonicity-noncollapsing}
  \dcref{Rmono}
  \group{Kleiner--Lott Long-Time Topology}
  \source{kleiner-lott}
 \label{Rmono}
If $R_{min}(t) \: \le \: 0$ for all $t$
(which we are assuming) then $\hat{R}(t)$ is nondecreasing.
\end{corollary}
\begin{proof}
  \uses{lem:kl-noncollapsed-pointed-limits-hyperbolic-noncollapsing-lemma}
If $\M$ is a nonsingular flow then the corollary follows from Lemma \ref{Rhatev}.
If there are surgeries then it only has the effect of decreasing $V(t)$, and so
possibly increasing $\hat{R}(t)$ (since $R_{min}(t) \le 0$).
\end{proof}

Put $\overline{R} \: = \: \lim_{t \rightarrow \infty} \hat{R}(t)$.

\begin{lemma}[noncollapsing lemma]
  \label{lem:kl-noncollapsed-pointed-limits-hyperbolic-noncollapsing-lemma-5}
  \dcref{RVlemma}
  \group{Kleiner--Lott Long-Time Topology}
  \source{kleiner-lott}
 \label{RVlemma}
If $\overline{V} \: > \: 0$ then
$\overline{R} \: \overline{V}^{-2/3} \: = \: - \: \frac32$.
\end{lemma}
\begin{proof}
  \uses{lem:kl-establishing-noncollapsing-presence-surgery-vanishing-noncollapsing}
Suppose that $\overline{V} \: > \: 0$.
Using Lemma \ref{Rev},
\begin{equation}
\overline{R} \: \overline{V}^{-2/3}  \: = \:
\lim_{t \rightarrow \infty} R_{min}(t) V(t)^{\frac23} \:
\left( V(t) \: \left(t + \frac14 \right)^{- \: \frac32} \right)^{- \: \frac23} \: = \:
\lim_{t \rightarrow \infty} \left( t + \frac14 \right) \: R_{min}(t) \: \ge \: - \: \frac32.
\end{equation}
In particular, there is a limit as $t \rightarrow \infty$ of
$\left( t + \frac14 \right) \: R_{min}(t)$.
Suppose that
\begin{equation}
\overline{R} \: \overline{V}^{-2/3} \: = \:
\lim_{t \rightarrow \infty} \left( t + \frac14 \right) \: R_{min}(t) \: = \:
c \: >
\: - \: \frac32.
\end{equation}
Combining this with (\ref{dVdt}) gives that for any $\mu > 0$,
$V(t) \: \le \: \const \:
t^{\mu-c}$ whenever $t$ is sufficiently large. Then
$V(t) t^{-3/2} \: \le \: \const \: t^{-(c+3/2 - \mu)}$.
Taking $\mu \: = \: \frac12 \: (c \: + \: \frac32)$, we
contradict the assumption that $\overline{V} \: > \: 0$.
\end{proof}

From the proof of Lemma \ref{RVlemma},
if $\overline{V} > 0$ then $R_{min}(t) \: \sim \: - \: \frac{3}{2t}$.

The next proposition shows that a long-time limit will
necessarily be hyperbolic.
\begin{proposition}[convergence for noncollapsing]
  \label{prop:kl-noncollapsed-pointed-limits-hyperbolic-convergence-noncollapsing}
  \dcref{propII.7.1}
  \group{Kleiner--Lott Long-Time Topology}
  \source{kleiner-lott}
 \label{propII.7.1}
Given the flow $\M$,
suppose that we have a sequence of parabolic neighborhoods
$P(x^\al, t^\al, r \sqrt{t^\al}, - r^2 t^\al)$, for
$t^\al \rightarrow \infty$ and some fixed $r \in (0,1)$, such that
the scalings of the parabolic neighborhoods with factor
$t^\al$ smoothly converge to some limit
solution $(\M_\infty, (\overline{x}, 1), g_\infty(\cdot))$
defined in a parabolic neighborhood
$P(\overline{x}, 1, r, -r^2)$. Then $g_\infty(t)$ has
constant sectional curvature $- \: \frac{1}{4t}$.
\end{proposition}
\begin{proof}
  \uses{lem:kl-noncollapsed-pointed-limits-hyperbolic-noncollapsing-lemma}
Suppose first that the flow is surgery-free. Because of the assumed existence of the limit
$(\M_\infty, (\overline{x}, 1), g_\infty(\cdot))$, the
original solution $\M$ has $\overline{V} > 0$.
We claim that the scalar curvature
on $P(\overline{x}, 1, r, -r^2)$ is spatially constant.
If not then there are numbers $c<0$ and $s_0, \mu > 0$ so that
\begin{equation}
\int_{B(\overline{x}, {s}, r)} (R^\prime_{min}({s})
\: - \: R(x, {s})) \: dV \: \le \: c
\end{equation}
whenever ${s} \in (s_0 - \mu, s_0 + \mu) \subset [1-r^2, 1]$,
where $R^\prime_{min}({s})$
is the minimum of $R$ over
$B(\overline{x}, {s}, r)$. Then for large $\alpha$,
\begin{equation}
\int_{B(x^\alpha, {s} t^\alpha, r \sqrt{t^\alpha})}
(R^\prime_{min}({s} t^\alpha)
\: - \: R(x, {s} t^\alpha)) \: dV \: < \: \frac{c}{2}
\: \sqrt{t^\alpha},
\end{equation}
where $R^\prime_{min}({s} t^\alpha)$ is now the minimum of $R$ over
$B(x^\alpha, {s} t^\alpha, r \sqrt{t^\alpha})$.
Thus
\begin{equation} \label{poscont}
\int_{B(x^\alpha, {s} t^\alpha, r \sqrt{t^\alpha})}
(R_{min}({s} t^\alpha) \: - \: R(x, {s} t^\alpha)) \:
dV \: < \: \frac{c}{2}
\: \sqrt{t^\alpha}.
\end{equation}

After passing to a subsequence, we can assume that
$\frac{t^{\alpha + 1}}{t^\al} \: > \: \frac{s_0 + \mu}{s_0 - \mu}$ for
all $\al$. From (\ref{Rhat}),
\begin{align} \label{Ralign}
\overline{R} \: - \: \hat{R}(0) \: & \ge \:
\frac23 \: \int_0^\infty \hat{R}(t) \: V(t)^{-1} \:
\int_M (R_{min}(t) - R(x,t)) \: dV(x) \: dt \\
& \ge \:
\frac23 \: \sum_{\al} t^\al
\int_{s_0- \mu}^{s_0 + \mu} \hat{R}(s t^\al) \: V(s t^\al)^{-1} \:
\int_M (R_{min}(s t^\al) - R(x,s t^\al)) \: dV(x) \: ds \notag \\
& \ge \:
\frac23 \: \sum_{\al} t^\al
\int_{s_0- \mu}^{s_0 + \mu} \hat{R}(s t^\al) \: V(s t^\al)^{-1} \:
\int_{B(x^\alpha, {s} t^\alpha, r \sqrt{t^\alpha})}
(R_{min}(s t^\al) - R(x,s t^\al)) \: dV(x) \: ds. \notag
\end{align}
Using the definitions of $\overline{V}$ and $\overline{R} \: = \:
- \: \frac32 \: \overline{V}^{\frac23}$, along with
(\ref{poscont}), it follows that
the right-hand side of (\ref{Ralign}) is infinite.
This contradicts the fact that $\overline{R} < \infty$.

Thus $R$ is spatially constant on
$P(\overline{x}, 1, r, -r^2)$.
As $R_{min}(t) \: \sim \: - \: \frac{3}{2t}$ on $M$, we know that the
scalar curvature $R$ at $(x,t) \in P(\overline{x}, 1, r, -r^2)$, which
only depends on $t$, satisfies
$R(t) \: \ge \: - \: \frac{3}{2t}$. It does not immediately follow that
the scalar curvature on $P(\overline{x}, 1, r, -r^2)$
equals $- \: \frac{3}{2t}$,  as $R_{min}(t)$ is the minimum of
the scalar curvature on all of $M$. However, if the scalar curvature is not identically
$- \: \frac{3}{2t}$ on $P(\overline{x}, 1, r, -r^2)$ then
again we can find $c<0$ and $s_0, \mu > 0$ so that
for large $\alpha$, (\ref{poscont}) holds for
$s \in (s_0-\mu, s_0 +\mu) \subset [1-r^2, 1]$.
Again we get a contradiction using (\ref{Rhat}).  Thus
$R(t) \: = \: - \: \frac{3}{2t}$ on $P(\overline{x}, 1, r, -r^2)$.
Then from (\ref{3dR}),
each time-slice of $P(\overline{x}, 1, r, -r^2)$
has an Einstein metric.  Thus the sectional curvature on
$P(\overline{x}, 1, r, -r^2)$ is $- \: \frac{1}{4t}$.


The argument goes through if one allows surgeries.
The main ingredient was the monotonicity formulas, which
still hold if there are surgeries. Note that for large
$\al$ there are no surgeries in $P(x^\al, t^\al, r \sqrt{t^\al}, - r^2 t^\al)$
by assumption.
\end{proof}

\section{II.7.2. Noncollapsed regions with a lower curvature bound are
almost hyperbolic on a large scale}

In this section it is shown that for fixed $A, r,w > 0$ and large time $t_0$, if
$B(x_0, t_0, r \sqrt{t_0}) \subset \M_{t_0}^+$ has volume at least
$wr^3 t_0^{\frac32}$ and sectional curvatures at least
$- \: r^{-2} t_0^{-1}$ then the Ricci flow on the parabolic neighborhood
$P(x_0, t_0, Ar \sqrt{t_0}, A r^2 t_0)$ is close to the flow on a
hyperbolic manifold.

We retain the assumptions of the previous section.

\begin{lemma}[noncollapsing lemma]
  \label{lem:kl-noncollapsed-regions-lower-curvature-bound-almost-hyperbolic-large-scale-noncollapsing-lemma}
  \dcref{lemII.7.2}
  \group{Kleiner--Lott Long-Time Topology}
  \source{kleiner-lott}
 \label{lemII.7.2} (cf. Lemma II.7.2) \\
(a) Given $w, r, \xi > 0$ one can find $T = T(w,r,\xi) < \infty$ such that
if the ball $B(x_0, t_0, r \sqrt{t_0}) \subset \M_{t_0}^+$ at some time $t_0 \ge T$ has
volume at least $wr^3 t_0^{\frac32}$ and sectional curvatures at least
$- \: r^{-2} t_0^{-1}$ then the curvature at $(x_0, t_0)$
satisfies
\begin{equation} \label{closeness}
|2tR_{ij}(x_0, t_0) + g_{ij}|^2 \: = \:
(2tR_{ij}(x_0, t_0) + g_{ij})(2tR^{ij}(x_0, t_0) + g^{ij}) \: < \: \xi^2.
\end{equation}
(b) Given in addition $A < \infty$ and allowing $T$ to depend on $A$, we
can ensure (\ref{closeness}) for all points in
$B(x_0, t_0, Ar \sqrt{t_0})$. \\
(c) The same is true for $P(x_0, t_0, Ar \sqrt{t_0}, A r^2 t_0)$.

Note that the time $T$ will depend on the initial metric.
\end{lemma}
\begin{proof}
  \uses{cor:kl-double-sided-curvature-bound-thick-part-existence-ricci-flow-surgery,prop:kl-noncollapsed-pointed-limits-hyperbolic-convergence-noncollapsing,prop:kl-proof-double-sided-curvature-bound-thick-part-modulo-two-propositions-existence-canonical-neighborhoods}
To prove (a), suppose that there is a sequence of points
$(x_0^\al, t_0^\al)$ with $t_0^\al \rightarrow \infty$
that provide a counterexample.  We wish to apply
Corollary \ref{II.6.8} with the parameter $r_0$ of the
corollary equal to $r \sqrt{t_0^\al}$.
Putting
\begin{equation}
\widehat{w} \: = \: \left( \min_{x \in [0,1]} \frac{\int_0^x \sinh^2(s) \:
ds}{x^3 \int_0^1 \sinh^2(s) \: ds} \right) \: w,
\end{equation}
if $r > \overline{r}(\widehat{w})$,
where $\overline{r}$ is from Corollary \ref{II.6.8},
then the hypotheses of the lemma will still be satisfied upon replacing
$w$ by $\widehat{w}$ and
$r$ by $\overline{r}(\widehat{w})$.
Thus
after redefining   $w$, if necessary,
we may assume that $r \le \overline{r}(w)$.
As the function
$h_{\max}(t)$
is nonincreasing, if $t_0^\al$ is sufficiently
large then
$\theta^{-1}(w)h_{\max}(t_0^\al) \: \le \: r \sqrt{t_0^\al}$.
Using Corollary \ref{II.6.8} with a redefinition of $w$,
we can take a convergent pointed subsequence
as $\al \rightarrow \infty$ of the $t_0^\al$-rescalings,
whose limit is defined in an abstract parabolic neighborhood.
From Proposition \ref{propII.7.1}
the limit will be hyperbolic, which is a contradiction.

For part (b), Corollary \ref{II.6.8} gives a bound $R \: \le \: K r_0^{-2}$ in
the unscathed parabolic neighborhood $P(x_0, t_0, r_0/4, -\tau r_0^2)$,
where
$r_0 = \min(r, \overline{r}(w^\prime)) \sqrt{t_0}$.
We apply Proposition \ref{propII.6.3} to the parabolic
neighborhood $P(x_0, t_0, r_0^\prime, - \: (r_0^\prime)^2)$ where
$K \: r_0^{-2} \: = \: (r_0^\prime)^{-2}$. By Proposition \ref{propII.6.3}(b), each
point $ y\in B(x_0, t_0, Ar\sqrt{t_0})$ with scalar curvature at
least $Q = K^\prime(A) r_0^{-2}$ has a canonical neighborhood.
Suppose that there is such a point.  From part (a) we have $R(x_0, t_0) < 0$, so
along a geodesic from $x_0$ to $y$ there will be some point
$x^\prime_0 \in B(x_0, t_0, Ar\sqrt{t_0})$
with scalar curvature $Q$. It also has a canonical neighborhood,
necessarily of type (a) or (b).
We can apply part (a) to a ball around $x^\prime_0$ with a radius
on the order of $(K^\prime(A))^{-1/2} r_0$, and with a value
of $w$ coming from the canonical neighborhood condition,
to get a contradiction
for large $t_0$.
(Note that $(K^\prime(A))^{-1/2} r_0$ is proportionate to
$\sqrt{t_0}$.)
Thus $R \: \le \: K^\prime(A) r_0^{-2}$
on $B(x_0, t_0, Ar\sqrt{t_0})$. If $T$ is large enough then the
$\Phi$-almost nonnegative curvature implies that $|\Rm| \: \le \:
K^\prime(A) r_0^{-2}$. Then the noncollapsing in
Proposition \ref{propII.6.3}(a) gives a lower local
volume bound. Hence we can apply part (a) of the lemma
to appropriate-sized balls in
$B(x_0, t_0, \frac{A}{2} r\sqrt{t_0})$. As $A$ is arbitrary, this
proves (b) of the lemma.

For part (c), without loss of generality we can take $\xi$ small.
Suppose that the claim is not true. Then there is a point $(x_0, t_0)$
that satisfies the hypotheses of the lemma but for which there is a point
$\left(x_1, t_1 \right) \in P(x_0, t_0, Ar\sqrt{t_0},
Ar^2 t_0)$ with $|2 t_1 R_{ij}(x_1, t_1) + g_{ij}| \ge \xi$.
Without loss of generality, we can take $(x_1, t_1)$ to be a first
such point in $P(x_0, t_0, Ar\sqrt{t_0}, Ar^2 t_0)$. By part (b),
$t_1 > t_0$. Then $|2 t R_{ij} + g_{ij}| \le \xi$ on
$P(x_0, t_0, Ar\sqrt{t_0}, t_1 -  t_0)$. If $\xi$ is small then this region has
negative sectional curvature and there are no
surgeries in the region. Using the length distortion estimates of
Section \ref{secI.8.3}, we can find
$r^\prime = r^\prime(r, A)>0$ so that the sectional curvature is bounded below
by $- \: (r^\prime)^{-2} \: t_1^{-1}$ on
$B(x_0, t_1, r^\prime \sqrt{t_1})$. Also, by the evolution of volume under
Ricci flow, there will be a $w^\prime =
w^\prime(r, w, \xi, A)$ so that
the volume of   $B(x_0, t_1, r^\prime \sqrt{t_1})$ is bounded below by
$w^\prime (r^\prime)^3 (t_1)^{\frac32}$. Thus for large $t_0$ we can apply
(b) to $B(x_0, t_1, A^\prime r^\prime \sqrt{t_1})$ with an appropriate choice of
$A^\prime$
to obtain a contradiction.
\end{proof}

\section{II.7.3. Thick-thin decomposition}

This section is concerned with the large-time decomposition of the
manifold in ``thick'' and ``thin'' parts.

\begin{definition}[curvature scale]
  \label{def:kl-thick-thin-decomposition-curvature-scale}
  \dcref{thick-thin-decomposition:definition-1}
  \group{Kleiner--Lott Long-Time Topology}
  \source{kleiner-lott}

For $x \in \M_t^+$,
let $\rho(x,t)$ be the unique number $\rho\in (0,\infty)$ such that
$\inf_{B^+(x,t, \rho)} \Rm \: = \: - \: \rho^{-2}$, if such a $\rho$ exists,
and put $\rho(x,t)=\infty$ otherwise.
\end{definition}
The function $\rho(x,t)$ is well-defined because $\M_t^+$ is a compact smooth
Riemannian manifold, so for fixed
$(x,t)\in\M_t^+$ the quantity
$\inf_{B^+(x,t,\rho)} \Rm$
is a continuous nonincreasing function of $\rho$ which is negative
for sufficiently large $\rho$ if and only if $(x,t)$ lies in a connected component
with negative sectional curvature somewhere; on the other hand the function $-\rho^{-2}$ is
continuous and  strictly increasing.  We note that when it is finite,
the quantity $\rho(x,t)$ may be larger than the diameter of the component
of $\M_t$ containing $(x,t)$.

As an example, if $\M$ is the flow on a manifold $M$ with spatially constant negative
curvature then for large $t$, $\rho(x,t) \sim 2 \sqrt{t}$ uniformly on $M$.
The ``thin'' part of $\M_t$, in the sense of hyperbolic geometry, can then be
characterized as the points $x$ so that $\vol(B(x,t,\rho(x,t)) \: < \: w \:
\rho^3(x,t)$, for an appropriate constant $w$.

\begin{lemma}[thick--thin geometry lemma]
  \label{lem:kl-thick-thin-decomposition-thick-thin-geometry-lemma}
  \dcref{bigenough}
  \group{Kleiner--Lott Long-Time Topology}
  \source{kleiner-lott}
 \label{bigenough}
For any $w > 0$ we can find $\overline{\rho} = \overline{\rho}(w) > 0$
and $\overline{T}=\overline{T}(w)$
such that if $t\geq \overline{T}$ and $\rho(x,t) < \overline{\rho} \sqrt{t}$ then
\begin{equation}
\vol(B(x,t,\rho(x,t))) \: < \: w \rho^3(x,t).
\end{equation}
\end{lemma}
\par\medskip\noindent\textit{Source argument (not certified as complete).}\ \begingroup\itshape
If the lemma is not true then there is a sequence
$(x^\al, t^\al)$ with $t^\al \rightarrow \infty$,
$\rho(x^\al, t^\al) \: (t^\al)^{-1/2} \rightarrow 0$ and
$\vol(B(x^\al, t^\al, \rho(x^\al, t^\al))) \: \ge \: w \rho^3(x^\al, t^\al)$.
The first step is to apply Corollary \ref{II.6.8},
but we need to know that for large $\al$
we have
$\rho(x^\al, t^\al) \: \ge \: \theta^{-1}(w) \: h_{\max}(t^\al)$, where
$\theta(w)$ and $h_{\max}$ are
from Corollary \ref{II.6.8}. Suppose that
this is not the case. Then after passing to a subsequence we have
$\rho(x^\al, t^\al) \: < \: \theta^{-1}(w) \:
h_{\max}(t^\al) \le r(t^{\al})$ for all $\al$,
where we used Remark \ref{hmax} in the last inequality.
There are points
$x^{\al,\prime} \in \overline{B(x^\al, t^\al, \rho(x^\al, t^\al))}$
with a sectional curvature
equal to  $- \: \rho^{-2}(x^\al, t^\al)$.
Applying the Hamilton-Ivey pinching estimate of
(\ref{estimate}) with $X^\al = \rho^{-2}(x^\al, t^\al)$, and using the fact
that
$\lim_{\al \rightarrow \infty} t^\al X^\al \: = \:  \infty$, gives
\begin{equation} \label{Rblowup}
\lim_{\al \rightarrow \infty}
R(x^{\al,\prime}, t^\al) \: \rho^2(x^\al, t^\al) \: = \: \infty.
\end{equation}

We claim that
the curvatures at the centers of the balls satisfy
\begin{equation}
\lim_{\al \rightarrow \infty} R(x^{\al}, t^\al) \: \rho^2(x^\al, t^\al) \: = \: \infty.
\end{equation}
Suppose not. Then there is some number $C \in (0, \infty)$ so that
after passing to a subsequence,
$R(x^{\al}, t^\al) \: \rho^2(x^\al, t^\al) \: \le \: C$ for all $\al$.
By continuity and (\ref{Rblowup}),
after passing to another subsequence we can assume that there is a point
$x^{\alpha\prime \prime}$ on a time-$t^\al$ geodesic segment between
$x^{\alpha}$ and $x^{\alpha \prime}$ so that
$R(x^{\al \prime \prime}, t^\al) \: \rho^2(x^\al, t^\al) \: = \: 2C$,
for all $\al$. We now apply Lemma \ref{claim2II.4.2} around
$(x^{\al \prime \prime}, t^\al)$ to get a contradiction to
(\ref{Rblowup}).
(More precisely, we apply a version of
Lemma \ref{claim2II.4.2} that applies along geodesics, as in Claim 2 of
II.4.2.) In applying Lemma \ref{claim2II.4.2}, we use the fact that
$\lim_{\al \rightarrow \infty} t^\al \: \rho^{-2}(x^\al, t^\al) \: = \:
\infty$ in order to say that
$\lim_{\al \rightarrow \infty}
\frac{\Phi(2C\rho^{-2}(x^\al, t^\al))}{2C\rho^{-2}(x^\al, t^\al)} \: = \: 0$,
with the notation of Lemma \ref{claim2II.4.2}.

Making a similar argument centered at other points in
$B(x^\al, t^\al, \rho(x^\al, t^\al))$, we deduce that
\begin{equation}
\lim_{\al \rightarrow \infty} \left( \inf R \big|_{B(x^\al, t^\al, \rho(x^\al, t^\al))} \right) \:
\rho^2(x^\al, t^\al) \: = \: \infty.
\end{equation}
In particular, since $\rho(x^\al, t^\al) \: \le \: r(t^\al)$, if $\al$ is
large then each point
$y^\al \in B(x^\al, t^\al, \rho(x^\al, t^\al))$ is the center of a canonical
neighborhood.
As $\al \rightarrow \infty$, the intrinsic scales
$R(y^{\al}, t^\al)^{-1/2}$ become arbitrarily small compared to $\rho(x^\al, t^\al)$.
However, by Lemma \ref{Alexlemma2} there is a subball $B^{\al, \prime}$ of
$B(x^\al,t^\al,\rho(x^\al,t^\al))$ with radius $\theta_0(w) \rho(x^\al,t^\al)$ so that every subball of
$B^{\al,\prime}$ has almost Euclidean volume. This contradicts the existence of a small
canonical neighborhood around each point of $B(x^\al,t^\al,\rho(x^\al,t^\al))$.


We now know that for large $\al$,
$\rho(x^\al, t^\al) \: \ge \: \theta^{-1}(w) \: h_{\max}(t^\al)$.
Thus we can apply Corollary \ref{II.6.8} to get an unscathed solution on the
parabolic neighborhood
$P(x^\al, t^\al, \rho(x^\al, t^\al)/4, - \tau \rho^2(x^\al, t^\al))$, with
$R \: < \: K_0 \: \rho(x^\al, t^\al)^{-2}$ there. Applying
Proposition \ref{propII.6.3}(c),
along with the fact that
$\lim_{\al \rightarrow \infty} t^\al \: \rho^{-2}(x^\al, t^\al) \: = \:
\infty$,
gives an estimate $R \: \le \: K_2 \: \rho(x^\al, t^\al)^{-2}$ on
$B(x^\al, t^\al, \rho(x^\al, t^\al))$.
But then for large $\al$, the Hamilton-Ivey
pinching gives $\Rm \: > \: - \: \frac12 \: \rho(x^\al, t^\al)^{-2}$ on
$B(x^\al, t^\al, \rho(x^\al, t^\al))$, which is a contradiction.
\endgroup\par\medskip

\begin{remark}[estimates for canonical neighborhoods]
  \label{rem:kl-thick-thin-decomposition-estimates-canonical-neighborhoods}
  \dcref{thick-thin-decomposition:remark-3}
  \group{Kleiner--Lott Long-Time Topology}
  \source{kleiner-lott}

Another approach to the above proof  would be to use the
canonical neighborhood at the center $(x^\al, t^\al)$ of the ball,
along with the Bishop-Gromov
inequality, to contradict the fact that
$\vol(B(x^\al, t^\al, \rho(x^\al, t^\al))) \: \ge \: w \rho^3(x^\al, t^\al)$. For this
to work we would have to know that the relative volume
$(\epsilon^{2} \: R(x^\al, t^\al))^{\frac32} \:  \vol(B(x^\al, t^\al, \epsilon^{-1}
R(x^\al, t^\al)^{- \: \frac12})$ of the canonical neighborhood (of
type (a) or (b)) around $(x^\al, t^\al)$ is small compared to $w$.
This will be the case if we take the constant $\epsilon$ to
be small enough, but in a $w$-dependent way.  Although $\epsilon$ is supposed
to be a universal constant, this approach will work because when
characterizing the graph manifold part in Section \ref{II.7.4}, $w$ can be taken to be a
small but fixed constant.
\end{remark}

\begin{definition}[thick--thin decomposition]
  \label{def:kl-thick-thin-decomposition-thick-thin-decomposition}
  \dcref{thick-thin-decomposition:definition-4}
  \group{Kleiner--Lott Long-Time Topology}
  \source{kleiner-lott}

The {\em $w$-thin part} $M^-(w, t) \subset \M_t^+$ is the set of points
$x \in M$ so that either $\rho(x,t)=\infty$ or
\begin{equation}
\vol(B(x, t, \rho(x, t))) \: < \: w \:
(\rho(x, t))^3.
\end{equation}
The {\em $w$-thick part} is $M^+(w, t) = \M_t^+  - M^-(w, t)$.
\end{definition}

\begin{lemma}[existence for thick--thin geometry]
  \label{lem:kl-thick-thin-decomposition-existence-thick-thin-geometry}
  \dcref{applies}
  \group{Kleiner--Lott Long-Time Topology}
  \source{kleiner-lott}
  \uses{lem:kl-thick-thin-decomposition-thick-thin-geometry-lemma}
 \label{applies}
Given $w > 0$, there are
$w^\prime = w^\prime(w) > 0$ and $T^\prime = T^\prime(w) < \infty$ so that
taking $r = \overline{\rho}(w)$
(with reference to Lemma \ref{bigenough}),
if $x_0 \in M^+(w,t)$ and $t_0 \ge T^\prime$ then
$B(x_0, t_0, r\sqrt{t_0})$ has volume at least $w^\prime r^3 t_0^{\frac32}$ and
sectional curvature at least $- \: r^{-2} \: t_0^{-1}$.
\end{lemma}
\begin{proof}
  \uses{lem:kl-thick-thin-decomposition-thick-thin-geometry-lemma}
Suppose that $x_0 \in M^+(w, t_0)$. From Lemma \ref{bigenough},
if $t_0$ is big enough (as a function of $w$) then $\rho(x_0, t_0)
\: \ge \: r \: \sqrt{t_0}$.  As $\Rm \: \ge \:
- \: \rho(x_0, t_0)^{-2}$ on $B(x_0, t_0, \rho(x_0, t_0))$, we have
$\Rm \: \ge \:
- \: r^{-2} \: t_0^{-1}$ on $B(x_0, t_0, r \sqrt{t_0})$.
As $\vol(B(x_0, t_0, \rho(x_0, t_0))) \: \ge \: w \:
(\rho(x_0, t_0))^3$, the Bishop-Gromov inequality gives a
lower bound on $\left( r \sqrt{t_0} \right)^{-3} \:
\vol(B(x_0, t_0, r \sqrt{t_0}))$ in terms of $w$.
\end{proof}

\section{Hyperbolic rigidity and stabilization of the thick part}
\label{hyprig}

Lemma \ref{applies} implies that Lemma \ref{lemII.7.2} applies to
$M^+(w,t)$ if $t$ is sufficiently large (as a function of $w$).
That is, if one takes a sequence of points
in the $w$-thick parts at a sequence of times tending
to infinity then the pointed time slices subconverge,
modulo rescaling the metrics by $t^{-1}$, to complete finite volume
hyperbolic manifolds with sectional curvatures equal to
$- \: \frac14$. (The $- \: \frac14$ comes from the Ricci flow
equation along with the equation $g(t) \: = \: t \: g(1)$ for the rescaled
limit, which implies that $g(1)$ has Einstein constant $- \: \frac12$.)
In what follows we will take the word ``hyperbolic'' for a $3$-manifold
to mean ``constant sectional curvature $- \: \frac14$''.
The next step, following Hamilton \cite{Hamiltonn},
is to show that
for large time the picture stabilizes, i.e. the limits
are unique in a strong sense.
\begin{proposition}[monotonicity for thick--thin geometry]
  \label{prop:kl-hyperbolic-rigidity-stabilization-thick-part-monotonicity-thick-thin-geometry}
  \dcref{slowisotopy}
  \group{Kleiner--Lott Long-Time Topology}
  \source{kleiner-lott}

\label{slowisotopy}
There exist a number $T_0<\infty$, a nonincreasing function
$\alpha:[T_0,\infty)\ra (0,\infty)$ with
$\lim_{t \rightarrow \infty} \alpha(t) \: = \: 0$,
a (possibly empty) collection
$\{(H_1,x_1),\ldots,(H_k,x_k)\}$
of complete connected pointed finite-volume hyperbolic $3$-manifolds
and a family of smooth maps
\begin{equation}
f(t):B_t= \bigcup_{i=1}^k \;B \left( x_i, \frac{1}{\alpha(t)} \right)\lra \M_t,
\end{equation}
 defined for $t\in [T_0,\infty)$, such that

1.  $f(t)$ is close to an isometry:
\begin{equation}
\|t^{-1}f(t)^*g_{\M_t}-g_{B_t}\|_{C^{[\frac{1}{\alpha(t)}]}}<\alpha(t),
\end{equation}

2. $f(t)$ defines a smooth family of maps which changes slowly
with time:
\begin{equation}
|\dot{f}(p,t)|<\alpha(t)t^{-\frac12}
\end{equation}
for all $p\in B_t$,
where $\dot{f}$ refers to the time derivative (as defined
with admissible curves),

and

3. $f(t)$ parametrizes more and more of the thick part:
 $M^+(\alpha(t),t)\subset \im(f(t))$ for all $t\geq T_0$.
\end{proposition}

\begin{remark}[characterization for thick--thin geometry]
  \label{rem:kl-hyperbolic-rigidity-stabilization-thick-part-characterization-thick-thin-geometry}
  \dcref{hyperbolic-rigidity-stabilization-thick-part:remark-2}
  \group{Kleiner--Lott Long-Time Topology}
  \source{kleiner-lott}
  \uses{prop:kl-hyperbolic-rigidity-stabilization-thick-part-monotonicity-thick-thin-geometry}

The analogous statement in \cite[Section 7.3]{Perelman2} is in terms of a fixed $w$.
That is, for a given $w$ one considers pointed limits of
$\{(\M^+_{t_j}, (x_j, t_j), t_j^{-1} g(t_j))\}_{j=1}^\infty$
with $\lim_{j \rightarrow \infty} t_j = \infty$ and the basepoint
satisfying
$x_j \in M^+(w, t_j) \subset \M^+_{t_j}$
for all $j$. Considering
the possible limit spaces in a certain order, as described below,
one extracts complete pointed finite-volume hyperbolic manifolds $\{(H_i, x_i)\}_{i=1}^k$
with $x_i$ in the $w$-thick part of $H_i$.
There is a number
$w_0 > 0$ so that as long as $w \le w_0$, the hyperbolic manifolds
$H_i$ are independent of $w$.
For any $w^\prime > 0$,
as time goes on the $w^\prime$-thick part of  $\bigcup_{i=1}^k H_i$
better approximates $M^+(w^\prime, t)$.
Hence the formulation of
Proposition \ref{slowisotopy} is equivalent to that of
\cite[Section 7.3]{Perelman2}.
\end{remark}

Rather than proving Proposition \ref{slowisotopy}
using harmonic maps
as in \cite{Hamiltonn},
we give a simple proof using
smooth compactness and a smoothing argument.  Roughly speaking, the idea is
to exploit a variant of Mostow rigidity to show that for large $t$, the components
of the $w$-thick part change slowly with time, and
are close to hyperbolic manifolds which are isolated (due to a refinement
of Mostow-Prasad rigidity). This forces them to eventually stabilize.

\begin{definition}[pointed smooth closeness]
  \label{def:kl-hyperbolic-rigidity-stabilization-thick-part-pointed-smooth-closeness}
  \dcref{fapprox}
  \group{Kleiner--Lott Long-Time Topology}
  \source{kleiner-lott}

If $(X,x)$ and $(Y,y)$ are pointed smooth Riemannian manifolds
and $\eps>0$ then $(X,x)$ is $\eps$-close to $(Y,y)$
if there is a pointed map $f:(X,x)\ra (Y,y)$ such that
\begin{equation}
f\restr_{\ol{B(x,\eps^{-1})}}:\ol{B(x,\eps^{-1})}\ra Y
\end{equation}
 is a diffeomorphism onto its image
and
\begin{equation}
\label{fapprox}
\|f^*g_Y-g_X\|_{C^{\eps^{-1}}}<\eps,
\end{equation}
where the norm is taken on $\ol{B(x,\eps^{-1})}$.
Note that nothing is required of $f$ on the complement of $\ol{B(x,\eps^{-1})}$.
Such a map $f$ is called an {\em $\eps$-approximation}.
\end{definition}

We will sometimes refer
to a partially defined map $f:(X,x)\supset (W,x)\ra (Y,y)$ as an
$\eps$-approximation provided that $W$ contains $\ol{B(x,\eps^{-1})}$ and
the conditions above are satisfied. By convention we will permit $X$ and $Y$
to be disconnnected, in which case $\eps$-closeness only says something about the
components containing the basepoints.
We say that two maps $f_1,f_2:(X,x)\ra Y$ (not necessarily basepoint-preserving)
are $\eps$-close if
\begin{equation}
\sup_{p\in \ol{B(x,\eps^{-1})}}d_Y(f_1(p),f_2(p))<\eps.
\end{equation}

We  recall some facts about hyperbolic manifolds.
There is a constant $\mu_0>0$, the Margulis constant,
such that if $X$ is a complete connected finite-volume
hyperbolic $3$-manifold (orientable, as usual), $\mu\leq \mu_0$, and
\begin{equation}
X_\mu=\{x\in X\mid \injrad(X,x)\geq \mu\}
\end{equation}
is the $\mu$-thick part of $X$, then $X_\mu$ is a nonempty
compact manifold-with-boundary whose complement $U$
is a finite union of components $U_1,\ldots,U_k$,
where each $U_i$ is isometric either to a geodesic tube
around a closed geodesic or to  a cusp.
In particular, $X_\mu$ is connected and there is a one-to-one
correspondence between the boundary components of $X_\mu$
and the ``thin'' components $U_i$.   For each $i$,
let $\rho_i:\ol{U_i}\ra\R$ denote either the distance function
from the core geodesic or a Busemann function, in the
tube and cusp cases respectively. In the latter case
we normalize $\rho_i$ so that $\rho_i^{-1}(0)=\D U_i$.  (The
Busemann function goes to $- \: \infty$ as one goes down
the cusp.)
The {\em radial direction}
is the direction field on $U_i -  \core(U_i)$
defined by $\nabla\rho_i$, where $\core(U_i)$ is the core
geodesic when $U_i$ is a geodesic tube and the empty set otherwise.

\begin{lemma}[existence for thick--thin geometry]
  \label{lem:kl-hyperbolic-rigidity-stabilization-thick-part-existence-thick-thin-geometry}
  \dcref{prasad}
  \group{Kleiner--Lott Long-Time Topology}
  \source{kleiner-lott}

\label{prasad}
Let $(X,x)$ be a pointed complete connected finite-volume hyperbolic
$3$-manifold.  Then for each $\zeta>0$ there exists $\xi>0$
such that if $X'$ is a complete finite-volume hyperbolic
manifold with at least as many cusps as $X$,
and $f:(X,x)\ra X'$ is a $\xi$-approximation, then there is an
isometry $\hat f: (X,x)\ra X'$ which
is $\zeta$-close to $f$.
\end{lemma}
This was stated as Theorem 8.1 in \cite{Hamiltonn} as going
back to the work of Mostow. We give a proof here.
The hypothesis about cusps is essential because every pointed noncompact
finite-volume hyperbolic $3$-manifold $(X,x)$ is a pointed limit of a sequence
$\{(X_i,x_i)\}_{i=1}^\infty$
of compact hyperbolic manifolds.
Hence for every $\xi > 0$, if $i$ is sufficiently large then there is a
$\xi$-approximation $f \: : \: (X, x) \rightarrow (X_i, x_i)$, but there is
no isometry from $X$ to $X_i$.

\begin{proof}
The main step is to show that for the fixed $(X, x)$, if
$\xi$ is sufficiently small then for any $\xi$-approximation
$f \: : \: (X, x) \rightarrow X^\prime$ satisfying the hypotheses of the lemma,
the manifolds $X$ and $X^\prime$ are
diffeomorphic. The proof of this will use the
Margulis thick-thin decomposition. The rest of the assertion then
follows readily from Mostow-Prasad rigidity
\cite{Mostow,Prasad}.

Pick $\mu_1\in (0,\mu_0)$ so that $X -  X_{\mu_1}$
consists only of cusps $U_1,\ldots, U_k$. The thick part
$X_{\mu_1}$ is compact and connected. Given $\xi > 0$, let
$f:(B(x,\xi^{-1}),x)\ra (X',x')$ be a $\xi$-approximation
as in (\ref{fapprox}).
The intuitive idea is that because of the compactness
of $X_{\mu_1}$, if $\xi$ is sufficiently small then
$f \big|_{X_{\mu_1}}$ is close to being an isometry
from $X_{\mu_1}$ to its image.
Then $f(X_{\mu_1})$ is close to a connected component of
the thick part $X^\prime_{\mu_1}$ of $X^\prime$. As
$f(X_{\mu_1})$ and $X^\prime_{\mu_1}$ are connected,
this means that $f(X_{\mu_1})$ is close to $X^\prime_{\mu_1}$.
We will show that in fact $X_{\mu_1}$ is diffeomorphic to
$X^\prime_{\mu_1}$. The boundary components of
$X_{\mu_1}$ correspond to the cusps of $X$ and the
boundary components of
$X^\prime_{\mu_1}$ correspond to the connected components of
$X^\prime - X^\prime_{\mu_1}$. As $X^\prime$ has at least
as many cusps as $X$ by assumption, it follows that the
connected components of $X^\prime - X^\prime_{\mu_1}$ are
all cusps. Hence $X$ and
$X^\prime$ are diffeomorphic.

In order to show that $X_{\mu_1}$ is diffeomorphic to
$X^\prime_{\mu_1}$, we will take a larger region
$W \supset X_{\mu_1}$ that is diffeomorphic to $X_{\mu_1}$
and show that $f(W)$ can be isotoped to $X^\prime_{\mu_1}$
by sliding it inward along the radial direction. More precisely,
in each cusp
$U_i$ put $V_i= \rho_i^{-1}([-3L,-L])$,
where $L\gg 1$ is large enough that  every cuspidal torus
$\rho_i^{-1}(s)$ with $s\in [-3L,-L]$ has diameter much less than one.
Let $W\subset X$ be the complement of the open horoballs
at height $-2L$, i.e.
\begin{equation}
W= X -  \bigcup_{i=1}^k  \rho_i^{-1}(-\infty, -2L).
\end{equation}
When $\xi$
is sufficiently small, $f$ will preserve injectivity radius
to within a factor close to $1$ for points $p\in B(x,\xi^{-1})$
with $d(p,\D B(x,\xi^{-1}))>2\injrad(X,p)$.  Therefore
when $\xi$ is small, $f$ will map each $V_i$ into
 $X' -  X'_{\mu_1}$, and hence into one of the
connected components $U'_{k_i}$ of $X' -  X'_{\mu_1}$.
Let $Z_i'$ be the image of $Z_i=\rho_i^{-1}(-2L)$ under $f$.   Note that
$d(\core(U'_{k_i}),Z_i')\gtrsim L$ (if $\core(U'_{k_i}) \neq
\emptyset$), for otherwise
$f^{-1}(\core(U'_{k_i}))$ would be a closed curve with
small diameter and curvature lying in $U_i$, which contradicts
the fact that the horospheres have principal curvatures
$- \: \frac12$ (because of our normalization that the sectional curvatures are
$- \: \frac14$). Thus
for each point $z'\in Z_i'$ there is a minimizing radial
geodesic segment $\ga'$ passing through $z'$ with
$d(z',\D\ga')\gtrsim L$.  The preimage of $\gamma^\prime$ under $f$ is
a curve $\ga\subset X$ with small curvature and length $\gtrsim L$
passing through $Z_i$. This forces the direction of $\ga$
to be nearly radial
and so transverse to $Z_i$.
Hence $Z_i'$ is transverse to the radial direction in
$U'_{k_i}$.  Combining this with the  fact that $Z_i'$ is
embedded  implies that $Z_i'$ is isotopic in $\ol{U'_{k_i}}$ to
$\D U'_{k_i}$.  It follows that $f(W)$
is isotopic to $X'_{\mu_1}$. Then by the preceding argument
involving counting the number of cusps,
$X$ and $X'$ are diffeomorphic.
We apply Mostow-Prasad rigidity \cite{Mostow,Prasad} to deduce that
$X$ is isometric to $X'$.

We now claim that for any $\zeta > 0$, if $\xi$
is sufficiently small then the map $f$ is $\zeta$-close to an
isometry from $X$ to $X^\prime$. Suppose not.  Then there are a number $\zeta > 0$ and
a sequence of $\frac{1}{i}$-approximations $f_i \: : \: (X, x) \rightarrow (X^\prime_i, x^\prime_i)$
so that none of the $f_i$'s are $\zeta$-close to
any
isometry from
$(X, x)$ to $(X^\prime_i, x^\prime_i)$. Taking a convergent subsequence of
the maps $f_i$ gives a limit isometry
$f_\infty \: : \: (X, x) \rightarrow (X^\prime_\infty, x^\prime_\infty)$.
From what has already been proven,
for large $i$ we know that $X^\prime_i$ is isometric to $X$, and so isometric
to $X^\prime_\infty$.
This is a contradiction.
\end{proof}

Recall the statement of Lemma \ref{lemII.7.2}.

\begin{definition}[hyperbolic pointed-limit space]
  \label{def:kl-hyperbolic-rigidity-stabilization-thick-part-hyperbolic-pointed-limit-space}
  \dcref{hyperbolic-rigidity-stabilization-thick-part:definition-5}
  \group{Kleiner--Lott Long-Time Topology}
  \source{kleiner-lott}

Given $w >0$, let $\Lambda_w$ be the space
of complete pointed finite-volume hyperbolic
$3$-manifolds that arise as pointed limits of
sequences $\{(\M^+_{t_i}, (x_i, t_i), t_i^{-1} g(t_i))\}_{i=1}^\infty$
with $\lim_{i \rightarrow \infty} t_i = \infty$ and the basepoint
$(x_i,t_i)$ satisfying
$(x_i,t_i) \in M^+(w, t_i) \subset \M^+_{t_i}$
for all $i$.
\end{definition}

The space $\La_w$ is compact in the smooth pointed topology.
Any element of $\La_w$ has volume at most $\overline{V}$,
the latter being defined in Definition \ref{VR}.

The next lemma summarizes the content of Lemma \ref{lemII.7.2}.

\bigskip
\begin{lemma}[existence for cap geometry]
  \label{lem:kl-hyperbolic-rigidity-stabilization-thick-part-existence-cap-geometry}
  \dcref{II.7.2again}
  \group{Kleiner--Lott Long-Time Topology}
  \source{kleiner-lott}

\label{II.7.2again}
Given $w>0$, there is a decreasing function
$\beta:[0,\infty)\ra (0,\infty]$
with $\lim_{s\ra\infty}\beta(s)=0$ such that
if
$(x,t)\in M^+(w,t)\subset\M_t^+$,
and  $Z_t$ denotes
the forward time slice
$\M_t^+$ rescaled by $t^{-1}$, then

1.  Some $(X,x)\in \La_w$
is $\beta(t)$-close to $(Z_t,(x,t))$.

2.
$B(x,t,\beta(t)^{-1}\sqrt{t})\subset \M_t^+$
is unscathed
on the interval $[t,2t]$ and if $\ga:[t,2t]\ra \M$ is a static
curve starting at $(x,t)$, $\bar t\in [t,2t]$, then the map
\begin{equation}
B(x,t,\beta(t)^{-1}\sqrt{t})\ra
P(x,t,\beta(t)^{-1}\sqrt{t},t)\cap\M_{\bar t}
\end{equation}
defined by following static curves induces a map
\begin{equation}
i_{t,\bar t}:(Z_t,(x,t))\supset (B(x,t,\beta(t)^{-1}),(x,t))
\ra (Z_{\bar t},\ga(\bar t))
\end{equation}
satisfying
\begin{equation}
\|\left(i_{t,\bar t}\right)^*g_{Z_{\bar
t}}-g_{Z_t}\|_{C^{\beta(t)^{-1}}}
<\beta(t).
\end{equation}
\end{lemma}
\begin{proof}
  \uses{lem:kl-noncollapsed-regions-lower-curvature-bound-almost-hyperbolic-large-scale-noncollapsing-lemma}
This follows immediately from Lemma \ref{lemII.7.2}.
\end{proof}

\bigskip
\noindent
{\em Proof of Proposition \ref{slowisotopy}.}
If for some $w>0$ we have $\La_w=\emptyset$ then
$M^+(w,t)=\emptyset$ for large $t$.  Thus if
$\La_w=\emptyset$ for all $w>0$,  we can take
the empty collection of pointed hyperbolic manifolds
and then 1 and 2 will be satisfied vacuously,
and $\alpha(t)$ may be chosen so that 3 holds.
So we assume that $\La_w\neq\emptyset$ for some $w>0$.

Since every complete finite-volume hyperbolic $3$-manifold
has a point with injectivity radius $\geq \mu_0$,
there is a $w_0>0$ such that the collections $\{\La_w\}_{w\leq w_0}$
contain the same sets of underlying hyperbolic manifolds
(although the basepoints have more freedom when $w$ is small).
We let $H_1$ be a hyperbolic manifold from this
collection with the fewest
cusps and we choose a basepoint $x_1\in H_1$ so that
$(H_1,x_1)\in \La_{w_0}$.   Put $w_1= \frac{w_0}{2}$.
Note that $x_1$ lies in the $w_1$-thick part of $H_1$.
In what follows we will use the fact that if $f$ is a
$\epsilon$-approximation from $H_1$, for sufficiently small
$\epsilon$, then $f(x_1)$ will lie in the $.9 w_0$-thick part of
the image.

The idea of the first step of the proof is to define a family $\{f_0(t)\}$
of $\de$-approximations $(H_1,x_1)\ra Z_t$,
for all $t$ sufficiently large, by taking
  a $\de$-approximation $(H_1,x_1)\ra Z_t$,
pushing it along static curves, and arguing using Lemma \ref{prasad}
that one can make small adjustments from time to time
to keep it a $\de$-approximation. The family $\{f_0(t)\}$
will not vary continuously with time, but it will have controlled
``jumps''.

More precisely, pick $T_0<\infty$ and let $\xi_1,\ldots,\xi_4>0$ be
parameters to be specified later.
We assume that $T_0$ is large enough so that
$
2\beta(T_0)<\xi_1,
$
where $\beta$ is from Lemma \ref{II.7.2again}.
By the definition of $\La_{w_1}$, we may pick $T_0$ so that there
is a point
$(\overline{x}_0,T_0)\in M^+(w_1,T_0)\subset\M_{T_0}^+$
and a $\xi_1$-approximation $f_0(T_0):(H_1,x_1)\ra (Z_{T_0},
\overline{x}_0)$.

To do the induction step, for a given $j \ge 0$ suppose that at time $2^j T_0$
there is a point
$(\overline{x}_j, 2^j T_0)\in M^+(w_1,2^j T_0)\subset\M_{2^j T_0}^+$
and a $\xi_1$-approximation $f_0(2^j T_0):(H_1,x_1)\ra
(Z_{2^j T_0},\overline{x}_j)$. As mentioned above, if
$\xi_1$ is small then in fact $\overline{x}_j \in M^+(.9 w_0,2^j T_0)$.

By part 2 of Lemma \ref{II.7.2again}, provided that $T_0$ is
sufficiently large we may define, for all $t\in [2^j T_0, 2^{j+1} T_0]$,
a  $2\xi_1$-approximation $f_0(t):(H_1,x_1)\ra Z_t$ by
moving $f_0(2^j T_0)$ along static curves.  Provided that
$\xi_1$ is sufficiently small we will have
$f_0(2^{j+1} T_0)(x_1)\in M^+(w_1,2^{j+1} T_0)$ and then part 1 of
Lemma \ref{II.7.2again} says there is some $(H',x')\in \La_{w_1}$
with a $\beta(2^{j+1} T_0)$-approximation
$\phi:(H',x')\ra (Z_{2^{j+1} T_0},f_0(2^{j+1} T_0)(x_1))$.
Provided that
$\beta(2^{j+1} T_0)$
and $\xi_1$ are sufficiently small, the partially defined map
$\phi^{-1}\circ f_0(2^{j+1} T_0)$ will define a $\xi_2$-approximation
from $(H_1,x_1)$ to $(H',x')$. Hence provided that $\xi_2$
is sufficiently small, by Lemma \ref{prasad}
the map will be $\xi_3$-close to an isometry
$\psi:(H_1,x_1)\ra H'$.
(In applying Lemma \ref{prasad} we use the fact that $H_1$ is
also a manifold with the fewest number of cusps in $\Lambda_{w_1}$.)
 Put $\phi_1= \phi\circ\psi$.
Provided that $\xi_3$ is sufficiently small,
$f_0(2^{j+1} T_0)$ and $\phi_1$ will be $\xi_4$-close as maps from
$(H_1,x_1)$ to $Z_{2^{j+1} T_0}$.
Since $\phi_1$ is a $\beta(2^{j+1} T_0)$-approximation
precomposed with an isometry which shifts basepoints a
distance at most $\xi_3$, it will be a $2\beta(2^{j+1} T_0)$-approximation
provided that $\xi_3<1$ and $\beta(2^{j+1} T_0)<\frac12$.  We now redefine
$f_0(2^{j+1} T_0)$ to be $\phi_1$ and let $\overline{x}_{j+1}$ be the
image of $x_1$ under $\phi_1$. This completes the induction step.

In this way we define a family of partially defined maps
$\{f_0(t):(H_1,x_1)\ra Z_t\}_{t\in [T_0,\infty)}$.
From the construction,
$f_0(2^jT_0)$ is a $2\beta(2^jT_0)$-approximation
for all $j\geq 0$.   Lemma \ref{II.7.2again} then implies
that there is a  function
$\alpha_1:[T_0,\infty)\ra(0,\infty)$ decreasing to
zero at infinity such that for all $t\in[T_0,\infty)$,
$f_0(t)$ is an $\alpha_1(t)$-approximation, and for every
$\bar t\in [t,2t]$ we may slide
$f_0(t)$
along
static curves to define an $\alpha_1(t)$-approximation
 $h(\bar t):(H_1,x_1)\ra Z_t$ which is $\alpha_1(t)$-close
to
$f_0(\bar t)$.

One may now employ a standard smoothing argument to convert the
family $\{f_0(t)\}_{t\in[T_0,\infty)}$ into  a family
$\{f_1(t)\}_{t\in [T_0,\infty)}$ which satisfies the first two
conditions of the proposition.   If condition 3 fails to hold then
we redefine the $\La_{w}$'s by considering  limits of only those
$\{(\M^+_{t_i}, (x_i, t_i), t_i^{-1} g(t_i))\}_{i=1}^\infty$
with $t_i \rightarrow \infty$ and
$x_i \in M^+(w,t_i)\subset \M_{t_i}^+$
not in
the image of $f_1(t_i)$.  Repeating the construction we
obtain a pointed hyperbolic manifold $(H_2,x_2)$ and a
family $\{f_2(t)\}$ defined for large $t$ satisfying conditions
1 and 2, where $\im(f_2(t))$ is disjoint from $\im(f_1(t))$ for large
$t$.  Iteration of this procedure must stop after  $k$
steps for some finite number $k$, in view of the fact that $\overline{V} < \infty$ and
the fact that there is a positive lower bound on the volumes of
complete hyperbolic $3$-manifolds. We get the
desired family $\{f(t)\}$ by taking the union of the maps
$f_1(t),\ldots,f_k(t)$.


\section{Incompressibility of cuspidal tori} \label{inctori}

By Proposition \ref{slowisotopy}, we know that for large times the thick
part of the manifold can
be parametrized by a collection of (truncated) finite volume hyperbolic
manifolds.  In this section we show that each cuspidal torus
maps to an embedded incompressible torus in $\M_t$.
(An alternative argument is given in Section \ref{II.8}.)
The strategy, due
to Hamilton, is to argue by contradiction. If such a torus were compressible
then there would be an embedded compressing disk of least area at each
time. By estimating the rate of change of the area of such disks one concludes
that the area must go to zero in finite time, which is absurd.

Let $T_0$, $\alpha$, $\{(H_1,x_1),\ldots,(H_k,x_k)\}$, $B_t$, and $f(t)$ be as
in Proposition \ref{slowisotopy}.  We will consider a fixed $H_i$, with $1\leq i\leq k$,  which
is noncompact.   Choose a  number
$a>0$ much smaller than the Margulis constant and let
$\{V_1,\ldots,V_l\}\subset H_i$
be  the cusp regions bounded by tori of diameter $a$.  Each
$V_j$ is an embedded $3$-dimensional submanifold (with boundary) of $H_i$
and is isometric to the quotient of a horoball in hyperbolic $3$-space $\mathbb{H}^3$
by the action of a copy of $\Z^2$ sitting in the stabilizer of the horoball.  The
boundary $\D V_j$ is a totally umbilic torus whose principal curvatures are
equal to $\frac12$ everywhere.   We let $Y\subset H_i$ be the closure of the complement
of $\bigcup_{j=1}^l V_j$ in $H_i$.

Let $T_a<\infty$ be large enough  that  $B_{T_a}$  (defined as in Proposition \ref{slowisotopy})
contains $Y$.

In order to focus on a given cusp, we now fix an integer $1\leq j \leq l$  and put
\begin{equation}
Z =  \partial V_{j},\quad\hat Z_t =  f(t)(Z)\\
\quad \hat Y_t =  f(t)(Y),\quad
\hat W_t =  \M_t^+ -  \Int(\hat Y_t)
\end{equation}
for every $t\geq T_a$.  The objective of this section is:

\begin{proposition}[incompressibility of cuspidal tori proposition]
  \label{prop:kl-incompressibility-cuspidal-tori-incompressibility-cuspidal-tori-proposition}
  \dcref{propincompressible}
  \group{Kleiner--Lott Long-Time Topology}
  \source{kleiner-lott}

\label{propincompressible}
The homomorphism
\begin{equation} \label{injeqn}
\pi_1(f(t)):\pi_1(Z,\star)\ra \pi_1(\M_t^+,f(t)(\star))
\end{equation}
is a monomorphism for all $t\geq T_a$.
\end{proposition}
\par\medskip\noindent\textit{Source argument (not certified as complete).}\ \begingroup\itshape
The proof will occupy the remainder of this section.
The first step is:
\begin{lemma}[incompressibility of cuspidal tori lemma]
  \label{lem:kl-incompressibility-cuspidal-tori-incompressibility-cuspidal-tori-lemma}
  \dcref{stablelemma}
  \group{Kleiner--Lott Long-Time Topology}
  \source{kleiner-lott}
 \label{stablelemma}
The kernels of the homomorphisms
\begin{equation}
\label{eqnpi1inclusion}
\pi_1(f(t)):\pi_1(Z,\star)\ra \pi_1(\M_t^+,f(t)(\star)),\quad
\pi_1(f(t)):\pi_1(Z,\star)\ra \pi_1(\hat W_t,f(t)(\star))
\end{equation}
are independent of $t$, for all $t\geq T_a$.
\end{lemma}
\par\medskip\noindent\textit{Source argument (not certified as complete).}\ \begingroup\itshape
We prove the assertion for the first homomorphism. The argument for the
second one is similar.

The kernel obviously remains constant on any time interval which is free of
singular times. Suppose that $t_0\geq T_a$ is a singular time.  Then the intersection
$
\M_{t_0}^+\cap\M_{t_0}^-
$
includes into $\M_{t_0}^+$ and, by using static curves, into $\M_t$ for $t \neq t_0$
close to $t_0$.  By Van Kampen's theorem, these inclusions induce monomorphisms of the
fundamental groups.  Therefore for $t$ close to $t_0$,
the kernel of (\ref{eqnpi1inclusion}) is the same as the kernel of
\begin{equation}
\pi_1(f(t)):\pi_1(Z,\star)\ra \pi_1(\M_{t_0}^+\cap\M_{t_0}^-),
\end{equation}
which is independent of $t$ for times $t$ close to $t_0$.
\endgroup\par\medskip

\bigskip
We now assume that the kernel of
\begin{equation}
\pi_1(f(t)):\pi_1(Z,\star)\ra \pi_1(\M_t^+,f(t)(\star))
\end{equation}
is nontrivial for some, and hence every, $t\geq T_a$.  By Van Kampen's
theorem and the fact that the cuspidal torus
$Z\subset Y$ is incompressible in $Y$, it follows that the kernel $K$
of
\begin{equation}
\pi_1(f(t)):\pi_1(Z,\star)\ra \pi_1(\hat W_t,f(t)(\star))
\end{equation}
in nontrivial for all $t\geq T_a$.   By Poincar\'e duality,
$\Image \left( \HH^1(\hat{W}_t; \R) \rightarrow \HH^1(\partial \hat{W}_t; \R) \right)$
is a Lagrangian subspace of $\HH^1(\partial \hat{W}_t; \R)$. In particular,
$\Image \left( \HH^1(\hat{W}_t; \R) \rightarrow \HH^1(Z; \R) \right)$
has
rank one. Dually,
$\Ker \left( \HH_1(Z; \R) \rightarrow \HH_1(\widehat{W}_t; \R) \right)$
has rank one and so $K$,
a subgroup of a rank-two free abelian group,
has rank one. We note that for all large $t$,
$\hat{Z}_t$ is a convex
boundary component of $\hat W_t$.
The main theorem of \cite{meeksyau}  implies that for
every such $t$, there is
is a least-area compressing disk
\begin{equation}
(N^2_t,\D N^2_t)\subset (\hat W_t,\hat Z_t).
\end{equation}
We  recall that a compressing disk is an embedded disk whose boundary curve is essential
in $\hat Z_t$. We note that by definition, $\hat W_t$ is a compact manifold
even when $t$ is a singular time.
The embedded curve $f(t)^{-1}(\D N_t)\subset Z$
represents a primitive element of $\pi_1(Z)$ which,
since $K$ has  rank one, must therefore generate $K$.
It follows that
modulo taking inverses, the homotopy class of $f(t)^{-1}(\D N_t)\subset Z$ is independent of $t$.

We define a function $A:[T_a,\infty)\ra (0,\infty)$
by letting $A(t)$ be the infimum of the areas of such embedded compressing disks.
We now show that the least-area compressing disks avoid the surgery regions.

\begin{lemma}[existence for Ricci flow with surgery]
  \label{lem:kl-incompressibility-cuspidal-tori-existence-ricci-flow-surgery}
  \dcref{lemavoids}
  \group{Kleiner--Lott Long-Time Topology}
  \source{kleiner-lott}

\label{lemavoids}
Let $\de(t)$ be the surgery parameter from
Section \ref{II.4.4}. There is a $T = T(a) < \infty$ so that whenever $t \ge T$,
no point in any area-minimizing compressing disk  $N_t\subset \hat W_t$ is in the
center of a $10 \delta(t)$-neck.
\end{lemma}
\begin{proof}
If the lemma were not true then
there would be a sequence of times $t_k\ra\infty$
and  for each $k$ an area-minimizing compressing disk $(N_{t_k},\D N_{t_k})
\subset (\hat W_{t_k},\hat Z_{t_k})$,
along with a point $x_k \in N_{t_k}$ that is in the center of a
$10 \delta(t_k)$-neck.
Note that the scalar curvature
near $\D N_{t_k}$ is comparable to $- \: \frac{3}{2t_k}$.
We now rescale by  $R(x_k, t_k)$, and consider the map of pointed
manifolds $f_k:(N_{t_k},\D N_{t_k},x_k)\hookrightarrow (\hat W_{t_k},\hat Z_{t_k},x_k)$
where the domain is equipped with the pullback Riemannian metric.
By \cite{schoen} and standard elliptic regularity, for all $\rho<\infty$
and every integer $j$, the $j^{th}$ covariant derivative of the
second fundamental form of $f_k$ is uniformly bounded on the ball $B(x_k,\rho)\subset N_{t_k}$,
for sufficiently large $k$.  Therefore the pointed Riemannian
manifolds $(N_{t_k},\D N_{t_k},x_k)$
subconverge in the smooth topology to a pointed, complete, connected, smooth manifold
$(N_\infty,x_\infty)$.  Using the same bounds on the derivatives
of the second fundamental form, we may extract a limit mapping
$\phi_{\infty} \: : \: N_{\infty} \ra \R \times S^2$ which
is a $2$-sided  isometric stable minimal immersion.
By  \cite[Theorem 2]{schoenyau}, $\phi_{\infty}$ is a totally geodesic
immersion whose normal vector field in $M$ has vanishing Ricci curvature.
It follows that $\phi_{\infty}$ is a cover of
a fiber $\{ \pt \}\times S^2$.  This contradicts the fact that $N_\infty$ is noncompact.
\end{proof}

We redefine $T_a$ if necessary so that $T_a$ is greater than the $T$ of
Lemma \ref{lemavoids}.


We can isotope the surface $Z$ by moving it down the cusp $V_j$. In doing so
we do not change the group $K$
but we can make the diameter of $Z$ as small as desired. The next lemma
refers to this isotopy freedom.


\begin{lemma}[existence for kappa solutions]
  \label{lem:kl-incompressibility-cuspidal-tori-existence-kappa-solutions}
  \dcref{Hamlemma}
  \group{Kleiner--Lott Long-Time Topology}
  \source{kleiner-lott}
 \label{Hamlemma}
Given $D > 0$, there is a number $a_0 > 0$
so that for any $a \in (0,a_0)$, if $\diam(Z) = a$ and $t$ is
sufficiently large then
$\int_{\D N_t}\kappa_{\D N_t}ds \: \le \: \frac{D}{2}$ and
$\length(\D N_t) \: \le \: \frac{D}{2} \: \sqrt{t}$,
where $\kappa_{\D N_t}$ is the geodesic curvature of
$\D N_t \subset N_t$.
\end{lemma}
\begin{proof}
  \uses{prop:kl-hyperbolic-rigidity-stabilization-thick-part-monotonicity-thick-thin-geometry}
This is proved in \cite[Sections 11 and 12]{Hamiltonn}. We just state the
main idea.  For the purposes of this proof, we give $\hat W_t$ the
metric $t^{-1} g(t)$. First,  $\D N_t$ is the intersection of $N_t$ with
$\hat Z_t$. Because $N_t$ is minimal with respect to free boundary
conditions (i.e. the only constraint is that $\D N_t$ is in the
right homotopy class in $\hat Z_t$), it follows that $N_t$ meets
$\hat Z_t$ orthogonally. Then $\kappa_{\D N_t} =
\Pi(v,v)$, where $\Pi$ is the second fundamental form of
$\hat Z_t$ in $\hat W_t$ and $v$ is the unit tangent vector of $\D N_t$.
Given $a>0$, let $Z$ be the horospherical torus
in $V_j$ of diameter
$a$.
By Proposition \ref{slowisotopy},
for large $t$ the map $f(t)$ is close to being an isometry of pairs
$(Y, Z) \rightarrow (\hat Y_t, \hat Z_t)$.
As $Z$ has principal
curvatures $\frac12$, we may assume that
$\Pi(v,v)$ is close to $\frac12$. This reduces the problem to showing that
with an appropriate choice of $a_0$, if $a \in (0, a_0)$ then
for large values of $t$ the length of
$\D N_t$ is guaranteed to be small. The intuition is that since
$\hat W_t$ is close to being the standard cusp $V_j$, a large piece of
the minimal disk
$N_t$ should be like a minimal surface $N_\infty$  in $V_j$
that intersects $Z$ in the given homotopy class.  Such a minimal
surface in $V_j$ essentially consists of a
geodesic curve in $Z$ going all the way down the cusp. The
length of the intersection of $N_\infty$ with the horospherical
torus of diameter $a$ is proportionate to $a$.  Hence if $a_0$ is small enough, one
would expect that if $a < a_0$ and if $t$ is large then
the length of $\D N_t$ is small.
In particular, the length of $\D N_t$ is uniformly bounded with
respect to $a$.
 A detailed proof appears in
\cite[Section 12]{Hamiltonn}.

Rescaling from the metric $t^{-1} g(t)$ to the original metric $g(t)$,
$\int_{\D N_t}\kappa_{\D N_t}ds$ is unchanged and
$\length(\D N_t)$ is multiplied by $\sqrt{t}$.
\end{proof}

\begin{lemma}[existence for diameter]
  \label{lem:kl-incompressibility-cuspidal-tori-existence-diameter}
  \dcref{lemsupport}
  \group{Kleiner--Lott Long-Time Topology}
  \source{kleiner-lott}
 \label{lemsupport}
For every $D>0$ there is a number $a_0>0$ with the following property.
Given $a \in (0, a_0)$, suppose that we take $Z$ to be the torus
cross-section
in $V_j$
of diameter $a$. Then
there is a number $T^\prime_a < \infty$ so that as long as
$t_0\geq T^\prime_a$, there is a smooth function $\bar A$ defined
on a neighborhood of $t_0$ such that $\bar  A(t_0)=A(t_0)$, $\bar A\geq A$
everywhere, and
\begin{equation} \label{dereqn}
\bar A'(t_0)< \frac{3}{4}\left(\frac{1}{t_0+\frac14}\right) A(t_0)-2\pi+D.
\end{equation}
\end{lemma}
\begin{proof}
  \uses{lem:kl-establishing-noncollapsing-presence-surgery-vanishing-noncollapsing,lem:kl-incompressibility-cuspidal-tori-existence-kappa-solutions,lem:kl-incompressibility-cuspidal-tori-existence-ricci-flow-surgery,prop:kl-hyperbolic-rigidity-stabilization-thick-part-monotonicity-thick-thin-geometry}
Take $a_0$ as in Lemma \ref{Hamlemma}.

For $t_0 > T_a$,
we begin with the minimizing compressing disk $N_{t_0}\subset\M_{t_0}^+$.
If $t_0$ is a surgery time and $N_{t_0}$ intersected the surgery region $\M_{t_0}^+ -
(\M_{t_0}^+ \cap \M_{t_0}^-)$ then $N_{t_0}$ would have to pass through a
$10\delta(t_0)$-neck, which is impossible by Lemma \ref{lemavoids}.
Thus $N_{t_0}$ avoids any parts added by surgery.

For $t$ close to $t_0$ we define an embedded compressing disk $S_t\subset\M_t^+$
as follows. We take $N_{t_0}$
and extend it slightly to a smooth surface $N_{t_0}'\subset \M_{t_0}^+$ which contains
$N_{t_0}$ in its interior.  The surface $N_{t_0}'$ will
be unscathed on some open time interval containing $t_0$. If we let
$S_t'\subset \M_t^+$ be the  surface obtained by moving $N_{t_0}'$ along
static curves then for some $b>0$, the surface $S_t'$ will
intersect $\D \hat W_t$ transversely for all $t\in (t_0-b,t_0+b)$.
Putting
\begin{equation}
S_t =  S_t'\cap \hat W_t
\end{equation}
defines a compressing disk for $\hat Z_t\subset \hat W_t$.

Define $\bar A:(t_0-b,t_0+b)\ra \R$ by
\begin{equation}
\bar A(t) =  \area(S_t).
\end{equation}
Clearly $\bar A(t_0)=A(t_0)$ and $\bar A\geq A$.

For the rest of the calculation, we will view $S_t$ as a surface sitting
in a fixed manifold $M$ (a fattening of $\hat W_t$)
 with a varying metric $g(\cdot)$, and
put $S =  S_{t_0}=N_{t_0}$.

By the first variation formula for area,
\begin{equation}
\label{eqninitial}
\bar A'(t_0)=\int_{\D S}\langle X,\nu_{\D S}\rangle\:ds
+\int_S\frac{d}{dt} \: \Big|_{t=t_0} \dvol_S,
\end{equation}
where $X$ denotes the variation vector field for $\hat Z_t$, viewed
as a surface moving in $M$, and
$\nu_{\D S}$ is the outward normal vector along $\D S$.
By Proposition \ref{slowisotopy}, there is an estimate $|X|\leq \alpha(t_0)t_0^{-\frac12}$,
where $\alpha(t_0)\ra 0$ as $t_0 \ra \infty$.
Therefore
\begin{equation}
\label{eqnbdyterm}
\left|\int_{\D S_{t_0}}\langle X,\nu_{\D S_{t_0}}\rangle\:ds\right| \:
\leq \: \alpha(t_0) \: t_0^{-\frac12} \: \length(\D N_{t_0}).
\end{equation}
By Lemma \ref{Hamlemma}, the right-hand side of (\ref{eqnbdyterm})
is bounded above by $\frac{D}{2}$ if $t_0$ is large.

We turn to the second term in (\ref{eqninitial}).  Pick $p\in S$
and let $e_1,e_2,e_3$ be an orthonormal basis for $T_pM$ with
$e_1$ and $e_2$ tangent to $S$.  Then
\begin{equation}
\frac{1}{\dvol_S} \: \frac{d}{dt} \: \Big|_{t=t_0} \dvol_S \: = \: \frac12
\sum_{i=1}^2 \frac{dg}{dt} \Big|_{t=t_0} (e_i, e_i)
= -\Ric(e_1,e_1)-\Ric(e_2,e_2).
\end{equation}
Now
\begin{align}
-\Ric(e_1,e_1)-\Ric(e_2,e_2)) & =-R+\Ric(e_3,e_3)
=-R+K(e_3,e_1)+K(e_3,e_2) \\
& =-\frac{R}{2}-K(e_1,e_2)
=-\frac{R}{2}-K_S+\gk_S, \notag
\end{align}
where $K_S$ denotes the Gauss curvature of $S$ and $\gk_S$
denotes the product of the principal curvatures.  Applying the Gauss-Bonnet formula
\begin{equation}
\int_{\D S}\kappa_{\D S}\:ds=2\pi-\int_S K_S\vol_S,
\end{equation}
the fact that $\gk_S\leq 0$
(since $S$ is time-$t_0$ minimal)
and the inequality
\begin{equation}
R_{\min}(t)\geq -\frac{3}{2}\left(\frac{1}{t+\frac14}\right)
\end{equation}
from Lemma \ref{Rev}, we obtain
\begin{equation} \label{evolve}
\int_S\:\frac{d}{dt} \: \Big|_{t=t_0} \dvol_S \: \le \:
\int_S\:\frac34\left(\frac{1}{t_0+\frac14}\right) \: \dvol_S+\int_{\D S}\kappa_{\D S}ds-2\pi.
\end{equation}
By Lemma \ref{Hamlemma}, if $a \in (0,a_0)$, $\diam(Z) = a$ and $t_0$ is
sufficiently large then
$\int_{\D S}\kappa_{\D S}ds \: \le \: \frac{D}{2}$.
Using (\ref{eqninitial}), (\ref{eqnbdyterm}) and (\ref{evolve}),
if $t_0$ is large then
\begin{equation}
\bar A'(t_0)< \frac34\left(\frac{1}{t_0+\frac14}\right)A(t_0)-2\pi+D.
\end{equation}
This proves the lemma.
\end{proof}
\noindent
{\em Proof of Proposition \ref{propincompressible}.}
Pick $D<2\pi$. Let $a<a_0$ and $T^\prime_a$ be as in Lemma
\ref{lemsupport}.

By Lemma \ref{lemsupport}, $A$ is bounded on compact subsets of
$[T^\prime_a, \infty)$.
By Lemma \ref{lemavoids}, for any
$t \in [T^\prime_a, \infty)$ we can find a compact set $K_t \subset \M_t^+$
so that for all $t^\prime \in [T^\prime_a, \infty)$ sufficiently close to $t$,
the compressing disk $N_{t^\prime}$ lies in $K_t$. If
$(K_t, g(t))$ and $(K_t, g(t^\prime))$ are $e^\sigma$-biLipschitz
equivalent then $A(t) \: \le \: e^{2\sigma} \area \left( N_{t^\prime} \right)
\: = \: e^{2 \sigma} \: A(t^\prime)$ and
$A(t^\prime) \: \le \: e^{2\sigma} \area \left( N_{t} \right)
\: = \: e^{2 \sigma} \: A(t)$. It follows that $A$ is continuous
on $[T^\prime_a, \infty)$.


For $t \ge T_a^\prime$, put
\begin{equation}
F(t) =
\left(t+\frac14\right)^{-\frac34} A(t) + 4 (2\pi-D) \left(t+\frac14\right)^{\frac14}.
\end{equation}
We claim that
$F(t) \le F(T_a^\prime)$ for all $t \ge T_a^\prime$. Suppose not.  Put
$t_0 = \inf \{t \ge T_a^\prime \: : \: F(t) > F(T_a^\prime) \}$.
By continuity, $F(t_0) = F(T_a^\prime)$.
Consider the function $\bar A$ of Lemma \ref{lemsupport}.
Put
\begin{equation}
\bar F(t) =  \left(t+\frac14\right)^{-\frac34} \bar A(t) + 4 (2\pi-D) \left(t+\frac14\right)^{\frac14}.
\end{equation}
Then $\bar F(t_0) \: = \: F(t_0)$ and
in a small interval around $t_0$, we have
$\bar F \ge F$. However, (\ref{dereqn}) implies that
$\bar F^\prime(t_0)
\: < \: 0$.
There is some $\sigma > 0$ so that
for $t \in (t_0, t_0 + \sigma)$, we have
\begin{equation}
F(t) \: \le \: \bar F (t) \: \le \: \bar F(t_0) \: + \: \frac12 \:
\bar F^\prime(t_0) \: (t - t_0) \:
< \: \bar F(t_0) \: = \: F(t_0) \: = \: F(T_a^\prime),
\end{equation}
which contradicts the definition of $t_0$.

Thus if $t \ge T_a^\prime$ then
$F(t) \le F(T_a^\prime)$. This implies that
$A(t)$ is negative for large $t$, which contradicts the fact
that an area is nonnegative.


We have shown that the homomorphism (\ref{injeqn}) is injective if $t$ is sufficiently large.
In view of Lemma \ref{stablelemma}, the same statement holds for all $t \ge T_a$.
\endgroup\par\medskip

\section{II.7.4. The thin part is a graph manifold} \label{II.7.4}

This section is concerned with showing that the thin part
$M^-(w,t)$ is a graph manifold.  We refer to Appendix \ref{geom}
for the definition of a graph manifold.
We remind the reader that this completes the proof the geometrization
conjecture.

The next two theorems are purely Riemannian. They say that if a $3$-manifold
is locally volume-collapsed, with sectional curvature bounded below, then it
is a graph manifold. They differ slightly in their hypotheses.

\begin{theorem}[existence for thick--thin geometry]
  \label{thm:kl-thin-part-graph-manifold-existence-thick-thin-geometry}
  \dcref{thmII.7.4}
  \group{Kleiner--Lott Long-Time Topology}
  \source{kleiner-lott}
 (cf. Theorem II.7.4) \label{thmII.7.4}
Suppose that $(M^\al, g^\al)$ is a sequence of compact oriented Riemannian
$3$-manifolds, closed or with convex boundary, and $w^\al \rightarrow 0$.
Assume that \\
(1) for each point $x \in M^\al$ there exists a radius $\rho = \rho^\al(x)$
not exceeding the diameter of $M^\al$
such that
the ball $B(x, \rho)$ in the metric $g^\al$ has volume at most
$w^\al \rho^3$ and sectional curvatures at least $- \: \rho^{-2}$.\\
(2) each component of the boundary of $M^\al$ has diameter at most
$w^\al$, and has a (topologically trivial) collar of length one, where
the sectional curvatures are between $- \: \frac14 \: - \: \epsilon$ and
$- \: \frac14 \: - \: \epsilon$.

Then for large $\alpha$, $M^\al$ is diffeomorphic to a graph manifold.
\end{theorem}

\begin{remark}[estimates for thick--thin geometry]
  \label{rem:kl-thin-part-graph-manifold-estimates-thick-thin-geometry}
  \dcref{thin-part-graph-manifold:remark-2}
  \group{Kleiner--Lott Long-Time Topology}
  \source{kleiner-lott}
  \uses{thm:kl-thin-part-graph-manifold-existence-thick-thin-geometry}

A proof of Theorem \ref{thmII.7.4} appears in \cite[Section 8]{Shioya-Yamaguchi}.
The proof in \cite{Shioya-Yamaguchi} is for closed manifolds, but in
view of condition (2) the method of proof clearly goes through to
manifolds-with-boundary as considered in Theorem \ref{thmII.7.4}.
The statement of the theorem in \cite[Theorem II.7.4]{Perelman2} also has a condition
$\rho < 1$, which seems to be unnecessary.
\end{remark}

\begin{theorem}[existence for thick--thin geometry]
  \label{thm:kl-thin-part-graph-manifold-existence-thick-thin-geometry-3}
  \dcref{thmII.7.4second}
  \group{Kleiner--Lott Long-Time Topology}
  \source{kleiner-lott}
 (cf. Theorem II.7.4) \label{thmII.7.4second}
Suppose that $(M^\al, g^\al)$ is a sequence of compact oriented Riemannian
$3$-manifolds, closed or with convex boundary, and $w^\al \rightarrow 0$.
Assume that \\
(1) for each point $x \in M^\al$ there exists a radius $\rho = \rho^\al(x)$
such that
the ball $B(x, \rho)$ in the metric $g^\al$ has volume at most
$w^\al \rho^3$ and sectional curvatures at least $- \: \rho^{-2}$.\\
(2) each component of the boundary of $M^\al$ has diameter at most
$w^\al$, and has a (topologically trivial) collar of length one, where
the sectional curvatures are between $- \: \frac14 \: - \: \epsilon$ and
$- \: \frac14 \: - \: \epsilon$. \\
(3) for every $w^\prime > 0$ and $m  \in \{0, 1, \ldots, [\epsilon^{-1}]\}$,
there exist $\overline{r} = \overline{r}(w^\prime) > 0$
and $K_m = K_m(w^\prime) <\infty$ such that for sufficiently
large $\alpha$, if $r \in (0, \overline{r}]$
and a ball $B(x,r)$ in the metric $g^\al$
has volume at least $w^\prime r^3$ and sectional curvatures at least
$- \: r^{-2}$ then $|\nabla^m \Rm|(x) \: \le \: K_m \: r^{-m-2}$.

Then for large $\alpha$, $M^\al$ is diffeomorphic to a graph manifold.
\end{theorem}

\begin{remark}[estimates for noncollapsing]
  \label{rem:kl-thin-part-graph-manifold-estimates-noncollapsing}
  \dcref{thin-part-graph-manifold:remark-4}
  \group{Kleiner--Lott Long-Time Topology}
  \source{kleiner-lott}

The statement of this theorem in \cite[Theorem II.7.4]{Perelman2} has the
stronger assumption that (3) holds for all $m \ge 0$. In the
application to the locally collapsing part of the Ricci flow,
it is not clear that this stronger condition holds.  However, one does get
a bound on a large number of derivatives, which is good enough.
\end{remark}

\begin{remark}[structural remark on thick--thin geometry]
  \label{rem:kl-thin-part-graph-manifold-structural-remark-thick-thin-geometry}
  \dcref{thin-part-graph-manifold:remark-5}
  \group{Kleiner--Lott Long-Time Topology}
  \source{kleiner-lott}
  \uses{thm:kl-thin-part-graph-manifold-existence-thick-thin-geometry-3}

As pointed out in \cite[Section 7.4]{Perelman2}, adding condition (3) simplifies
the proof
and allows one to avoid both Alexandrov spaces and Perelman's
stability theorem.
(A proof of Perelman's stability theorem appears in \cite{Kapovitch}).
Proofs of Theorem
\ref{thmII.7.4second} are in \cite{Besson},
\cite{Kleiner-Lott} and \cite{Morgan-Tian2}.
\end{remark}

\begin{remark}[positivity for thick--thin geometry]
  \label{rem:kl-thin-part-graph-manifold-positivity-thick-thin-geometry}
  \dcref{thin-part-graph-manifold:remark-6}
  \group{Kleiner--Lott Long-Time Topology}
  \source{kleiner-lott}
  \uses{thm:kl-thin-part-graph-manifold-existence-thick-thin-geometry,thm:kl-thin-part-graph-manifold-existence-thick-thin-geometry-3}

Comparing Theorems \ref{thmII.7.4} and \ref{thmII.7.4second},
Theorem \ref{thmII.7.4} has the extra assumption that
$\rho^\al(x)$ does not exceed the diameter of $M^\al$.
Without this extra assumption, the Alexandrov space arguments could give
that for large $\al$, $M^\al$ is homeomorphic to a nonnegatively curved
Alexandrov space \cite[Theorem 1.1(2)]{Shioya-Yamaguchi}. This does not
immediately imply that $M^\al$ is a graph manifold.
\end{remark}


\begin{remark}[existence for thick--thin geometry]
  \label{rem:kl-thin-part-graph-manifold-existence-thick-thin-geometry}
  \dcref{thin-part-graph-manifold:remark-7}
  \group{Kleiner--Lott Long-Time Topology}
  \source{kleiner-lott}
  \uses{thm:kl-no-local-collapsing-theorem-ii-existence-kappa-solutions,thm:kl-thin-part-graph-manifold-existence-thick-thin-geometry,thm:kl-thin-part-graph-manifold-existence-thick-thin-geometry-3}

We give some simple examples where the collapsing theorems apply.
Let $(\Sigma, g_{hyp})$ be a closed surface with the hyperbolic
metric.  Let $S^1(\mu)$ be a circle of length $\mu$ and consider
the Ricci flow on $S^1 \times \Sigma$ with the initial metric
$S^1(\mu) \times (\Sigma, c_0 \: g_{hyp})$. The Ricci flow solution at
time $t$ is $S^1(\mu) \times (\Sigma, (c_0 + 2t) g_{hyp})$.
In the rest of this example we consider
the rescaled metric $t^{-1} g(t)$.
Its diameter goes like
$O(t^0)$ and its sectional curvatures  go like $O(t^0)$. If we
take $\rho$ to be a small constant then for large $t$,
$B(x, t, \rho)$ is approximately a circle bundle over a ball in a hyperbolic surface
of constant sectional curvature $- \: \frac12$, with circle lengths
that go like $t^{-1/2}$. The sectional curvature on
$B(x, t, \rho)$ is bounded below by $- \: \rho^{-2}$, and
$\rho^{-3} \: \vol(B(x, t, \rho)) \: \sim \: t^{-1/2}$.
Theorems \ref{thmII.7.4}
and \ref{thmII.7.4second} both apply.

Next, consider a compact
$3$-dimensional nilmanifold that evolves under the Ricci flow.
Let $\theta_1, \theta_2, \theta_3$ be affine-parallel $1$-forms on $M$
which lift to Maurer-Cartan forms on the Heisenberg group, with
$d\theta_1 \: = \: d\theta_2 \: = \: 0$ and $d\theta_3 \: = \:
\theta_1 \: \wedge \: \theta_2$.
Consider the metric
\begin{equation}
g(t) \: = \: \alpha^2(t) \: \theta_1^2 \: + \:
\beta^2(t) \: \theta_2^2 \: + \:
\gamma^2(t) \: \theta_3^2.
\end{equation}
Its sectional curvatures are
$R_{1212} \: = \: - \: \frac34 \: \frac{\gamma^2}{\alpha^2 \beta^2}$ and
$R_{1313} \: = \: R_{2323}
\: = \: \frac14 \: \frac{\gamma^2}{\alpha^2 \beta^2}$.
The Ricci tensor is
\begin{equation}
\Ric \: = \: \frac12 \: \frac{\gamma^2}{\alpha^2 \beta^2} \:
\left( - \: \alpha^2 \: \theta_1^2 \: - \: \beta^2 \: \theta_2^2 \:
+ \: \gamma^2 \: \theta_3^2 \right).
\end{equation}
The general solution to the Ricci flow equation is of the form
\begin{align}
\alpha^2(t) \: & = \: A_0 (t + t_0)^{1/3}, \\
\beta^2(t) \: & = \: B_0 (t + t_0)^{1/3}, \notag \\
\gamma^2(t) \: & = \: \frac{A_0 B_0}{3} (t + t_0)^{-1/3}. \notag
\end{align}
In the rest of this example we consider
the rescaled metric $t^{-1} g(t)$.
Its diameter goes like
$t^{-1/3}$, its volume goes like $t^{-4/3}$
and its sectional curvatures
go like $t^0$. If we take
$\rho(x) \: = \: \diam$ then
$\rho^{-3} \vol(B(x, \rho)) \sim t^{-1/3}$, so both Theorem \ref{thmII.7.4}
and Theorem \ref{thmII.7.4second} apply.
We could also take $\rho(x)$ to be a small constant
$c > 0$, in which case Theorem \ref{thmII.7.4second} applies.


In general, among the eight maximal homogeneous geometries, the
rescaled solution for a compact $3$-manifold with geometry
$H^2 \times \R$ or $\widetilde{\SL_2(\R)}$
will collapse to a hyperbolic surface of constant
sectional curvature $- \: \frac12$.
The rescaled solution for a $\Sol$ geometry will collapse to a circle.
The rescaled solution for an $\R^3$ or $\Nil$ geometry will collapse to
a point.

We remark that although these homogeneous solutions
are collapsing in the sense of Theorem
\ref{thmII.7.4}, there is no contradiction with the
no local collapsing result of Theorem \ref{nolocalcollapse},
which only rules out
local collapsing on a finite time interval.
\end{remark}

Returning to our Ricci flow with surgery, recall the statement of
Proposition \ref{slowisotopy}. If the collection
$\{H_1,\ldots, H_k\}$ of Proposition \ref{slowisotopy} is nonempty then
for large $t$, let $\widehat{H}_i(t)$ be the result of removing from $H_i$ the horoballs whose
boundaries are at
distance approximately $\frac12 \: \alpha(t)$ from the basepoint $x_i$.
(If there are no such horoballs then $\widehat{H}_i(t) = H_i$.)
Put
\begin{equation}
M_{thin}(t) \: = \: \M_t^+ - f(t) (\widehat{H}_1(t) \cup \ldots \widehat{H}_k(t)).
\end{equation}

\begin{proposition}[thick--thin geometry proposition]
  \label{prop:kl-thin-part-graph-manifold-thick-thin-geometry-proposition}
  \dcref{thingraph}
  \group{Kleiner--Lott Long-Time Topology}
  \source{kleiner-lott}
 \label{thingraph}
For large $t$, $M_{thin}(t)$ is a graph manifold.
\end{proposition}
\begin{proof}
  \uses{cor:kl-double-sided-curvature-bound-thick-part-existence-ricci-flow-surgery,lem:kl-curvature-bounds-from-priori-assumptions-canonical-neighborhoods-lemma,lem:kl-curvature-bounds-from-priori-assumptions-estimates-curvature,lem:kl-locating-small-balls-whose-subballs-almost-euclidean-volume-existence-sectional-curvature,lem:kl-thin-part-graph-manifold-thick-thin-geometry-lemma,rem:kl-earlier-scalar-curvature-bounds-smaller-balls-from-lower-curvature-bound-structural-remark-scalar-curvature,thm:kl-thin-part-graph-manifold-existence-thick-thin-geometry,thm:kl-thin-part-graph-manifold-existence-thick-thin-geometry-3}
We give two closely related proofs, one using Theorem \ref{thmII.7.4second} and
one using Theorem \ref{thmII.7.4}.


If the proposition is not true then there is a sequence $t^\al \rightarrow \infty$
so that for each $\al$, $M_{thin}(t^\al)$ is not a graph manifold.
Let $M^\al$ be the manifold obtained from  $M_{thin}(t^\al)$
by throwing away connected components which are closed and admit metrics
of nonnegative sectional curvature, and put $g^\al \: = \: (t^\al)^{-1} g(t^\al)$.
Since any closed manifold of nonnegative sectional curvature is a graph
manifold by \cite{Hamiltonnnnn}, for each $\al$ the manifold $M^\al$ is not a graph
manifold.

We first show that the assumptions of Theorem \ref{thmII.7.4second} are verified.

\begin{lemma}[thick--thin geometry lemma]
  \label{lem:kl-thin-part-graph-manifold-thick-thin-geometry-lemma}
  \dcref{condition}
  \group{Kleiner--Lott Long-Time Topology}
  \source{kleiner-lott}
  \uses{thm:kl-thin-part-graph-manifold-existence-thick-thin-geometry-3}
 \label{condition}
Condition (3) in Theorem \ref{thmII.7.4second} holds for the $M^\al$'s.
\end{lemma}
\begin{proof}
  \uses{cor:kl-double-sided-curvature-bound-thick-part-existence-ricci-flow-surgery,lem:kl-curvature-bounds-from-priori-assumptions-canonical-neighborhoods-lemma,lem:kl-curvature-bounds-from-priori-assumptions-estimates-curvature,lem:kl-locating-small-balls-whose-subballs-almost-euclidean-volume-existence-sectional-curvature,rem:kl-earlier-scalar-curvature-bounds-smaller-balls-from-lower-curvature-bound-structural-remark-scalar-curvature,thm:kl-thin-part-graph-manifold-existence-thick-thin-geometry-3}
With $w^\prime$ being a parameter as in
Condition (3) in Theorem \ref{thmII.7.4second},
let $\overline{r} = \overline{r}(w^\prime)$ be the parameter of Corollary \ref{II.6.8}.
It is enough to show that for large $\al$, if $r \in (0, \overline{r} \sqrt{t^\al}]$,
$x^\al \in \M_{t^\al}^+$, and $B(x^\al, t^\al, r)$
has volume at least $w^\prime r^3$ and sectional curvatures
bounded below by $- \: r^{-2}$, then $|\nabla^m \Rm|(x^\al, t^\al) \: \le \:
K_m \: r^{-m-2}$ for an appropriate choice of constants $K_m$.

To prove this by contradiction, we  assume that after passing
to a subsequence if necessary, there are  $r^\al\in (0,\overline{r} \sqrt{t^\al}]$
and $x^\al \in \M_{t^\al}^+$ such that $B(x^\al, t^\al,r^\al)$ has
volume
at least $w^\prime (r^\al)^3$ and sectional curvature at least
$-(r^\al)^{-2}$, but $\lim_{\al \rightarrow \infty} \:
(r^\al)^{m+2} \: |\nabla^m \Rm|(x^\al, t^\al) \: = \:
\infty$
for some $m \le [\epsilon^{-1}]$.

In the notation of Corollary \ref{II.6.8}, if
$r^\al \ge \theta^{-1}(w^\prime) h_{\max}(t^\al)$ for infinitely many $\al$
then for these $\al$, Corollary \ref{II.6.8} gives a curvature bound on an unscathed
parabolic neighborhood
$P \left( x^\al, t^\al, r^\al/4, - \tau (r^\al)^2 \right)$ and hence,
by Appendix \ref{applocalder}, derivative bounds at $(x^\al, t^\al)$.
This is a contradiction.  Therefore we may assume that
$r^\al < \theta^{-1}(w^\prime) h_{\max}(t^\al) \le r(t^\al)$ for all $\al$,
where we used Remark \ref{hmax} for the last inequality.

Suppose first that
$R(x^\al, t^\al) \le \left( r(t^\al) \right)^{-2}$.
By Lemma \ref{claim1II.4.2}, there is an estimate
$R \: \le \: 16 \: \left( r(t^\al) \right)^{-2}$ on the parabolic neighborhood
$P \left( x^\al, t^\al, \frac14 \eta^{-1} r(t^\al),
- \: \frac{1}{16} \eta^{-1} \left( r(t^\al) \right)^2 \right)$.
A surgery in this neighborhood could only occur where
$R \ge h_{\max}(t^\al)^{-2}$. For large $\al$,
$h_{\max}(t^\al)^{-2} >> r(t^\al)^{-2}$ by Remark \ref{hmax}.
Hence this neighborhood is unscathed.
Appendix \ref{applocalder} now gives bounds of the form
$|\nabla^m \Rm|(x^\al, t^\al) \: \le \: \const \:
\left( r(t^\al) \right)^{-m-2} \: \le \: \const \:
(r^\al)^{-m-2}$, which is a contradiction..

Suppose now that $R(x^\al, t^\al) > \left( r(t^\al) \right)^{-2}$.
Then $(x^\al, t^\al)$ is in the center of a canonical neighborhood and there
are universal estimates $|\nabla^m \Rm|(x^\al, t^\al) \: \le \: \const(m) \:
R(x^\al, t^\al)^{\frac{m+2}{2}}$ for all $m \le [\epsilon^{-1}]$. Hence
in this case, it suffices to show that $R(x^\al, t^\al)$ is bounded above by
a constant times $(r^\al)^{-2}$, i.e. it suffices to get a contradiction
just to the assumption that
$\lim_{\al \rightarrow \infty} \:
(r^\al)^{2} \: R(x^\al, t^\al) \: = \: \infty$.

So suppose that $\lim_{\al \rightarrow \infty} \:
(r^\al)^{2} \: R(x^\al, t^\al) \: = \: \infty$.
We claim that $\lim_{\al \rightarrow \infty} \:
(r^\al)^{2} \: \inf R \Big|_{B \left( x^\al, t^\al,
r^\al \right)} \: = \: \infty$.
Suppose not.  Then there is some $C \in (0, \infty)$ so that
after passing to a subsequence, there are points
$x^{\al \prime} \in B \left( x^\al, t^\al, r^\al
 \right)$ with
$(r^\al)^{2} \: R(x^{\al \prime}, t^\al) < C$. Considering points
along the time-$t^\al$ geodesic segment from
$x^{\al \prime}$ to $x^\al$,
for large $\al$ we can find points
$x^{\al \prime \prime} \in B \left( x^\al, t^\al,
r^\al \right)$ with
$(r^\al)^{2} \: R(x^{\al \prime\prime}, t^\al) \: = \: 2C$.
Applying Lemma \ref{claim2II.4.2} at $(x^{\al \prime \prime}, t^\al)$,
or more precisely a version that applies along geodesics as in
Claim 2 of II.4.2, we obtain a contradiction to the assumption
that $\lim_{\al \rightarrow \infty} \:
(r^\al)^{2} \: R(x^\al, t^\al) \: = \: \infty$. In applying
Lemma \ref{claim2II.4.2} we use that
$r^\al < r(t^\al)$ and $\lim_{\al \rightarrow \infty} r(t^\al) \: = \: 0$,
giving
$\lim_{\al \rightarrow \infty} \left( r^\al \right)^{-2} \: t^\al \: = \:
\infty$, in order to say that
$\lim_{\al \rightarrow \infty}
\frac{\Phi \left( 2C (r^\al)^{-2} \right)}{ 2C (r^\al)^{-2}} \: = \: 0$.

Hence for large $\al$, every point $x \in B \left( x^\al, t^\al,
r^\al \right)$ is in the
center of a canonical neighborhood of size comparable to
$R(x, t^\al)^{- \: \frac12}$, which is small compared to $r^\al$.
On the other hand,
from Lemma \ref{Alexlemma2},
there is a ball $B^\prime$ of radius $\theta_0(w^{\prime}) \:
r^\al$ in
$B \left( x^\al, t^\al, r^\al \right)$ so that every subball of $B^\prime$ has
almost-Euclidean volume. This is a contradiction.

This proves the lemma.
\end{proof}

We continue with the proof of Proposition \ref{thingraph}.
By construction there is a sequence $w^\al \rightarrow 0$ so that
conditions (1) and (2) of Theorem \ref{thmII.7.4second} hold
with $\rho^\al(x^\al) = (t^\al)^{- \frac12} \: \rho(x^\al, t^\al)$.
Hence for large $\al$, $M^\al$ is
diffeomorphic to a graph manifold.  This is a contradiction to the choice of
the $M^\al$'s and proves the theorem.

We now give a proof that instead uses Theorem \ref{thmII.7.4}.
Let $d^\al$ denote the diameter of $M^\al$.
If we take $\rho^\al(x^\al) = (t^\al)^{- \frac12} \: \rho(x^\al, t^\al)$
then we can apply Theorem \ref{thmII.7.4} as long as
that the diameter statement in condition (1) of Theorem \ref{thmII.7.4} is
satisfied.
If it is not satisfied then there is some point $x^\al \in M^\al$ with
$\rho^\al(x^\al) \: > \:  d^\al$.
The sectional curvatures of $M^\al$ are bounded below by
$- \: \frac{1}{\rho^\al(x^\al)^2}$, and so are bounded below by
$- \: \frac{1}{(d^\al)^2}$.
If there is a subsequence  with
$\frac{\rho^\al(x^\al)}{d^\al} \le C < \infty$  then
\begin{equation}
\vol(M^\al) \: = \: \vol(B(x^\al, \rho^\al(x^\al))) \: \le \: w^\al \:
\rho^\al(x^\al)^3 \: \le \: w^\al \: C^3 \: (d^\al)^3
\end{equation}
and we can apply Theorem \ref{thmII.7.4} with $\rho^\al \: = \: d^\al$,
after redefining $w^\al$.
Thus we may assume that
$\lim_{\al \rightarrow \infty} \frac{\rho^\al(x^\al)}{d^\al} = \infty$.
If there is a subsequence with
$\frac{\vol(M^\al)}{(d^\al)^3}  \rightarrow 0$ then we
can apply Theorem \ref{thmII.7.4} with $\rho^\al \: = \: d^\al$.
Thus we may assume that $\frac{\vol(M^\al)}{(d^\al)^3}$ is bounded away from zero.
After rescaling the metric to make the diameter one, we are in a
noncollapsing situation
with the lower sectional curvature bound going to zero.
By the argument in the proof of Lemma \ref{condition} (which used
Corollary \ref{II.6.8}) there are uniform $L^\infty$-bounds on $\Rm(M^\al)$ and
its covariant derivatives.  After passing to a subsequence, there is a
limit $(M_\infty, g_\infty)$ in the smooth topology which is diffeomorphic
to $M^\al$ for large $\al$, and carries a metric of nonnegative sectional curvature.
As any boundary component of $M_\infty$ would have to have a neighborhood
of negative sectional curvature (see
the definition of $M_{thin}$ and
condition (2) of Theorem \ref{thmII.7.4}),
$M_\infty$ is closed. However, by construction $M^\al$ has no connected components
which are closed and admit metrics of nonnegative sectional curvature.
 This contradiction shows that the diameter statement in condition
(1) of Theorem \ref{thmII.7.4} is satisfied. The other conditions of
Theorem \ref{thmII.7.4} are satisfied as before.
\end{proof}

Thus for large $t$,  $\M_t^+$ has a decomposition into
a piece  $f(t) (\widehat{H}_1(t) \cup \ldots \widehat{H}_k(t))$, whose interior admits
a complete finite-volume hyperbolic metric, and the complement, which is a
graph manifold.

In addition, by Section \ref{inctori}
the cuspidal tori are incompressible in $\M_t^+$.
By Lemma \ref{reconstruct2},
the initial (connected) manifold $\M_0$ is diffeomorphic to a connected sum of
the connected components of $\M_t$, along with some possible additional connected sums
with a finite number of $S^1 \times S^2$'s and quotients of the round $S^3$.
This proves the geometrization conjecture of Appendix \ref{geom}.

\section{II.8. Alternative proof of cusp incompressibility} \label{II.8}

The goal of this section is to prove Perelman's Proposition II.8.2, which
gives a numerical characterization of
the geometric type of a compact $3$-manifold. It also contains an
independent proof of the incompressibility of the cuspidal ends
of the hyperbolic piece in the geometric decomposition.

We recall from Section \ref{I.2.2}
that $\lambda(g)$ is the first eigenvalue of
$-4 \triangle \: + \: R$, and can also be expressed as
\begin{equation} \label{Rayleigh}
\lambda(g) \: = \: \inf_{\Phi \in C^\infty(M) \: : \:
\Phi \neq 0} \frac{\int_M \left( 4 |\nabla \Phi|^2 \: + \:
R \: \Phi^2 \right) \: dV}{\int_M \Phi^2 \: dV}.
\end{equation}

From Lemma \ref{2/3}, if $g(\cdot)$ is a Ricci flow and
$\lambda(t) \: = \: \lambda(g(t))$ then
\begin{equation} \label{ll}
\frac{d}{dt} \lambda(t) \: \ge \: \frac23 \:
\lambda^2(t).
\end{equation}
From Lemma \ref{expandingprop},
$\lambda(t) \: V(t)^{\frac23}$ is nondecreasing
when it is nonpositive.

For any metric $g$, there are inequalities
\begin{equation} \label{test}
\min R \: \le \: \lambda(g) \: \le \: \frac{\int_M R \: dV}{\vol(M, g)},
\end{equation}
where the first inequality follows directly from (\ref{Rayleigh})
and the second inequality comes from using $1$ as a
test function in (\ref{Rayleigh}).


Perelman's proof of his Proposition II.8.2 uses the functional
$\lambda(g) V(g)^{\frac23}$. The functional
$R_{min} V(g)^{\frac23}$ plays a similar role.  For example, from
Corollary \ref{Rmono},
$\widehat{R}(t) = R_{min}(t) V(t)^{\frac23}$ is nondecreasing
when it is nonpositive. We first give a proof of an analog of
Proposition II.8.2 that uses $R_{min} V(g)^{\frac23}$ instead
of $\lambda(g) V(g)^{\frac23}$. The technical simplification is that
when
$R_{min} V(g)^{\frac23}$ is nonpositive, it is nondecreasing under a surgery,
as surgeries are only done in regions of large positive scalar curvature,
so $R_{min}$ doesn't change,
and a surgery reduces volume. (A possible extinction of a component
clearly doesn't change $R_{min} V(g)^{\frac23}$.)
We show that a minimal-volume hyperbolic submanifold of $M$
has incompressible tori, which gives a different approach to
Section \ref{inctori}.

Perelman's alternative approach to
Section \ref{inctori} uses the functional $\lambda(g) V(g)^{\frac23}$
instead of $R_{min} V(g)^{\frac23}$. Our use of $R_{min} V(g)^{\frac23}$ and
the sigma-invariant
$\sigma(M)$, instead of $\lambda(g) V(g)^{\frac23}$ and $\overline{\lambda}$,
is inspired by \cite{Anderson2}.

We then give the arguments using $\lambda(g) V(g)^{\frac23}$,
thereby proving Perelman's Proposition II.8.2.
The main technical difficulty is to
control how $\lambda(g) V(g)^{\frac23}$ changes under a surgery.

\subsection{The approach using the $\sigma$-invariant}

We first give some well-known results about the sigma-invariant.
We recall that the sigma-invariant of a closed connected
manifold $M$ of dimension $n \ge 3$ is given by
\begin{equation} \label{yamabe}
\sigma(M) \: = \: \sup_{{\mathcal C}} \inf_{g \in {\mathcal C}}
\frac{\int_M R(g) \: \dvol(g)}{\vol(M, g)^{\frac{n-2}{n}}},
\end{equation}
where ${\mathcal C}$ runs over the conformal classes of Riemannian
metrics on $M$.
From the solution to the Yamabe problem,
the infimum in (\ref{yamabe}) is realized by a metric of constant
scalar curvature in the given conformal class.
It follows that if $\sigma(M) > 0$ then $M$ admits a metric with
positive scalar curvature. Conversely, suppose that $M$ admits
a metric $g_0$ with positive scalar curvature. Let ${\mathcal C}$
be the conformal class containing $g_0$.  Then
\begin{equation}
\inf_{g \in {\mathcal C}}
\frac{\int_M R(g) \: \dvol(g)}{\vol(M, g)^{\frac{n-2}{n}}} \: = \:
\inf_{u > 0} \frac{
\int_M \left(
\frac{4(n-1)}{n-2} \:  |\nabla u|^2 \: + \:
R(g_0) \: u^2 \right) \dvol_M(g_0)
}{
\left( \int_M u^{\frac{2n}{n-2}} \: \dvol_M(g_0) \right)^{\frac{n-2}{n}}}
\end{equation}
is positive, in view of the Sobolev embedding theorem, and so $\sigma(M) > 0$.

We claim that if $\sigma(M) \le 0$ then
\begin{equation} \label{testt}
\sigma(M) \: = \: \sup_{g} R_{min}(g) \: V(g)^{\frac2n}.
\end{equation}
To see this, as
the infimum in (\ref{yamabe}) is realized by a metric of constant
scalar curvature in the given conformal class, it follows that
$\sigma(M) \: \le \: \sup_{g} R_{min}(g) \: V(g)^{\frac2n}$.
Now given a Riemannian metric $g$, the infimum in (\ref{yamabe})
within the corresponding conformal class ${\mathcal C}$
equals $\widetilde{R} \: V(\widetilde{g})^{\frac{2}{n}}$
for a metric $\widetilde{g} \: = \: u^{\frac{4}{n-2}} g$ with
constant scalar curvature $\widetilde{R}$. Then
\begin{equation}
\widetilde{R} \: V(\widetilde{g})^{\frac{2}{n}} \: = \:
\inf_{u > 0} \frac{
\int_M \left(
\frac{4(n-1)}{n-2} \:  |\nabla u|^2 \: + \:
R \: u^2 \right) \dvol_M
}{
\left( \int_M u^{\frac{2n}{n-2}} \: \dvol_M \right)^{\frac{n-2}{n}}}.
\end{equation}
As
\begin{equation}
\frac{ \int_M \left(
\frac{4(n-1)}{n-2} \:  |\nabla u|^2 \: + \:
R \: u^2 \right) \dvol_M
}{
\left( \int_M u^{\frac{2n}{n-2}} \: \dvol_M \right)^{\frac{n-2}{n}}}
\: \ge \: R_{min}(g) \: \frac{
\int_M u^2 \: \dvol_M
}{
\left( \int_M u^{\frac{2n}{n-2}} \: \dvol_M \right)^{\frac{n-2}{n}}}
\end{equation}
and $R_{min}(g) \le 0$,
Holder's inequality implies that
$\widetilde{R} \: V(\widetilde{g})^{\frac{2}{n}} \: \ge \:
{R}_{min}(g) \: V({g})^{\frac{2}{n}}$.
It follows that
$\sigma(M) \: \ge \: \sup_{g} R_{min}(g) \: V(g)^{\frac2n}$.

The next proposition
answers conjectures of Anderson \cite{Anderson3}.

\begin{proposition}[characterization for graph manifolds]
  \label{prop:kl-alternative-proof-cusp-incompressibility-characterization-graph-manifolds}
  \dcref{propYamabe}
  \group{Kleiner--Lott Long-Time Topology}
  \source{kleiner-lott}
 \label{propYamabe}
Let $M$ be a closed connected oriented $3$-manifold. \\
(a) If $\sigma(M) > 0$ then $M$ is
diffeomorphic to a connected sum of a finite number of
$S^1 \times S^2$'s and metric quotients of the round $S^3$.
Conversely, each such manifold has $\sigma(M) > 0$. \\
(b) $M$ is a graph manifold if and only if $\sigma(M)
\ge 0$. \\
(c) If $\sigma(M) < 0$
 then
$\left( - \: \frac23 \: \sigma(M)
\right)^{\frac32}$ is the minimum
of the numbers
$V$ with the following property : $M$ can be decomposed as a connected sum
of a finite collection of $S^1 \times S^2$'s, metric quotients of the round
$S^3$ and some other components, the union of which is denoted by
$M^\prime$, and there exists a (possibly disconnected) complete
finite-volume
manifold $N$ with constant sectional curvature $- \: \frac14$ and
volume $V$ which can be embedded in $M^\prime$ so that the
complement $M^\prime - N$ (if nonempty) is a graph manifold.

Moreover, if
$\vol(N) \: = \: \left( - \: \frac23 \: \sigma(M)
\right)^{\frac32}$
then the
cusps of $N$ (if any) are incompressible in $M^\prime$.
\end{proposition}

\begin{proof}
  \uses{lem:kl-double-sided-curvature-bound-thick-part-compactness-ricci-flow-surgery,lem:kl-double-sided-curvature-bound-thick-part-compactness-thick-thin-geometry}
If $\sigma(M) > 0$ then $M$ has a metric $g$ of positive scalar curvature.
From Lemmas \ref{pscempty} and \ref{finitetime}, $M$ is a connected sum of $S^1 \times S^2$'s and
metric quotients of the round $S^3$.
Conversely, if $M$ is a connected sum of $S^1 \times S^2$'s and
metric quotients of the round $S^3$ then $M$ admits a metric $g$ of
positive scalar curvature and so $\sigma(M) > 0$.

Now suppose that $\sigma(M) \le 0$.
If $M$ is a graph manifold then $M$ volume-collapses with bounded
curvature, so (\ref{testt}) implies that $\sigma(M) = 0$.

Suppose that $M$ is not a graph manifold. Suppose that we have
a given decomposition of $M$
as a connected sum
of a finite collection of $S^1 \times S^2$'s, metric quotients of the round
$S^3$ and some other components, the union of which is denoted by
$M^\prime$, and there exists a (possibly disconnected) finite-volume complete
manifold $N$ with constant sectional curvature $- \: \frac14$
which can be embedded in $M^\prime$ so that the
complement (if nonempty) is a graph manifold.
Let
$V_{hyp}$ denote the hyperbolic volume of $N$.  We do not
assume that the cusps of $N$ are incompressible in $M^\prime$.
For any $\epsilon > 0$, we claim that
there is a metric $g_\epsilon$ on $M$ with
$R \: \ge \: - \: 6 \cdot \frac14 \: - \: \epsilon$ and volume
$V(g_\epsilon) \: \le \: V_{hyp} + \epsilon$. This comes from
collapsing the graph manifold pieces, along with the
fact that the connected sum operation can be performed
while  decreasing the scalar curvature arbitrarily little and increasing
the volume arbitrarily little.
Then
$R_{min}(g_\epsilon) \: V(g_\epsilon)^{\frac23}
 \: \ge \: - \: \frac32 \: V_{hyp}^{2/3} \: - \: \const \:
\epsilon$. Thus
$\sigma(M) \: \ge \:
- \: \frac32 \: V_{hyp}^{2/3}$.

Let
$\widehat{V}$
denote the minimum of $V_{hyp}$ over
all such decompositions of $M$. (As the set of
volumes of complete finite-volume $3$-manifolds with constant curvature
$- \: \frac14$ is well-ordered, there is a minimum.)
Then
$\sigma(M) \: \ge \:
- \: \frac32 \: \widehat{V}^{2/3}$.

Next, take an arbitrary metric
$g_0$ on $M$ and consider the Ricci flow $g(t)$ with initial metric $g_0$.
From Sections \ref{hyprig} and  \ref{II.7.4}, there is a nonempty
manifold $N$ with a complete finite-volume metric
of constant curvature $- \: \frac14$
so that for large $t$, there is a decomposition $\M^+_t = M_1(t) \cup M_2(t)$ of the time-$t$ manifold,
where $M_1(t)$ is a graph manifold and $(M_2(t), \frac{1}{t} \: g(t) \big|_{M_2(t)})$
is close to a large
piece of $N$. In terms of condition (c) of Proposition \ref{propYamabe}, we will think
of $M^\prime$ as being $\M^+_t$.
Because of the presence of $N$, we know that $t \: R_{min}(t) \: \le \: - \:
\frac32 \: + \: \epsilon(t)$ and
$V(t) \: \ge \:
t^{2/3} \: V_{hyp}(N) \: - \: \epsilon(t)$ for a function $\epsilon(t)$
with $\lim_{t \rightarrow \infty}
\epsilon(t) \: = \: 0$. The monotonicity of
$R_{min}(t) \: V^{2/3}(t)$, even through surgeries, implies that
\begin{equation}
R_{min}(g_0) \: V^{2/3}(g_0) \: \le \: - \: \frac32 \: V_{hyp}(N)^{2/3}
 \: \le \: - \: \frac32 \: \widehat{V}^{2/3}.
 \end{equation}
Thus $\sigma(M) \: \le \: - \: \frac32 \: \widehat{V}^{2/3}$.

This shows that
$\sigma(M) \: = \: - \: \frac32 \: \widehat{V}^{2/3}$.
Now take a decomposition of $M$ as in
condition (c) of Proposition \ref{propYamabe}, with
$V_{hyp}(N) \: = \: \widehat{V}$. We claim that the cuspidal $2$-tori of $N$
are incompressible in $M^\prime$. If not
then there would be a metric  $g$ on $M$ with $R(g) \: \ge \: - \: \frac32$ and
$\vol(g) \: < \: V_{hyp}(N)$ \cite[Pf. of Theorem 2.9]{Anderson}.
This would contradict the fact that
$\sigma(M) \: = \: - \: \frac32 \: V_{hyp}(N)^{2/3}$.
\end{proof}

\subsection{The approach using the $\overline{\lambda}$-invariant}

\begin{proposition}[characterization for graph manifolds]
  \label{prop:kl-alternative-proof-cusp-incompressibility-characterization-graph-manifolds-2}
  \dcref{propII.8.2}
  \group{Kleiner--Lott Long-Time Topology}
  \source{kleiner-lott}
 (cf. II.8.2) \label{propII.8.2}
Let $M$ be a closed connected oriented $3$-manifold. \\
(a) If $M$ admits a metric $g$ with $\lambda(g) > 0$ then it is
diffeomorphic to a connected sum of a finite number of
$S^1 \times S^2$'s and metric quotients of the round $S^3$.
Conversely, each such manifold admits a metric $g$ with
$\lambda(g) > 0$. \\
(b) Suppose that $M$ does not admit any metric $g$ with
$\lambda(g) > 0$. Let $\overline{\lambda}$ denote the
supremum of $\lambda(g) V(g)^{\frac23}$ over all metrics $g$ on $M$.
Then $M$ is a graph manifold if and only if $\overline{\lambda}
= 0$. \\
(c) Suppose that $M$ does not admit any metric $g$ with
$\lambda(g) > 0$, and $\overline{\lambda} < 0$.
Then
$\left( - \: \frac23 \: \overline{\lambda}
\right)^{\frac32}$ is the minimum
of the numbers
$V$ with the following property : $M$ can be decomposed as a connected sum
of a finite collection of $S^1 \times S^2$'s, metric quotients of the round
$S^3$ and some other components, the union of which is denoted by
$M^\prime$, and there exists a (possibly disconnected) complete
manifold $N$ with constant sectional curvature $- \: \frac14$ and
volume $V$ which can be embedded in $M^\prime$ so that the
complement $M^\prime - N$ (if nonempty) is a graph manifold.

Moreover, if
$\vol(N) \: = \: \left( - \: \frac23 \: \overline{\lambda}
\right)^{\frac32}$
then the cusps of $N$ (if any) are incompressible in $M^\prime$.
\end{proposition}


\begin{proof}
  \uses{lem:kl-alternative-proof-cusp-incompressibility-estimates-spectral-geometry,lem:kl-alternative-proof-cusp-incompressibility-positivity-entropy,lem:kl-double-sided-curvature-bound-thick-part-compactness-ricci-flow-surgery,lem:kl-establishing-noncollapsing-presence-surgery-vanishing-noncollapsing,prop:kl-statement-existence-theorem-ricci-flow-surgery-monotonicity-kappa-solutions}
We first give the argument for Proposition \ref{propII.8.2} under the
pretense that all Ricci flows are smooth, except for possible
extinction of components. (Of course this is not the
case, but it will allow us to present the main idea of the proof.)

If $\lambda(g) > 0$ for some metric $g$ then from (\ref{ll}), the Ricci flow
starting from $g$ will become extinct within time
$\frac{3}{2\lambda(g)}$. Hence Lemma \ref{finitetime} applies.
Conversely, if $M$ is a connected sum of $S^1 \times S^2$'s and
metric quotients of the round $S^3$ then $M$ admits a metric $g$ of
positive scalar curvature.  From (\ref{test}), $\lambda(g) > 0$.

Now suppose that $M$ does
not admit any metric $g$ with $\lambda(g) > 0$.
If $M$ is a graph manifold then $M$ volume-collapses with bounded
curvature, so (\ref{test}) implies that $\overline{\lambda} = 0$.

Suppose that $M$ is not a graph manifold. Suppose that we have
a given decomposition of $M$
as a connected sum
of a finite collection of $S^1 \times S^2$'s, metric quotients of the round
$S^3$ and some other components, the union of which is denoted by
$M^\prime$, and there exists a (possibly disconnected) complete
manifold $N$ with constant sectional curvature $- \: \frac14$
which can be embedded in $M^\prime$ so that the
complement (if nonempty) is a graph manifold.
Let
$V_{hyp}$ denote the hyperbolic volume of $N$.  We do not
assume that the cusps of $N$ are incompressible in $M^\prime$.
For any $\epsilon > 0$, we claim that
there is a metric $g_\epsilon$ on $M$ with
$R \: \ge \: - \: 6 \cdot \frac14 \: - \: \epsilon$ and volume
$V(g_\epsilon) \: \le \: V_{hyp} + \epsilon$. This comes from
collapsing the graph manifold pieces, along with the
fact that the connected sum operation can be performed
while  decreasing the scalar curvature arbitrarily little and increasing
the volume arbitrarily little.
Then (\ref{test}) implies that
$\lambda(g_\epsilon) \: V(g_\epsilon)^{\frac23}
 \: \ge \: - \: \frac32 \: V_{hyp}^{2/3} \: - \: \const \:
\epsilon$. Thus
$\overline{\lambda} \: \ge \:
- \: \frac32 \: V_{hyp}^{2/3}$.

Let $\widehat{V}$ denote the minimum of $V_{hyp}$ over
all such decompositions of $M$. (As the set of
volumes of complete finite-volume $3$-manifolds with constant curvature
$- \: \frac14$ is well-ordered, there is a minimum.)
Then
$\overline{\lambda} \: \ge \:
- \: \frac32 \: \widehat{V}^{2/3}$.

Next, take an arbitrary metric
$g_0$ on $M$ and consider the Ricci flow $g(t)$ with initial metric $g_0$.
From Sections \ref{hyprig} and \ref{II.7.4}, there is a
nonempty manifold $N$ with a finite-volume complete metric
of constant curvature $- \: \frac14$
so that for large $t$, there is a decomposition of the time-$t$ manifold
$\M^+_t = M_1(t) \cup M_2(t)$
where $M_1(t)$ is a graph manifold and $(M_2(t), \frac{1}{t} \: g(t) \big|_{M_2(t)})$
is close to a large
piece of $N$.
As $N$ has finite volume, a constant function on $N$ is square-integrable
and so  $\inf \spec(- \triangle_N) \: = \: 0$.
Equivalently,
\begin{equation}
\inf_{f \in C^\infty_c(N), f \neq 0}
\frac{\int_N |\nabla f|^2 \: \dvol_N}{\int_N f^2 \: \dvol_N} \: = \: 0.
\end{equation}
Taking an appropriate test function $\Phi$ on $M$ with compact
support in $M_2(t)$
gives $t \: \lambda(t) \: \le \: - \:
\frac32 \: + \: \epsilon_1(t)$, with $\lim_{t \rightarrow \infty}
\epsilon_1(t) \: = \: 0$. In terms of condition (c) of Proposition \ref{propII.8.2},
we will think of $M^\prime$ as being $\M^+_t$.
From the presence of $N$, we know that $V(t) \: \ge \:
t^{2/3} \: V_{hyp}(N) \: - \: \epsilon_2(t)$, with $\lim_{t \rightarrow \infty}
\epsilon_2(t) \: = \: 0$. As we are assuming that the Ricci flow is nonsingular,
the monotonicity of
$\lambda(t) \: V^{2/3}(t)$ implies that
\begin{equation}
\lambda(g_0) \: V^{2/3}(g_0) \: \le \: - \: \frac32 \: V_{hyp}(N)^{2/3}
 \: \le \: - \: \frac32 \: \widehat{V}^{2/3}.
 \end{equation}
Thus $\overline{\lambda} \: \le \: - \: \frac32 \: \widehat{V}^{2/3}$.

This shows that
$\overline{\lambda} \: = \: - \: \frac32 \: \widehat{V}^{2/3}$.
Now take a decomposition of $M$ as in condition (c) of Proposition \ref{propII.8.2},
with
$V_{hyp}(N) \: = \: \widehat{V}$. We claim that the cuspidal $2$-tori of $N$ are
incompressible in $M^\prime$. If not
then there would be a metric $g$ on $M$ with $R(g) \: \ge \: - \: \frac32$ and
$\vol(g) \: < \: V_{hyp}(N)$ \cite[Pf. of Theorem 2.9]{Anderson}.
Using (\ref{test}), one
would obtain a contradiction to the fact that
$\overline{\lambda} \: = \: - \: \frac32 \: V_{hyp}(N)^{2/3}$.

To handle the behaviour of $\lambda(t) \: V^{\frac23}(t)$ under
Ricci flows with surgery, we
first state a couple of general facts about
Schr\"odinger operators.

\begin{lemma}[estimates for spectral geometry]
  \label{lem:kl-alternative-proof-cusp-incompressibility-estimates-spectral-geometry}
  \dcref{newprop1}
  \group{Kleiner--Lott Long-Time Topology}
  \source{kleiner-lott}
 \label{newprop1}
Given a closed Riemannian manifold $M$, let $X$ be a codimension-$0$
submanifold-with-boundary of $M$. Given $R \in C^\infty(M)$, let
$\lambda_M$ be the lowest eigenvalue of $-4\triangle \: + \: R$
on $M$, with corresponding eigenfunction $\psi$.
Let $\lambda_X$ be the lowest eigenvalue of the corresponding
operator on $X$, with Dirichlet boundary conditions, and similarly
for $\lambda_{M - \Int(X)}$. Then for all
$\eta \in C^\infty_c(\Int(X))$, we have
\begin{equation} \label{eig1}
\lambda_M \: \le \: \min(\lambda_X, \lambda_{M - \Int(X)})
\end{equation}
and
\begin{equation} \label{eig2}
\lambda_X \: \le \: \lambda_M \: + \: 4 \:
\frac{\int_M |\nabla \eta|^2 \: \psi^2 \: dV}{\int_M \eta^2 \:
\psi^2 \: dV}.
\end{equation}
\end{lemma}
\begin{proof}
Equation (\ref{eig1}) follows from Dirichlet-Neumann bracketing
\cite[Chapter XIII.15]{Reed-Simon}. To prove (\ref{eig2}),
$\eta \psi$ is supported in $\Int(X)$ and so
\begin{align}
\lambda_X \: & \le \: \frac{\int_M \left( 4 |\nabla(\eta \psi)|^2 \: + \:
R \: \eta^2 \psi^2 \right) \: dV}{\int_M \eta^2 \psi^2 \: dV} \\
& = \: \frac{\int_M \left( 4 |\nabla \eta|^2 \psi^2 \: + \:
8 \langle \nabla\eta, \nabla \psi \rangle \eta \psi \: + \:
4 \eta^2 |\nabla \psi|^2 \: + \:
R \: \eta^2 \psi^2 \right) \: dV}{\int_M \eta^2 \psi^2 \: dV}. \notag
\end{align}
As $- 4 \triangle \psi \: + \: R \psi \: = \: \lambda_M \psi$, we have

\begin{align}
\lambda_M \int_M \eta^2 \psi^2 dV \: & = \:
- \: 4 \: \int_M \eta^2 \psi \triangle \psi \: dV \: + \:
\int_M R \eta^2 \psi^2 \: dV \\
& = \:
4 \: \int_M \langle \nabla(\eta^2 \psi), \nabla \psi \rangle \: dV \:
+ \: \int_M R \eta^2 \psi^2 \: dV \notag \\
& = \: 4 \: \int_M \eta^2 \: |\nabla \psi|^2 \: dV \: \: + \:
8 \int_M \langle \nabla\eta, \nabla \psi \rangle \eta \psi \: dV
\: + \: \int_M R \eta^2 \psi^2 \: dV. \notag
\end{align}
Equation (\ref{eig2}) follows.
\end{proof}

The next result is an Agmon-type estimate.

\begin{lemma}[positivity for entropy]
  \label{lem:kl-alternative-proof-cusp-incompressibility-positivity-entropy}
  \dcref{newprop2}
  \group{Kleiner--Lott Long-Time Topology}
  \source{kleiner-lott}
  \uses{lem:kl-alternative-proof-cusp-incompressibility-estimates-spectral-geometry}
 \label{newprop2}
With the notation of Lemma \ref{newprop1},
given a nonnegative function $\phi \in C^\infty(M)$,
suppose that $f \in C^\infty(M)$ satisfies
\begin{equation}
4 |\nabla f|^2 \: \le \: R \: - \: \lambda_M \: - \: c
\end{equation}
on $\supp(\phi)$, for some $c > 0$. Then
\begin{equation} \label{eig3}
\|e^f \phi \psi\|_2 \: \le \: 4 \: c^{-1} \: \left( \|e^f \triangle \phi\|_\infty
\: + \: \|e^f \nabla \phi\|_\infty (\lambda_M \: - \: \min R)^{1/2} \right)
\: \|\psi\|_2.
\end{equation}
\end{lemma}
\begin{proof}
Put $H \: = \: - \: 4 \triangle \: + \: R$.
By assumption,
\begin{equation}
\phi (R \: - \: 4 |\nabla f|^2 \: - \: \lambda_M) \: \phi \: \ge \: c \: \phi^2
\end{equation}
and so there is an
 inequality of operators on $L^2(M)$ :
\begin{equation}
\phi (H \: - \: 4 |\nabla f|^2 \: - \: \lambda_M) \: \phi \: = \:
4 \phi d^* d \phi \: + \:
\phi (R \: - \: 4 |\nabla f|^2 \: - \: \lambda_M) \: \phi \:
\ge \: c \: \phi^2.
\end{equation}
In particular,
\begin{equation} \label{part1}
\int_M e^f \psi\phi \: (H \: - \: 4 |\nabla f|^2 \: - \: \lambda_M) \: \phi
e^f \psi dV \: \ge \: c \int_M \phi^2 e^{2f} \psi^2 \: dV.
\end{equation}

For $\rho \in C^\infty(M)$,
\begin{equation} \label{Hf}
e^f H (e^{-f} \rho) \: = \: H \rho \: + \: 4 \: \nabla \cdot
((\nabla f) \rho) \: + \: 4 \: \langle \nabla f, \nabla \rho \rangle
\: - \: 4 \: |\nabla f|^2 \rho
\end{equation}
and so
\begin{equation}
\int_M \rho e^f H(e^{-f} \rho) \: dV \: = \:
\int_M \rho (H \: - \: 4 \: |\nabla f|^2) \rho \: dV.
\end{equation}
Taking $\rho \: = \: e^f \phi \psi$ gives
\begin{equation} \label{part2}
\int_M e^{2f} \phi \psi H(\phi \psi) \: dV \: = \:
\int_M e^f \phi \psi (H \: - \: 4 \:
|\nabla f|^2)  e^{f} \phi \psi \: dV.
\end{equation}
From (\ref{part1}) and (\ref{part2}),
\begin{equation}
c \|e^f \phi \psi \|_2^2 \: \le \:
\int_M e^{2f} \phi \psi (H \: - \: \lambda_M) (\phi \psi) \: dV \: = \:
\int_M e^{2f} \phi \psi [H, \phi] \psi \: dV \: = \:
\langle e^{f} \phi \psi, e^{f} [H, \phi] \psi \rangle_2.
\end{equation}
Thus
\begin{equation}
c \|e^f \phi \psi \|_2 \: \le \:
\|e^{f} [H, \phi] \psi\|_2.
\end{equation}
Now
\begin{equation}
e^{f} [H, \phi] \psi \: = \: - \: 4 \: e^f (\triangle \phi) \psi \: - \:
8 \: e^f \langle \nabla\phi, \nabla \psi \rangle.
\end{equation}
Then
\begin{equation}
\|e^{f} [H, \phi] \psi\|_2 \: \le \:
4 \: \|e^f \triangle \phi\|_\infty \: \|\psi\|_2 \: + \: 8 \:
\|e^f \nabla\phi\|_\infty \: \|\nabla \psi\|_2.
\end{equation}
Finally,
\begin{equation}
4 \: \|\nabla \psi\|_2^2 \: = \: \int_M (\lambda_M \: - \: R) \: \psi^2 \:
dV
\end{equation}
and so
\begin{equation}
2 \|\nabla \psi\|_2 \: \le \: (\lambda_M \: - \: \min R)^{1/2} \:
\|\psi\|_2.
\end{equation}
This proves the lemma.
\end{proof}

Clearly Lemma \ref{newprop2} is also true if
$f$ is just assumed to be Lipschitz-regular.

We now apply Lemmas \ref{newprop1} and \ref{newprop2}
to a Ricci flow with surgery.
A singularity caused by extinction of a component will not be a
problem, so
let $T_0$ be a surgery time  and let
$M_+ = \M_{T_0}^+$ be the postsurgery manifold.  We will write
$\lambda^+$ instead of $\lambda_{M_+}$. Let $M_{cap} =
\M_{T_0}^+ - (\M_{T_0}^+ \cap \M_{T_0}^-)$ be the added
caps and put $X \: = \: \overline{M_+ - M_{cap}} =
\overline{\M_{T_0}^+ \cap \M_{T_0}^-}$.
For simplicity,
let us assume that $M_{cap}$ has a single component; the argument in
the general case is similar.
From the nature of the surgery procedure, the surgery is done in an
$\epsilon$-horn extending from $\Omega_\rho$, where
$\rho \: = \: \delta(T_0) r(T_0)$.
In fact, because of the canonical
neighborhood assumption, we can extend the $\epsilon$-horn inward
until $R \sim r(T_0)^{-2}$. Appplying (\ref{Rayleigh}) with a test function
supported in an $\epsilon$-tube near this inner boundary, it follows that
$\lambda^+ \: \le \: c^\prime \: r(T_0)^{-2}$ for some universal constant
$c^\prime >> 1$.


In what follows we take $\delta(T_0)$ to be small.
As $R$ is much greater than $r(T_0)^{-2}$ on $M_{cap}$, it follows that
$\lambda_{M - \Int(X)}$ is much greater than $r(T_0)^{-2}$. Then from
(\ref{eig1}), $\lambda^+ \: \le \: \lambda_X$. We can apply
(\ref{eig2}) to get an inequality the other way.
We take the function $\eta$ to interpolate from being $1$ outside of the
$h(T_0)$-neighborhood
$N_h M_{cap}$ of $M_{cap}$, to being $0$ on $M_{cap}$.
In terms of the normalized eigenfunction $\psi$ on $M_+$,
this gives a bound of the form
\begin{equation} \label{bound11}
\lambda_X \: \le \: \lambda^+ \: + \: \const \: h(T_0)^{-2} \:
\frac{\int_{N_h M_{cap}} \psi^2 \: dV}{1 - \int_{N_h M_{cap}} \psi^2 \: dV}.
\end{equation}

We now wish to show that $\int_{N_h M_{cap}} \psi^2 \: dV$ is small.
For this we apply Lemma \ref{newprop2} with
$c \: = \: c^\prime \: r(T_0)^{-2}$. Take an $\epsilon$-tube $U$,
in the $\epsilon$-horn,
whose center has scalar curvature roughly $200 \: c^\prime \: r(T_0)^{-2}$
and which is the closest tube to the cap with this property.
Let $x \: : \: U \rightarrow
(- \: \epsilon^{-1}, \epsilon^{-1})$
be the longitudinal parametrization of the tube,
which we take to be increasing in the direction of the surgery cap.
Let $\Phi \: : \: (-1, 1) \rightarrow [0,1]$
be a fixed nondecreasing smooth function which is zero on $(-1, 1/4)$ and one on
$(1/2, 1)$. Put $\phi \: = \: \Phi \circ x$ on $U$.
Extend $\phi$ to $M_+$ by making it zero to the left of $U$ and one to the right of $U$,
where ``right of $U$'' means the connected component of $M_+ - U$ containing the
surgery cap.
Dimensionally, $|\nabla \phi|_\infty \: \le \: \const \:
r(T_0)^{-1}$ and $|\triangle \phi|_\infty \: \le \: \const \:  r(T_0)^{-2}$.
Define a function $f$ to the right of $x^{-1}(0)$ by setting it to be the distance from $x^{-1}(0)$
with respect to the metric $\frac14 \: (R - \lambda^+ - c) \: g_{M_+}$.
(Note that to the right of $x^{-1}(0)$, we have
$R \: \ge \: 200 c^\prime r(T_0)^{-2} \: \ge \: \lambda^+ + c$.)
Then
equation (\ref{eig3}) gives
$|e^f \phi \psi|_2 \: \le \: \const (T_0)$. The point is that $\const (T_0)$ is independent
of the (small) surgery parameter $\delta(T_0)$.

Hence
\begin{equation} \label{bound12}
\int_{N_h M_{cap}} \psi^2 \: dV \: \le \:
\left( \sup_{N_h M_{cap}} e^{-2f} \right) \:
\int_{N_h M_{cap}} \: e^{2f} \psi^2 \: dV \: \le \:
\const \: \sup_{N_h M_{cap}} e^{-2f}.
\end{equation}
To estimate $\sup_{N_h M_{cap}} e^{-2f}$, we use the fact that
the $\epsilon$-horn consists of a
sequence of $\epsilon$-tubes stacked together.  In the
region of $M_+$ from $x^{-1}(0)$ to the surgery cap, the scalar curvature
ranges from roughly $200 \: c^\prime \: r(T_0)^{-2}$
to $h(T_0)^{-2}$.
On a given
$\epsilon$-tube, if $\epsilon$ is sufficiently small then
the ratio of the scalar curvatures between the
two ends is bounded by $e$. Hence in going
from $x^{-1}(0)$ to the surgery cap,
one must cross at least $N$ disjoint $\epsilon$-tubes,
with $e^N \: = \: \frac{1}{200 c^\prime} \: r(T_0)^2 \: h(T_0)^{-2}$.
Traversing a given $\epsilon$-tube (say of
radius $r^\prime$) in going towards the surgery cap, $f$ increases by
roughly $\const \: \int_{- \epsilon^{-1}r^\prime}^{\epsilon^{-1}r^\prime}
(r^\prime)^{-1} \: ds$, which is $\const \epsilon^{-1}$. Hence near
the surgery cap, we have
\begin{equation}
\sup_{N_h M_{cap}} e^{-2f} \: \le \: \const \: e^{-\const N \epsilon^{-1}}
\: = \: \const \: (r(T_0)^2 \: h(T_0)^{-2})^{-\const \epsilon^{-1}}.
\end{equation}
Combining this with (\ref{bound11}) and (\ref{bound12}), we obtain
\begin{equation}
\lambda_X \: \le \: \lambda^+ \: + \: \const \: h(T_0)^{-2} \:
(r(T_0)^2 \: h(T_0)^{-2})^{- \const \epsilon^{-1}}.
\end{equation}


By making a single redefinition of $\epsilon$,
we can ensure that
$\lambda_X \: \le \: \lambda^+ \: + \: \const \: h(T_0)^{4}$. The last
constant will depend on $r(T_0)$ but is independent of $\delta(T_0)$.
Thus if $\delta(T_0)$ is small enough, we can ensure that
$|\lambda_X \: - \: \lambda^+|$ is small in comparison to the
volume change $V^-(T_0) - V^+(T_0)$, which is comparable to $h(T_0)^3$.


If $\lambda^-$ is the smallest eigenvalue of $- \:4 \: \triangle \: + \: R$
on the presurgery manifold $\M_t$, for $t$ slightly less than $T_0$,
then we can estimate $|\lambda_X \: - \: \lambda^-|$ in a similar way.
Hence for an arbitrary positive continuous function $\xi(t)$,
we can make the parameters $\overline{\delta}_j$ of
Proposition \ref{surgeryflowexists} small enough to ensure that
\begin{equation}
|\lambda^+(T_0) - \lambda^-(T_0)| \: \le \: \xi(T_0) \:
(V^-(T_0) - V^+(T_0))
\end{equation}
for a surgery at time $T_0$.

We now redo the argument for the proposition, as given above in the
surgery-free case, in the presence of surgeries.
Suppose first that $\lambda(g_0) > 0$ for some metric $g_0$ on $M$.
After possible rescaling, we can assume that $g_0$ is the initial
condition for a Ricci flow with surgery $(\M, g(\cdot))$, with normalized initial condition.
Using the lower scalar curvature bound of Lemma \ref{Rev} and the Ricci
flow equation,
the volume on the time interval $[0, \frac{3}{\lambda(0)}]$ has an {\it a priori} upper
bound of the form $\const V(0)$. As a surgery at time $T_0$ removes a
volume comparable to $h(T_0)^3$, we have
$\sum h(T_0)^3 \: \le \: \const V(0)$, where the sum is over the
surgeries and $T_0$ denotes the
surgery time.  From the above discussion,
the change in $\lambda$ due to the
surgeries is bounded below by $-\: \const \sum_{T_0} h(T_0)^4$.
Then the decrease in $\lambda$ due to surgeries on the time interval
$[0, \frac{3}{\lambda(0)}]$ is bounded above by
\begin{equation}
\const \sum_{T_0 \in [0, \frac{3}{\lambda(0)}]} h(T_0)^4 \: \le \:  \const \:
\left( \sup_{t \in [0, \frac{3}{\lambda(0)}]} h(t) \right) \: V(0).
\end{equation}
By choosing the function $\delta(t)$ to be sufficiently small,
the decrease in $\lambda$ due to surgeries
is not enough to prevent the blowup of $\lambda$ on the time
interval $[0, \frac{3}{\lambda(0)}]$ coming from the increase of
$\lambda$ between the surgeries. Hence the solution goes extinct.

Now suppose that $M$ does not admit a metric $g$ with $\lambda(g) > 0$.
Again, if $M$ is a graph manifold then $\overline{\lambda} \: = \: 0$.

Suppose that $M$ is not a graph manifold.
As before, $\overline{\lambda} \: \ge \: - \: \frac32 \: \widehat{V}^{2/3}$.
Given an initial metric $g_0$,
we wish to show that by choosing the function $\delta(t)$
small enough we can make the function $\lambda(t) V^{2/3}(t)$
arbitrarily close to being nondecreasing. To see this, we consider
the effect of a surgery on $\lambda(t) V^{2/3}(t)$. Upon performing a surgery
the volume decreases, which in itself cannot decrease
$\lambda(t) V^{2/3}(t)$. (We are using the fact that $\lambda(t)$ is
nonpositive.)
Then from the above discussion, the
change in $\lambda(t) V^{2/3}(t)$ due from a surgery at time $T_0$,
is bounded below by $- \: \xi(T_0) \: (V^-(T_0) - V^+(T_0)) \:
V^-(T_0)^{2/3}$.
With normalized initial conditions,
we have an {\it a priori} upper bound on $V(t)$ in terms of $V(0)$ and $t$.
Over any time interval $[T_1, T_2]$, we must have
\begin{equation}
\sum_{T_0 \in [T_1, T_2]} \left( V^-(T_0) - V^+(T_0) \right) \: \le \:
\sup_{t \in [T_1,T_2]} V(t),
\end{equation}
where the sum is over the surgeries in the
interval $[T_1,T_2]$.
Then along with the monotonicity of $\lambda(t) V^{2/3}(t)$ in between the
 surgery times,  by choosing the function
$\delta(t)$ appropriately we can ensure that
for any
$\sigma > 0$
there is a Ricci flow with $(r, \delta)$-cutoff
starting from $g_0$ so that
$\lambda(g_0) \: V^{2/3}(g_0) \: \le \: \lambda(t) \: V^{2/3}(t) \: + \:
\sigma$ for all $t$.
It follows that
$\overline{\lambda} \: \le \: - \: \frac32 \: \widehat{V}^{2/3}$.

This shows that $\overline{\lambda} \: = \: - \: \frac32 \: \widehat{V}^{2/3}$.
The same argument as before shows that if we have a decomposition
with the hyperbolic volume of $N$ equal to  $\widehat{V}$ then the cusps of $N$
(if any) are incompressible in $M^\prime$.
\end{proof}

\begin{remark}[positivity for scalar curvature]
  \label{rem:kl-alternative-proof-cusp-incompressibility-positivity-scalar-curvature}
  \dcref{alternative-proof-cusp-incompressibility:remark-5}
  \group{Kleiner--Lott Long-Time Topology}
  \source{kleiner-lott}

It follows that if the three-manifold $M$ does not admit a metric of positive scalar curvature then
$\sigma(M) \: = \: \overline{\lambda}$. In fact, this is true in any dimension
$n \ge 3$ \cite{Akutagawa-Ishida-Lebrun}.
\end{remark}
